%============================================================
%  Bd. 06 - Axiomatischer Aufbau der Halbverbände und Verbände
%============================================================

\documentclass[main.tex]{subfiles}

\ifSubfilesClassLoaded{
  \BandSetupStandalone{B06}
}{
}

\title{Bd. 06 - Axiomatischer Aufbau der Halbverbände und Verbände}
\author{Martin Kunze}
\date{}
\setcounter{file}{6}

\begin{document}

\title{Bd. 06 - Axiomatischer Aufbau der Halbverbände und Verbände}
\author{Martin Kunze}
\date{}
\hypersetup{pageanchor=false}
\maketitle
\hypersetup{pageanchor=true}
\setcounter{tocdepth}{3}
\tableofcontents
\setcounter{file}{6}
\renewcommand{\FormulaBandID}{6}

\chapter{Algebraische Halbverbände}

Band~03 entwickelt partielle Ordnungen sowie obere und untere Schranken,
Infima und Suprema. Der vorliegende Band geht den umgekehrten Weg: Eine
einzige binäre Operation wird axiomatisch festgelegt, und aus ihren
algebraischen Gesetzen werden zwei zueinander duale partielle Ordnungen
gewonnen.

Die Bezeichnung \emph{Halbgitter} wird synonym zu \emph{Halbverband}
verwendet. Algebraisch gibt es keinen Unterschied zwischen einem
Supremums- und einem Infimums-Halbverband: Beide sind assoziative,
kommutative und idempotente binäre Operationen. Der Unterschied entsteht erst
durch die Orientierung der induzierten Relation. Für die allgemeine Operation
verwenden wir \(\NeutralSemiOp\); in der supremalen Lesart schreiben wir
\(\JoinOp\), in der infimalen Lesart \(\MeetOp\).

\section{Allgemeine axiomatische Definition}

Wie in den Bänden~03 und~04 werden die Strukturaxiome zunächst einzeln
registriert und anschließend in einer Definition gesammelt. In allen vier
Axiomen bezeichnet \(\NeutralSemiOp\) dieselbe Operation.

\FormulaAxiomDeltaK[Abgeschlossenheit einer Halbverbandsoperation]{%
  x,y\in A\vdash x\NeutralSemiOp{}y\in A%
}{SemilatticeClosureLaw}{%
  \DeltaRow{Mengen}{A\dsep x\dsep y}
  \DeltaRow{Funktionen}{\NeutralSemiOp}
}
\par\smallskip

\FormulaAxiomDeltaK[Assoziativität einer Halbverbandsoperation]{%
  x,y,z\in A\vdash
  \bigl(x\NeutralSemiOp{}y\bigr)\NeutralSemiOp{}z
    =x\NeutralSemiOp\bigl(y\NeutralSemiOp{}z\bigr)%
}{SemilatticeAssociativityLaw}{%
  \DeltaRow{Mengen}{A\dsep x\dsep y\dsep z}
  \DeltaRow{Funktionen}{\NeutralSemiOp}
}
\par\smallskip

\FormulaAxiomDeltaK[Kommutativität einer Halbverbandsoperation]{%
  x,y\in A\vdash x\NeutralSemiOp{}y=y\NeutralSemiOp{}x%
}{SemilatticeCommutativityLaw}{%
  \DeltaRow{Mengen}{A\dsep x\dsep y}
  \DeltaRow{Funktionen}{\NeutralSemiOp}
}
\par\smallskip

\FormulaAxiomDeltaK[Idempotenz einer Halbverbandsoperation]{%
  x\in A\vdash x\NeutralSemiOp{}x=x%
}{SemilatticeIdempotenceLaw}{%
  \DeltaRow{Mengen}{A\dsep x}
  \DeltaRow{Funktionen}{\NeutralSemiOp}
}

\FormulaDefDeltaK[Begriff des Halbverbandes]{%
  \Semi{\NeutralSemiOp,A}%
}{JoinSemilattice}{%
  \DeltaRow{Mengen}{A\dsep x\dsep y\dsep z}
  \DeltaRow{Funktionen}{\NeutralSemiOp}
  \DeltaRow{\textbf{Axiome}}{}
  \DeltaRow{Abgeschlossenheit}{%
    x,y\in A\vdash x\NeutralSemiOp{}y\in A}
    [\FormulaRefAuto{SemilatticeClosureLaw}[ax]]
  \DeltaRow{Assoziativität}{%
    \begin{aligned}[t]
      &x,y,z\in A\vdash{}\\
      &\bigl(x\NeutralSemiOp{}y\bigr)\NeutralSemiOp{}z
        =x\NeutralSemiOp\bigl(y\NeutralSemiOp{}z\bigr)
    \end{aligned}}
    [\FormulaRefAuto{SemilatticeAssociativityLaw}[ax]]
  \DeltaRow{Kommutativität}{%
    x,y\in A\vdash x\NeutralSemiOp{}y=y\NeutralSemiOp{}x}
    [\FormulaRefAuto{SemilatticeCommutativityLaw}[ax]]
  \DeltaRow{Idempotenz}{%
    x\in A\vdash x\NeutralSemiOp{}x=x}
    [\FormulaRefAuto{SemilatticeIdempotenceLaw}[ax]]
  \DeltaRow{\textbf{Neues Symbol}}{}
  \DeltaPrem{Halbverbände}{\Semi{\NeutralSemiOp,A}}
}

Die folgenden Zugriffsaxiome verbinden das neue Strukturprädikat mit den
vier soeben registrierten Gesetzen. Spätere Sätze zitieren jeweils genau das
benötigte Zugriffsaxiom und nicht die zusammenfassende Definition.

\FormulaAxiomDeltaK[Zugriff auf die Abgeschlossenheit eines Halbverbandes]{%
  \Semi{\NeutralSemiOp,A}\dsep x,y\in A
  \vdash x\NeutralSemiOp{}y\in A%
}{SemilatticeClosureAxiom}{%
  \DeltaRow{Mengen}{A\dsep x\dsep y}
  \DeltaRow{Funktionen}{\NeutralSemiOp}
}
\par\smallskip

\FormulaAxiomDeltaK[Zugriff auf die Assoziativität eines Halbverbandes]{%
  \Semi{\NeutralSemiOp,A}\dsep x,y,z\in A\vdash
  \bigl(x\NeutralSemiOp{}y\bigr)\NeutralSemiOp{}z
    =x\NeutralSemiOp\bigl(y\NeutralSemiOp{}z\bigr)%
}{SemilatticeAssociativityAxiom}{%
  \DeltaRow{Mengen}{A\dsep x\dsep y\dsep z}
  \DeltaRow{Funktionen}{\NeutralSemiOp}
}
\par\smallskip

\FormulaAxiomDeltaK[Zugriff auf die Kommutativität eines Halbverbandes]{%
  \Semi{\NeutralSemiOp,A}\dsep x,y\in A
  \vdash x\NeutralSemiOp{}y=y\NeutralSemiOp{}x%
}{SemilatticeCommutativityAxiom}{%
  \DeltaRow{Mengen}{A\dsep x\dsep y}
  \DeltaRow{Funktionen}{\NeutralSemiOp}
}
\par\smallskip

\FormulaAxiomDeltaK[Zugriff auf die Idempotenz eines Halbverbandes]{%
  \Semi{\NeutralSemiOp,A}\dsep x\in A
  \vdash x\NeutralSemiOp{}x=x%
}{SemilatticeIdempotenceAxiom}{%
  \DeltaRow{Mengen}{A\dsep x}
  \DeltaRow{Funktionen}{\NeutralSemiOp}
}

Die Abgeschlossenheit ist die auf den Träger eingeschränkte Form der in
Band~03 eingeführten Funktionsauswertung. Die übrigen drei Zeilen sind die
eigentlichen Halbverbandsgesetze. Beim Nachweis eines konkreten
Halbverbandes werden alle vier Eigenschaften einzeln gezeigt; erst die letzte
Beweiszeile verwendet die zusammenfassende Definition.

\section{Die supremal induzierte Relation}

\subsection{Axiomatische Definition und Schreibweise}

\FormulaDefDeltaK[Supremal induzierte Relation und Schreibweise]{%
  x\JoinLeq{\JoinOp}{A}y
  \quad:\Longleftrightarrow\quad
  (x,y)\in\JoinOrd{\JoinOp}{A}%
}{JoinOrder}{%
  \DeltaRow{Mengen}{A\dsep x\dsep y}
  \DeltaRow{Funktionen}{\JoinOp}
  \DeltaRow{Relationen}{\JoinOrd{\JoinOp}{A}}
  \DeltaRow{\textbf{Axiome}}{}
  \DeltaRow{Relation auf \(A\)}{%
    \JoinOrd{\JoinOp}{A}\subseteq A\times A}
    [\FormulaRefAuto{F \subseteq A \times B}[def]]
  \DeltaRow{Auswertung}{%
    \begin{aligned}[t]
      &\Semi{\JoinOp,A}\dsep x,y\in A\vdash{}\\
      &(x,y)\in\JoinOrd{\JoinOp}{A}
        \leftrightarrow x\JoinOp{}y=y
    \end{aligned}}
    [\FormulaRefAuto{JoinOrderEvaluationAxiom}[ax]]
  \DeltaRow{\textbf{Neue Symbole}}{}
  \DeltaPrem{Supremale Relation}{\JoinOrd{\JoinOp}{A}}
  \DeltaPrem{Infixschreibweise}{\JoinLeq{\JoinOp}{A}}
}

\FormulaAxiomDeltaK[Auswertung der supremal induzierten Relation]{%
  \Semi{\JoinOp,A}\dsep x,y\in A\vdash
  \bigl((x,y)\in\JoinOrd{\JoinOp}{A}
    \leftrightarrow x\JoinOp{}y=y\bigr)%
}{JoinOrderEvaluationAxiom}{%
  \DeltaRow{Mengen}{A\dsep x\dsep y}
  \DeltaRow{Funktionen}{\JoinOp}
  \DeltaRow{Relationen}{\JoinOrd{\JoinOp}{A}}
}

Damit darf in späteren Sätzen gleichwertig
\[
  x\JoinLeq{\JoinOp}{A}y,
  \qquad (x,y)\in\JoinOrd{\JoinOp}{A},
  \qquad x\JoinOp{}y=y
\]
geschrieben werden, sobald \(x,y\in A\) feststehen.

\subsection{Träger und Auswertung}

\FormulaThmDeltaK[Erste Komponente der supremal induzierten Relation]{%
  \Semi{\JoinOp,A}\dsep (x,y)\in\JoinOrd{\JoinOp}{A}\vdash x\in A
}{JoinOrderFirstCarrier}{%
  \DeltaRow{Mengen}{A\dsep x\dsep y}
  \DeltaRow{Funktionen}{\JoinOp}
  \DeltaRow{Relationen}{\JoinOrd{\JoinOp}{A}}
}
\begin{tabproof}
  \proofstep{1}{\Semi{\JoinOp,A}}{\rA}
  \proofstep{2}{(x,y)\in\JoinOrd{\JoinOp}{A}}{\rA}
  \proofstep{}{\JoinOrd{\JoinOp}{A}\subseteq A\times A}{%
    \FormulaRefAuto{F \subseteq A \times B}[def]}
  \proofstep{2}{(x,y)\in A\times A}{%
    \FormulaRefAuto{A\subseteq B,\, x\in A \vdash x\in B}{3,2}}
  \proofstep{2}{x\in A}{%
    \FormulaRefAuto{(a,b)\in A\times B\vdash a\in A}{4}}
\end{tabproof}

\FormulaThmDeltaK[Zweite Komponente der supremal induzierten Relation]{%
  \Semi{\JoinOp,A}\dsep (x,y)\in\JoinOrd{\JoinOp}{A}\vdash y\in A
}{JoinOrderSecondCarrier}{%
  \DeltaRow{Mengen}{A\dsep x\dsep y}
  \DeltaRow{Funktionen}{\JoinOp}
  \DeltaRow{Relationen}{\JoinOrd{\JoinOp}{A}}
}
\begin{tabproof}
  \proofstep{1}{\Semi{\JoinOp,A}}{\rA}
  \proofstep{2}{(x,y)\in\JoinOrd{\JoinOp}{A}}{\rA}
  \proofstep{}{\JoinOrd{\JoinOp}{A}\subseteq A\times A}{%
    \FormulaRefAuto{F \subseteq A \times B}[def]}
  \proofstep{2}{(x,y)\in A\times A}{%
    \FormulaRefAuto{A\subseteq B,\, x\in A \vdash x\in B}{3,2}}
  \proofstep{2}{y\in A}{%
    \FormulaRefAuto{(a,b)\in A\times B\vdash b\in B}{4}}
\end{tabproof}

\FormulaThmDeltaK[Auswertung der supremal induzierten Relation]{%
  \Semi{\JoinOp,A}\dsep (x,y)\in\JoinOrd{\JoinOp}{A}
  \vdash x\JoinOp{}y=y
}{JoinOrderEvaluation}{%
  \DeltaRow{Mengen}{A\dsep x\dsep y}
  \DeltaRow{Funktionen}{\JoinOp}
  \DeltaRow{Relationen}{\JoinOrd{\JoinOp}{A}}
}
\begin{tabproof}
  \proofstep{1}{\Semi{\JoinOp,A}}{\rA}
  \proofstep{2}{(x,y)\in\JoinOrd{\JoinOp}{A}}{\rA}
  \proofstep{1,2}{x\in A}{\FormulaRefAuto{JoinOrderFirstCarrier}[thm]{1,2}}
  \proofstep{1,2}{y\in A}{\FormulaRefAuto{JoinOrderSecondCarrier}[thm]{1,2}}
  \proofstep{1,2}{%
    (x,y)\in\JoinOrd{\JoinOp}{A}\leftrightarrow x\JoinOp{}y=y}{%
    \FormulaRefAuto{JoinOrderEvaluationAxiom}[ax]{1,3,4}}
  \proofstep{1,2}{x\JoinOp{}y=y}{%
    \FormulaRefAuto{P\leftrightarrow Q, P\vdash Q}{5,2}}
\end{tabproof}

\subsection{Die induzierte partielle Ordnung}

\FormulaThmDeltaK[Reflexivität der supremal induzierten Relation]{%
  \Semi{\JoinOp,A}\dsep x\in A\vdash (x,x)\in\JoinOrd{\JoinOp}{A}
}{JoinOrderReflexive}{%
  \DeltaRow{Mengen}{A\dsep x}
  \DeltaRow{Funktionen}{\JoinOp}
}
\begin{tabproof}
  \proofstep{1}{\Semi{\JoinOp,A}}{\rA}
  \proofstep{2}{x\in A}{\rA}
  \proofstep{1,2}{x\JoinOp{}x=x}{%
    \FormulaRefAuto{SemilatticeIdempotenceAxiom}[ax]{1,2}}
  \proofstep{1,2}{%
    (x,x)\in\JoinOrd{\JoinOp}{A}\leftrightarrow x\JoinOp{}x=x}{%
    \FormulaRefAuto{JoinOrderEvaluationAxiom}[ax]{1,2,2}}
  \proofstep{1,2}{(x,x)\in\JoinOrd{\JoinOp}{A}}{%
    \FormulaRefAuto{P\leftrightarrow Q, Q\vdash P}{4,3}}
\end{tabproof}

\FormulaThmDeltaK[Transitivität der supremal induzierten Relation]{%
  \Semi{\JoinOp,A}\dsep (x,y)\in\JoinOrd{\JoinOp}{A}
  \dsep (y,z)\in\JoinOrd{\JoinOp}{A}
  \vdash (x,z)\in\JoinOrd{\JoinOp}{A}
}{JoinOrderTransitive}{%
  \DeltaRow{Mengen}{A\dsep x\dsep y\dsep z}
  \DeltaRow{Funktionen}{\JoinOp}
}
\begin{tabproof}
  \proofstep{1}{\Semi{\JoinOp,A}}{\rA}
  \proofstep{2}{(x,y)\in\JoinOrd{\JoinOp}{A}}{\rA}
  \proofstep{3}{(y,z)\in\JoinOrd{\JoinOp}{A}}{\rA}
  \proofstep{1,2}{x\in A}{\FormulaRefAuto{JoinOrderFirstCarrier}[thm]{1,2}}
  \proofstep{1,2}{y\in A}{\FormulaRefAuto{JoinOrderSecondCarrier}[thm]{1,2}}
  \proofstep{1,3}{z\in A}{\FormulaRefAuto{JoinOrderSecondCarrier}[thm]{1,3}}
  \proofstep{1,2}{x\JoinOp{}y=y}{\FormulaRefAuto{JoinOrderEvaluation}[thm]{1,2}}
  \proofstep{1,3}{y\JoinOp{}z=z}{\FormulaRefAuto{JoinOrderEvaluation}[thm]{1,3}}
  \proofstep{1,2,3}{\bigl(x\JoinOp{}y\bigr)\JoinOp{}z
    =x\JoinOp\bigl(y\JoinOp{}z\bigr)}{%
    \FormulaRefAuto{SemilatticeAssociativityAxiom}[ax]{1,4,5,6}}
  \proofstep{1,2,3}{\bigl(x\JoinOp{}y\bigr)\JoinOp{}z=z}{%
    \rIE{7,8}}
  \proofstep{1,2,3}{x\JoinOp\bigl(y\JoinOp{}z\bigr)=x\JoinOp{}z}{%
    \rIE{8,\rII}}
  \proofstep{1,2,3}{z=\bigl(x\JoinOp{}y\bigr)\JoinOp{}z}{%
    \FormulaRefAuto{a=b\vdash b=a}{10}}
  \proofstep{1,2,3}{z=x\JoinOp\bigl(y\JoinOp{}z\bigr)}{%
    \FormulaRefAuto{a=b,\,b=c\vdash a=c}{12,9}}
  \proofstep{1,2,3}{x\JoinOp\bigl(y\JoinOp{}z\bigr)=z}{%
    \FormulaRefAuto{a=b\vdash b=a}{13}}
  \proofstep{1,2,3}{x\JoinOp{}z=z}{%
    \FormulaRefAuto{a=b,\,a=c\vdash b=c}{11,14}}
  \proofstep{1,2,3}{%
    (x,z)\in\JoinOrd{\JoinOp}{A}\leftrightarrow x\JoinOp{}z=z}{%
    \FormulaRefAuto{JoinOrderEvaluationAxiom}[ax]{1,4,6}}
  \proofstep{1,2,3}{(x,z)\in\JoinOrd{\JoinOp}{A}}{%
    \FormulaRefAuto{P\leftrightarrow Q, Q\vdash P}{16,15}}
\end{tabproof}

\FormulaThmDeltaK[Antisymmetrie der supremal induzierten Relation]{%
  \Semi{\JoinOp,A}\dsep (x,y)\in\JoinOrd{\JoinOp}{A}
  \dsep (y,x)\in\JoinOrd{\JoinOp}{A}\vdash x=y
}{JoinOrderAntisymmetric}{%
  \DeltaRow{Mengen}{A\dsep x\dsep y}
  \DeltaRow{Funktionen}{\JoinOp}
}
\begin{tabproof}
  \proofstep{1}{\Semi{\JoinOp,A}}{\rA}
  \proofstep{2}{(x,y)\in\JoinOrd{\JoinOp}{A}}{\rA}
  \proofstep{3}{(y,x)\in\JoinOrd{\JoinOp}{A}}{\rA}
  \proofstep{1,2}{x\in A}{\FormulaRefAuto{JoinOrderFirstCarrier}[thm]{1,2}}
  \proofstep{1,2}{y\in A}{\FormulaRefAuto{JoinOrderSecondCarrier}[thm]{1,2}}
  \proofstep{1,2}{x\JoinOp{}y=y}{\FormulaRefAuto{JoinOrderEvaluation}[thm]{1,2}}
  \proofstep{1,3}{y\JoinOp{}x=x}{\FormulaRefAuto{JoinOrderEvaluation}[thm]{1,3}}
  \proofstep{1,2}{x\JoinOp{}y=y\JoinOp{}x}{%
    \FormulaRefAuto{SemilatticeCommutativityAxiom}[ax]{1,4,5}}
  \proofstep{1,2,3}{x\JoinOp{}y=x}{\rIE{8,7}}
  \proofstep{1,2,3}{y=x}{%
    \FormulaRefAuto{a=b,\,a=c\vdash b=c}{6,9}}
  \proofstep{1,2,3}{x=y}{\FormulaRefAuto{a=b\vdash b=a}{10}}
\end{tabproof}

\FormulaThmDeltaK[Ein supremal gelesener Halbverband induziert eine partielle Ordnung]{%
  \Semi{\JoinOp,A}\vdash \PartOrd{\JoinOrd{\JoinOp}{A},A}
}{JoinSemilatticeInducesOrder}{%
  \DeltaRow{Mengen}{A\dsep x\dsep y\dsep z}
  \DeltaRow{Funktionen}{\JoinOp}
  \DeltaRow{Relationen}{\JoinOrd{\JoinOp}{A}}
}
\begin{tabproofsplitwide}
  \proofpartwideindR[Relation auf dem Träger]{%
    \JoinOrd{\JoinOp}{A}\subseteq A\times A
  }{\JoinOrd{\JoinOp}{A}\subseteq A\times A}[JoinOrderCarrierPart]
    \proofstepwidestar[]{\JoinOrd{\JoinOp}{A}\subseteq A\times A}{%
      \FormulaRefAuto{F \subseteq A \times B}[def]}
  \closeproofpartwideindR

  \proofpartwideindR[Reflexivität]{%
    \Semi{\JoinOp,A}\vdash
    x\in A\vdash (x,x)\in\JoinOrd{\JoinOp}{A}
  }{\Semi{\JoinOp,A}\vdash
    x\in A\vdash (x,x)\in\JoinOrd{\JoinOp}{A}}[JoinOrderReflexivePart]
    \proofstepwidestar[1]{\Semi{\JoinOp,A}}{\rA}
    \proofstepwidestar[1]{x\in A\vdash (x,x)\in\JoinOrd{\JoinOp}{A}}{%
      \FormulaRefAuto{JoinOrderReflexive}[thm]{1}}
  \closeproofpartwideindR

  \proofpartwideindR[Transitivität]{%
    \begin{aligned}[t]
      &\Semi{\JoinOp,A}\vdash{}\\[-2pt]
      &(x,y)\in\JoinOrd{\JoinOp}{A}\dsep
        (y,z)\in\JoinOrd{\JoinOp}{A}
        \vdash (x,z)\in\JoinOrd{\JoinOp}{A}
    \end{aligned}
  }{\Semi{\JoinOp,A}\vdash
    (x,y)\in\JoinOrd{\JoinOp}{A}\dsep
    (y,z)\in\JoinOrd{\JoinOp}{A}\vdash
    (x,z)\in\JoinOrd{\JoinOp}{A}}[JoinOrderTransitivePart]
    \proofstepwidestar[1]{\Semi{\JoinOp,A}}{\rA}
    \proofstepwidestar[1]{%
      (x,y)\in\JoinOrd{\JoinOp}{A}\dsep
      (y,z)\in\JoinOrd{\JoinOp}{A}\vdash
      (x,z)\in\JoinOrd{\JoinOp}{A}}{%
      \FormulaRefAuto{JoinOrderTransitive}[thm]{1}}
  \closeproofpartwideindR

  \proofpartwideindR[Antisymmetrie]{%
    \begin{aligned}[t]
      &\Semi{\JoinOp,A}\vdash{}\\[-2pt]
      &(x,y)\in\JoinOrd{\JoinOp}{A}\dsep
        (y,x)\in\JoinOrd{\JoinOp}{A}\vdash x=y
    \end{aligned}
  }{\Semi{\JoinOp,A}\vdash
    (x,y)\in\JoinOrd{\JoinOp}{A}\dsep
    (y,x)\in\JoinOrd{\JoinOp}{A}\vdash x=y}[JoinOrderAntisymmetricPart]
    \proofstepwidestar[1]{\Semi{\JoinOp,A}}{\rA}
    \proofstepwidestar[1]{%
      (x,y)\in\JoinOrd{\JoinOp}{A}\dsep
      (y,x)\in\JoinOrd{\JoinOp}{A}\vdash x=y}{%
      \FormulaRefAuto{JoinOrderAntisymmetric}[thm]{1}}
  \closeproofpartwideindR

  \proofpartwide{Schluss}
    \proofstepwidestar[1]{\Semi{\JoinOp,A}}{\rA}
    \proofstepwidestar[]{\JoinOrd{\JoinOp}{A}\subseteq A\times A}{%
      \FormulaRefAuto{JoinOrderCarrierPart}[thm]}
    \proofstepwidestar[1]{x\in A\vdash (x,x)\in\JoinOrd{\JoinOp}{A}}{%
      \FormulaRefAuto{JoinOrderReflexivePart}[thm]{1}}
    \proofstepwidestar[1]{%
      (x,y)\in\JoinOrd{\JoinOp}{A}\dsep
      (y,z)\in\JoinOrd{\JoinOp}{A}\vdash
      (x,z)\in\JoinOrd{\JoinOp}{A}}{%
      \FormulaRefAuto{JoinOrderTransitivePart}[thm]{1}}
    \proofstepwidestar[1]{%
      (x,y)\in\JoinOrd{\JoinOp}{A}\dsep
      (y,x)\in\JoinOrd{\JoinOp}{A}\vdash x=y}{%
      \FormulaRefAuto{JoinOrderAntisymmetricPart}[thm]{1}}
    \proofstepwidestar[1]{\PartOrd{\JoinOrd{\JoinOp}{A},A}}{%
      \FormulaRefAuto{Partielle Ordnung}[def]{2,3,4,5}}
  \closeproofpartwide
\end{tabproofsplitwide}

\subsection{Universelle Eigenschaft der supremalen Lesart}

\FormulaDefDeltaK[Supremum eines Elementpaares]{%
  \PairSup{R}{A}{x}{y}{u}%
}{PairSupremum}{%
  \DeltaRow{Mengen}{A\dsep x\dsep y\dsep u\dsep v}
  \DeltaRow{Relationen}{R}
  \DeltaRow{Elemente}{x,y,u\in A}
  \DeltaRow{Obere Schranken}{(x,u)\in R\dsep (y,u)\in R}
  \DeltaRow{Kleinste obere Schranke}{%
    v\in A\dsep (x,v)\in R\dsep (y,v)\in R\vdash (u,v)\in R}
  \DeltaRow{\textbf{Neues Symbol}}{}
  \DeltaPrem{Paarsupremum}{\PairSup{R}{A}{x}{y}{u}}
}

\FormulaThmDeltaK[Der erste Operand liegt unter dem algebraischen Supremum]{%
  \Semi{\JoinOp,A}\dsep x,y\in A
  \vdash (x,x\JoinOp{}y)\in\JoinOrd{\JoinOp}{A}
}{JoinUpperLeft}{%
  \DeltaRow{Mengen}{A\dsep x\dsep y}
  \DeltaRow{Funktionen}{\JoinOp}
}
\begin{tabproof}
  \proofstep{1}{\Semi{\JoinOp,A}}{\rA}
  \proofstep{2}{x\in A}{\rA}
  \proofstep{3}{y\in A}{\rA}
  \proofstep{1,2,3}{x\JoinOp{}y\in A}{%
    \FormulaRefAuto{SemilatticeClosureAxiom}[ax]{1,2,3}}
  \proofstep{1,2}{x\JoinOp{}x=x}{\FormulaRefAuto{SemilatticeIdempotenceAxiom}[ax]{1,2}}
  \proofstep{1,2,3}{\bigl(x\JoinOp{}x\bigr)\JoinOp{}y=x\JoinOp{}y}{\rIE{5,\rII}}
  \proofstep{1,2,3}{\bigl(x\JoinOp{}x\bigr)\JoinOp{}y=x\JoinOp\bigl(x\JoinOp{}y\bigr)}{%
    \FormulaRefAuto{SemilatticeAssociativityAxiom}[ax]{1,2,2,3}}
  \proofstep{1,2,3}{x\JoinOp\bigl(x\JoinOp{}y\bigr)=x\JoinOp{}y}{%
    \FormulaRefAuto{a=b,\,a=c\vdash b=c}{7,6}}
  \proofstep{1,2,3}{%
    (x,x\JoinOp{}y)\in\JoinOrd{\JoinOp}{A}
    \leftrightarrow x\JoinOp\bigl(x\JoinOp{}y\bigr)=x\JoinOp{}y}{%
    \FormulaRefAuto{JoinOrderEvaluationAxiom}[ax]{1,2,4}}
  \proofstep{1,2,3}{(x,x\JoinOp{}y)\in\JoinOrd{\JoinOp}{A}}{%
    \FormulaRefAuto{P\leftrightarrow Q, Q\vdash P}{9,8}}
\end{tabproof}

\FormulaThmDeltaK[Der zweite Operand liegt unter dem algebraischen Supremum]{%
  \Semi{\JoinOp,A}\dsep x,y\in A
  \vdash (y,x\JoinOp{}y)\in\JoinOrd{\JoinOp}{A}
}{JoinUpperRight}{%
  \DeltaRow{Mengen}{A\dsep x\dsep y}
  \DeltaRow{Funktionen}{\JoinOp}
}
\begin{tabproof}
  \proofstep{1}{\Semi{\JoinOp,A}}{\rA}
  \proofstep{2}{x\in A}{\rA}
  \proofstep{3}{y\in A}{\rA}
  \proofstep{1,2,3}{(y,y\JoinOp{}x)\in\JoinOrd{\JoinOp}{A}}{%
    \FormulaRefAuto{JoinUpperLeft}[thm]{1,3,2}}
  \proofstep{1,2,3}{y\JoinOp{}x=x\JoinOp{}y}{%
    \FormulaRefAuto{SemilatticeCommutativityAxiom}[ax]{1,3,2}}
  \proofstep{1,2,3}{(y,x\JoinOp{}y)\in\JoinOrd{\JoinOp}{A}}{\rIE{5,4}}
\end{tabproof}

\FormulaThmDeltaK[Das algebraische Supremum ist die kleinste obere Schranke]{%
  \begin{aligned}[t]
    &\Semi{\JoinOp,A}\dsep x,y,u\in A\dsep
      (x,u)\in\JoinOrd{\JoinOp}{A}\dsep (y,u)\in\JoinOrd{\JoinOp}{A}\vdash{}\\
    &(x\JoinOp{}y,u)\in\JoinOrd{\JoinOp}{A}
  \end{aligned}
}{JoinLeastUpper}{%
  \DeltaRow{Mengen}{A\dsep x\dsep y\dsep u}
  \DeltaRow{Funktionen}{\JoinOp}
}
\begin{tabproof}
  \proofstep{1}{\Semi{\JoinOp,A}}{\rA}
  \proofstep{2}{x\in A}{\rA}
  \proofstep{3}{y\in A}{\rA}
  \proofstep{4}{u\in A}{\rA}
  \proofstep{5}{(x,u)\in\JoinOrd{\JoinOp}{A}}{\rA}
  \proofstep{6}{(y,u)\in\JoinOrd{\JoinOp}{A}}{\rA}
  \proofstep{1,5}{x\JoinOp{}u=u}{%
    \FormulaRefAuto{JoinOrderEvaluation}[thm]{1,5}}
  \proofstep{1,6}{y\JoinOp{}u=u}{%
    \FormulaRefAuto{JoinOrderEvaluation}[thm]{1,6}}
  \proofstep{1,2,3}{x\JoinOp{}y\in A}{%
    \FormulaRefAuto{SemilatticeClosureAxiom}[ax]{1,2,3}}
  \proofstep{1,2,3,4}{\bigl(x\JoinOp{}y\bigr)\JoinOp{}u=x\JoinOp\bigl(y\JoinOp{}u\bigr)}{%
    \FormulaRefAuto{SemilatticeAssociativityAxiom}[ax]{1,2,3,4}}
  \proofstep{1,2,3,4,6}{x\JoinOp\bigl(y\JoinOp{}u\bigr)=x\JoinOp{}u}{\rIE{8,\rII}}
  \proofstep{1,2,3,4,5,6}{\bigl(x\JoinOp{}y\bigr)\JoinOp{}u=u}{%
    \FormulaRefAuto{a=b,\,b=c\vdash a=c}{10,\rIE{11,7}}}
  \proofstep{1,2,3,4}{%
    (x\JoinOp{}y,u)\in\JoinOrd{\JoinOp}{A}
    \leftrightarrow \bigl(x\JoinOp{}y\bigr)\JoinOp{}u=u}{%
    \FormulaRefAuto{JoinOrderEvaluationAxiom}[ax]{1,9,4}}
  \proofstep{1,2,3,4,5,6}{(x\JoinOp{}y,u)\in\JoinOrd{\JoinOp}{A}}{%
    \FormulaRefAuto{P\leftrightarrow Q, Q\vdash P}{13,12}}
\end{tabproof}

\FormulaThmDeltaK[Die Halbverbandsoperation ist das Paarsupremum]{%
  \Semi{\JoinOp,A}\dsep x,y\in A
  \vdash \PairSup{\JoinOrd{\JoinOp}{A}}{A}{x}{y}{x\JoinOp{}y}
}{JoinIsPairSupremum}{%
  \DeltaRow{Mengen}{A\dsep x\dsep y\dsep u}
  \DeltaRow{Funktionen}{\JoinOp}
}
\begin{tabproof}
  \proofstep{1}{\Semi{\JoinOp,A}}{\rA}
  \proofstep{2}{x\in A}{\rA}
  \proofstep{3}{y\in A}{\rA}
  \proofstep{1,2,3}{x\JoinOp{}y\in A}{%
    \FormulaRefAuto{SemilatticeClosureAxiom}[ax]{1,2,3}}
  \proofstep{1,2,3}{(x,x\JoinOp{}y)\in\JoinOrd{\JoinOp}{A}}{%
    \FormulaRefAuto{JoinUpperLeft}[thm]{1,2,3}}
  \proofstep{1,2,3}{(y,x\JoinOp{}y)\in\JoinOrd{\JoinOp}{A}}{%
    \FormulaRefAuto{JoinUpperRight}[thm]{1,2,3}}
  \proofstep{1}{%
    \begin{aligned}[t]
      &u\in A\dsep (x,u)\in\JoinOrd{\JoinOp}{A}
       \dsep (y,u)\in\JoinOrd{\JoinOp}{A}\vdash{}\\
      &(x\JoinOp{}y,u)\in\JoinOrd{\JoinOp}{A}
    \end{aligned}}{%
    \FormulaRefAuto{JoinLeastUpper}[thm]{1,2,3}}
  \proofstep{1,2,3}{\PairSup{\JoinOrd{\JoinOp}{A}}{A}{x}{y}{x\JoinOp{}y}}{%
    \FormulaRefAuto{PairSupremum}[def]{\rAI{\rAI{2,3},4},5,6,7}}
\end{tabproof}

\section{Die ordnungstheoretische Charakterisierung}

\FormulaDefDeltaK[Durch Paarsuprema bestimmte Operation]{%
  \PairJoinOp{\JoinOp,R,A}%
}{PairJoinOperation}{%
  \DeltaRow{Mengen}{A\dsep x\dsep y}
  \DeltaRow{Relationen}{R}
  \DeltaRow{Funktionen}{\JoinOp}
  \DeltaRow{Abgeschlossenheit}{%
    x,y\in A\vdash x\JoinOp{}y\in A}
  \DeltaRow{Paarsuprema}{%
    x,y\in A\vdash \PairSup{R}{A}{x}{y}{x\JoinOp{}y}}
  \DeltaRow{\textbf{Neues Symbol}}{}
  \DeltaPrem{Paarsupremumsoperation}{\PairJoinOp{\JoinOp,R,A}}
}

\FormulaThmDeltaK[Eindeutigkeit des Paarsupremums]{%
  \begin{aligned}[t]
    &\PartOrd{R,A}\dsep \PairSup{R}{A}{x}{y}{u}
      \dsep \PairSup{R}{A}{x}{y}{v}\vdash u=v
  \end{aligned}
}{PairSupremumUnique}{%
  \DeltaRow{Mengen}{A\dsep x\dsep y\dsep u\dsep v}
  \DeltaRow{Relationen}{R}
}
\begin{tabproof}
  \proofstep{1}{\PartOrd{R,A}}{\rA}
  \proofstep{2}{\PairSup{R}{A}{x}{y}{u}}{\rA}
  \proofstep{3}{\PairSup{R}{A}{x}{y}{v}}{\rA}
  \proofstep{2}{u\in A}{\FormulaRefAuto{PairSupremum}[def]{2}}
  \proofstep{3}{v\in A}{\FormulaRefAuto{PairSupremum}[def]{3}}
  \proofstep{3}{(x,v)\in R}{\FormulaRefAuto{PairSupremum}[def]{3}}
  \proofstep{3}{(y,v)\in R}{\FormulaRefAuto{PairSupremum}[def]{3}}
  \proofstep{2,3}{(u,v)\in R}{%
    \FormulaRefAuto{PairSupremum}[def]{2,5,6,7}}
  \proofstep{2}{(x,u)\in R}{\FormulaRefAuto{PairSupremum}[def]{2}}
  \proofstep{2}{(y,u)\in R}{\FormulaRefAuto{PairSupremum}[def]{2}}
  \proofstep{2,3}{(v,u)\in R}{%
    \FormulaRefAuto{PairSupremum}[def]{3,4,9,10}}
  \proofstep{1,2,3}{u=v}{%
    \FormulaRefAuto{(x,y)\in R \dsep (y,x)\in R \vdash x=y}[ax]{1,8,11}}
\end{tabproof}

\FormulaThmDeltaK[Idempotenz einer Paarsupremumsoperation]{%
  \PartOrd{R,A}\dsep \PairJoinOp{\JoinOp,R,A}\dsep x\in A
  \vdash x\JoinOp{}x=x
}{PairJoinIdempotent}{%
  \DeltaRow{Mengen}{A\dsep x}
  \DeltaRow{Relationen}{R}
  \DeltaRow{Funktionen}{\JoinOp}
}
\begin{tabproof}
  \proofstep{1}{\PartOrd{R,A}}{\rA}
  \proofstep{2}{\PairJoinOp{\JoinOp,R,A}}{\rA}
  \proofstep{3}{x\in A}{\rA}
  \proofstep{2,3}{\PairSup{R}{A}{x}{x}{x\JoinOp{}x}}{%
    \FormulaRefAuto{PairJoinOperation}[def]{2,3,3}}
  \proofstep{1,3}{(x,x)\in R}{%
    \FormulaRefAuto{EqRelReflexivity}[ax]{1,3}}
  \proofstep{}{%
    v\in A\dsep (x,v)\in R\dsep (x,v)\in R\vdash (x,v)\in R}{\rA}
  \proofstep{1,3}{\PairSup{R}{A}{x}{x}{x}}{%
    \FormulaRefAuto{PairSupremum}[def]{\rAI{\rAI{3,3},3},5,5,6}}
  \proofstep{1,2,3}{x\JoinOp{}x=x}{%
    \FormulaRefAuto{PairSupremumUnique}[thm]{1,4,7}}
\end{tabproof}

\FormulaThmDeltaK[Kommutativität einer Paarsupremumsoperation]{%
  \PartOrd{R,A}\dsep \PairJoinOp{\JoinOp,R,A}\dsep x,y\in A
  \vdash x\JoinOp{}y=y\JoinOp{}x
}{PairJoinCommutative}{%
  \DeltaRow{Mengen}{A\dsep x\dsep y}
  \DeltaRow{Relationen}{R}
  \DeltaRow{Funktionen}{\JoinOp}
}
\begin{tabproof}
  \proofstep{1}{\PartOrd{R,A}}{\rA}
  \proofstep{2}{\PairJoinOp{\JoinOp,R,A}}{\rA}
  \proofstep{3}{x\in A}{\rA}
  \proofstep{4}{y\in A}{\rA}
  \proofstep{2,3,4}{\PairSup{R}{A}{x}{y}{x\JoinOp{}y}}{%
    \FormulaRefAuto{PairJoinOperation}[def]{2,3,4}}
  \proofstep{2,3,4}{\PairSup{R}{A}{y}{x}{y\JoinOp{}x}}{%
    \FormulaRefAuto{PairJoinOperation}[def]{2,4,3}}
  \proofstep{2,3,4}{(y,y\JoinOp{}x)\in R}{\FormulaRefAuto{PairSupremum}[def]{6}}
  \proofstep{2,3,4}{(x,y\JoinOp{}x)\in R}{\FormulaRefAuto{PairSupremum}[def]{6}}
  \proofstep{2,3,4}{y\JoinOp{}x\in A}{\FormulaRefAuto{PairSupremum}[def]{6}}
  \proofstep{2,3,4}{%
    u\in A\dsep (x,u)\in R\dsep (y,u)\in R
    \vdash (y\JoinOp{}x,u)\in R}{%
    \FormulaRefAuto{PairSupremum}[def]{6}}
  \proofstep{2,3,4}{\PairSup{R}{A}{x}{y}{y\JoinOp{}x}}{%
    \FormulaRefAuto{PairSupremum}[def]{%
      \rAI{\rAI{3,4},9},8,7,10}}
  \proofstep{1,2,3,4}{x\JoinOp{}y=y\JoinOp{}x}{%
    \FormulaRefAuto{PairSupremumUnique}[thm]{1,5,10}}
\end{tabproof}

\FormulaThmDeltaK[Assoziativität einer Paarsupremumsoperation]{%
  \begin{aligned}[t]
    &\PartOrd{R,A}\dsep \PairJoinOp{\JoinOp,R,A}\dsep x,y,z\in A\vdash{}\\
    &\bigl(x\JoinOp{}y\bigr)\JoinOp{}z=x\JoinOp\bigl(y\JoinOp{}z\bigr)
  \end{aligned}
}{PairJoinAssociative}{%
  \DeltaRow{Mengen}{A\dsep x\dsep y\dsep z}
  \DeltaRow{Relationen}{R}
  \DeltaRow{Funktionen}{\JoinOp}
}
\begin{tabproofwide}
  \proofstepwidestar[1]{\PartOrd{R,A}}{\rA}
  \proofstepwidestar[2]{\PairJoinOp{\JoinOp,R,A}}{\rA}
  \proofstepwidestar[3]{x\in A}{\rA}
  \proofstepwidestar[4]{y\in A}{\rA}
  \proofstepwidestar[5]{z\in A}{\rA}
  \proofstepwidestar[2,3,4]{\PairSup{R}{A}{x}{y}{x\JoinOp{}y}}{%
    \FormulaRefAuto{PairJoinOperation}[def]{2,3,4}}
  \proofstepwidestar[2,4,5]{\PairSup{R}{A}{y}{z}{y\JoinOp{}z}}{%
    \FormulaRefAuto{PairJoinOperation}[def]{2,4,5}}
  \proofstepwidestar[2,3,4]{x\JoinOp{}y\in A}{\FormulaRefAuto{PairSupremum}[def]{6}}
  \proofstepwidestar[2,4,5]{y\JoinOp{}z\in A}{\FormulaRefAuto{PairSupremum}[def]{7}}
  \proofstepwidestar[2,3,4,5]{\PairSup{R}{A}{x\JoinOp{}y}{z}{\bigl(x\JoinOp{}y\bigr)\JoinOp{}z}}{%
    \FormulaRefAuto{PairJoinOperation}[def]{2,8,5}}
  \proofstepwidestar[2,3,4,5]{\PairSup{R}{A}{x}{y\JoinOp{}z}{x\JoinOp\bigl(y\JoinOp{}z\bigr)}}{%
    \FormulaRefAuto{PairJoinOperation}[def]{2,3,9}}
  \proofstepwidestar[2,3,4,5]{x\JoinOp\bigl(y\JoinOp{}z\bigr)\in A}{%
    \FormulaRefAuto{PairSupremum}[def]{11}}
  \proofstepwidestar[2,3,4,5]{\bigl(x\JoinOp{}y\bigr)\JoinOp{}z\in A}{%
    \FormulaRefAuto{PairSupremum}[def]{10}}
  \proofstepwidestar[2,3,4,5]{(x,x\JoinOp\bigl(y\JoinOp{}z\bigr))\in R}{%
    \FormulaRefAuto{PairSupremum}[def]{11}}
  \proofstepwidestar[2,3,4,5]{(y\JoinOp{}z,x\JoinOp\bigl(y\JoinOp{}z\bigr))\in R}{%
    \FormulaRefAuto{PairSupremum}[def]{11}}
  \proofstepwidestar[2,3,4,5]{(y,y\JoinOp{}z)\in R}{%
    \FormulaRefAuto{PairSupremum}[def]{7}}
  \proofstepwidestar[2,3,4,5]{(z,y\JoinOp{}z)\in R}{%
    \FormulaRefAuto{PairSupremum}[def]{7}}
  \proofstepwidestar[1,2,3,4,5]{(y,x\JoinOp\bigl(y\JoinOp{}z\bigr))\in R}{%
    \FormulaRefAuto{EqRelTransitivity}[ax]{1,16,15}}
  \proofstepwidestar[1,2,3,4,5]{(z,x\JoinOp\bigl(y\JoinOp{}z\bigr))\in R}{%
    \FormulaRefAuto{EqRelTransitivity}[ax]{1,17,15}}
  \proofstepwidestar[1,2,3,4,5]{(x\JoinOp{}y,x\JoinOp\bigl(y\JoinOp{}z\bigr))\in R}{%
    \FormulaRefAuto{PairSupremum}[def]{6,12,14,18}}
  \proofstepwidestar[1,2,3,4,5]{(\bigl(x\JoinOp{}y\bigr)\JoinOp{}z,x\JoinOp\bigl(y\JoinOp{}z\bigr))\in R}{%
    \FormulaRefAuto{PairSupremum}[def]{10,12,20,19}}
  \proofstepwidestar[2,3,4,5]{(x\JoinOp{}y,\bigl(x\JoinOp{}y\bigr)\JoinOp{}z)\in R}{%
    \FormulaRefAuto{PairSupremum}[def]{10}}
  \proofstepwidestar[2,3,4,5]{(z,\bigl(x\JoinOp{}y\bigr)\JoinOp{}z)\in R}{%
    \FormulaRefAuto{PairSupremum}[def]{10}}
  \proofstepwidestar[2,3,4,5]{(x,x\JoinOp{}y)\in R}{%
    \FormulaRefAuto{PairSupremum}[def]{6}}
  \proofstepwidestar[2,3,4,5]{(y,x\JoinOp{}y)\in R}{%
    \FormulaRefAuto{PairSupremum}[def]{6}}
  \proofstepwidestar[1,2,3,4,5]{(x,\bigl(x\JoinOp{}y\bigr)\JoinOp{}z)\in R}{%
    \FormulaRefAuto{EqRelTransitivity}[ax]{1,24,22}}
  \proofstepwidestar[1,2,3,4,5]{(y,\bigl(x\JoinOp{}y\bigr)\JoinOp{}z)\in R}{%
    \FormulaRefAuto{EqRelTransitivity}[ax]{1,25,22}}
  \proofstepwidestar[1,2,3,4,5]{(y\JoinOp{}z,\bigl(x\JoinOp{}y\bigr)\JoinOp{}z)\in R}{%
    \FormulaRefAuto{PairSupremum}[def]{7,13,27,23}}
  \proofstepwidestar[1,2,3,4,5]{(x\JoinOp\bigl(y\JoinOp{}z\bigr),\bigl(x\JoinOp{}y\bigr)\JoinOp{}z)\in R}{%
    \FormulaRefAuto{PairSupremum}[def]{11,13,26,28}}
  \proofstepwidestar[1,2,3,4,5]{\bigl(x\JoinOp{}y\bigr)\JoinOp{}z=x\JoinOp\bigl(y\JoinOp{}z\bigr)}{%
    \FormulaRefAuto{(x,y)\in R \dsep (y,x)\in R \vdash x=y}[ax]{1,21,29}}
\end{tabproofwide}

\FormulaThmDeltaK[Charakterisierung durch Paarsuprema]{%
  \PartOrd{R,A}\dsep \PairJoinOp{\JoinOp,R,A}\vdash \Semi{\JoinOp,A}
}{PairJoinCharacterization}{%
  \DeltaRow{Mengen}{A\dsep x\dsep y\dsep z}
  \DeltaRow{Relationen}{R}
  \DeltaRow{Funktionen}{\JoinOp}
}
\begin{tabproof}
  \proofstep{1}{\PartOrd{R,A}}{\rA}
  \proofstep{2}{\PairJoinOp{\JoinOp,R,A}}{\rA}
  \proofstep{2}{x,y\in A\vdash x\JoinOp{}y\in A}{%
    \FormulaRefAuto{PairJoinOperation}[def]{2}}
  \proofstep{1,2}{%
    x,y,z\in A\vdash \bigl(x\JoinOp{}y\bigr)\JoinOp{}z=x\JoinOp\bigl(y\JoinOp{}z\bigr)}{%
    \FormulaRefAuto{PairJoinAssociative}[thm]{1,2}}
  \proofstep{1,2}{x,y\in A\vdash x\JoinOp{}y=y\JoinOp{}x}{%
    \FormulaRefAuto{PairJoinCommutative}[thm]{1,2}}
  \proofstep{1,2}{x\in A\vdash x\JoinOp{}x=x}{%
    \FormulaRefAuto{PairJoinIdempotent}[thm]{1,2}}
  \proofstep{1,2}{\Semi{\JoinOp,A}}{%
    \FormulaRefAuto{JoinSemilattice}[def]{3,4,5,6}}
\end{tabproof}

\FormulaThmDeltaK[Die Paarsupremumsoperation gewinnt die Ordnung zurück]{%
  \begin{aligned}[t]
    &\PartOrd{R,A}\dsep \PairJoinOp{\JoinOp,R,A}\dsep x,y\in A\vdash{}\\
    &(x,y)\in R\leftrightarrow (x,y)\in\JoinOrd{\JoinOp}{A}
  \end{aligned}
}{PairJoinRecoversOrder}{%
  \DeltaRow{Mengen}{A\dsep x\dsep y}
  \DeltaRow{Relationen}{R\dsep \JoinOrd{\JoinOp}{A}}
  \DeltaRow{Funktionen}{\JoinOp}
}
\begin{tabproofsplitwide}
  \proofpartwide{\(\vdash\)}
    \proofstepwidestar[1]{\PartOrd{R,A}}{\rA}
    \proofstepwidestar[2]{\PairJoinOp{\JoinOp,R,A}}{\rA}
    \proofstepwidestar[3]{x\in A}{\rA}
    \proofstepwidestar[4]{y\in A}{\rA}
    \proofstepwidestar[5]{(x,y)\in R}{\rA}
    \proofstepwidestar[1,4]{(y,y)\in R}{%
      \FormulaRefAuto{EqRelReflexivity}[ax]{1,4}}
    \proofstepwidestar[2,3,4]{\PairSup{R}{A}{x}{y}{x\JoinOp{}y}}{%
      \FormulaRefAuto{PairJoinOperation}[def]{2,3,4}}
    \proofstepwidestar[2,3,4]{x\JoinOp{}y\in A}{%
      \FormulaRefAuto{PairSupremum}[def]{7}}
    \proofstepwidestar[1,2,3,4,5]{(x\JoinOp{}y,y)\in R}{%
      \FormulaRefAuto{PairSupremum}[def]{7,4,5,6}}
    \proofstepwidestar[2,3,4]{(y,x\JoinOp{}y)\in R}{%
      \FormulaRefAuto{PairSupremum}[def]{7}}
    \proofstepwidestar[1,2,3,4,5]{x\JoinOp{}y=y}{%
      \FormulaRefAuto{(x,y)\in R \dsep (y,x)\in R \vdash x=y}[ax]{1,9,10}}
    \proofstepwidestar[1,2,3,4,5]{(x,y)\in\JoinOrd{\JoinOp}{A}}{%
      \FormulaRefAuto{P\leftrightarrow Q, Q\vdash P}{%
        \FormulaRefAuto{JoinOrderEvaluationAxiom}[ax]{%
          \FormulaRefAuto{PairJoinCharacterization}[thm]{1,2},3,4},11}}
  \closeproofpartwide

  \proofpartwide{\(\dashv\)}
    \proofstepwidestar[1]{\PartOrd{R,A}}{\rA}
    \proofstepwidestar[2]{\PairJoinOp{\JoinOp,R,A}}{\rA}
    \proofstepwidestar[3]{x\in A}{\rA}
    \proofstepwidestar[4]{y\in A}{\rA}
    \proofstepwidestar[5]{(x,y)\in\JoinOrd{\JoinOp}{A}}{\rA}
    \proofstepwidestar[1,2,3,4,5]{x\JoinOp{}y=y}{%
      \FormulaRefAuto{JoinOrderEvaluation}[thm]{%
        \FormulaRefAuto{PairJoinCharacterization}[thm]{1,2},5}}
    \proofstepwidestar[2,3,4]{\PairSup{R}{A}{x}{y}{x\JoinOp{}y}}{%
      \FormulaRefAuto{PairJoinOperation}[def]{2,3,4}}
    \proofstepwidestar[2,3,4]{(x,x\JoinOp{}y)\in R}{%
      \FormulaRefAuto{PairSupremum}[def]{7}}
    \proofstepwidestar[2,3,4,5]{(x,y)\in R}{\rIE{6,8}}
  \closeproofpartwide
\end{tabproofsplitwide}

Damit ist die zentrale Äquivalenz bewiesen:
\[
  \boxed{
    \begin{aligned}
      &\text{assoziative, kommutative und idempotente binäre Operation}\\
      &\qquad\Longleftrightarrow\\
      &\text{partielle Ordnung mit Suprema aller Elementpaare.}
    \end{aligned}}
\]

\section{Die infimal induzierte Relation}

Die Operation \(\MeetOp\) erfüllt keine anderen algebraischen Axiome als
\(\JoinOp\). Sie ist lediglich eine zweite Schreibweise für eine
Halbverbandsoperation, deren Relation in der entgegengesetzten Richtung
gelesen wird.

\subsection{Axiomatische Definition und Schreibweise}

\FormulaDefDeltaK[Infimal induzierte Relation und Schreibweise]{%
  x\MeetLeq{\MeetOp}{A}y
  \quad:\Longleftrightarrow\quad
  (x,y)\in\MeetOrd{\MeetOp}{A}%
}{MeetOrder}{%
  \DeltaRow{Mengen}{A\dsep x\dsep y}
  \DeltaRow{Funktionen}{\MeetOp}
  \DeltaRow{Relationen}{\MeetOrd{\MeetOp}{A}}
  \DeltaRow{\textbf{Axiome}}{}
  \DeltaRow{Relation auf \(A\)}{%
    \MeetOrd{\MeetOp}{A}\subseteq A\times A}
    [\FormulaRefAuto{F \subseteq A \times B}[def]]
  \DeltaRow{Auswertung}{%
    \begin{aligned}[t]
      &\Semi{\MeetOp,A}\dsep x,y\in A\vdash{}\\
      &(x,y)\in\MeetOrd{\MeetOp}{A}
        \leftrightarrow x\MeetOp{}y=x
    \end{aligned}}
    [\FormulaRefAuto{MeetOrderEvaluationAxiom}[ax]]
  \DeltaRow{\textbf{Neue Symbole}}{}
  \DeltaPrem{Infimale Relation}{\MeetOrd{\MeetOp}{A}}
  \DeltaPrem{Infixschreibweise}{\MeetLeq{\MeetOp}{A}}
}

\FormulaAxiomDeltaK[Auswertung der infimal induzierten Relation]{%
  \Semi{\MeetOp,A}\dsep x,y\in A\vdash
  \bigl((x,y)\in\MeetOrd{\MeetOp}{A}
    \leftrightarrow x\MeetOp{}y=x\bigr)%
}{MeetOrderEvaluationAxiom}{%
  \DeltaRow{Mengen}{A\dsep x\dsep y}
  \DeltaRow{Funktionen}{\MeetOp}
  \DeltaRow{Relationen}{\MeetOrd{\MeetOp}{A}}
}

Somit sind
\[
  x\MeetLeq{\MeetOp}{A}y,
  \qquad (x,y)\in\MeetOrd{\MeetOp}{A},
  \qquad x\MeetOp{}y=x
\]
gleichwertige Schreibweisen für dieselbe infimale Vergleichsaussage.

\subsection{Träger und Auswertung}

\FormulaThmDeltaK[Erste Komponente der infimal induzierten Relation]{%
  \Semi{\MeetOp,A}\dsep (x,y)\in\MeetOrd{\MeetOp}{A}\vdash x\in A
}{MeetOrderFirstCarrier}{%
  \DeltaRow{Mengen}{A\dsep x\dsep y}
  \DeltaRow{Funktionen}{\MeetOp}
  \DeltaRow{Relationen}{\MeetOrd{\MeetOp}{A}}
}
\begin{tabproof}
  \proofstep{1}{\Semi{\MeetOp,A}}{\rA}
  \proofstep{2}{(x,y)\in\MeetOrd{\MeetOp}{A}}{\rA}
  \proofstep{}{\MeetOrd{\MeetOp}{A}\subseteq A\times A}{%
    \FormulaRefAuto{F \subseteq A \times B}[def]}
  \proofstep{2}{(x,y)\in A\times A}{%
    \FormulaRefAuto{A\subseteq B,\, x\in A \vdash x\in B}{3,2}}
  \proofstep{2}{x\in A}{%
    \FormulaRefAuto{(a,b)\in A\times B\vdash a\in A}{4}}
\end{tabproof}

\FormulaThmDeltaK[Zweite Komponente der infimal induzierten Relation]{%
  \Semi{\MeetOp,A}\dsep (x,y)\in\MeetOrd{\MeetOp}{A}\vdash y\in A
}{MeetOrderSecondCarrier}{%
  \DeltaRow{Mengen}{A\dsep x\dsep y}
  \DeltaRow{Funktionen}{\MeetOp}
  \DeltaRow{Relationen}{\MeetOrd{\MeetOp}{A}}
}
\begin{tabproof}
  \proofstep{1}{\Semi{\MeetOp,A}}{\rA}
  \proofstep{2}{(x,y)\in\MeetOrd{\MeetOp}{A}}{\rA}
  \proofstep{}{\MeetOrd{\MeetOp}{A}\subseteq A\times A}{%
    \FormulaRefAuto{F \subseteq A \times B}[def]}
  \proofstep{2}{(x,y)\in A\times A}{%
    \FormulaRefAuto{A\subseteq B,\, x\in A \vdash x\in B}{3,2}}
  \proofstep{2}{y\in A}{%
    \FormulaRefAuto{(a,b)\in A\times B\vdash b\in B}{4}}
\end{tabproof}

\FormulaThmDeltaK[Auswertung der infimal induzierten Relation]{%
  \Semi{\MeetOp,A}\dsep (x,y)\in\MeetOrd{\MeetOp}{A}
  \vdash x\MeetOp{}y=x
}{MeetOrderEvaluation}{%
  \DeltaRow{Mengen}{A\dsep x\dsep y}
  \DeltaRow{Funktionen}{\MeetOp}
  \DeltaRow{Relationen}{\MeetOrd{\MeetOp}{A}}
}
\begin{tabproof}
  \proofstep{1}{\Semi{\MeetOp,A}}{\rA}
  \proofstep{2}{(x,y)\in\MeetOrd{\MeetOp}{A}}{\rA}
  \proofstep{1,2}{x\in A}{\FormulaRefAuto{MeetOrderFirstCarrier}[thm]{1,2}}
  \proofstep{1,2}{y\in A}{\FormulaRefAuto{MeetOrderSecondCarrier}[thm]{1,2}}
  \proofstep{1,2}{%
    (x,y)\in\MeetOrd{\MeetOp}{A}\leftrightarrow x\MeetOp{}y=x}{%
    \FormulaRefAuto{MeetOrderEvaluationAxiom}[ax]{1,3,4}}
  \proofstep{1,2}{x\MeetOp{}y=x}{%
    \FormulaRefAuto{P\leftrightarrow Q, P\vdash Q}{5,2}}
\end{tabproof}

\subsection{Die induzierte partielle Ordnung}

\FormulaThmDeltaK[Reflexivität der infimal induzierten Relation]{%
  \Semi{\MeetOp,A}\dsep x\in A\vdash (x,x)\in\MeetOrd{\MeetOp}{A}
}{MeetOrderReflexive}{%
  \DeltaRow{Mengen}{A\dsep x}
  \DeltaRow{Funktionen}{\MeetOp}
}
\begin{tabproof}
  \proofstep{1}{\Semi{\MeetOp,A}}{\rA}
  \proofstep{2}{x\in A}{\rA}
  \proofstep{1,2}{x\MeetOp{}x=x}{%
    \FormulaRefAuto{SemilatticeIdempotenceAxiom}[ax]{1,2}}
  \proofstep{1,2}{%
    (x,x)\in\MeetOrd{\MeetOp}{A}\leftrightarrow x\MeetOp{}x=x}{%
    \FormulaRefAuto{MeetOrderEvaluationAxiom}[ax]{1,2,2}}
  \proofstep{1,2}{(x,x)\in\MeetOrd{\MeetOp}{A}}{%
    \FormulaRefAuto{P\leftrightarrow Q, Q\vdash P}{4,3}}
\end{tabproof}

\FormulaThmDeltaK[Transitivität der infimal induzierten Relation]{%
  \Semi{\MeetOp,A}\dsep (x,y)\in\MeetOrd{\MeetOp}{A}
  \dsep (y,z)\in\MeetOrd{\MeetOp}{A}
  \vdash (x,z)\in\MeetOrd{\MeetOp}{A}
}{MeetOrderTransitive}{%
  \DeltaRow{Mengen}{A\dsep x\dsep y\dsep z}
  \DeltaRow{Funktionen}{\MeetOp}
}
\begin{tabproof}
  \proofstep{1}{\Semi{\MeetOp,A}}{\rA}
  \proofstep{2}{(x,y)\in\MeetOrd{\MeetOp}{A}}{\rA}
  \proofstep{3}{(y,z)\in\MeetOrd{\MeetOp}{A}}{\rA}
  \proofstep{1,2}{x\in A}{\FormulaRefAuto{MeetOrderFirstCarrier}[thm]{1,2}}
  \proofstep{1,2}{y\in A}{\FormulaRefAuto{MeetOrderSecondCarrier}[thm]{1,2}}
  \proofstep{1,3}{z\in A}{\FormulaRefAuto{MeetOrderSecondCarrier}[thm]{1,3}}
  \proofstep{1,2}{x\MeetOp{}y=x}{\FormulaRefAuto{MeetOrderEvaluation}[thm]{1,2}}
  \proofstep{1,3}{y\MeetOp{}z=y}{\FormulaRefAuto{MeetOrderEvaluation}[thm]{1,3}}
  \proofstep{1,2,3}{\bigl(x\MeetOp{}y\bigr)\MeetOp{}z
    =x\MeetOp\bigl(y\MeetOp{}z\bigr)}{%
    \FormulaRefAuto{SemilatticeAssociativityAxiom}[ax]{1,4,5,6}}
  \proofstep{1,2,3}{\bigl(x\MeetOp{}y\bigr)\MeetOp{}z=x\MeetOp{}z}{%
    \rIE{7,\rII}}
  \proofstep{1,2,3}{x\MeetOp\bigl(y\MeetOp{}z\bigr)=x\MeetOp{}y}{%
    \rIE{8,\rII}}
  \proofstep{1,2,3}{x\MeetOp{}z=\bigl(x\MeetOp{}y\bigr)\MeetOp{}z}{%
    \FormulaRefAuto{a=b\vdash b=a}{10}}
  \proofstep{1,2,3}{x\MeetOp{}z=x\MeetOp\bigl(y\MeetOp{}z\bigr)}{%
    \FormulaRefAuto{a=b,\,b=c\vdash a=c}{12,9}}
  \proofstep{1,2,3}{x\MeetOp{}z=x\MeetOp{}y}{%
    \FormulaRefAuto{a=b,\,b=c\vdash a=c}{13,11}}
  \proofstep{1,2,3}{x\MeetOp{}z=x}{%
    \FormulaRefAuto{a=b,\,b=c\vdash a=c}{14,7}}
  \proofstep{1,2,3}{%
    (x,z)\in\MeetOrd{\MeetOp}{A}\leftrightarrow x\MeetOp{}z=x}{%
    \FormulaRefAuto{MeetOrderEvaluationAxiom}[ax]{1,4,6}}
  \proofstep{1,2,3}{(x,z)\in\MeetOrd{\MeetOp}{A}}{%
    \FormulaRefAuto{P\leftrightarrow Q, Q\vdash P}{16,15}}
\end{tabproof}

\FormulaThmDeltaK[Antisymmetrie der infimal induzierten Relation]{%
  \Semi{\MeetOp,A}\dsep (x,y)\in\MeetOrd{\MeetOp}{A}
  \dsep (y,x)\in\MeetOrd{\MeetOp}{A}\vdash x=y
}{MeetOrderAntisymmetric}{%
  \DeltaRow{Mengen}{A\dsep x\dsep y}
  \DeltaRow{Funktionen}{\MeetOp}
}
\begin{tabproof}
  \proofstep{1}{\Semi{\MeetOp,A}}{\rA}
  \proofstep{2}{(x,y)\in\MeetOrd{\MeetOp}{A}}{\rA}
  \proofstep{3}{(y,x)\in\MeetOrd{\MeetOp}{A}}{\rA}
  \proofstep{1,2}{x\in A}{\FormulaRefAuto{MeetOrderFirstCarrier}[thm]{1,2}}
  \proofstep{1,2}{y\in A}{\FormulaRefAuto{MeetOrderSecondCarrier}[thm]{1,2}}
  \proofstep{1,2}{x\MeetOp{}y=x}{\FormulaRefAuto{MeetOrderEvaluation}[thm]{1,2}}
  \proofstep{1,3}{y\MeetOp{}x=y}{\FormulaRefAuto{MeetOrderEvaluation}[thm]{1,3}}
  \proofstep{1,2}{x\MeetOp{}y=y\MeetOp{}x}{%
    \FormulaRefAuto{SemilatticeCommutativityAxiom}[ax]{1,4,5}}
  \proofstep{1,2,3}{x\MeetOp{}y=y}{\rIE{8,7}}
  \proofstep{1,2,3}{x=y}{%
    \FormulaRefAuto{a=b,\,a=c\vdash b=c}{6,9}}
\end{tabproof}

\FormulaThmDeltaK[Ein infimal gelesener Halbverband induziert eine partielle Ordnung]{%
  \Semi{\MeetOp,A}\vdash \PartOrd{\MeetOrd{\MeetOp}{A},A}
}{MeetSemilatticeInducesOrder}{%
  \DeltaRow{Mengen}{A\dsep x\dsep y\dsep z}
  \DeltaRow{Funktionen}{\MeetOp}
  \DeltaRow{Relationen}{\MeetOrd{\MeetOp}{A}}
}
\begin{tabproofsplitwide}
  \proofpartwideindR[Relation auf dem Träger]{%
    \MeetOrd{\MeetOp}{A}\subseteq A\times A
  }{\MeetOrd{\MeetOp}{A}\subseteq A\times A}[MeetOrderCarrierPart]
    \proofstepwidestar[]{\MeetOrd{\MeetOp}{A}\subseteq A\times A}{%
      \FormulaRefAuto{F \subseteq A \times B}[def]}
  \closeproofpartwideindR

  \proofpartwideindR[Reflexivität]{%
    \Semi{\MeetOp,A}\vdash
    x\in A\vdash (x,x)\in\MeetOrd{\MeetOp}{A}
  }{\Semi{\MeetOp,A}\vdash
    x\in A\vdash (x,x)\in\MeetOrd{\MeetOp}{A}}[MeetOrderReflexivePart]
    \proofstepwidestar[1]{\Semi{\MeetOp,A}}{\rA}
    \proofstepwidestar[1]{x\in A\vdash (x,x)\in\MeetOrd{\MeetOp}{A}}{%
      \FormulaRefAuto{MeetOrderReflexive}[thm]{1}}
  \closeproofpartwideindR

  \proofpartwideindR[Transitivität]{%
    \begin{aligned}[t]
      &\Semi{\MeetOp,A}\vdash{}\\[-2pt]
      &(x,y)\in\MeetOrd{\MeetOp}{A}\dsep
        (y,z)\in\MeetOrd{\MeetOp}{A}
        \vdash (x,z)\in\MeetOrd{\MeetOp}{A}
    \end{aligned}
  }{\Semi{\MeetOp,A}\vdash
    (x,y)\in\MeetOrd{\MeetOp}{A}\dsep
    (y,z)\in\MeetOrd{\MeetOp}{A}\vdash
    (x,z)\in\MeetOrd{\MeetOp}{A}}[MeetOrderTransitivePart]
    \proofstepwidestar[1]{\Semi{\MeetOp,A}}{\rA}
    \proofstepwidestar[1]{%
      (x,y)\in\MeetOrd{\MeetOp}{A}\dsep
      (y,z)\in\MeetOrd{\MeetOp}{A}\vdash
      (x,z)\in\MeetOrd{\MeetOp}{A}}{%
      \FormulaRefAuto{MeetOrderTransitive}[thm]{1}}
  \closeproofpartwideindR

  \proofpartwideindR[Antisymmetrie]{%
    \begin{aligned}[t]
      &\Semi{\MeetOp,A}\vdash{}\\[-2pt]
      &(x,y)\in\MeetOrd{\MeetOp}{A}\dsep
        (y,x)\in\MeetOrd{\MeetOp}{A}\vdash x=y
    \end{aligned}
  }{\Semi{\MeetOp,A}\vdash
    (x,y)\in\MeetOrd{\MeetOp}{A}\dsep
    (y,x)\in\MeetOrd{\MeetOp}{A}\vdash x=y}[MeetOrderAntisymmetricPart]
    \proofstepwidestar[1]{\Semi{\MeetOp,A}}{\rA}
    \proofstepwidestar[1]{%
      (x,y)\in\MeetOrd{\MeetOp}{A}\dsep
      (y,x)\in\MeetOrd{\MeetOp}{A}\vdash x=y}{%
      \FormulaRefAuto{MeetOrderAntisymmetric}[thm]{1}}
  \closeproofpartwideindR

  \proofpartwide{Schluss}
    \proofstepwidestar[1]{\Semi{\MeetOp,A}}{\rA}
    \proofstepwidestar[]{\MeetOrd{\MeetOp}{A}\subseteq A\times A}{%
      \FormulaRefAuto{MeetOrderCarrierPart}[thm]}
    \proofstepwidestar[1]{x\in A\vdash (x,x)\in\MeetOrd{\MeetOp}{A}}{%
      \FormulaRefAuto{MeetOrderReflexivePart}[thm]{1}}
    \proofstepwidestar[1]{%
      (x,y)\in\MeetOrd{\MeetOp}{A}\dsep
      (y,z)\in\MeetOrd{\MeetOp}{A}\vdash
      (x,z)\in\MeetOrd{\MeetOp}{A}}{%
      \FormulaRefAuto{MeetOrderTransitivePart}[thm]{1}}
    \proofstepwidestar[1]{%
      (x,y)\in\MeetOrd{\MeetOp}{A}\dsep
      (y,x)\in\MeetOrd{\MeetOp}{A}\vdash x=y}{%
      \FormulaRefAuto{MeetOrderAntisymmetricPart}[thm]{1}}
    \proofstepwidestar[1]{\PartOrd{\MeetOrd{\MeetOp}{A},A}}{%
      \FormulaRefAuto{Partielle Ordnung}[def]{2,3,4,5}}
  \closeproofpartwide
\end{tabproofsplitwide}

\subsection{Dualität der beiden Orientierungen}

\FormulaThmDeltaK[Die beiden induzierten Relationen sind dual]{%
  \Semi{\MeetOp,A}\dsep x,y\in A\vdash
  \bigl((x,y)\in\MeetOrd{\MeetOp}{A}
    \leftrightarrow (y,x)\in\JoinOrd{\MeetOp}{A}\bigr)
}{MeetJoinOrderDual}{%
  \DeltaRow{Mengen}{A\dsep x\dsep y}
  \DeltaRow{Funktionen}{\MeetOp}
  \DeltaRow{Relationen}{\MeetOrd{\MeetOp}{A}\dsep\JoinOrd{\MeetOp}{A}}
}
\begin{tabproofwide}
  \proofstepwidestar[1]{\Semi{\MeetOp,A}}{\rA}
  \proofstepwidestar[2]{x\in A}{\rA}
  \proofstepwidestar[3]{y\in A}{\rA}
  \proofstepwide[1,2,3]{(x,y)\in\MeetOrd{\MeetOp}{A}}{\leftrightarrow}{x\MeetOp{}y=x}{%
    \FormulaRefAuto{MeetOrderEvaluationAxiom}[ax]{1,2,3}}
  \proofstepwide[1,2,3]{}{\leftrightarrow}{y\MeetOp{}x=x}{%
    \rLRS{\FormulaRefAuto{SemilatticeCommutativityAxiom}[ax]{1,2,3},4}}
  \proofstepwide[1,2,3]{}{\leftrightarrow}{(y,x)\in\JoinOrd{\MeetOp}{A}}{%
    \FormulaRefAuto{JoinOrderEvaluationAxiom}[ax]{1,3,2}}
\end{tabproofwide}

\subsection{Universelle Eigenschaft der infimalen Lesart}

\FormulaDefDeltaK[Infimum eines Elementpaares]{%
  \PairInf{R}{A}{x}{y}{u}%
}{PairInfimum}{%
  \DeltaRow{Mengen}{A\dsep x\dsep y\dsep u\dsep v}
  \DeltaRow{Relationen}{R}
  \DeltaRow{Elemente}{x,y,u\in A}
  \DeltaRow{Untere Schranken}{(u,x)\in R\dsep (u,y)\in R}
  \DeltaRow{Größte untere Schranke}{%
    v\in A\dsep (v,x)\in R\dsep (v,y)\in R\vdash (v,u)\in R}
  \DeltaRow{\textbf{Neues Symbol}}{}
  \DeltaPrem{Paarinfimum}{\PairInf{R}{A}{x}{y}{u}}
}

\FormulaThmDeltaK[Die Halbverbandsoperation ist das Paarinfimum]{%
  \Semi{\MeetOp,A}\dsep x,y\in A
  \vdash \PairInf{\MeetOrd{\MeetOp}{A}}{A}{x}{y}{x\MeetOp{}y}
}{MeetIsPairInfimum}{%
  \DeltaRow{Mengen}{A\dsep x\dsep y\dsep u}
  \DeltaRow{Funktionen}{\MeetOp}
}
\begin{tabproofsplitwide}
  \proofpartwideindR[Erste untere Schranke]{%
    \Semi{\MeetOp,A}\dsep x,y\in A\vdash
    (x\MeetOp{}y,x)\in\MeetOrd{\MeetOp}{A}
  }{\Semi{\MeetOp,A}\dsep x,y\in A\vdash
    (x\MeetOp{}y,x)\in\MeetOrd{\MeetOp}{A}}[MeetLowerLeft]
    \proofstepwidestar[1]{\Semi{\MeetOp,A}}{\rA}
    \proofstepwidestar[2]{x\in A}{\rA}
    \proofstepwidestar[3]{y\in A}{\rA}
    \proofstepwidestar[1,2,3]{x\MeetOp{}y\in A}{%
      \FormulaRefAuto{SemilatticeClosureAxiom}[ax]{1,2,3}}
    \proofstepwidestar[1,2,3]{(x,x\MeetOp{}y)\in\JoinOrd{\MeetOp}{A}}{%
      \FormulaRefAuto{JoinUpperLeft}[thm]{1,2,3}}
    \proofstepwidestar[1,2,3]{(x\MeetOp{}y,x)\in\MeetOrd{\MeetOp}{A}}{%
      \FormulaRefAuto{MeetJoinOrderDual}[thm]{1,4,2,5}}
  \closeproofpartwideindR

  \proofpartwideindR[Zweite untere Schranke]{%
    \Semi{\MeetOp,A}\dsep x,y\in A\vdash
    (x\MeetOp{}y,y)\in\MeetOrd{\MeetOp}{A}
  }{\Semi{\MeetOp,A}\dsep x,y\in A\vdash
    (x\MeetOp{}y,y)\in\MeetOrd{\MeetOp}{A}}[MeetLowerRight]
    \proofstepwidestar[1]{\Semi{\MeetOp,A}}{\rA}
    \proofstepwidestar[2]{x\in A}{\rA}
    \proofstepwidestar[3]{y\in A}{\rA}
    \proofstepwidestar[1,2,3]{x\MeetOp{}y\in A}{%
      \FormulaRefAuto{SemilatticeClosureAxiom}[ax]{1,2,3}}
    \proofstepwidestar[1,2,3]{(y,x\MeetOp{}y)\in\JoinOrd{\MeetOp}{A}}{%
      \FormulaRefAuto{JoinUpperRight}[thm]{1,2,3}}
    \proofstepwidestar[1,2,3]{(x\MeetOp{}y,y)\in\MeetOrd{\MeetOp}{A}}{%
      \FormulaRefAuto{MeetJoinOrderDual}[thm]{1,4,3,5}}
  \closeproofpartwideindR

  \proofpartwideindR[Größte untere Schranke]{%
    \begin{aligned}[t]
      &\Semi{\MeetOp,A}\dsep x,y,u\in A\dsep
        (u,x)\in\MeetOrd{\MeetOp}{A}\dsep{}\\[-2pt]
      &(u,y)\in\MeetOrd{\MeetOp}{A}\vdash
        (u,x\MeetOp{}y)\in\MeetOrd{\MeetOp}{A}
    \end{aligned}
  }{\Semi{\MeetOp,A}\dsep x,y,u\in A\dsep
    (u,x)\in\MeetOrd{\MeetOp}{A}\dsep
    (u,y)\in\MeetOrd{\MeetOp}{A}\vdash
    (u,x\MeetOp{}y)\in\MeetOrd{\MeetOp}{A}}[MeetGreatestLower]
    \proofstepwidestar[1]{\Semi{\MeetOp,A}}{\rA}
    \proofstepwidestar[2]{x\in A}{\rA}
    \proofstepwidestar[3]{y\in A}{\rA}
    \proofstepwidestar[4]{u\in A}{\rA}
    \proofstepwidestar[5]{(u,x)\in\MeetOrd{\MeetOp}{A}}{\rA}
    \proofstepwidestar[6]{(u,y)\in\MeetOrd{\MeetOp}{A}}{\rA}
    \proofstepwidestar[1,2,4,5]{(x,u)\in\JoinOrd{\MeetOp}{A}}{%
      \FormulaRefAuto{MeetJoinOrderDual}[thm]{1,4,2,5}}
    \proofstepwidestar[1,3,4,6]{(y,u)\in\JoinOrd{\MeetOp}{A}}{%
      \FormulaRefAuto{MeetJoinOrderDual}[thm]{1,4,3,6}}
    \proofstepwidestar[1,2,3,4,5,6]{(x\MeetOp{}y,u)\in\JoinOrd{\MeetOp}{A}}{%
      \FormulaRefAuto{JoinLeastUpper}[thm]{1,2,3,4,7,8}}
    \proofstepwidestar[1,2,3]{x\MeetOp{}y\in A}{%
      \FormulaRefAuto{SemilatticeClosureAxiom}[ax]{1,2,3}}
    \proofstepwidestar[1,2,3,4,5,6]{(u,x\MeetOp{}y)\in\MeetOrd{\MeetOp}{A}}{%
      \FormulaRefAuto{MeetJoinOrderDual}[thm]{1,4,10,9}}
  \closeproofpartwideindR

  \proofpartwide{Schluss}
    \proofstepwidestar[1]{\Semi{\MeetOp,A}}{\rA}
    \proofstepwidestar[2]{x\in A}{\rA}
    \proofstepwidestar[3]{y\in A}{\rA}
    \proofstepwidestar[1,2,3]{x\MeetOp{}y\in A}{%
      \FormulaRefAuto{SemilatticeClosureAxiom}[ax]{1,2,3}}
    \proofstepwidestar[1,2,3]{(x\MeetOp{}y,x)\in\MeetOrd{\MeetOp}{A}}{%
      \FormulaRefAuto{MeetLowerLeft}[thm]{1,2,3}}
    \proofstepwidestar[1,2,3]{(x\MeetOp{}y,y)\in\MeetOrd{\MeetOp}{A}}{%
      \FormulaRefAuto{MeetLowerRight}[thm]{1,2,3}}
    \proofstepwidestar[1]{%
      u\in A\dsep (u,x)\in\MeetOrd{\MeetOp}{A}\dsep
      (u,y)\in\MeetOrd{\MeetOp}{A}\vdash
      (u,x\MeetOp{}y)\in\MeetOrd{\MeetOp}{A}}{%
      \FormulaRefAuto{MeetGreatestLower}[thm]{1,2,3}}
    \proofstepwidestar[1,2,3]{%
      \PairInf{\MeetOrd{\MeetOp}{A}}{A}{x}{y}{x\MeetOp{}y}}{%
      \FormulaRefAuto{PairInfimum}[def]{\rAI{\rAI{2,3},4},5,6,7}}
  \closeproofpartwide
\end{tabproofsplitwide}

\section{Vereinigung und Durchschnitt als Halbverbände}

Die beiden grundlegenden Mengenoperationen aus Band~03 liefern die ersten
konkreten Beispiele der allgemeinen Axiomatik. Wir zeigen zunächst jede
Halbverbandsstruktur für sich und bestimmen anschließend ihre induzierte
Relation.

\FormulaThmDeltaK[Abgeschlossenheit der Vereinigung auf einer Potenzmenge]{%
  X,Y\in\powerset(S)\vdash X\cup Y\in\powerset(S)
}{PowerSetUnionClosure}{%
  \DeltaRow{Mengen}{S\dsep X\dsep Y}
}
\begin{tabproof}
  \proofstep{1}{X\in\powerset(S)}{\rA}
  \proofstep{2}{Y\in\powerset(S)}{\rA}
  \proofstep{1}{X\subseteq S}{%
    \FormulaRefAuto{x\in\powerset(A)\eqvdash x\subseteq A}{1}}
  \proofstep{2}{Y\subseteq S}{%
    \FormulaRefAuto{x\in\powerset(A)\eqvdash x\subseteq A}{2}}
  \proofstep{1,2}{X\cup Y\subseteq S}{%
    \FormulaRefAuto{A\subseteq M\dsep B\subseteq M \vdash A\cup B\subseteq M}{3,4}}
  \proofstep{1,2}{X\cup Y\in\powerset(S)}{%
    \FormulaRefAuto{x\in\powerset(A)\eqvdash x\subseteq A}{5}}
\end{tabproof}

\FormulaThmDeltaK[Vereinigung ist eine Halbverbandsoperation]{%
  \Semi{\cup,\powerset(S)}
}{PowerSetUnionSemilattice}{%
  \DeltaRow{Mengen}{S\dsep X\dsep Y\dsep Z}
  \DeltaRow{Funktionen}{\cup}
}
\begin{tabproofwide}
  \proofstepwidestar[]{X,Y\in\powerset(S)\vdash X\cup Y\in\powerset(S)}{%
    \FormulaRefAuto{PowerSetUnionClosure}[thm]}
  \proofstepwidestar[]{X,Y,Z\in\powerset(S)\vdash
    (X\cup Y)\cup Z=X\cup(Y\cup Z)}{%
    \FormulaRefAuto{(A \cup B) \cup C = A \cup (B \cup C)}}
  \proofstepwidestar[]{X,Y\in\powerset(S)\vdash X\cup Y=Y\cup X}{%
    \FormulaRefAuto{A \cup B = B \cup A}}
  \proofstepwidestar[]{X\in\powerset(S)\vdash X\cup X=X}{%
    \FormulaRefAuto{A \cup A = A}}
  \proofstepwidestar[]{\Semi{\cup,\powerset(S)}}{%
    \FormulaRefAuto{JoinSemilattice}[def]{1,2,3,4}}
\end{tabproofwide}

\FormulaThmDeltaK[Abgeschlossenheit des Durchschnitts auf einer Potenzmenge]{%
  X,Y\in\powerset(S)\vdash X\cap Y\in\powerset(S)
}{PowerSetIntersectionClosure}{%
  \DeltaRow{Mengen}{S\dsep X\dsep Y}
}
\begin{tabproof}
  \proofstep{1}{X\in\powerset(S)}{\rA}
  \proofstep{2}{Y\in\powerset(S)}{\rA}
  \proofstep{1}{X\subseteq S}{%
    \FormulaRefAuto{x\in\powerset(A)\eqvdash x\subseteq A}{1}}
  \proofstep{1}{X\cap Y\subseteq S}{%
    \FormulaRefAuto{A\subseteq M \vdash A\cap B\subseteq M}{3}}
  \proofstep{1,2}{X\cap Y\in\powerset(S)}{%
    \FormulaRefAuto{x\in\powerset(A)\eqvdash x\subseteq A}{4}}
\end{tabproof}

\FormulaThmDeltaK[Durchschnitt ist eine Halbverbandsoperation]{%
  \Semi{\cap,\powerset(S)}
}{PowerSetIntersectionSemilattice}{%
  \DeltaRow{Mengen}{S\dsep X\dsep Y\dsep Z}
  \DeltaRow{Funktionen}{\cap}
}
\begin{tabproofwide}
  \proofstepwidestar[]{X,Y\in\powerset(S)\vdash X\cap Y\in\powerset(S)}{%
    \FormulaRefAuto{PowerSetIntersectionClosure}[thm]}
  \proofstepwidestar[]{X,Y,Z\in\powerset(S)\vdash
    (X\cap Y)\cap Z=X\cap(Y\cap Z)}{%
    \FormulaRefAuto{(A \cap B) \cap C = A \cap (B \cap C)}}
  \proofstepwidestar[]{X,Y\in\powerset(S)\vdash X\cap Y=Y\cap X}{%
    \FormulaRefAuto{A \cap B = B \cap A}}
  \proofstepwidestar[]{X\in\powerset(S)\vdash X\cap X=X}{%
    \FormulaRefAuto{a=b\vdash b=a}{\FormulaRefAuto{A = A \cap A}}}
  \proofstepwidestar[]{\Semi{\cap,\powerset(S)}}{%
    \FormulaRefAuto{JoinSemilattice}[def]{1,2,3,4}}
\end{tabproofwide}

\FormulaThmDeltaK[Inklusion durch Vereinigung]{%
  X,Y\in\powerset(S)\vdash
  \bigl(X\subseteq Y\leftrightarrow X\cup Y=Y\bigr)
}{SubsetIffUnionEqualsRight}{%
  \DeltaRow{Mengen}{S\dsep X\dsep Y}
}
\begin{tabproofsplitwide}
  \proofpartwide{\(\vdash\)}
    \proofstepwidestar[1]{X\in\powerset(S)}{\rA}
    \proofstepwidestar[2]{Y\in\powerset(S)}{\rA}
    \proofstepwidestar[3]{X\subseteq Y}{\rA}
    \proofstepwidestar[]{Y\subseteq Y}{\FormulaRefAuto{A\subseteq A}}
    \proofstepwidestar[3]{X\cup Y\subseteq Y}{%
      \FormulaRefAuto{A\subseteq C,\, B\subseteq C \vdash A\cup B\subseteq C}{3,4}}
    \proofstepwidestar[]{Y\subseteq Y\cup X}{%
      \FormulaRefAuto{A \subseteq A \cup B}}
    \proofstepwidestar[]{Y\cup X=X\cup Y}{%
      \FormulaRefAuto{A \cup B = B \cup A}}
    \proofstepwidestar[]{Y\subseteq X\cup Y}{\rIE{7,6}}
    \proofstepwidestar[3]{X\cup Y=Y}{%
      \FormulaRefAuto{A \subseteq B \land B \subseteq A \eqvdash A = B}{5,8}}
  \closeproofpartwide

  \proofpartwide{\(\dashv\)}
    \proofstepwidestar[1]{X\in\powerset(S)}{\rA}
    \proofstepwidestar[2]{Y\in\powerset(S)}{\rA}
    \proofstepwidestar[3]{X\cup Y=Y}{\rA}
    \proofstepwidestar[]{X\subseteq X\cup Y}{%
      \FormulaRefAuto{A \subseteq A \cup B}}
    \proofstepwidestar[3]{X\subseteq Y}{\rIE{3,4}}
  \closeproofpartwide
\end{tabproofsplitwide}

\FormulaThmDeltaK[Die Vereinigungsrelation ist die Inklusion]{%
  X,Y\in\powerset(S)\vdash
  \bigl((X,Y)\in\JoinOrd{\cup}{\powerset(S)}
    \leftrightarrow X\subseteq Y\bigr)
}{PowerSetUnionOrderIsInclusion}{%
  \DeltaRow{Mengen}{S\dsep X\dsep Y}
  \DeltaRow{Relationen}{\JoinOrd{\cup}{\powerset(S)}}
}
\begin{tabproofwide}
  \proofstepwidestar[1]{X\in\powerset(S)}{\rA}
  \proofstepwidestar[2]{Y\in\powerset(S)}{\rA}
  \proofstepwidestar[]{\Semi{\cup,\powerset(S)}}{%
    \FormulaRefAuto{PowerSetUnionSemilattice}[thm]}
  \proofstepwide[1,2]{(X,Y)\in\JoinOrd{\cup}{\powerset(S)}}{\leftrightarrow}{X\cup Y=Y}{%
    \FormulaRefAuto{JoinOrderEvaluationAxiom}[ax]{3,1,2}}
  \proofstepwide[1,2]{}{\leftrightarrow}{X\subseteq Y}{%
    \FormulaRefAuto{SubsetIffUnionEqualsRight}[thm]{1,2}}
\end{tabproofwide}

\FormulaThmDeltaK[Die Durchschnittsrelation ist die Inklusion]{%
  X,Y\in\powerset(S)\vdash
  \bigl((X,Y)\in\MeetOrd{\cap}{\powerset(S)}
    \leftrightarrow X\subseteq Y\bigr)
}{PowerSetIntersectionOrderIsInclusion}{%
  \DeltaRow{Mengen}{S\dsep X\dsep Y}
  \DeltaRow{Relationen}{\MeetOrd{\cap}{\powerset(S)}}
}
\begin{tabproofwide}
  \proofstepwidestar[1]{X\in\powerset(S)}{\rA}
  \proofstepwidestar[2]{Y\in\powerset(S)}{\rA}
  \proofstepwidestar[]{\Semi{\cap,\powerset(S)}}{%
    \FormulaRefAuto{PowerSetIntersectionSemilattice}[thm]}
  \proofstepwide[1,2]{(X,Y)\in\MeetOrd{\cap}{\powerset(S)}}{\leftrightarrow}{X\cap Y=X}{%
    \FormulaRefAuto{MeetOrderEvaluationAxiom}[ax]{3,1,2}}
  \proofstepwide[1,2]{}{\leftrightarrow}{X\subseteq Y}{%
    \FormulaRefAuto{A\subseteq B \eqvdash A \cap B = A}}
\end{tabproofwide}

\chapter{Verbände}

\section{Axiomatische Definition}

Zwei beliebige Halbverbandsoperationen auf derselben Menge müssen nicht
dieselbe Ordnung erzeugen. Erst die beiden Absorptionsgesetze verbinden sie zu
einem Verband. Die folgende Axiomtafel führt alle acht Halbverbandseigenschaften
und beide Absorptionseigenschaften ausdrücklich auf.

\FormulaAxiomDeltaK[Absorption des Infimums]{%
  x,y\in A\vdash x\MeetOp\bigl(x\JoinOp{}y\bigr)=x%
}{LatticeMeetAbsorptionLaw}{%
  \DeltaRow{Mengen}{A\dsep x\dsep y}
  \DeltaRow{Funktionen}{\MeetOp\dsep \JoinOp}
}
\par\smallskip

\FormulaAxiomDeltaK[Absorption des Supremums]{%
  x,y\in A\vdash x\JoinOp\bigl(x\MeetOp{}y\bigr)=x%
}{LatticeJoinAbsorptionLaw}{%
  \DeltaRow{Mengen}{A\dsep x\dsep y}
  \DeltaRow{Funktionen}{\MeetOp\dsep \JoinOp}
}

\FormulaDefDeltaK[Begriff des algebraischen Verbandes]{%
  \LatticeAlg{\MeetOp,\JoinOp,A}%
}{AlgebraicLattice}{%
  \DeltaRow{Mengen}{A\dsep x\dsep y\dsep z}
  \DeltaRow{Funktionen}{\MeetOp\dsep \JoinOp}
  \DeltaRow{\textbf{Axiome für \(\MeetOp\)}}{}
  \DeltaRow{Abgeschlossenheit}{%
    x,y\in A\vdash x\MeetOp{}y\in A}
    [\FormulaRefAuto{SemilatticeClosureLaw}[ax]]
  \DeltaRow{Assoziativität}{%
    \begin{aligned}[t]
      &x,y,z\in A\vdash{}\\[-2pt]
      &(x\MeetOp{}y)\MeetOp{}z=x\MeetOp(y\MeetOp{}z)
    \end{aligned}}
    [\FormulaRefAuto{SemilatticeAssociativityLaw}[ax]]
  \DeltaRow{Kommutativität}{%
    x,y\in A\vdash x\MeetOp{}y=y\MeetOp{}x}
    [\FormulaRefAuto{SemilatticeCommutativityLaw}[ax]]
  \DeltaRow{Idempotenz}{%
    x\in A\vdash x\MeetOp{}x=x}
    [\FormulaRefAuto{SemilatticeIdempotenceLaw}[ax]]
  \DeltaRow{\textbf{Axiome für \(\JoinOp\)}}{}
  \DeltaRow{Abgeschlossenheit}{%
    x,y\in A\vdash x\JoinOp{}y\in A}
    [\FormulaRefAuto{SemilatticeClosureLaw}[ax]]
  \DeltaRow{Assoziativität}{%
    \begin{aligned}[t]
      &x,y,z\in A\vdash{}\\[-2pt]
      &(x\JoinOp{}y)\JoinOp{}z=x\JoinOp(y\JoinOp{}z)
    \end{aligned}}
    [\FormulaRefAuto{SemilatticeAssociativityLaw}[ax]]
  \DeltaRow{Kommutativität}{%
    x,y\in A\vdash x\JoinOp{}y=y\JoinOp{}x}
    [\FormulaRefAuto{SemilatticeCommutativityLaw}[ax]]
  \DeltaRow{Idempotenz}{%
    x\in A\vdash x\JoinOp{}x=x}
    [\FormulaRefAuto{SemilatticeIdempotenceLaw}[ax]]
  \DeltaRow{\textbf{Verträglichkeitsaxiome}}{}
  \DeltaRow{\(\MeetOp\)-Absorption}{%
    x,y\in A\vdash x\MeetOp(x\JoinOp{}y)=x}
    [\FormulaRefAuto{LatticeMeetAbsorptionLaw}[ax]]
  \DeltaRow{\(\JoinOp\)-Absorption}{%
    x,y\in A\vdash x\JoinOp(x\MeetOp{}y)=x}
    [\FormulaRefAuto{LatticeJoinAbsorptionLaw}[ax]]
  \DeltaRow{\textbf{Neues Symbol}}{}
  \DeltaPrem{Verbände}{\LatticeAlg{\MeetOp,\JoinOp,A}}
}

\FormulaAxiomDeltaK[Zugriff auf die infimale Halbverbandsstruktur eines Verbandes]{%
  \LatticeAlg{\MeetOp,\JoinOp,A}\vdash \Semi{\MeetOp,A}%
}{LatticeMeetSemilatticeAxiom}{%
  \DeltaRow{Mengen}{A}
  \DeltaRow{Funktionen}{\MeetOp\dsep\JoinOp}
}
\par\smallskip

\FormulaAxiomDeltaK[Zugriff auf die supremale Halbverbandsstruktur eines Verbandes]{%
  \LatticeAlg{\MeetOp,\JoinOp,A}\vdash \Semi{\JoinOp,A}%
}{LatticeJoinSemilatticeAxiom}{%
  \DeltaRow{Mengen}{A}
  \DeltaRow{Funktionen}{\MeetOp\dsep\JoinOp}
}
\par\smallskip

\FormulaAxiomDeltaK[Zugriff auf das Absorptionsaxiom für \(\MeetOp\)]{%
  \LatticeAlg{\MeetOp,\JoinOp,A}\dsep x,y\in A
  \vdash x\MeetOp\bigl(x\JoinOp{}y\bigr)=x%
}{LatticeMeetAbsorptionAxiom}{%
  \DeltaRow{Mengen}{A\dsep x\dsep y}
  \DeltaRow{Funktionen}{\MeetOp\dsep \JoinOp}
}
\par\smallskip

\FormulaAxiomDeltaK[Zugriff auf das Absorptionsaxiom für \(\JoinOp\)]{%
  \LatticeAlg{\MeetOp,\JoinOp,A}\dsep x,y\in A
  \vdash x\JoinOp\bigl(x\MeetOp{}y\bigr)=x%
}{LatticeJoinAbsorptionAxiom}{%
  \DeltaRow{Mengen}{A\dsep x\dsep y}
  \DeltaRow{Funktionen}{\MeetOp\dsep \JoinOp}
}

\FormulaThmDeltaK[Übereinstimmung der beiden Verbandsordnungen]{%
  \begin{aligned}[t]
    &\LatticeAlg{\MeetOp,\JoinOp,A}\dsep x,y\in A\vdash{}\\
    &(x,y)\in\MeetOrd{\MeetOp}{A}
      \leftrightarrow (x,y)\in\JoinOrd{\JoinOp}{A}
  \end{aligned}
}{LatticeOrdersCoincide}{%
  \DeltaRow{Mengen}{A\dsep x\dsep y}
  \DeltaRow{Funktionen}{\MeetOp\dsep \JoinOp}
}
\begin{tabproofsplitwide}
  \proofpartwide{\(\vdash\)}
    \proofstepwidestar[1]{\LatticeAlg{\MeetOp,\JoinOp,A}}{\rA}
    \proofstepwidestar[2]{x\in A}{\rA}
    \proofstepwidestar[3]{y\in A}{\rA}
    \proofstepwidestar[4]{(x,y)\in\MeetOrd{\MeetOp}{A}}{\rA}
    \proofstepwidestar[1]{\Semi{\MeetOp,A}}{%
      \FormulaRefAuto{LatticeMeetSemilatticeAxiom}[ax]{1}}
    \proofstepwidestar[1]{\Semi{\JoinOp,A}}{%
      \FormulaRefAuto{LatticeJoinSemilatticeAxiom}[ax]{1}}
    \proofstepwidestar[1,2,3,4]{x\MeetOp{}y=x}{%
      \FormulaRefAuto{MeetOrderEvaluation}[thm]{5,4}}
    \proofstepwidestar[1,2,3]{y\MeetOp{}x=x\MeetOp{}y}{%
      \FormulaRefAuto{SemilatticeCommutativityAxiom}[ax]{5,3,2}}
    \proofstepwidestar[1,2,3,4]{y\MeetOp{}x=x}{\rIE{7,8}}
    \proofstepwidestar[1,2,3]{y\JoinOp\bigl(y\MeetOp{}x\bigr)=y}{%
      \FormulaRefAuto{LatticeJoinAbsorptionAxiom}[ax]{1,3,2}}
    \proofstepwidestar[1,2,3,4]{y\JoinOp{}x=y}{\rIE{9,10}}
    \proofstepwidestar[1,2,3]{y\JoinOp{}x=x\JoinOp{}y}{%
      \FormulaRefAuto{SemilatticeCommutativityAxiom}[ax]{6,3,2}}
    \proofstepwidestar[1,2,3,4]{x\JoinOp{}y=y}{%
      \FormulaRefAuto{a=b,\,a=c\vdash b=c}{12,11}}
    \proofstepwidestar[1,2,3,4]{(x,y)\in\JoinOrd{\JoinOp}{A}}{%
      \FormulaRefAuto{P\leftrightarrow Q, Q\vdash P}{%
        \FormulaRefAuto{JoinOrderEvaluationAxiom}[ax]{6,2,3},13}}
  \closeproofpartwide

  \proofpartwide{\(\dashv\)}
    \proofstepwidestar[1]{\LatticeAlg{\MeetOp,\JoinOp,A}}{\rA}
    \proofstepwidestar[2]{x\in A}{\rA}
    \proofstepwidestar[3]{y\in A}{\rA}
    \proofstepwidestar[4]{(x,y)\in\JoinOrd{\JoinOp}{A}}{\rA}
    \proofstepwidestar[1]{\Semi{\MeetOp,A}}{%
      \FormulaRefAuto{LatticeMeetSemilatticeAxiom}[ax]{1}}
    \proofstepwidestar[1]{\Semi{\JoinOp,A}}{%
      \FormulaRefAuto{LatticeJoinSemilatticeAxiom}[ax]{1}}
    \proofstepwidestar[1,2,3,4]{x\JoinOp{}y=y}{%
      \FormulaRefAuto{JoinOrderEvaluation}[thm]{6,4}}
    \proofstepwidestar[1,2,3]{x\MeetOp\bigl(x\JoinOp{}y\bigr)=x}{%
      \FormulaRefAuto{LatticeMeetAbsorptionAxiom}[ax]{1,2,3}}
    \proofstepwidestar[1,2,3,4]{x\MeetOp{}y=x}{\rIE{7,8}}
    \proofstepwidestar[1,2,3,4]{(x,y)\in\MeetOrd{\MeetOp}{A}}{%
      \FormulaRefAuto{P\leftrightarrow Q, Q\vdash P}{%
        \FormulaRefAuto{MeetOrderEvaluationAxiom}[ax]{5,2,3},9}}
  \closeproofpartwide
\end{tabproofsplitwide}

Somit gilt in jedem Verband die charakteristische Dreifachäquivalenz
\[
  x\MeetLeq{\MeetOp}{A}y
  \quad\Longleftrightarrow\quad
  x\MeetOp{}y=x
  \quad\Longleftrightarrow\quad
  x\JoinOp{}y=y
  \quad\Longleftrightarrow\quad
  x\JoinLeq{\JoinOp}{A}y.
\]

\section{Das Potenzmengengitter}

Die beiden Halbverbandsbeispiele aus Kapitel~1 werden nun durch die getrennt
zu beweisenden Absorptionsgesetze zu einem Verband verbunden.

\FormulaThmDeltaK[Absorption des Durchschnitts über der Vereinigung]{%
  X,Y\in\powerset(S)\vdash X\cap(X\cup Y)=X
}{PowerSetMeetAbsorption}{%
  \DeltaRow{Mengen}{S\dsep X\dsep Y}
}
\begin{tabproof}
  \proofstep{1}{X\in\powerset(S)}{\rA}
  \proofstep{2}{Y\in\powerset(S)}{\rA}
  \proofstep{}{X\subseteq X\cup Y}{\FormulaRefAuto{A \subseteq A \cup B}}
  \proofstep{}{X\cap(X\cup Y)=X}{%
    \FormulaRefAuto{A \subseteq B \eqvdash A \cap B = A}{3}}
\end{tabproof}

\FormulaThmDeltaK[Absorption der Vereinigung über dem Durchschnitt]{%
  X,Y\in\powerset(S)\vdash X\cup(X\cap Y)=X
}{PowerSetJoinAbsorption}{%
  \DeltaRow{Mengen}{S\dsep X\dsep Y}
}
\begin{tabproof}
  \proofstep{1}{X\in\powerset(S)}{\rA}
  \proofstep{2}{Y\in\powerset(S)}{\rA}
  \proofstep{}{X\cap Y\subseteq X}{\FormulaRefAuto{A \cap B \subseteq A}}
  \proofstep{}{X\subseteq X}{\FormulaRefAuto{A\subseteq A}}
  \proofstep{}{X\cup(X\cap Y)\subseteq X}{%
    \FormulaRefAuto{A \subseteq C,\, B \subseteq C \vdash A \cup B \subseteq C}{4,3}}
  \proofstep{}{X\subseteq X\cup(X\cap Y)}{%
    \FormulaRefAuto{A \subseteq A \cup B}}
  \proofstep{}{X\cup(X\cap Y)=X}{%
    \FormulaRefAuto{A \subseteq B \land B \subseteq A \eqvdash A = B}{5,6}}
\end{tabproof}

\FormulaThmDeltaK[Potenzmengen bilden Verbände]{%
  \LatticeAlg{\cap,\cup,\powerset(S)}
}{PowerSetLattice}{%
  \DeltaRow{Mengen}{S\dsep X\dsep Y\dsep Z}
  \DeltaRow{Funktionen}{\cap\dsep\cup}
}
\begin{tabproofsplitwide}
  \proofpartwide{Axiome}
    \proofstepwidestar[]{\Semi{\cap,\powerset(S)}}{%
      \FormulaRefAuto{PowerSetIntersectionSemilattice}[thm]}
    \proofstepwidestar[]{X,Y\in\powerset(S)\vdash X\cap Y\in\powerset(S)}{%
      \FormulaRefAuto{SemilatticeClosureAxiom}[ax]{1}}
    \proofstepwidestar[]{X,Y,Z\in\powerset(S)\vdash
      (X\cap Y)\cap Z=X\cap(Y\cap Z)}{%
      \FormulaRefAuto{SemilatticeAssociativityAxiom}[ax]{1}}
    \proofstepwidestar[]{X,Y\in\powerset(S)\vdash X\cap Y=Y\cap X}{%
      \FormulaRefAuto{SemilatticeCommutativityAxiom}[ax]{1}}
    \proofstepwidestar[]{X\in\powerset(S)\vdash X\cap X=X}{%
      \FormulaRefAuto{SemilatticeIdempotenceAxiom}[ax]{1}}
    \proofstepwidestar[]{\Semi{\cup,\powerset(S)}}{%
      \FormulaRefAuto{PowerSetUnionSemilattice}[thm]}
    \proofstepwidestar[]{X,Y\in\powerset(S)\vdash X\cup Y\in\powerset(S)}{%
      \FormulaRefAuto{SemilatticeClosureAxiom}[ax]{6}}
    \proofstepwidestar[]{X,Y,Z\in\powerset(S)\vdash
      (X\cup Y)\cup Z=X\cup(Y\cup Z)}{%
      \FormulaRefAuto{SemilatticeAssociativityAxiom}[ax]{6}}
    \proofstepwidestar[]{X,Y\in\powerset(S)\vdash X\cup Y=Y\cup X}{%
      \FormulaRefAuto{SemilatticeCommutativityAxiom}[ax]{6}}
    \proofstepwidestar[]{X\in\powerset(S)\vdash X\cup X=X}{%
      \FormulaRefAuto{SemilatticeIdempotenceAxiom}[ax]{6}}
    \proofstepwidestar[]{X,Y\in\powerset(S)\vdash X\cap(X\cup Y)=X}{%
      \FormulaRefAuto{PowerSetMeetAbsorption}[thm]}
    \proofstepwidestar[]{X,Y\in\powerset(S)\vdash X\cup(X\cap Y)=X}{%
      \FormulaRefAuto{PowerSetJoinAbsorption}[thm]}
    \proofstepwidestar[]{\LatticeAlg{\cap,\cup,\powerset(S)}}{%
      \FormulaRefAuto{AlgebraicLattice}[def]{2,3,4,5,7,8,9,10,11,12}}
  \closeproofpartwide
\end{tabproofsplitwide}

Nach den beiden bereits bewiesenen Ordnungscharakterisierungen ist die
gemeinsame Verbandsordnung genau die Inklusion:
\[
  X\subseteq Y
  \quad\Longleftrightarrow\quad
  X\cap Y=X
  \quad\Longleftrightarrow\quad
  X\cup Y=Y.
\]

\chapter{Frankls Vermutung in Halbverbandsform}

\section{Endliche Supremums-Halbverbände mit Nullelement}

Die in Band~05 eingeführte Frankl-Vermutung besitzt eine rein
halbverbandstheoretische Form. Damit die Formulierung exakt äquivalent ist,
muss der Halbverband ein Nullelement besitzen. Dieses entspricht in einer
vereinigungsabgeschlossenen Familie der leeren Menge.

Zusätzlich zu den Halbverbandsaxiomen werden die Trägerzugehörigkeit des
Nullelements und sein Neutralitätsaxiom als eigene Strukturaxiome festgehalten.

\FormulaAxiomDeltaK[Halbverbandsaxiom einer Nullelementstruktur]{%
  \Semi{\JoinOp,A}%
}{ZeroJoinBaseSemilatticeLaw}{%
  \DeltaRow{Mengen}{A\dsep 0}
  \DeltaRow{Funktionen}{\JoinOp}
}
\par\smallskip

\FormulaAxiomDeltaK[Trägeraxiom eines Nullelements]{%
  0\in A%
}{ZeroJoinCarrierLaw}{%
  \DeltaRow{Mengen}{A\dsep 0}
  \DeltaRow{Funktionen}{\JoinOp}
}
\par\smallskip

\FormulaAxiomDeltaK[Neutralitätsaxiom eines Nullelements]{%
  x\in A\vdash 0\JoinOp{}x=x%
}{ZeroJoinIdentityLaw}{%
  \DeltaRow{Mengen}{A\dsep x\dsep 0}
  \DeltaRow{Funktionen}{\JoinOp}
}

\FormulaDefDeltaK[Supremums-Halbverband mit Nullelement]{%
  \ZeroJoinSemi{\JoinOp,A,0}%
}{ZeroJoinSemilattice}{%
  \DeltaRow{Mengen}{A\dsep x\dsep 0}
  \DeltaRow{Funktionen}{\JoinOp}
  \DeltaRow{\textbf{Axiome}}{}
  \DeltaRow{Halbverband}
           {\Semi{\JoinOp,A}}
           [\FormulaRefAuto{ZeroJoinBaseSemilatticeLaw}[ax]]
  \DeltaRow{Nullelement}
           {0\in A}
           [\FormulaRefAuto{ZeroJoinCarrierLaw}[ax]]
  \DeltaRow{Neutralitätsaxiom}
           {x\in A\vdash 0\JoinOp{}x=x}
           [\FormulaRefAuto{ZeroJoinIdentityLaw}[ax]]
  \DeltaRow{\textbf{Neues Symbol}}{}
  \DeltaPrem{\makecell[l]{Supremums-Halbverbände\\mit Nullelement}}{%
    \ZeroJoinSemi{\JoinOp,A,0}}
}

\newpage
\FormulaAxiomDeltaK[Zugriff auf das Halbverbandsaxiom der Nullelementstruktur]{%
  \ZeroJoinSemi{\JoinOp,A,0}\vdash \Semi{\JoinOp,A}%
}{ZeroJoinBaseSemilatticeAxiom}{%
  \DeltaRow{Mengen}{A\dsep 0}
  \DeltaRow{Funktionen}{\JoinOp}
}
\par\smallskip

\FormulaAxiomDeltaK[Zugriff auf das Nullelement im Träger]{%
  \ZeroJoinSemi{\JoinOp,A,0}\vdash 0\in A%
}{ZeroJoinCarrierAxiom}{%
  \DeltaRow{Mengen}{A\dsep 0}
  \DeltaRow{Funktionen}{\JoinOp}
}

\FormulaAxiomDeltaK[Zugriff auf das Neutralitätsaxiom des Nullelements]{%
  \ZeroJoinSemi{\JoinOp,A,0}\dsep x\in A\vdash 0\JoinOp{}x=x%
}{ZeroJoinIdentityAxiom}{%
  \DeltaRow{Mengen}{A\dsep x\dsep 0}
  \DeltaRow{Funktionen}{\JoinOp}
}
\leavevmode\par\medskip

\FormulaThmDeltaK[Das Nullelement ist das kleinste Element]{%
  \ZeroJoinSemi{\JoinOp,A,0}\dsep x\in A
  \vdash (0,x)\in\JoinOrd{\JoinOp}{A}
}{ZeroJoinLeast}{%
  \DeltaRow{Mengen}{A\dsep x\dsep 0}
  \DeltaRow{Funktionen}{\JoinOp}
}
\begin{tabproof}
  \proofstep{1}{\ZeroJoinSemi{\JoinOp,A,0}}{\rA}
  \proofstep{2}{x\in A}{\rA}
  \proofstep{1}{0\in A}{\FormulaRefAuto{ZeroJoinCarrierAxiom}[ax]{1}}
  \proofstep{1,2}{0\JoinOp{}x=x}{\FormulaRefAuto{ZeroJoinIdentityAxiom}[ax]{1,2}}
  \proofstep{1,2}{%
    (0,x)\in\JoinOrd{\JoinOp}{A}\leftrightarrow 0\JoinOp{}x=x}{%
    \FormulaRefAuto{JoinOrderEvaluationAxiom}[ax]{%
      \FormulaRefAuto{ZeroJoinBaseSemilatticeAxiom}[ax]{1},3,2}}
  \proofstep{1,2}{(0,x)\in\JoinOrd{\JoinOp}{A}}{%
    \FormulaRefAuto{P\leftrightarrow Q, Q\vdash P}{5,4}}
\end{tabproof}

\FormulaDefDeltaK[\(\JoinOp\)-irreduzibles Element]{%
  \JoinIrr{\JoinOp,A,0}{j}%
}{JoinIrreducible}{%
  \DeltaRow{Mengen}{A\dsep 0\dsep j\dsep x\dsep y}
  \DeltaRow{Funktionen}{\JoinOp}
  \DeltaPrem{Halbverband mit Nullelement}{\ZeroJoinSemi{\JoinOp,A,0}}
  \DeltaRow{Von Null verschieden}{j\in A\dsep j\neq 0}
  \DeltaRow{Irreduzibilität}{%
    x,y\in A\dsep x\JoinOp{}y=j\vdash x=j\lor y=j}
  \DeltaRow{\textbf{Neues Symbol}}{}
  \DeltaPrem{\(\JoinOp\)-Irreduzible}{\JoinIrr{\JoinOp,A,0}{j}}
}
\leavevmode\par\medskip

\FormulaDefDeltaK[Hauptfilter eines Elements]{%
  \PrincipalFilter{R,A}{j}
  \coloneqq \{x\in A\mid (j,x)\in R\}%
}{PrincipalFilter}{%
  \DeltaRow{Mengen}{A\dsep j\dsep x}
  \DeltaRow{Relationen}{R}
  \DeltaRow{\textbf{Neues Symbol}}{}
  \DeltaPrem{Hauptfilter}{\PrincipalFilter{R,A}{j}}
}
\leavevmode\par\medskip

\FormulaThmDeltaK[Charakterisierung des Hauptfilters im Halbverband]{%
  \begin{aligned}[t]
    &\Semi{\JoinOp,A}\dsep j,x\in A\vdash{}\\
    &x\in\PrincipalFilter{\JoinOrd{\JoinOp}{A},A}{j}
      \leftrightarrow j\JoinOp{}x=x
  \end{aligned}
}{JoinPrincipalFilterChar}{%
  \DeltaRow{Mengen}{A\dsep j\dsep x}
  \DeltaRow{Funktionen}{\JoinOp}
}
\begin{tabproofwide}
  \proofstepwidestar[1]{\Semi{\JoinOp,A}}{\rA}
  \proofstepwidestar[2]{j\in A}{\rA}
  \proofstepwidestar[3]{x\in A}{\rA}
  \proofstepwide[2,3]{x\in\PrincipalFilter{\JoinOrd{\JoinOp}{A},A}{j}}{\leftrightarrow}{%
    (j,x)\in\JoinOrd{\JoinOp}{A}}{\FormulaRefAuto{PrincipalFilter}[def]}
  \proofstepwide[2,3]{}{\leftrightarrow}{j\JoinOp{}x=x}{%
    \FormulaRefAuto{JoinOrderEvaluationAxiom}[ax]{1,2,3}}
\end{tabproofwide}

\FormulaDefDeltaK[Injektive Form von höchstens der Hälfte]{%
  B\subseteq A\vdash
  \AtMostHalf{B}{A}\coloneqq
  \exists F\colon B\inj A\setminus B
}{AtMostHalf}{%
  \DeltaRow{Mengen}{A\dsep B}
  \DeltaRow{Funktionen}{F}
  \DeltaRow{\textbf{Neues Symbol}}{}
  \DeltaPrem{Höchstens die Hälfte}{\AtMostHalf{B}{A}}
}

\FormulaDefDeltaK[Seltenes Element einer endlichen Familie]{%
  \RareOcc{a}{M}\coloneqq
  a\in\bigcup M\land
  \exists F\colon \FamContaining{M}{a}\inj\FamAvoiding{M}{a}
}{RareOccurrence}{%
  \DeltaRow{Mengen}{M\dsep a}
  \DeltaRow{Funktionen}{F}
  \DeltaRow{Teilmengenfamilien}{%
    \FamContaining{M}{a}\dsep\FamAvoiding{M}{a}}
  \DeltaRow{\textbf{Neues Symbol}}{}
  \DeltaPrem{Seltene Elemente}{\RareOcc{a}{M}}
}

\FormulaDefDeltaK[Frankl-Eigenschaft]{%
  \FranklSemi{\JoinOp,A,0}\coloneqq
  \exists j\Bigl(
    \JoinIrr{\JoinOp,A,0}{j}\land
    \AtMostHalf{\PrincipalFilter{\JoinOrd{\JoinOp}{A},A}{j}}{A}
  \Bigr)
}{FranklSemilatticeProperty}{%
  \DeltaRow{Mengen}{A\dsep 0\dsep j}
  \DeltaRow{Funktionen}{\JoinOp\dsep F}
  \DeltaPrem{Halbverband mit Nullelement}{\ZeroJoinSemi{\JoinOp,A,0}}
  \DeltaRow{Hauptfilter}{\PrincipalFilter{\JoinOrd{\JoinOp}{A},A}{j}}
  \DeltaRow{\textbf{Neues Symbol}}{}
  \DeltaPrem{Frankl-Halbverbände}{\FranklSemi{\JoinOp,A,0}}
}

Die rein halbverbandstheoretische Frankl-Vermutung lautet damit:

\begin{samepage}
\begin{center}
\fbox{\begin{minipage}{0.88\linewidth}
\textbf{Frankl-Vermutung für Supremums-Halbverbände.}
Ist \((A,\JoinOp,0)\) ein endlicher Supremums-Halbverband mit Nullelement und besitzt
\(A\) ein von \(0\) verschiedenes Element, dann gilt
\[
  \FranklSemi{\JoinOp,A,0}.
\]
Äquivalent existieren ein \(\JoinOp\)-irreduzibles \(j\) und eine Injektion
\[
  \PrincipalFilter{\JoinOrd{\JoinOp}{A},A}{j}
  \longrightarrow
  A\setminus\PrincipalFilter{\JoinOrd{\JoinOp}{A},A}{j}.
\]
\end{minipage}}
\end{center}
\end{samepage}

Diese Aussage wird bewusst weder als Axiom noch als Theorem registriert. Sie
ist mit Stand vom 18.~Juli~2026 weiterhin offen. Bewiesen werden im Folgenden
ihre exakte Äquivalenz zur Mengenform aus Band~05 und einige Spezialfälle.

\begin{remark}
Das Nullelement ist wesentlich. Betrachte den dreielementigen Supremums-Halbverband
\[
  A=\{a,b,1\},\qquad a\JoinOp{}b=1.
\]
Hier sind \(a\) und \(b\) die einzigen \(\JoinOp\)-irreduziblen Elemente. Beide
Hauptfilter haben zwei Elemente, ihre
Komplemente aber nur eines. Es kann daher keine Injektion vom Hauptfilter in
sein Komplement geben. Nach Adjunktion eines Nullelements entsteht das
viergliedrige boolesche Gitter; dort haben die entsprechenden Hauptfilter
genau zwei Elemente.
\end{remark}

\section{Zwei vollständig bewiesene Spezialfälle}

Zunächst isolieren wir eine elementare Abbildung, die auch außerhalb der
Frankl-Frage nützlich ist.

\FormulaThmDeltaK[Eine Einermenge injiziert in jede nichtleere Menge]{%
  b\in B\vdash \exists F\colon\{a\}\inj B
}{SingletonInjectionIntoNonempty}{%
  \DeltaRow{Mengen}{A\dsep B\dsep a\dsep b\dsep x\dsep y}
  \DeltaRow{Funktionen}{F\dsep\ConstFun{\{a\}}{B}{b}}
}
\begin{tabproof}
  \proofstep{1}{b\in B}{\rA}
  \proofstep{1}{\ConstFun{\{a\}}{B}{b}\colon\{a\}\to B}{%
    \FormulaRefAuto{b\in B\vdash \ConstFun{A}{B}{b}\colon A\to B}{1}}
  \proofstep{3}{x\in\{a\}}{\rA}
  \proofstep{4}{y\in\{a\}}{\rA}
  \proofstep{5}{\ConstFun{\{a\}}{B}{b}(x)=
    \ConstFun{\{a\}}{B}{b}(y)}{\rA}
  \proofstep{3}{x=a}{\FormulaRefAuto{x\in\{a\}\eqvdash x=a}{3}}
  \proofstep{4}{y=a}{\FormulaRefAuto{x\in\{a\}\eqvdash x=a}{4}}
  \proofstep{3,4}{x=y}{%
    \FormulaRefAuto{a=b,\,c=b\vdash a=c}{6,7}}
  \proofstep{1}{\ConstFun{\{a\}}{B}{b}\colon\{a\}\inj B}{%
    \FormulaRefAuto{Injektive Funktion}[def]{2,3,4,5,8}}
  \proofstep{1}{\exists F\colon\{a\}\inj B}{\rEI{9}}
\end{tabproof}

\FormulaThmDeltaK[\(\JoinOp\)-irreduzibles größtes Element erfüllt Frankl]{%
  \begin{aligned}[t]
    &\ZeroJoinSemi{\JoinOp,A,0}\dsep \JoinIrr{\JoinOp,A,0}{t}\dsep
      \forall x\in A\,(x,t)\in\JoinOrd{\JoinOp}{A}\vdash{}\\
    &\FranklSemi{\JoinOp,A,0}
  \end{aligned}
}{JoinIrreducibleTopFrankl}{%
  \DeltaRow{Mengen}{A\dsep 0\dsep t\dsep x}
  \DeltaRow{Funktionen}{\JoinOp\dsep F}
}
\begin{tabproofsplitwide}
  \proofpartwideindR[Der Hauptfilter ist eine Einermenge]{%
    \begin{aligned}[t]
      &\ZeroJoinSemi{\JoinOp,A,0}\dsep \JoinIrr{\JoinOp,A,0}{t}\dsep
        \forall x\in A\,(x,t)\in\JoinOrd{\JoinOp}{A}\vdash{}\\[-2pt]
      &\PrincipalFilter{\JoinOrd{\JoinOp}{A},A}{t}=\{t\}
    \end{aligned}
  }{\ZeroJoinSemi{\JoinOp,A,0}\dsep \JoinIrr{\JoinOp,A,0}{t}\dsep
    \forall x\in A\,(x,t)\in\JoinOrd{\JoinOp}{A}\vdash
    \PrincipalFilter{\JoinOrd{\JoinOp}{A},A}{t}=\{t\}}%
  [JoinIrreducibleTopFilterSingleton]
    \proofstepwidestar[1]{\ZeroJoinSemi{\JoinOp,A,0}}{\rA}
    \proofstepwidestar[2]{\JoinIrr{\JoinOp,A,0}{t}}{\rA}
    \proofstepwidestar[3]{\forall x\in A\,(x,t)\in\JoinOrd{\JoinOp}{A}}{\rA}
    \proofstepwidestar[2]{t\in A}{\FormulaRefAuto{JoinIrreducible}[def]{2}}
    \proofstepwidestar[1]{\Semi{\JoinOp,A}}{%
      \FormulaRefAuto{ZeroJoinBaseSemilatticeAxiom}[ax]{1}}
    \proofstepwidestar[4,5]{(t,t)\in\JoinOrd{\JoinOp}{A}}{%
      \FormulaRefAuto{JoinOrderReflexive}[thm]{5,4}}
    \proofstepwidestar[4,6]{t\in\PrincipalFilter{\JoinOrd{\JoinOp}{A},A}{t}}{%
      \FormulaRefAuto{PrincipalFilter}[def]{4,6}}
    \proofstepwidestar[8]{x\in\PrincipalFilter{\JoinOrd{\JoinOp}{A},A}{t}}{\rA}
    \proofstepwidestar[8]{x\in A\land(t,x)\in\JoinOrd{\JoinOp}{A}}{%
      \FormulaRefAuto{PrincipalFilter}[def]{8}}
    \proofstepwidestar[8]{x\in A}{\rAEa{9}}
    \proofstepwidestar[8]{(t,x)\in\JoinOrd{\JoinOp}{A}}{\rAEb{9}}
    \proofstepwidestar[3,8]{(x,t)\in\JoinOrd{\JoinOp}{A}}{\rRE{10,\rUE{3}}}
    \proofstepwidestar[1,3,8]{x=t}{%
      \FormulaRefAuto{JoinOrderAntisymmetric}[thm]{5,12,11}}
    \proofstepwidestar[1,3,8]{x\in\{t\}}{%
      \FormulaRefAuto{x\in\{a\}\eqvdash x=a}{13}}
    \proofstepwidestar[1,2,3]{%
      \PrincipalFilter{\JoinOrd{\JoinOp}{A},A}{t}\subseteq\{t\}}{%
      \FormulaRefAuto{A\subseteq B \coloneq \forall x\,(x\in A\rightarrow x\in B)}{%
        \rUI{\rRI{8,14}}}}
    \proofstepwidestar[1,2,3]{%
      \{t\}\subseteq\PrincipalFilter{\JoinOrd{\JoinOp}{A},A}{t}}{%
      \FormulaRefAuto{a\in A\vdash\{a\}\subseteq A}{7}}
    \proofstepwidestar[1,2,3]{%
      \PrincipalFilter{\JoinOrd{\JoinOp}{A},A}{t}=\{t\}}{%
      \FormulaRefAuto{A \subseteq B \land B \subseteq A \eqvdash A = B}{15,16}}
  \closeproofpartwideindR

  \proofpartwideindR[Injektion in das Komplement]{%
    \begin{aligned}[t]
      &\ZeroJoinSemi{\JoinOp,A,0}\dsep \JoinIrr{\JoinOp,A,0}{t}\dsep
        \forall x\in A\,(x,t)\in\JoinOrd{\JoinOp}{A}\vdash{}\\[-2pt]
      &\AtMostHalf{\PrincipalFilter{\JoinOrd{\JoinOp}{A},A}{t}}{A}
    \end{aligned}
  }{\ZeroJoinSemi{\JoinOp,A,0}\dsep \JoinIrr{\JoinOp,A,0}{t}\dsep
    \forall x\in A\,(x,t)\in\JoinOrd{\JoinOp}{A}\vdash
    \AtMostHalf{\PrincipalFilter{\JoinOrd{\JoinOp}{A},A}{t}}{A}}%
  [JoinIrreducibleTopAtMostHalf]
    \proofstepwidestar[1]{\ZeroJoinSemi{\JoinOp,A,0}}{\rA}
    \proofstepwidestar[2]{\JoinIrr{\JoinOp,A,0}{t}}{\rA}
    \proofstepwidestar[3]{\forall x\in A\,(x,t)\in\JoinOrd{\JoinOp}{A}}{\rA}
    \proofstepwidestar[1]{0\in A}{\FormulaRefAuto{ZeroJoinCarrierAxiom}[ax]{1}}
    \proofstepwidestar[2]{t\neq0}{\FormulaRefAuto{JoinIrreducible}[def]{2}}
    \proofstepwidestar[1,2]{0\in A\setminus\{t\}}{%
      \FormulaRefAuto{x\in A\setminus B\eqvdash x\in A\land x\notin B}{%
        \rAI{4,\FormulaRefAuto{x\notin\{a\}\eqvdash x\neq a}{%
          \FormulaRefAuto{a=b\vdash b=a}{5}}}}}
    \proofstepwidestar[1,2]{\exists F\colon\{t\}\inj A\setminus\{t\}}{%
      \FormulaRefAuto{SingletonInjectionIntoNonempty}[thm]{6}}
    \proofstepwidestar[1,2,3]{%
      \PrincipalFilter{\JoinOrd{\JoinOp}{A},A}{t}=\{t\}}{%
      \FormulaRefAuto{JoinIrreducibleTopFilterSingleton}[thm]{1,2,3}}
    \proofstepwidestar[]{\{t\}\subseteq A}{%
      \FormulaRefAuto{a\in A\vdash\{a\}\subseteq A}{%
        \FormulaRefAuto{JoinIrreducible}[def]{2}}}
    \proofstepwidestar[1,2,3]{%
      \AtMostHalf{\PrincipalFilter{\JoinOrd{\JoinOp}{A},A}{t}}{A}}{%
      \FormulaRefAuto{AtMostHalf}[def]{\rIE{8,9},\rIE{8,7}}}
  \closeproofpartwideindR

  \proofpartwide{Schluss}
    \proofstepwidestar[1]{\ZeroJoinSemi{\JoinOp,A,0}}{\rA}
    \proofstepwidestar[2]{\JoinIrr{\JoinOp,A,0}{t}}{\rA}
    \proofstepwidestar[3]{\forall x\in A\,(x,t)\in\JoinOrd{\JoinOp}{A}}{\rA}
    \proofstepwidestar[1,2,3]{%
      \AtMostHalf{\PrincipalFilter{\JoinOrd{\JoinOp}{A},A}{t}}{A}}{%
      \FormulaRefAuto{JoinIrreducibleTopAtMostHalf}[thm]{1,2,3}}
    \proofstepwidestar[1,2,3]{\FranklSemi{\JoinOp,A,0}}{%
      \FormulaRefAuto{FranklSemilatticeProperty}[def]{%
        \rEI{\rAI{2,4}}}}
  \closeproofpartwide
\end{tabproofsplitwide}

\FormulaThmDeltaK[Endliche Ketten erfüllen die Halbverbandsform]{%
  \begin{aligned}[t]
    &\ZeroJoinSemi{\JoinOp,A,0}\dsep \Finite A\dsep
      \TotOrd{\JoinOrd{\JoinOp}{A},A}\dsep
      t\in A\dsep{}\\
    &\forall x\in A\,(x,t)\in\JoinOrd{\JoinOp}{A}\dsep
      t\neq0\vdash{}\\
    &\FranklSemi{\JoinOp,A,0}
  \end{aligned}
}{FranklJoinChains}{%
  \DeltaRow{Mengen}{A\dsep 0\dsep t\dsep x\dsep y}
  \DeltaRow{Funktionen}{\JoinOp}
}
\begin{tabproofsplitwide}
  \proofpartwideindR[Das größte Kettenelement ist irreduzibel]{%
    \begin{aligned}[t]
      &\ZeroJoinSemi{\JoinOp,A,0}\dsep
        \TotOrd{\JoinOrd{\JoinOp}{A},A}\dsep t\in A\dsep{}\\[-2pt]
      &\forall x\in A\,(x,t)\in\JoinOrd{\JoinOp}{A}\dsep t\neq0
        \vdash \JoinIrr{\JoinOp,A,0}{t}
    \end{aligned}
  }{\ZeroJoinSemi{\JoinOp,A,0}\dsep
    \TotOrd{\JoinOrd{\JoinOp}{A},A}\dsep t\in A\dsep
    \forall x\in A\,(x,t)\in\JoinOrd{\JoinOp}{A}\dsep t\neq0
    \vdash \JoinIrr{\JoinOp,A,0}{t}}[FranklJoinChainTopIrreducible]
    \proofstepwidestar[1]{\ZeroJoinSemi{\JoinOp,A,0}}{\rA}
    \proofstepwidestar[2]{\TotOrd{\JoinOrd{\JoinOp}{A},A}}{\rA}
    \proofstepwidestar[3]{t\in A}{\rA}
    \proofstepwidestar[4]{\forall x\in A\,(x,t)\in\JoinOrd{\JoinOp}{A}}{\rA}
    \proofstepwidestar[5]{t\neq0}{\rA}
    \proofstepwidestar[1]{\Semi{\JoinOp,A}}{%
      \FormulaRefAuto{ZeroJoinBaseSemilatticeAxiom}[ax]{1}}
    \proofstepwidestar[7]{x\in A}{\rA}
    \proofstepwidestar[8]{y\in A}{\rA}
    \proofstepwidestar[9]{x\JoinOp{}y=t}{\rA}
    \proofstepwidestar[2,7,8]{%
      (x,y)\in\JoinOrd{\JoinOp}{A}\lor(y,x)\in\JoinOrd{\JoinOp}{A}}{%
      \FormulaRefAuto{x\in A \dsep y\in A \vdash
        (x,y)\in R \lor (y,x)\in R}[ax]{2,7,8}}
    \proofstepwidestar[11]{(x,y)\in\JoinOrd{\JoinOp}{A}}{\rA}
    \proofstepwidestar[1,7,8,11]{x\JoinOp{}y=y}{%
      \FormulaRefAuto{JoinOrderEvaluation}[thm]{6,11}}
    \proofstepwidestar[7,8,9,11]{y=t}{%
      \FormulaRefAuto{a=b,\,a=c\vdash b=c}{12,9}}
    \proofstepwidestar[7,8,9,11]{x=t\lor y=t}{\rOIb{13}}
    \proofstepwidestar[15]{(y,x)\in\JoinOrd{\JoinOp}{A}}{\rA}
    \proofstepwidestar[1,7,8,15]{y\JoinOp{}x=x}{%
      \FormulaRefAuto{JoinOrderEvaluation}[thm]{6,15}}
    \proofstepwidestar[1,7,8]{y\JoinOp{}x=x\JoinOp{}y}{%
      \FormulaRefAuto{SemilatticeCommutativityAxiom}[ax]{6,8,7}}
    \proofstepwidestar[1,7,8,15]{x\JoinOp{}y=x}{%
      \rIE{17,16}}
    \proofstepwidestar[7,8,9,15]{x=t}{%
      \FormulaRefAuto{a=b,\,a=c\vdash b=c}{18,9}}
    \proofstepwidestar[7,8,9,15]{x=t\lor y=t}{\rOIa{19}}
    \proofstepwidestar[2,7,8,9]{x=t\lor y=t}{%
      \rOE{10,11,14,15,20}}
    \proofstepwidestar[1,2,3,4,5]{\JoinIrr{\JoinOp,A,0}{t}}{%
      \FormulaRefAuto{JoinIrreducible}[def]{3,5,7,8,9,21}}
  \closeproofpartwideindR

  \proofpartwide{Schluss}
    \proofstepwidestar[1]{\ZeroJoinSemi{\JoinOp,A,0}}{\rA}
    \proofstepwidestar[2]{\Finite A}{\rA}
    \proofstepwidestar[3]{\TotOrd{\JoinOrd{\JoinOp}{A},A}}{\rA}
    \proofstepwidestar[4]{t\in A}{\rA}
    \proofstepwidestar[5]{\forall x\in A\,(x,t)\in\JoinOrd{\JoinOp}{A}}{\rA}
    \proofstepwidestar[6]{t\neq0}{\rA}
    \proofstepwidestar[1,3,4,5,6]{\JoinIrr{\JoinOp,A,0}{t}}{%
      \FormulaRefAuto{FranklJoinChainTopIrreducible}[thm]{1,3,4,5,6}}
    \proofstepwidestar[1,3,4,5,6]{\FranklSemi{\JoinOp,A,0}}{%
      \FormulaRefAuto{JoinIrreducibleTopFrankl}[thm]{1,7,5}}
  \closeproofpartwide
\end{tabproofsplitwide}

Der Endlichkeitsparameter wird im letzten Beweis nur benötigt, um die
Kettenaussage als Spezialfall der endlichen Frankl-Vermutung einzuordnen. Ist
ein größtes Element bereits gegeben, funktioniert die Injektion unabhängig
von der Endlichkeit.

\section{Endliche Ordnungslemmata}

Für die Übersetzung zwischen Mengenfamilien und Halbverbänden werden zwei
elementare Endlichkeitslemmata benötigt. Die im Induktionsbeweis verwendete
Eigenschaft wird zunächst als eigenes Prädikat festgelegt.

\FormulaDefDeltaK[Induktionseigenschaft für maximale Elemente]{%
  \operatorname{HatMax}_{R,A}(C)
  \coloneqq
  \Bigl(
    C\subseteq A\land C\neq\varnothing
    \to
    \exists m\in C\,\forall b\in C\,
      \bigl((m,b)\in R\to b=m\bigr)
  \Bigr)
}{FinitePosetMaximalProperty}{%
  \DeltaRow{Mengen}{A\dsep C\dsep b\dsep m}
  \DeltaRow{Relationen}{R}
  \DeltaRow{\textbf{Neues Symbol}}{}
  \DeltaPrem{Maximalelementeigenschaft}{\operatorname{HatMax}_{R,A}(C)}
}

\FormulaThmDeltaK[Maximales Element einer endlichen Teilordnung]{%
  \begin{aligned}[t]
    &\PartOrd{R,A}\dsep \Finite B\dsep
      B\neq\varnothing\dsep B\subseteq A\vdash{}\\
    &\exists m\in B\,\forall b\in B\,
      \bigl((m,b)\in R\to b=m\bigr)
  \end{aligned}
}{FinitePosetMaximalElement}{%
  \DeltaRow{Mengen}{A\dsep B\dsep C\dsep b\dsep c\dsep m}
  \DeltaRow{Relationen}{R}
}

\begin{tabproofsplitwide}
  \proofpartwideindR[Einermengen besitzen ein maximales Element]{%
    c\in A\vdash
    \exists q\in\{c\}\,\forall b\in\{c\}\,
      \bigl((q,b)\in R\to b=q\bigr)
  }{%
    c\in A\vdash
    \exists q\in\{c\}\,\forall b\in\{c\}\,
      \bigl((q,b)\in R\to b=q\bigr)
  }[SingletonPosetMaximalElement]
    \proofstepwidestar[1]{c\in A}{\rA}
    \proofstepwidestar[2]{b\in\{c\}}{\rA}
    \proofstepwidestar[3]{(c,b)\in R}{\rA}
    \proofstepwidestar[2]{b=c}{%
      \FormulaRefAuto{x\in\{a\}\eqvdash x=a}{2}}
    \proofstepwidestar[2]{(c,b)\in R\to b=c}{\rRI{3,4}}
    \proofstepwidestar[]{%
      \forall b\in\{c\}\,\bigl((c,b)\in R\to b=c\bigr)}{%
      \rUI{\rRI{2,5}}}
    \proofstepwidestar[]{c\in\{c\}}{\FormulaRefAuto{a\in\{a\}}}
    \proofstepwidestar[1]{%
      \exists q\in\{c\}\,\forall b\in\{c\}\,
        \bigl((q,b)\in R\to b=q\bigr)}{%
      \rEI{\rAI{7,6}}}
  \closeproofpartwideindR

  \proofpartwideindR[Der neue Kandidat ist maximal]{%
    \begin{aligned}[t]
      &\PartOrd{R,A}\dsep C\subseteq A\dsep c\in A\dsep
        m\in C\dsep{}\\
      &\forall b\in C\,\bigl((m,b)\in R\to b=m\bigr)
        \dsep (m,c)\in R\vdash{}\\
      &\exists q\in C\cup\{c\}\,\forall b\in C\cup\{c\}\,
        \bigl((q,b)\in R\to b=q\bigr)
    \end{aligned}
  }{%
    \begin{aligned}[t]
      &\PartOrd{R,A}\dsep C\subseteq A\dsep c\in A\dsep
        m\in C\dsep{}\\
      &\forall b\in C\,\bigl((m,b)\in R\to b=m\bigr)
        \dsep (m,c)\in R\vdash{}\\
      &\exists q\in C\cup\{c\}\,\forall b\in C\cup\{c\}\,
        \bigl((q,b)\in R\to b=q\bigr)
    \end{aligned}
  }[FinitePosetMaximalAdjoinAbove]
    \proofstepwidestar[1]{\PartOrd{R,A}}{\rA}
    \proofstepwidestar[2]{C\subseteq A}{\rA}
    \proofstepwidestar[3]{c\in A}{\rA}
    \proofstepwidestar[4]{m\in C}{\rA}
    \proofstepwidestar[5]{%
      \forall b\in C\,\bigl((m,b)\in R\to b=m\bigr)}{\rA}
    \proofstepwidestar[6]{(m,c)\in R}{\rA}
    \proofstepwidestar[7]{b\in C\cup\{c\}}{\rA}
    \proofstepwidestar[8]{(c,b)\in R}{\rA}
    \proofstepwidestar[7]{b\in C\lor b=c}{%
      \FormulaRefAuto{x\in A\cup\{a\} \vdash x\in A\lor x=a}{7}}
    \proofstepwidestar[10]{b\in C}{\rA}
    \proofstepwidestar[6,8]{(m,b)\in R}{%
      \FormulaRefAuto{(x,y)\in R \dsep (y,z)\in R
        \vdash (x,z)\in R}{6,8}}
    \proofstepwidestar[5,6,8,10]{b=m}{%
      \rRE{\rRE{10,\rUE{5}},11}}
    \proofstepwidestar[5,6,8,10]{m=b}{%
      \FormulaRefAuto{a=b\vdash b=a}{12}}
    \proofstepwidestar[5,6,8,10]{(b,c)\in R}{\rIE{13,6}}
    \proofstepwidestar[5,6,8,10]{b=c}{%
      \FormulaRefAuto{(x,y)\in R \dsep (y,x)\in R
        \vdash x=y}{14,8}}
    \proofstepwidestar[16]{b=c}{\rA}
    \proofstepwidestar[1,5,6,7,8]{b=c}{%
      \rOE{9,10,15,16,16}}
    \proofstepwidestar[1,5,6,7]{(c,b)\in R\to b=c}{\rRI{8,17}}
    \proofstepwidestar[1,5,6]{%
      \forall b\in C\cup\{c\}\,\bigl((c,b)\in R\to b=c\bigr)}{%
      \rUI{\rRI{7,18}}}
    \proofstepwidestar[]{c\in C\cup\{c\}}{\FormulaRefAuto{a\in A\cup\{a\}}}
    \proofstepwidestar[1,2,3,4,5,6]{%
      \exists q\in C\cup\{c\}\,\forall b\in C\cup\{c\}\,
        \bigl((q,b)\in R\to b=q\bigr)}{%
      \rEI{\rAI{20,19}}}
  \closeproofpartwideindR

  \proofpartwideindR[Der bisherige Kandidat bleibt maximal]{%
    \begin{aligned}[t]
      &\PartOrd{R,A}\dsep C\subseteq A\dsep c\in A\dsep
        m\in C\dsep{}\\
      &\forall b\in C\,\bigl((m,b)\in R\to b=m\bigr)
        \dsep (m,c)\notin R\vdash{}\\
      &\exists q\in C\cup\{c\}\,\forall b\in C\cup\{c\}\,
        \bigl((q,b)\in R\to b=q\bigr)
    \end{aligned}
  }{%
    \begin{aligned}[t]
      &\PartOrd{R,A}\dsep C\subseteq A\dsep c\in A\dsep
        m\in C\dsep{}\\
      &\forall b\in C\,\bigl((m,b)\in R\to b=m\bigr)
        \dsep (m,c)\notin R\vdash{}\\
      &\exists q\in C\cup\{c\}\,\forall b\in C\cup\{c\}\,
        \bigl((q,b)\in R\to b=q\bigr)
    \end{aligned}
  }[FinitePosetMaximalAdjoinNotAbove]
    \proofstepwidestar[1]{\PartOrd{R,A}}{\rA}
    \proofstepwidestar[2]{C\subseteq A}{\rA}
    \proofstepwidestar[3]{c\in A}{\rA}
    \proofstepwidestar[4]{m\in C}{\rA}
    \proofstepwidestar[5]{%
      \forall b\in C\,\bigl((m,b)\in R\to b=m\bigr)}{\rA}
    \proofstepwidestar[6]{(m,c)\notin R}{\rA}
    \proofstepwidestar[7]{b\in C\cup\{c\}}{\rA}
    \proofstepwidestar[8]{(m,b)\in R}{\rA}
    \proofstepwidestar[7]{b\in C\lor b=c}{%
      \FormulaRefAuto{x\in A\cup\{a\} \vdash x\in A\lor x=a}{7}}
    \proofstepwidestar[10]{b\in C}{\rA}
    \proofstepwidestar[5,8,10]{b=m}{\rRE{\rRE{10,\rUE{5}},8}}
    \proofstepwidestar[12]{b=c}{\rA}
    \proofstepwidestar[8,12]{(m,c)\in R}{\rIE{12,8}}
    \proofstepwidestar[6,8,12]{\bot}{\rBI{6,13}}
    \proofstepwidestar[6,8,12]{b=m}{\FormulaRefAuto{\bot\vdash P}{14}}
    \proofstepwidestar[5,6,7,8]{b=m}{%
      \rOE{9,10,11,12,15}}
    \proofstepwidestar[5,6,7]{(m,b)\in R\to b=m}{\rRI{8,16}}
    \proofstepwidestar[5,6]{%
      \forall b\in C\cup\{c\}\,\bigl((m,b)\in R\to b=m\bigr)}{%
      \rUI{\rRI{7,17}}}
    \proofstepwidestar[4]{m\in C\cup\{c\}}{%
      \FormulaRefAuto{z\in A\vdash z\in A\cup B}[thm]{4}}
    \proofstepwidestar[1,2,3,4,5,6]{%
      \exists q\in C\cup\{c\}\,\forall b\in C\cup\{c\}\,
        \bigl((q,b)\in R\to b=q\bigr)}{%
      \rEI{\rAI{19,18}}}
  \closeproofpartwideindR

  \proofpartwideindR[Induktionsanfang]{%
    \operatorname{HatMax}_{R,A}(\varnothing)
  }{%
    \operatorname{HatMax}_{R,A}(\varnothing)
  }[FinitePosetMaximalElementBase]
    \proofstepwidestar[1]{%
      \varnothing\subseteq A\land\varnothing\neq\varnothing}{\rA}
    \proofstepwidestar[]{\varnothing=\varnothing}{\rII}
    \proofstepwidestar[1]{\bot}{\rCE{2,\rAEb{1}}}
    \proofstepwidestar[1]{%
      \exists m\in\varnothing\,\forall b\in\varnothing\,
        \bigl((m,b)\in R\to b=m\bigr)}{%
      \FormulaRefAuto{\bot\vdash P}{3}}
    \proofstepwidestar[]{%
      \begin{aligned}[t]
        &\varnothing\subseteq A\land\varnothing\neq\varnothing\to{}\\
        &\exists m\in\varnothing\,\forall b\in\varnothing\,
          \bigl((m,b)\in R\to b=m\bigr)
      \end{aligned}}{\rRI{1,4}}
    \proofstepwidestar[]{\operatorname{HatMax}_{R,A}(\varnothing)}{%
      \FormulaRefAuto{FinitePosetMaximalProperty}[def]{5}}
  \closeproofpartwideindR

  \proofpartwideindR[Induktionsschritt]{%
    \PartOrd{R,A}\dsep\Finite C\dsep c\notin C\dsep
    \operatorname{HatMax}_{R,A}(C)
    \vdash \operatorname{HatMax}_{R,A}(C\cup\{c\})
  }{%
    \PartOrd{R,A}\dsep\Finite C\dsep c\notin C\dsep
    \operatorname{HatMax}_{R,A}(C)
    \vdash \operatorname{HatMax}_{R,A}(C\cup\{c\})
  }[FinitePosetMaximalElementStep]
    \proofstepwidestar[1]{\PartOrd{R,A}}{\rA}
    \proofstepwidestar[2]{\Finite C}{\rA}
    \proofstepwidestar[3]{c\notin C}{\rA}
    \proofstepwidestar[4]{\operatorname{HatMax}_{R,A}(C)}{\rA}
    \proofstepwidestar[5]{%
      C\cup\{c\}\subseteq A\land C\cup\{c\}\neq\varnothing}{\rA}
    \proofstepwidestar[5]{C\cup\{c\}\subseteq A}{\rAEa{5}}
    \proofstepwidestar[]{C\subseteq C\cup\{c\}}{%
      \FormulaRefAuto{A\subseteq A\cup B}}
    \proofstepwidestar[5]{C\subseteq A}{%
      \FormulaRefAuto{A\subseteq B\dsep B\subseteq C\vdash A\subseteq C}{7,6}}
    \proofstepwidestar[]{c\in C\cup\{c\}}{\FormulaRefAuto{a\in A\cup\{a\}}}
    \proofstepwidestar[5]{c\in A}{%
      \FormulaRefAuto{A\subseteq B,\, x\in A \vdash x\in B}{6,9}}
    \proofstepwidestar[]{C=\varnothing\lor C\neq\varnothing}{%
      \FormulaRefAuto{P\lor\neg P}}
    \proofstepwidestar[12]{C=\varnothing}{\rA}
    \proofstepwidestar[12]{C\cup\{c\}=\{c\}}{%
      \rIE{\FormulaRefAuto{a=b\vdash b=a}{12},
        \FormulaRefAuto{\varnothing \cup A = A}}}
    \proofstepwidestar[5]{%
      \exists q\in\{c\}\,\forall b\in\{c\}\,
        \bigl((q,b)\in R\to b=q\bigr)}{%
      \FormulaRefAuto{SingletonPosetMaximalElement}[thm]{10}}
    \proofstepwidestar[5,12]{%
      \exists q\in C\cup\{c\}\,\forall b\in C\cup\{c\}\,
        \bigl((q,b)\in R\to b=q\bigr)}{%
      \rIE{\FormulaRefAuto{a=b\vdash b=a}{13},14}}
    \proofstepwidestar[16]{C\neq\varnothing}{\rA}
    \proofstepwidestar[4]{%
      \begin{aligned}[t]
        &C\subseteq A\land C\neq\varnothing\to{}\\
        &\exists m\in C\,\forall b\in C\,
          \bigl((m,b)\in R\to b=m\bigr)
      \end{aligned}}{%
      \FormulaRefAuto{FinitePosetMaximalProperty}[def]{4}}
    \proofstepwidestar[4,5,16]{%
      \exists m\in C\,\forall b\in C\,
        \bigl((m,b)\in R\to b=m\bigr)}{%
      \rRE{17,\rAI{8,16}}}
    \proofstepwidestar[19]{%
      m\in C\land
      \forall b\in C\,\bigl((m,b)\in R\to b=m\bigr)}{\rA}
    \proofstepwidestar[]{(m,c)\in R\lor(m,c)\notin R}{%
      \FormulaRefAuto{P\lor\neg P}}
    \proofstepwidestar[21]{(m,c)\in R}{\rA}
    \proofstepwidestar[1,5,19,21]{%
      \exists q\in C\cup\{c\}\,\forall b\in C\cup\{c\}\,
        \bigl((q,b)\in R\to b=q\bigr)}{%
      \FormulaRefAuto{FinitePosetMaximalAdjoinAbove}[thm]{%
        1,8,10,\rAEa{19},\rAEb{19},21}}
    \proofstepwidestar[23]{(m,c)\notin R}{\rA}
    \proofstepwidestar[1,5,19,23]{%
      \exists q\in C\cup\{c\}\,\forall b\in C\cup\{c\}\,
        \bigl((q,b)\in R\to b=q\bigr)}{%
      \FormulaRefAuto{FinitePosetMaximalAdjoinNotAbove}[thm]{%
        1,8,10,\rAEa{19},\rAEb{19},23}}
    \proofstepwidestar[1,5,19]{%
      \exists q\in C\cup\{c\}\,\forall b\in C\cup\{c\}\,
        \bigl((q,b)\in R\to b=q\bigr)}{%
      \rOE{20,21,22,23,24}}
    \proofstepwidestar[1,4,5,16]{%
      \exists q\in C\cup\{c\}\,\forall b\in C\cup\{c\}\,
        \bigl((q,b)\in R\to b=q\bigr)}{%
      \rEE{18,19,25}}
    \proofstepwidestar[1,4,5]{%
      \exists q\in C\cup\{c\}\,\forall b\in C\cup\{c\}\,
        \bigl((q,b)\in R\to b=q\bigr)}{%
      \rOE{11,12,15,16,26}}
    \proofstepwidestar[1,4]{%
      \begin{aligned}[t]
        &C\cup\{c\}\subseteq A\land C\cup\{c\}\neq\varnothing\to{}\\
        &\exists q\in C\cup\{c\}\,\forall b\in C\cup\{c\}\,
          \bigl((q,b)\in R\to b=q\bigr)
      \end{aligned}}{\rRI{5,27}}
    \proofstepwidestar[1,4]{\operatorname{HatMax}_{R,A}(C\cup\{c\})}{%
      \FormulaRefAuto{FinitePosetMaximalProperty}[def]{28}}
  \closeproofpartwideindR

  \proofpartwide{Induktionsschluss}
    \proofstepwidestar[1]{\PartOrd{R,A}}{\rA}
    \proofstepwidestar[2]{\Finite B}{\rA}
    \proofstepwidestar[3]{B\neq\varnothing}{\rA}
    \proofstepwidestar[4]{B\subseteq A}{\rA}
    \proofstepwidestar[1,2]{\operatorname{HatMax}_{R,A}(B)}{%
      \FormulaRefAuto{FiniteInductionRule}[ax]{%
        \FormulaRefAuto{FinitePosetMaximalElementBase}[thm],
        \FormulaRefAuto{FinitePosetMaximalElementStep}[thm]{1},2}}
    \proofstepwidestar[1,2]{%
      \begin{aligned}[t]
        &B\subseteq A\land B\neq\varnothing\to{}\\
        &\exists m\in B\,\forall b\in B\,
          \bigl((m,b)\in R\to b=m\bigr)
      \end{aligned}}{%
      \FormulaRefAuto{FinitePosetMaximalProperty}[def]{5}}
    \proofstepwidestar[1,2,3,4]{%
      \exists m\in B\,\forall b\in B\,((m,b)\in R\to b=m)}{%
      \rRE{6,\rAI{4,3}}}
  \closeproofpartwide
\end{tabproofsplitwide}

\FormulaDefDeltaK[Duale Relation auf einem Träger]{%
  \begin{aligned}[t]
    &R_A^{\mathrm{op}}\subseteq A\times A,\\
    &x,y\in A\vdash
      \bigl((x,y)\in R_A^{\mathrm{op}}\leftrightarrow (y,x)\in R\bigr)
  \end{aligned}
}{OrderDualRelation}{%
  \DeltaRow{Mengen}{A\dsep x\dsep y}
  \DeltaRow{Relationen}{R\dsep R_A^{\mathrm{op}}}
  \DeltaRow{Träger}{R_A^{\mathrm{op}}\subseteq A\times A}
  \DeltaRow{Auswertung}{%
    x,y\in A\vdash
    \bigl((x,y)\in R_A^{\mathrm{op}}\leftrightarrow (y,x)\in R\bigr)}
  \DeltaRow{\textbf{Neues Symbol}}{}
  \DeltaPrem{Duale Relation}{R_A^{\mathrm{op}}}
}

\FormulaThmDeltaK[Die duale Relation einer partiellen Ordnung ist partiell geordnet]{%
  \PartOrd{R,A}\vdash\PartOrd{R_A^{\mathrm{op}},A}
}{OrderDualPartialOrder}{%
  \DeltaRow{Mengen}{A\dsep x\dsep y\dsep z}
  \DeltaRow{Relationen}{R\dsep R_A^{\mathrm{op}}}
}
\begin{tabproofsplitwide}
  \proofpartwideindR[Reflexivität]{%
    \PartOrd{R,A}\dsep x\in A\vdash(x,x)\in R_A^{\mathrm{op}}
  }{%
    \PartOrd{R,A}\dsep x\in A\vdash(x,x)\in R_A^{\mathrm{op}}
  }[OrderDualReflexive]
    \proofstepwidestar[1]{\PartOrd{R,A}}{\rA}
    \proofstepwidestar[2]{x\in A}{\rA}
    \proofstepwidestar[2]{(x,x)\in R}{%
      \FormulaRefAuto{x\in A \vdash (x,x)\in R}{2}}
    \proofstepwidestar[2]{%
      (x,x)\in R_A^{\mathrm{op}}\leftrightarrow(x,x)\in R}{%
      \FormulaRefAuto{OrderDualRelation}[def]{2,2}}
    \proofstepwidestar[1,2]{(x,x)\in R_A^{\mathrm{op}}}{%
      \FormulaRefAuto{P\leftrightarrow Q, Q\vdash P}{4,3}}
  \closeproofpartwideindR

  \proofpartwideindR[Transitivität]{%
    \begin{aligned}[t]
      &\PartOrd{R,A}\dsep(x,y)\in R_A^{\mathrm{op}}
        \dsep(y,z)\in R_A^{\mathrm{op}}\vdash{}\\
      &(x,z)\in R_A^{\mathrm{op}}
    \end{aligned}
  }{%
    \begin{aligned}[t]
      &\PartOrd{R,A}\dsep(x,y)\in R_A^{\mathrm{op}}
        \dsep(y,z)\in R_A^{\mathrm{op}}\vdash{}\\
      &(x,z)\in R_A^{\mathrm{op}}
    \end{aligned}
  }[OrderDualTransitive]
    \proofstepwidestar[1]{\PartOrd{R,A}}{\rA}
    \proofstepwidestar[2]{(x,y)\in R_A^{\mathrm{op}}}{\rA}
    \proofstepwidestar[3]{(y,z)\in R_A^{\mathrm{op}}}{\rA}
    \proofstepwidestar[]{R_A^{\mathrm{op}}\subseteq A\times A}{%
      \FormulaRefAuto{OrderDualRelation}[def]}
    \proofstepwidestar[2]{x\in A\land y\in A}{%
      \FormulaRefAuto{(a,b)\in A\times B \eqvdash a\in A \land b\in B}{%
        \FormulaRefAuto{A\subseteq B,\,x\in A\vdash x\in B}{4,2}}}
    \proofstepwidestar[3]{y\in A\land z\in A}{%
      \FormulaRefAuto{(a,b)\in A\times B \eqvdash a\in A \land b\in B}{%
        \FormulaRefAuto{A\subseteq B,\,x\in A\vdash x\in B}{4,3}}}
    \proofstepwidestar[2]{(y,x)\in R}{%
      \FormulaRefAuto{P\leftrightarrow Q, P\vdash Q}{%
        \FormulaRefAuto{OrderDualRelation}[def]{\rAEa{5},\rAEb{5}},2}}
    \proofstepwidestar[3]{(z,y)\in R}{%
      \FormulaRefAuto{P\leftrightarrow Q, P\vdash Q}{%
        \FormulaRefAuto{OrderDualRelation}[def]{\rAEa{6},\rAEb{6}},3}}
    \proofstepwidestar[2,3]{(z,x)\in R}{%
      \FormulaRefAuto{(x,y)\in R \dsep (y,z)\in R
        \vdash (x,z)\in R}{8,7}}
    \proofstepwidestar[2,3]{%
      (x,z)\in R_A^{\mathrm{op}}\leftrightarrow(z,x)\in R}{%
      \FormulaRefAuto{OrderDualRelation}[def]{\rAEa{5},\rAEb{6}}}
    \proofstepwidestar[1,2,3]{(x,z)\in R_A^{\mathrm{op}}}{%
      \FormulaRefAuto{P\leftrightarrow Q, Q\vdash P}{10,9}}
  \closeproofpartwideindR

  \proofpartwideindR[Antisymmetrie]{%
    \begin{aligned}[t]
      &\PartOrd{R,A}\dsep(x,y)\in R_A^{\mathrm{op}}
        \dsep(y,x)\in R_A^{\mathrm{op}}\vdash x=y
    \end{aligned}
  }{%
    \begin{aligned}[t]
      &\PartOrd{R,A}\dsep(x,y)\in R_A^{\mathrm{op}}
        \dsep(y,x)\in R_A^{\mathrm{op}}\vdash x=y
    \end{aligned}
  }[OrderDualAntisymmetric]
    \proofstepwidestar[1]{\PartOrd{R,A}}{\rA}
    \proofstepwidestar[2]{(x,y)\in R_A^{\mathrm{op}}}{\rA}
    \proofstepwidestar[3]{(y,x)\in R_A^{\mathrm{op}}}{\rA}
    \proofstepwidestar[]{R_A^{\mathrm{op}}\subseteq A\times A}{%
      \FormulaRefAuto{OrderDualRelation}[def]}
    \proofstepwidestar[2]{x\in A\land y\in A}{%
      \FormulaRefAuto{(a,b)\in A\times B \eqvdash a\in A \land b\in B}{%
        \FormulaRefAuto{A\subseteq B,\,x\in A\vdash x\in B}{4,2}}}
    \proofstepwidestar[2]{(y,x)\in R}{%
      \FormulaRefAuto{P\leftrightarrow Q, P\vdash Q}{%
        \FormulaRefAuto{OrderDualRelation}[def]{\rAEa{5},\rAEb{5}},2}}
    \proofstepwidestar[3]{(x,y)\in R}{%
      \FormulaRefAuto{P\leftrightarrow Q, P\vdash Q}{%
        \FormulaRefAuto{OrderDualRelation}[def]{\rAEb{5},\rAEa{5}},3}}
    \proofstepwidestar[2,3]{x=y}{%
      \FormulaRefAuto{(x,y)\in R \dsep (y,x)\in R
        \vdash x=y}{7,6}}
  \closeproofpartwideindR

  \proofpartwide{Schluss}
    \proofstepwidestar[1]{\PartOrd{R,A}}{\rA}
    \proofstepwidestar[]{R_A^{\mathrm{op}}\subseteq A\times A}{%
      \FormulaRefAuto{OrderDualRelation}[def]}
    \proofstepwidestar[1]{x\in A\vdash(x,x)\in R_A^{\mathrm{op}}}{%
      \FormulaRefAuto{OrderDualReflexive}[thm]{1}}
    \proofstepwidestar[1]{%
      (x,y)\in R_A^{\mathrm{op}}\dsep(y,z)\in R_A^{\mathrm{op}}
      \vdash(x,z)\in R_A^{\mathrm{op}}}{%
      \FormulaRefAuto{OrderDualTransitive}[thm]{1}}
    \proofstepwidestar[1]{%
      (x,y)\in R_A^{\mathrm{op}}\dsep(y,x)\in R_A^{\mathrm{op}}
      \vdash x=y}{%
      \FormulaRefAuto{OrderDualAntisymmetric}[thm]{1}}
    \proofstepwidestar[1]{\PartOrd{R_A^{\mathrm{op}},A}}{%
      \FormulaRefAuto{Partielle Ordnung}[def]{2,3,4,5}}
  \closeproofpartwide
\end{tabproofsplitwide}

\FormulaThmDeltaK[Minimales Element einer endlichen Teilordnung]{%
  \begin{aligned}[t]
    &\PartOrd{R,A}\dsep \Finite B\dsep
      B\neq\varnothing\dsep B\subseteq A\vdash{}\\
    &\exists m\in B\,\forall b\in B\,
      \bigl((b,m)\in R\to b=m\bigr)
  \end{aligned}
}{FinitePosetMinimalElement}{%
  \DeltaRow{Mengen}{A\dsep B\dsep b\dsep m}
  \DeltaRow{Relationen}{R\dsep R_A^{\mathrm{op}}}
}
\begin{tabproofwide}
  \proofstepwidestar[1]{\PartOrd{R,A}}{\rA}
  \proofstepwidestar[2]{\Finite B}{\rA}
  \proofstepwidestar[3]{B\neq\varnothing}{\rA}
  \proofstepwidestar[4]{B\subseteq A}{\rA}
  \proofstepwidestar[1]{\PartOrd{R_A^{\mathrm{op}},A}}{%
    \FormulaRefAuto{OrderDualPartialOrder}[thm]{1}}
  \proofstepwidestar[1,2,3,4]{%
    \exists m\in B\,\forall b\in B\,
      \bigl((m,b)\in R_A^{\mathrm{op}}\to b=m\bigr)}{%
    \FormulaRefAuto{FinitePosetMaximalElement}[thm]{5,2,3,4}}
  \proofstepwidestar[7]{%
    m\in B\land
    \forall b\in B\,\bigl((m,b)\in R_A^{\mathrm{op}}\to b=m\bigr)}{\rA}
  \proofstepwidestar[8]{b\in B}{\rA}
  \proofstepwidestar[9]{(b,m)\in R}{\rA}
  \proofstepwidestar[4,8,9]{%
    (m,b)\in R_A^{\mathrm{op}}\leftrightarrow(b,m)\in R}{%
    \FormulaRefAuto{OrderDualRelation}[def]{%
      \FormulaRefAuto{A\subseteq B,\,x\in A\vdash x\in B}{4,\rAEa{7}},
      \FormulaRefAuto{A\subseteq B,\,x\in A\vdash x\in B}{4,8}}}
  \proofstepwidestar[4,8,9]{(m,b)\in R_A^{\mathrm{op}}}{%
    \FormulaRefAuto{P\leftrightarrow Q, Q\vdash P}{10,9}}
  \proofstepwidestar[4,7,8,9]{b=m}{%
    \rRE{\rRE{8,\rUE{\rAEb{7}}},11}}
  \proofstepwidestar[4,7]{%
    \forall b\in B\,\bigl((b,m)\in R\to b=m\bigr)}{%
    \rUI{\rRI{8,\rRI{9,12}}}}
  \proofstepwidestar[4,7]{%
    \begin{aligned}[t]
      &\exists m\in B\,\forall b\in B\\
      &\qquad ((b,m)\in R\to b=m)
    \end{aligned}}{%
    \rEI{\rAI{\rAEa{7},13}}}
  \proofstepwidestar[1,2,3,4]{%
    \exists m\in B\,\forall b\in B\,\bigl((b,m)\in R\to b=m\bigr)}{%
    \rEE{6,7,14}}
\end{tabproofwide}

\FormulaDefDeltaK[Menge der \(\JoinOp\)-Irreduziblen]{%
  \JoinIrrSet{\JoinOp,A,0}
  \coloneqq\{j\in A\mid\JoinIrr{\JoinOp,A,0}{j}\}
}{JoinIrreducibleSet}{%
  \DeltaRow{Mengen}{A\dsep j\dsep 0}
  \DeltaRow{Funktionen}{\JoinOp}
  \DeltaRow{\textbf{Neues Symbol}}{}
  \DeltaPrem{\(\JoinOp\)-Irreduzible}{\JoinIrrSet{\JoinOp,A,0}}
}

\FormulaDefDeltaK[Gemeinsame untere Schranken eines Elementpaares]{%
  \operatorname{US}_{\JoinOp,A}(x,y)
  \coloneqq
  \bigl\{d\in A\mid
    (d,x)\in\JoinOrd{\JoinOp}{A}\land
    (d,y)\in\JoinOrd{\JoinOp}{A}
  \bigr\}
}{PairLowerBoundSet}{%
  \DeltaRow{Mengen}{A\dsep d\dsep x\dsep y}
  \DeltaRow{Funktionen}{\JoinOp}
  \DeltaRow{Relationen}{\JoinOrd{\JoinOp}{A}}
  \DeltaRow{\textbf{Neues Symbol}}{}
  \DeltaPrem{Gemeinsame untere Schranken}{%
    \operatorname{US}_{\JoinOp,A}(x,y)}
}

Dies ist die in Band~03 definierte untere Schrankenmenge für die
Zweiermenge \(\{x,y\}\), nun in der vom Halbverband induzierten Ordnung.

\FormulaThmDeltaK[Charakterisierung der gemeinsamen unteren Schranken]{%
  \begin{aligned}[t]
    d\in\operatorname{US}_{\JoinOp,A}(x,y)
    \quad\eqvdash\quad{}&d\in A\land
      (d,x)\in\JoinOrd{\JoinOp}{A}\land{}\\
    &(d,y)\in\JoinOrd{\JoinOp}{A}
  \end{aligned}
}{PairLowerBoundSetCharacterization}{%
  \DeltaRow{Mengen}{A\dsep d\dsep x\dsep y}
  \DeltaRow{Funktionen}{\JoinOp}
  \DeltaRow{Relationen}{\JoinOrd{\JoinOp}{A}}
}
\begin{tabproofwide}
  \proofstepwide{d\in\operatorname{US}_{\JoinOp,A}(x,y)}{\leftrightarrow}{%
    d\in A\land(d,x)\in\JoinOrd{\JoinOp}{A}\land
    (d,y)\in\JoinOrd{\JoinOp}{A}}{%
    \FormulaRefAuto{PairLowerBoundSet}[def]}
\end{tabproofwide}

\FormulaThmDeltaK[Gemeinsame untere Schranken liegen im Träger]{%
  \operatorname{US}_{\JoinOp,A}(x,y)\subseteq A
}{PairLowerBoundSetSubset}{%
  \DeltaRow{Mengen}{A\dsep d\dsep x\dsep y}
  \DeltaRow{Funktionen}{\JoinOp}
}
\begin{tabproof}
  \proofstep{1}{d\in\operatorname{US}_{\JoinOp,A}(x,y)}{\rA}
  \proofstep{1}{%
    d\in A\land(d,x)\in\JoinOrd{\JoinOp}{A}\land
    (d,y)\in\JoinOrd{\JoinOp}{A}}{%
    \FormulaRefAuto{P\leftrightarrow Q, P\vdash Q}{%
      \FormulaRefAuto{PairLowerBoundSetCharacterization}[thm],1}}
  \proofstep{1}{d\in A}{\rAEa{\rAEa{2}}}
  \proofstep{}{\operatorname{US}_{\JoinOp,A}(x,y)\subseteq A}{%
    \FormulaRefAuto{A \subseteq B \coloneq \forall x\,(x\in A \rightarrow x\in B)}[def]{%
      \rUI{\rRI{1,3}}}}
\end{tabproof}

\FormulaThmDeltaK[Endlichkeit der gemeinsamen unteren Schranken]{%
  \Finite A\vdash\Finite{\operatorname{US}_{\JoinOp,A}(x,y)}
}{PairLowerBoundSetFinite}{%
  \DeltaRow{Mengen}{A\dsep x\dsep y}
  \DeltaRow{Funktionen}{\JoinOp}
}
\begin{tabproof}
  \proofstep{1}{\Finite A}{\rA}
  \proofstep{}{\operatorname{US}_{\JoinOp,A}(x,y)\subseteq A}{%
    \FormulaRefAuto{PairLowerBoundSetSubset}[thm]}
  \proofstep{1}{\Finite{\operatorname{US}_{\JoinOp,A}(x,y)}}{%
    \FormulaRefAuto{FiniteSubset}[thm]{2,1}}
\end{tabproof}

\FormulaThmDeltaK[Das Nullelement ist eine gemeinsame untere Schranke]{%
  \ZeroJoinSemi{\JoinOp,A,0}\dsep x,y\in A
  \vdash 0\in\operatorname{US}_{\JoinOp,A}(x,y)
}{PairLowerBoundSetContainsZero}{%
  \DeltaRow{Mengen}{A\dsep 0\dsep x\dsep y}
  \DeltaRow{Funktionen}{\JoinOp}
}
\begin{tabproof}
  \proofstep{1}{\ZeroJoinSemi{\JoinOp,A,0}}{\rA}
  \proofstep{2}{x\in A}{\rA}
  \proofstep{3}{y\in A}{\rA}
  \proofstep{1}{0\in A}{\FormulaRefAuto{ZeroJoinCarrierAxiom}[ax]{1}}
  \proofstep{1,2}{(0,x)\in\JoinOrd{\JoinOp}{A}}{%
    \FormulaRefAuto{ZeroJoinLeast}[thm]{1,2}}
  \proofstep{1,3}{(0,y)\in\JoinOrd{\JoinOp}{A}}{%
    \FormulaRefAuto{ZeroJoinLeast}[thm]{1,3}}
  \proofstep{1,2,3}{%
    0\in A\land(0,x)\in\JoinOrd{\JoinOp}{A}\land
    (0,y)\in\JoinOrd{\JoinOp}{A}}{\rAI{\rAI{4,5},6}}
  \proofstep{1,2,3}{0\in\operatorname{US}_{\JoinOp,A}(x,y)}{%
    \FormulaRefAuto{P\leftrightarrow Q, Q\vdash P}{%
      \FormulaRefAuto{PairLowerBoundSetCharacterization}[thm],7}}
\end{tabproof}

\FormulaThmDeltaK[Abschluss der gemeinsamen unteren Schranken]{%
  \begin{aligned}[t]
    &\ZeroJoinSemi{\JoinOp,A,0}\dsep x,y\in A\dsep
      d,e\in\operatorname{US}_{\JoinOp,A}(x,y)\vdash{}\\
    &d\JoinOp{}e\in\operatorname{US}_{\JoinOp,A}(x,y)
  \end{aligned}
}{PairLowerBoundSetJoinClosed}{%
  \DeltaRow{Mengen}{A\dsep 0\dsep d\dsep e\dsep x\dsep y}
  \DeltaRow{Funktionen}{\JoinOp}
}
\begin{tabproofwide}
  \proofstepwidestar[1]{\ZeroJoinSemi{\JoinOp,A,0}}{\rA}
  \proofstepwidestar[2]{x\in A}{\rA}
  \proofstepwidestar[3]{y\in A}{\rA}
  \proofstepwidestar[4]{d\in\operatorname{US}_{\JoinOp,A}(x,y)}{\rA}
  \proofstepwidestar[5]{e\in\operatorname{US}_{\JoinOp,A}(x,y)}{\rA}
  \proofstepwidestar[1]{\Semi{\JoinOp,A}}{%
    \FormulaRefAuto{ZeroJoinBaseSemilatticeAxiom}[ax]{1}}
  \proofstepwidestar[4]{%
    d\in A\land(d,x)\in\JoinOrd{\JoinOp}{A}\land
    (d,y)\in\JoinOrd{\JoinOp}{A}}{%
    \FormulaRefAuto{P\leftrightarrow Q, P\vdash Q}{%
      \FormulaRefAuto{PairLowerBoundSetCharacterization}[thm],4}}
  \proofstepwidestar[5]{%
    e\in A\land(e,x)\in\JoinOrd{\JoinOp}{A}\land
    (e,y)\in\JoinOrd{\JoinOp}{A}}{%
    \FormulaRefAuto{P\leftrightarrow Q, P\vdash Q}{%
      \FormulaRefAuto{PairLowerBoundSetCharacterization}[thm],5}}
  \proofstepwidestar[4]{d\in A}{\rAEa{\rAEa{7}}}
  \proofstepwidestar[5]{e\in A}{\rAEa{\rAEa{8}}}
  \proofstepwidestar[4]{(d,x)\in\JoinOrd{\JoinOp}{A}}{%
    \FormulaRefAuto{P\land Q\land R\vdash Q}{7}}
  \proofstepwidestar[5]{(e,x)\in\JoinOrd{\JoinOp}{A}}{%
    \FormulaRefAuto{P\land Q\land R\vdash Q}{8}}
  \proofstepwidestar[4]{(d,y)\in\JoinOrd{\JoinOp}{A}}{%
    \FormulaRefAuto{P\land Q\land R\vdash R}{7}}
  \proofstepwidestar[5]{(e,y)\in\JoinOrd{\JoinOp}{A}}{%
    \FormulaRefAuto{P\land Q\land R\vdash R}{8}}
  \proofstepwidestar[1,4,5]{d\JoinOp{}e\in A}{%
    \FormulaRefAuto{SemilatticeClosureAxiom}[ax]{6,9,10}}
  \proofstepwidestar[1,2,4,5]{%
    (d\JoinOp{}e,x)\in\JoinOrd{\JoinOp}{A}}{%
    \FormulaRefAuto{JoinLeastUpper}[thm]{6,9,10,2,11,12}}
  \proofstepwidestar[1,3,4,5]{%
    (d\JoinOp{}e,y)\in\JoinOrd{\JoinOp}{A}}{%
    \FormulaRefAuto{JoinLeastUpper}[thm]{6,9,10,3,13,14}}
  \proofstepwidestar[1,2,3,4,5]{%
    d\JoinOp{}e\in\operatorname{US}_{\JoinOp,A}(x,y)}{%
    \FormulaRefAuto{P\leftrightarrow Q, Q\vdash P}{%
      \FormulaRefAuto{PairLowerBoundSetCharacterization}[thm],
      \rAI{\rAI{15,16},17}}}
\end{tabproofwide}

\FormulaThmDeltaK[Ein maximales Element einer abgeschlossenen Teilmenge ist größtes Element]{%
  \begin{aligned}[t]
    &\Semi{\JoinOp,A}\dsep D\subseteq A\dsep m\in D\dsep
      \forall q\in D\,\bigl((m,q)\in\JoinOrd{\JoinOp}{A}\to q=m\bigr)\dsep{}\\
    &\forall d,e\in D\,(d\JoinOp{}e\in D)
      \vdash\forall d\in D\,(d,m)\in\JoinOrd{\JoinOp}{A}
  \end{aligned}
}{JoinClosedMaximalIsGreatest}{%
  \DeltaRow{Mengen}{A\dsep D\dsep d\dsep e\dsep m\dsep q}
  \DeltaRow{Funktionen}{\JoinOp}
}
\begin{tabproofwide}
  \proofstepwidestar[1]{\Semi{\JoinOp,A}}{\rA}
  \proofstepwidestar[2]{D\subseteq A}{\rA}
  \proofstepwidestar[3]{m\in D}{\rA}
  \proofstepwidestar[4]{%
    \forall q\in D\,\bigl((m,q)\in\JoinOrd{\JoinOp}{A}\to q=m\bigr)}{\rA}
  \proofstepwidestar[5]{\forall d,e\in D\,(d\JoinOp{}e\in D)}{\rA}
  \proofstepwidestar[6]{d\in D}{\rA}
  \proofstepwidestar[2,3]{m\in A}{%
    \FormulaRefAuto{A\subseteq B,\,x\in A\vdash x\in B}{2,3}}
  \proofstepwidestar[2,6]{d\in A}{%
    \FormulaRefAuto{A\subseteq B,\,x\in A\vdash x\in B}{2,6}}
  \proofstepwidestar[3,5,6]{m\JoinOp{}d\in D}{%
    \rRE{6,\rRE{3,\rUE{\rUE{5}}}}}
  \proofstepwidestar[1,2,3,6]{%
    (m,m\JoinOp{}d)\in\JoinOrd{\JoinOp}{A}}{%
    \FormulaRefAuto{JoinUpperLeft}[thm]{1,7,8}}
  \proofstepwidestar[3,4,5,6]{m\JoinOp{}d=m}{%
    \rRE{\rRE{9,\rUE{4}},10}}
  \proofstepwidestar[1,2,3,6]{%
    (d,m\JoinOp{}d)\in\JoinOrd{\JoinOp}{A}}{%
    \FormulaRefAuto{JoinUpperRight}[thm]{1,7,8}}
  \proofstepwidestar[1,2,3,4,5,6]{%
    (d,m)\in\JoinOrd{\JoinOp}{A}}{\rIE{11,12}}
  \proofstepwidestar[1,2,3,4,5]{%
    \forall d\in D\,(d,m)\in\JoinOrd{\JoinOp}{A}}{%
    \rUI{\rRI{6,13}}}
\end{tabproofwide}

\FormulaThmDeltaK[Ein größtes gemeinsames unteres Element ist das Paarinfimum]{%
  \begin{aligned}[t]
    &\Semi{\JoinOp,A}\dsep x,y\in A\dsep
      m\in\operatorname{US}_{\JoinOp,A}(x,y)\dsep{}\\
    &\forall d\in\operatorname{US}_{\JoinOp,A}(x,y)\,
      (d,m)\in\JoinOrd{\JoinOp}{A}\vdash{}\\
    &\PairInf{\JoinOrd{\JoinOp}{A}}{A}{x}{y}{m}
  \end{aligned}
}{PairLowerBoundGreatestIsInfimum}{%
  \DeltaRow{Mengen}{A\dsep d\dsep m\dsep v\dsep x\dsep y}
  \DeltaRow{Funktionen}{\JoinOp}
}
\begin{tabproofwide}
  \proofstepwidestar[1]{\Semi{\JoinOp,A}}{\rA}
  \proofstepwidestar[2]{x\in A}{\rA}
  \proofstepwidestar[3]{y\in A}{\rA}
  \proofstepwidestar[4]{m\in\operatorname{US}_{\JoinOp,A}(x,y)}{\rA}
  \proofstepwidestar[5]{%
    \forall d\in\operatorname{US}_{\JoinOp,A}(x,y)\,
      (d,m)\in\JoinOrd{\JoinOp}{A}}{\rA}
  \proofstepwidestar[4]{%
    m\in A\land(m,x)\in\JoinOrd{\JoinOp}{A}\land
    (m,y)\in\JoinOrd{\JoinOp}{A}}{%
    \FormulaRefAuto{P\leftrightarrow Q, P\vdash Q}{%
      \FormulaRefAuto{PairLowerBoundSetCharacterization}[thm],4}}
  \proofstepwidestar[4]{m\in A}{\rAEa{\rAEa{6}}}
  \proofstepwidestar[4]{(m,x)\in\JoinOrd{\JoinOp}{A}}{%
    \FormulaRefAuto{P\land Q\land R\vdash Q}{6}}
  \proofstepwidestar[4]{(m,y)\in\JoinOrd{\JoinOp}{A}}{%
    \FormulaRefAuto{P\land Q\land R\vdash R}{6}}
  \proofstepwidestar[10]{v\in A}{\rA}
  \proofstepwidestar[11]{(v,x)\in\JoinOrd{\JoinOp}{A}}{\rA}
  \proofstepwidestar[12]{(v,y)\in\JoinOrd{\JoinOp}{A}}{\rA}
  \proofstepwidestar[10,11,12]{v\in\operatorname{US}_{\JoinOp,A}(x,y)}{%
    \FormulaRefAuto{P\leftrightarrow Q, Q\vdash P}{%
      \FormulaRefAuto{PairLowerBoundSetCharacterization}[thm],
      \rAI{\rAI{10,11},12}}}
  \proofstepwidestar[5,10,11,12]{(v,m)\in\JoinOrd{\JoinOp}{A}}{%
    \rRE{13,\rUE{5}}}
  \proofstepwidestar[1,2,3,4,5]{%
    \PairInf{\JoinOrd{\JoinOp}{A}}{A}{x}{y}{m}}{%
    \FormulaRefAuto{PairInfimum}[def]{%
      \rAI{\rAI{2,3},7},8,9,14}}
\end{tabproofwide}

\FormulaThmDeltaK[Paarinfima in endlichen Halbverbänden mit Nullelement]{%
  \begin{aligned}[t]
    &\Finite A\dsep\ZeroJoinSemi{\JoinOp,A,0}\dsep x,y\in A\vdash{}\\
    &\exists m\in A\,
      \PairInf{\JoinOrd{\JoinOp}{A}}{A}{x}{y}{m}
  \end{aligned}
}{FiniteZeroJoinSemilatticePairInfima}{%
  \DeltaRow{Mengen}{A\dsep d\dsep e\dsep m\dsep x\dsep y\dsep 0}
  \DeltaRow{Funktionen}{\JoinOp}
}
\begin{tabproofwide}
  \proofstepwidestar[1]{\Finite A}{\rA}
  \proofstepwidestar[2]{\ZeroJoinSemi{\JoinOp,A,0}}{\rA}
  \proofstepwidestar[3]{x\in A}{\rA}
  \proofstepwidestar[4]{y\in A}{\rA}
  \proofstepwidestar[2]{\Semi{\JoinOp,A}}{%
    \FormulaRefAuto{ZeroJoinBaseSemilatticeAxiom}[ax]{2}}
  \proofstepwidestar[2]{\PartOrd{\JoinOrd{\JoinOp}{A},A}}{%
    \FormulaRefAuto{JoinSemilatticeInducesOrder}[thm]{5}}
  \proofstepwidestar[]{\operatorname{US}_{\JoinOp,A}(x,y)\subseteq A}{%
    \FormulaRefAuto{PairLowerBoundSetSubset}[thm]}
  \proofstepwidestar[1]{\Finite{\operatorname{US}_{\JoinOp,A}(x,y)}}{%
    \FormulaRefAuto{PairLowerBoundSetFinite}[thm]{1}}
  \proofstepwidestar[2,3,4]{%
    0\in\operatorname{US}_{\JoinOp,A}(x,y)}{%
    \FormulaRefAuto{PairLowerBoundSetContainsZero}[thm]{2,3,4}}
  \proofstepwidestar[2,3,4]{%
    \operatorname{US}_{\JoinOp,A}(x,y)\neq\varnothing}{%
    \FormulaRefAuto{A\neq\varnothing \eqvdash \exists x\,(x \in A)}{%
      \rEI{9}}}
  \proofstepwidestar[2,3,4]{%
    d,e\in\operatorname{US}_{\JoinOp,A}(x,y)
    \vdash d\JoinOp{}e\in\operatorname{US}_{\JoinOp,A}(x,y)}{%
    \FormulaRefAuto{PairLowerBoundSetJoinClosed}[thm]{2,3,4}}
  \proofstepwidestar[1,2,3,4]{%
    \begin{aligned}[t]
      &\exists m\in\operatorname{US}_{\JoinOp,A}(x,y)\,
        \forall d\in\operatorname{US}_{\JoinOp,A}(x,y)\\
      &\qquad ((m,d)\in\JoinOrd{\JoinOp}{A}\to d=m)
    \end{aligned}}{%
    \FormulaRefAuto{FinitePosetMaximalElement}[thm]{6,8,10,7}}
  \proofstepwidestar[12]{%
    \begin{aligned}[t]
      &m\in\operatorname{US}_{\JoinOp,A}(x,y)\land
        \forall d\in\operatorname{US}_{\JoinOp,A}(x,y)\\
      &\qquad ((m,d)\in\JoinOrd{\JoinOp}{A}\to d=m)
    \end{aligned}}{\rA}
  \proofstepwidestar[2,3,4,13]{%
    \forall d\in\operatorname{US}_{\JoinOp,A}(x,y)\,
      (d,m)\in\JoinOrd{\JoinOp}{A}}{%
    \FormulaRefAuto{JoinClosedMaximalIsGreatest}[thm]{%
      5,7,\rAEa{13},\rAEb{13},11}}
  \proofstepwidestar[2,3,4,13]{%
    \PairInf{\JoinOrd{\JoinOp}{A}}{A}{x}{y}{m}}{%
    \FormulaRefAuto{PairLowerBoundGreatestIsInfimum}[thm]{%
      5,3,4,\rAEa{13},14}}
  \proofstepwidestar[13]{m\in A}{%
    \rAEa{\rAEa{%
      \FormulaRefAuto{P\leftrightarrow Q, P\vdash Q}{%
        \FormulaRefAuto{PairLowerBoundSetCharacterization}[thm],
        \rAEa{13}}}}}
  \proofstepwidestar[2,3,4,13]{%
    m\in A\land\PairInf{\JoinOrd{\JoinOp}{A}}{A}{x}{y}{m}}{%
    \rAI{16,15}}
  \proofstepwidestar[1,2,3,4]{%
    \exists m\in A\,\PairInf{\JoinOrd{\JoinOp}{A}}{A}{x}{y}{m}}{%
    \rEE{12,13,\rEI{17}}}
\end{tabproofwide}

\FormulaDefDeltaK[Trennende untere Elemente]{%
  \operatorname{Sep}_{\JoinOp,A}(x,y)
  \coloneqq
  \bigl\{d\in A\mid
    (d,x)\in\JoinOrd{\JoinOp}{A}\land
    (d,y)\notin\JoinOrd{\JoinOp}{A}
  \bigr\}
}{JoinSeparatorSet}{%
  \DeltaRow{Mengen}{A\dsep d\dsep x\dsep y}
  \DeltaRow{Funktionen}{\JoinOp}
  \DeltaRow{Relationen}{\JoinOrd{\JoinOp}{A}}
  \DeltaRow{\textbf{Neues Symbol}}{}
  \DeltaPrem{Trennende untere Elemente}{%
    \operatorname{Sep}_{\JoinOp,A}(x,y)}
}

\FormulaThmDeltaK[Charakterisierung der trennenden unteren Elemente]{%
  \begin{aligned}[t]
    d\in\operatorname{Sep}_{\JoinOp,A}(x,y)
    \quad\eqvdash\quad{}&d\in A\land
      (d,x)\in\JoinOrd{\JoinOp}{A}\land{}\\
    &(d,y)\notin\JoinOrd{\JoinOp}{A}
  \end{aligned}
}{JoinSeparatorSetCharacterization}{%
  \DeltaRow{Mengen}{A\dsep d\dsep x\dsep y}
  \DeltaRow{Funktionen}{\JoinOp}
}
\begin{tabproofwide}
  \proofstepwide{d\in\operatorname{Sep}_{\JoinOp,A}(x,y)}{\leftrightarrow}{%
    d\in A\land(d,x)\in\JoinOrd{\JoinOp}{A}\land
    (d,y)\notin\JoinOrd{\JoinOp}{A}}{%
    \FormulaRefAuto{JoinSeparatorSet}[def]}
\end{tabproofwide}

\FormulaThmDeltaK[Trennende untere Elemente liegen im Träger]{%
  \operatorname{Sep}_{\JoinOp,A}(x,y)\subseteq A
}{JoinSeparatorSetSubset}{%
  \DeltaRow{Mengen}{A\dsep d\dsep x\dsep y}
  \DeltaRow{Funktionen}{\JoinOp}
}
\begin{tabproof}
  \proofstep{1}{d\in\operatorname{Sep}_{\JoinOp,A}(x,y)}{\rA}
  \proofstep{1}{%
    d\in A\land(d,x)\in\JoinOrd{\JoinOp}{A}\land
    (d,y)\notin\JoinOrd{\JoinOp}{A}}{%
    \FormulaRefAuto{P\leftrightarrow Q, P\vdash Q}{%
      \FormulaRefAuto{JoinSeparatorSetCharacterization}[thm],1}}
  \proofstep{1}{d\in A}{\rAEa{\rAEa{2}}}
  \proofstep{}{\operatorname{Sep}_{\JoinOp,A}(x,y)\subseteq A}{%
    \FormulaRefAuto{A \subseteq B \coloneq \forall x\,(x\in A \rightarrow x\in B)}[def]{%
      \rUI{\rRI{1,3}}}}
\end{tabproof}

\FormulaThmDeltaK[Das linke Element trennt sich selbst vom rechten]{%
  \begin{aligned}[t]
    &\ZeroJoinSemi{\JoinOp,A,0}\dsep x,y\in A\dsep
      (x,y)\notin\JoinOrd{\JoinOp}{A}\vdash{}\\
    &x\in\operatorname{Sep}_{\JoinOp,A}(x,y)
  \end{aligned}
}{JoinSeparatorSetContainsLeft}{%
  \DeltaRow{Mengen}{A\dsep 0\dsep x\dsep y}
  \DeltaRow{Funktionen}{\JoinOp}
}
\begin{tabproofwide}
  \proofstepwidestar[1]{\ZeroJoinSemi{\JoinOp,A,0}}{\rA}
  \proofstepwidestar[2]{x\in A}{\rA}
  \proofstepwidestar[3]{y\in A}{\rA}
  \proofstepwidestar[4]{(x,y)\notin\JoinOrd{\JoinOp}{A}}{\rA}
  \proofstepwidestar[1]{\Semi{\JoinOp,A}}{%
    \FormulaRefAuto{ZeroJoinBaseSemilatticeAxiom}[ax]{1}}
  \proofstepwidestar[1,2]{(x,x)\in\JoinOrd{\JoinOp}{A}}{%
    \FormulaRefAuto{JoinOrderReflexive}[thm]{5,2}}
  \proofstepwidestar[1,2,3,4]{%
    x\in A\land(x,x)\in\JoinOrd{\JoinOp}{A}\land
    (x,y)\notin\JoinOrd{\JoinOp}{A}}{\rAI{\rAI{2,6},4}}
  \proofstepwidestar[1,2,3,4]{%
    x\in\operatorname{Sep}_{\JoinOp,A}(x,y)}{%
    \FormulaRefAuto{P\leftrightarrow Q, Q\vdash P}{%
      \FormulaRefAuto{JoinSeparatorSetCharacterization}[thm],7}}
\end{tabproofwide}

\FormulaThmDeltaK[Endlichkeit der trennenden unteren Elemente]{%
  \Finite A\vdash\Finite{\operatorname{Sep}_{\JoinOp,A}(x,y)}
}{JoinSeparatorSetFinite}{%
  \DeltaRow{Mengen}{A\dsep x\dsep y}
  \DeltaRow{Funktionen}{\JoinOp}
}
\begin{tabproof}
  \proofstep{1}{\Finite A}{\rA}
  \proofstep{}{\operatorname{Sep}_{\JoinOp,A}(x,y)\subseteq A}{%
    \FormulaRefAuto{JoinSeparatorSetSubset}[thm]}
  \proofstep{1}{\Finite{\operatorname{Sep}_{\JoinOp,A}(x,y)}}{%
    \FormulaRefAuto{FiniteSubset}[thm]{2,1}}
\end{tabproof}

\FormulaThmDeltaK[Echte untere Trenner liegen unter dem rechten Element]{%
  \begin{aligned}[t]
    &\Semi{\JoinOp,A}\dsep x,y,u\in A\dsep
      j\in\operatorname{Sep}_{\JoinOp,A}(x,y)\dsep{}\\
    &\forall d\in\operatorname{Sep}_{\JoinOp,A}(x,y)\,
      \bigl((d,j)\in\JoinOrd{\JoinOp}{A}\to d=j\bigr)\dsep{}\\
    &(u,j)\in\JoinOrd{\JoinOp}{A}\dsep u\neq j
      \vdash (u,y)\in\JoinOrd{\JoinOp}{A}
  \end{aligned}
}{StrictLowerSeparatorBelowRight}{%
  \DeltaRow{Mengen}{A\dsep d\dsep j\dsep u\dsep x\dsep y}
  \DeltaRow{Funktionen}{\JoinOp}
}
\begin{tabproofwide}
  \proofstepwidestar[1]{\Semi{\JoinOp,A}}{\rA}
  \proofstepwidestar[2]{x\in A}{\rA}
  \proofstepwidestar[3]{y\in A}{\rA}
  \proofstepwidestar[4]{u\in A}{\rA}
  \proofstepwidestar[5]{j\in\operatorname{Sep}_{\JoinOp,A}(x,y)}{\rA}
  \proofstepwidestar[6]{%
    \forall d\in\operatorname{Sep}_{\JoinOp,A}(x,y)\,
      \bigl((d,j)\in\JoinOrd{\JoinOp}{A}\to d=j\bigr)}{\rA}
  \proofstepwidestar[7]{(u,j)\in\JoinOrd{\JoinOp}{A}}{\rA}
  \proofstepwidestar[8]{u\neq j}{\rA}
  \proofstepwidestar[5]{%
    j\in A\land(j,x)\in\JoinOrd{\JoinOp}{A}\land
    (j,y)\notin\JoinOrd{\JoinOp}{A}}{%
    \FormulaRefAuto{P\leftrightarrow Q, P\vdash Q}{%
      \FormulaRefAuto{JoinSeparatorSetCharacterization}[thm],5}}
  \proofstepwidestar[5]{(j,x)\in\JoinOrd{\JoinOp}{A}}{%
    \FormulaRefAuto{P\land Q\land R\vdash Q}{9}}
  \proofstepwidestar[5,7]{(u,x)\in\JoinOrd{\JoinOp}{A}}{%
    \FormulaRefAuto{(x,y)\in R \dsep (y,z)\in R
      \vdash (x,z)\in R}{7,10}}
  \proofstepwidestar[12]{(u,y)\notin\JoinOrd{\JoinOp}{A}}{\rA}
  \proofstepwidestar[4,5,7,12]{%
    u\in\operatorname{Sep}_{\JoinOp,A}(x,y)}{%
    \FormulaRefAuto{P\leftrightarrow Q, Q\vdash P}{%
      \FormulaRefAuto{JoinSeparatorSetCharacterization}[thm],
      \rAI{\rAI{4,11},12}}}
  \proofstepwidestar[5,6,7,12]{u=j}{%
    \rRE{\rRE{13,\rUE{6}},7}}
  \proofstepwidestar[5,6,7,8,12]{\bot}{\rBI{8,14}}
  \proofstepwidestar[1,2,3,4,5,6,7,8]{%
    \neg\bigl((u,y)\notin\JoinOrd{\JoinOp}{A}\bigr)}{\rCI{12,15}}
  \proofstepwidestar[1,2,3,4,5,6,7,8]{%
    (u,y)\in\JoinOrd{\JoinOp}{A}}{\FormulaRefAuto{rDN}{16}}
\end{tabproofwide}

\FormulaThmDeltaK[Ein minimaler Trenner ist \(\JoinOp\)-irreduzibel]{%
  \begin{aligned}[t]
    &\ZeroJoinSemi{\JoinOp,A,0}\dsep x,y\in A\dsep
      j\in\operatorname{Sep}_{\JoinOp,A}(x,y)\dsep{}\\
    &\forall d\in\operatorname{Sep}_{\JoinOp,A}(x,y)\,
      \bigl((d,j)\in\JoinOrd{\JoinOp}{A}\to d=j\bigr)
      \vdash\JoinIrr{\JoinOp,A,0}{j}
  \end{aligned}
}{MinimalSeparatorJoinIrreducible}{%
  \DeltaRow{Mengen}{A\dsep 0\dsep d\dsep j\dsep u\dsep v\dsep x\dsep y}
  \DeltaRow{Funktionen}{\JoinOp}
}
\begin{tabproofwide}
  \proofstepwidestar[1]{\ZeroJoinSemi{\JoinOp,A,0}}{\rA}
  \proofstepwidestar[2]{x\in A}{\rA}
  \proofstepwidestar[3]{y\in A}{\rA}
  \proofstepwidestar[4]{j\in\operatorname{Sep}_{\JoinOp,A}(x,y)}{\rA}
  \proofstepwidestar[5]{%
    \forall d\in\operatorname{Sep}_{\JoinOp,A}(x,y)\,
      \bigl((d,j)\in\JoinOrd{\JoinOp}{A}\to d=j\bigr)}{\rA}
  \proofstepwidestar[1]{\Semi{\JoinOp,A}}{%
    \FormulaRefAuto{ZeroJoinBaseSemilatticeAxiom}[ax]{1}}
  \proofstepwidestar[4]{%
    j\in A\land(j,x)\in\JoinOrd{\JoinOp}{A}\land
    (j,y)\notin\JoinOrd{\JoinOp}{A}}{%
    \FormulaRefAuto{P\leftrightarrow Q, P\vdash Q}{%
      \FormulaRefAuto{JoinSeparatorSetCharacterization}[thm],4}}
  \proofstepwidestar[4]{j\in A}{\rAEa{\rAEa{7}}}
  \proofstepwidestar[4]{(j,y)\notin\JoinOrd{\JoinOp}{A}}{%
    \FormulaRefAuto{P\land Q\land R\vdash R}{7}}
  \proofstepwidestar[1,3]{(0,y)\in\JoinOrd{\JoinOp}{A}}{%
    \FormulaRefAuto{ZeroJoinLeast}[thm]{1,3}}
  \proofstepwidestar[11]{j=0}{\rA}
  \proofstepwidestar[11]{0=j}{\FormulaRefAuto{a=b\vdash b=a}{11}}
  \proofstepwidestar[1,3,11]{(j,y)\in\JoinOrd{\JoinOp}{A}}{\rIE{12,10}}
  \proofstepwidestar[1,3,4,11]{\bot}{\rBI{9,13}}
  \proofstepwidestar[1,3,4]{j\neq0}{\rCI{11,14}}
  \proofstepwidestar[16]{u\in A}{\rA}
  \proofstepwidestar[17]{v\in A}{\rA}
  \proofstepwidestar[18]{u\JoinOp{}v=j}{\rA}
  \proofstepwidestar[6,16,17]{%
    (u,u\JoinOp{}v)\in\JoinOrd{\JoinOp}{A}}{%
    \FormulaRefAuto{JoinUpperLeft}[thm]{6,16,17}}
  \proofstepwidestar[6,16,17,18]{(u,j)\in\JoinOrd{\JoinOp}{A}}{%
    \rIE{18,19}}
  \proofstepwidestar[6,16,17]{%
    (v,u\JoinOp{}v)\in\JoinOrd{\JoinOp}{A}}{%
    \FormulaRefAuto{JoinUpperRight}[thm]{6,16,17}}
  \proofstepwidestar[6,16,17,18]{(v,j)\in\JoinOrd{\JoinOp}{A}}{%
    \rIE{18,21}}
  \proofstepwidestar[]{u=j\lor u\neq j}{\FormulaRefAuto{P\lor\neg P}}
  \proofstepwidestar[24]{u=j}{\rA}
  \proofstepwidestar[24]{u=j\lor v=j}{\rOIa{24}}
  \proofstepwidestar[26]{u\neq j}{\rA}
  \proofstepwidestar[1,2,3,4,5,16,20,26]{%
    (u,y)\in\JoinOrd{\JoinOp}{A}}{%
    \FormulaRefAuto{StrictLowerSeparatorBelowRight}[thm]{%
      6,2,3,16,4,5,20,26}}
  \proofstepwidestar[]{v=j\lor v\neq j}{\FormulaRefAuto{P\lor\neg P}}
  \proofstepwidestar[29]{v=j}{\rA}
  \proofstepwidestar[29]{u=j\lor v=j}{\rOIb{29}}
  \proofstepwidestar[31]{v\neq j}{\rA}
  \proofstepwidestar[1,2,3,4,5,17,22,31]{%
    (v,y)\in\JoinOrd{\JoinOp}{A}}{%
    \FormulaRefAuto{StrictLowerSeparatorBelowRight}[thm]{%
      6,2,3,17,4,5,22,31}}
  \proofstepwidestar[1,2,3,16,17,26,31]{%
    (u\JoinOp{}v,y)\in\JoinOrd{\JoinOp}{A}}{%
    \FormulaRefAuto{JoinLeastUpper}[thm]{6,16,17,3,27,32}}
  \proofstepwidestar[1,2,3,4,16,17,18,26,31]{%
    (j,y)\in\JoinOrd{\JoinOp}{A}}{\rIE{18,33}}
  \proofstepwidestar[1,2,3,4,16,17,18,26,31]{\bot}{\rBI{9,34}}
  \proofstepwidestar[1,2,3,4,16,17,18,26,31]{u=j\lor v=j}{%
    \FormulaRefAuto{\bot\vdash P}{35}}
  \proofstepwidestar[1,2,3,4,5,16,17,18,26]{u=j\lor v=j}{%
    \rOE{28,29,30,31,36}}
  \proofstepwidestar[1,2,3,4,5,16,17,18]{u=j\lor v=j}{%
    \rOE{23,24,25,26,37}}
  \proofstepwidestar[1,2,3,4,5]{\JoinIrr{\JoinOp,A,0}{j}}{%
    \FormulaRefAuto{JoinIrreducible}[def]{8,15,16,17,18,38}}
\end{tabproofwide}

\FormulaThmDeltaK[\(\JoinOp\)-Irreduzible trennen Elemente]{%
  \begin{aligned}[t]
    &\Finite A\dsep\ZeroJoinSemi{\JoinOp,A,0}\dsep x,y\in A\dsep
      (x,y)\notin\JoinOrd{\JoinOp}{A}\vdash{}\\
    &\exists j\in\JoinIrrSet{\JoinOp,A,0}\,
      \bigl((j,x)\in\JoinOrd{\JoinOp}{A}\land
            (j,y)\notin\JoinOrd{\JoinOp}{A}\bigr)
  \end{aligned}
}{JoinIrreduciblesSeparate}{%
  \DeltaRow{Mengen}{A\dsep d\dsep j\dsep x\dsep y\dsep 0}
  \DeltaRow{Funktionen}{\JoinOp}
}
\begin{tabproofwide}
  \proofstepwidestar[1]{\Finite A}{\rA}
  \proofstepwidestar[2]{\ZeroJoinSemi{\JoinOp,A,0}}{\rA}
  \proofstepwidestar[3]{x\in A}{\rA}
  \proofstepwidestar[4]{y\in A}{\rA}
  \proofstepwidestar[5]{(x,y)\notin\JoinOrd{\JoinOp}{A}}{\rA}
  \proofstepwidestar[2]{\Semi{\JoinOp,A}}{%
    \FormulaRefAuto{ZeroJoinBaseSemilatticeAxiom}[ax]{2}}
  \proofstepwidestar[2]{\PartOrd{\JoinOrd{\JoinOp}{A},A}}{%
    \FormulaRefAuto{JoinSemilatticeInducesOrder}[thm]{6}}
  \proofstepwidestar[]{\operatorname{Sep}_{\JoinOp,A}(x,y)\subseteq A}{%
    \FormulaRefAuto{JoinSeparatorSetSubset}[thm]}
  \proofstepwidestar[1]{\Finite{\operatorname{Sep}_{\JoinOp,A}(x,y)}}{%
    \FormulaRefAuto{JoinSeparatorSetFinite}[thm]{1}}
  \proofstepwidestar[2,3,4,5]{%
    x\in\operatorname{Sep}_{\JoinOp,A}(x,y)}{%
    \FormulaRefAuto{JoinSeparatorSetContainsLeft}[thm]{2,3,4,5}}
  \proofstepwidestar[2,3,4,5]{%
    \operatorname{Sep}_{\JoinOp,A}(x,y)\neq\varnothing}{%
    \FormulaRefAuto{A\neq\varnothing \eqvdash \exists x\,(x \in A)}{%
      \rEI{10}}}
  \proofstepwidestar[1,2,3,4]{%
    \exists j\in\operatorname{Sep}_{\JoinOp,A}(x,y)\,
      \forall d\in\operatorname{Sep}_{\JoinOp,A}(x,y)\,
      ((d,j)\in\JoinOrd{\JoinOp}{A}\to d=j)}{%
    \FormulaRefAuto{FinitePosetMinimalElement}[thm]{7,9,11,8}}
  \proofstepwidestar[12]{%
    j\in\operatorname{Sep}_{\JoinOp,A}(x,y)\land
    \forall d\in\operatorname{Sep}_{\JoinOp,A}(x,y)\,
      ((d,j)\in\JoinOrd{\JoinOp}{A}\to d=j)}{\rA}
  \proofstepwidestar[2,3,4,13]{\JoinIrr{\JoinOp,A,0}{j}}{%
    \FormulaRefAuto{MinimalSeparatorJoinIrreducible}[thm]{%
      2,3,4,\rAEa{13},\rAEb{13}}}
  \proofstepwidestar[13]{%
    j\in A\land(j,x)\in\JoinOrd{\JoinOp}{A}\land
    (j,y)\notin\JoinOrd{\JoinOp}{A}}{%
    \FormulaRefAuto{P\leftrightarrow Q, P\vdash Q}{%
      \FormulaRefAuto{JoinSeparatorSetCharacterization}[thm],
      \rAEa{13}}}
  \proofstepwidestar[2,3,4,13]{j\in\JoinIrrSet{\JoinOp,A,0}}{%
    \FormulaRefAuto{JoinIrreducibleSet}[def]{14}}
  \proofstepwidestar[13]{%
    (j,x)\in\JoinOrd{\JoinOp}{A}\land
    (j,y)\notin\JoinOrd{\JoinOp}{A}}{%
    \rAI{\FormulaRefAuto{P\land Q\land R\vdash Q}{15},
      \FormulaRefAuto{P\land Q\land R\vdash R}{15}}}
  \proofstepwidestar[2,3,4,13]{%
    \begin{aligned}[t]
      &\exists j\in\JoinIrrSet{\JoinOp,A,0}\,\bigl(
        (j,x)\in\JoinOrd{\JoinOp}{A}\land{}\\
      &\qquad (j,y)\notin\JoinOrd{\JoinOp}{A}\bigr)
    \end{aligned}}{\rEI{\rAI{16,17}}}
  \proofstepwidestar[1,2,3,4,5]{%
    \begin{aligned}[t]
      &\exists j\in\JoinIrrSet{\JoinOp,A,0}\,\bigl(
        (j,x)\in\JoinOrd{\JoinOp}{A}\land{}\\
      &\qquad (j,y)\notin\JoinOrd{\JoinOp}{A}\bigr)
    \end{aligned}}{\rEE{12,13,18}}
\end{tabproofwide}

\section{Vom Halbverband zur Vereinigungsfamilie}

\FormulaDefDeltaK[\(\JoinOp\)-Irreduzible unter einem Element]{%
  \IrrBelow{\JoinOp,A,0}{x}
  \coloneqq
  \{j\in\JoinIrrSet{\JoinOp,A,0}\mid(j,x)\in\JoinOrd{\JoinOp}{A}\}
}{IrreduciblesBelow}{%
  \DeltaRow{Mengen}{A\dsep j\dsep x\dsep 0}
  \DeltaRow{Funktionen}{\JoinOp}
  \DeltaRow{\textbf{Neues Symbol}}{}
  \DeltaPrem{Irreduziblenprofil}{\IrrBelow{\JoinOp,A,0}{x}}
}

\FormulaDefDeltaK[Kanonische Profilfamilien]{%
  \begin{aligned}[t]
    \CanonicalMeetFamily{\JoinOp,A,0}
      &\coloneqq\{\IrrBelow{\JoinOp,A,0}{x}\mid x\in A\},\\
    \mathcal U_{\JoinOp,A,0}
      &\coloneqq
      \{\JoinIrrSet{\JoinOp,A,0}\setminus\IrrBelow{\JoinOp,A,0}{x}\mid x\in A\}
  \end{aligned}
}{CanonicalProfileFamilies}{%
  \DeltaRow{Mengen}{A\dsep x\dsep 0}
  \DeltaRow{Funktionen}{\JoinOp}
  \DeltaRow{\textbf{Neue Symbole}}{}
  \DeltaPrem{Schnittprofilfamilie}{\CanonicalMeetFamily{\JoinOp,A,0}}
  \DeltaPrem{Komplementprofilfamilie}{\mathcal U_{\JoinOp,A,0}}
}

Die erste Profilfamilie ist schnittabgeschlossen. Ist \(m\) das Paarinfimum
von \(x\) und \(y\), so folgt unmittelbar aus dessen universeller Eigenschaft
\[
  \IrrBelow{\JoinOp,A,0}{m}
  =
  \IrrBelow{\JoinOp,A,0}{x}\cap\IrrBelow{\JoinOp,A,0}{y}.
\]
Die Komplementprofilfamilie ist daher nach de~Morgan
vereinigungsabgeschlossen. Entscheidend ist außerdem, dass die Profile alle
Elemente unterscheiden.

\FormulaThmDeltaK[Charakterisierung der Komplementprofilfamilie]{%
  \begin{aligned}[t]
    Y\in\mathcal U_{\JoinOp,A,0}
    \quad\Longleftrightarrow\quad
    \exists x\in A\,\bigl(
      Y={}&\JoinIrrSet{\JoinOp,A,0}\setminus
      \IrrBelow{\JoinOp,A,0}{x}\bigr)
  \end{aligned}
}{CanonicalProfileFamilyMembership}{%
  \DeltaRow{Mengen}{A\dsep Y\dsep x\dsep 0}
  \DeltaRow{Funktionen}{\JoinOp}
}
\begin{tabproofwide}
  \proofstepwide{Y\in\mathcal U_{\JoinOp,A,0}}{\leftrightarrow}{%
    \exists x\in A\,\bigl(
      Y=\JoinIrrSet{\JoinOp,A,0}\setminus
      \IrrBelow{\JoinOp,A,0}{x}\bigr)}{%
    \FormulaRefAuto{CanonicalProfileFamilies}[def]}
\end{tabproofwide}

\FormulaDefDeltaK[Funktionsvorschrift der kanonischen Profilabbildung]{%
  \ZeroJoinSemi{\JoinOp,A,0}\dsep x\in A\vdash
  \widehat c_{\JoinOp,A,0}(x)\coloneqq
  \JoinIrrSet{\JoinOp,A,0}\setminus\IrrBelow{\JoinOp,A,0}{x}
}{CanonicalProfileRule}{%
  \DeltaRow{Mengen}{A\dsep x\dsep 0}
  \DeltaRow{Funktionen}{\JoinOp}
  \DeltaRow{Funktionensymbol}{\widehat c_{\JoinOp,A,0}}
}

\FormulaThmDeltaK[Typisierung der kanonischen Profilvorschrift]{%
  \ZeroJoinSemi{\JoinOp,A,0}\dsep x\in A\vdash
  \widehat c_{\JoinOp,A,0}(x)\in\mathcal U_{\JoinOp,A,0}
}{CanonicalProfileRuleTyped}{%
  \DeltaRow{Mengen}{A\dsep x\dsep 0}
  \DeltaRow{Funktionen}{\JoinOp}
  \DeltaRow{Funktionensymbol}{\widehat c_{\JoinOp,A,0}}
}
\begin{tabproofwide}
  \proofstepwidestar[1]{\ZeroJoinSemi{\JoinOp,A,0}}{\rA}
  \proofstepwidestar[2]{x\in A}{\rA}
  \proofstepwidestar[1,2]{%
    \widehat c_{\JoinOp,A,0}(x)=
    \JoinIrrSet{\JoinOp,A,0}\setminus\IrrBelow{\JoinOp,A,0}{x}}{%
    \FormulaRefAuto{CanonicalProfileRule}[def]{1,2}}
  \proofstepwidestar[2]{%
    \exists u\in A\,\bigl(
      \JoinIrrSet{\JoinOp,A,0}\setminus\IrrBelow{\JoinOp,A,0}{x}
      =\JoinIrrSet{\JoinOp,A,0}\setminus\IrrBelow{\JoinOp,A,0}{u}\bigr)}{%
    \rEI{\rAI{2,\rII}}}
  \proofstepwidestar[2]{%
    \JoinIrrSet{\JoinOp,A,0}\setminus\IrrBelow{\JoinOp,A,0}{x}
    \in\mathcal U_{\JoinOp,A,0}}{%
    \FormulaRefAuto{CanonicalProfileFamilyMembership}[thm]{4}}
  \proofstepwidestar[1,2]{%
    \widehat c_{\JoinOp,A,0}(x)\in\mathcal U_{\JoinOp,A,0}}{%
    \rIE{3,5}}
\end{tabproofwide}

\FormulaThmDeltaK[Universelle Typisierung der kanonischen Profilvorschrift]{%
  \ZeroJoinSemi{\JoinOp,A,0}\vdash
  \forall x\in A\,
    \widehat c_{\JoinOp,A,0}(x)\in\mathcal U_{\JoinOp,A,0}
}{CanonicalProfileRuleTypedUniversal}{%
  \DeltaRow{Mengen}{A\dsep x\dsep 0}
  \DeltaRow{Funktionen}{\JoinOp}
  \DeltaRow{Funktionensymbol}{\widehat c_{\JoinOp,A,0}}
}
\begin{tabproof}
  \proofstep{1}{\ZeroJoinSemi{\JoinOp,A,0}}{\rA}
  \proofstep{2}{x\in A}{\rA}
  \proofstep{1,2}{%
    \widehat c_{\JoinOp,A,0}(x)\in\mathcal U_{\JoinOp,A,0}}{%
    \FormulaRefAuto{CanonicalProfileRuleTyped}[thm]{1,2}}
  \proofstep{1}{%
    \forall x\in A\,
      \widehat c_{\JoinOp,A,0}(x)\in\mathcal U_{\JoinOp,A,0}}{%
    \rUI{\rRI{2,3}}}
\end{tabproof}

\FormulaDefDeltaK[Kanonische Profilabbildung]{%
  \ZeroJoinSemi{\JoinOp,A,0}\vdash
  \CanonicalProfileMap{\JoinOp,A,0}
    \coloneqq\Graph(\widehat c_{\JoinOp,A,0})
}{CanonicalProfileMapDefinition}{%
  \DeltaRow{Mengen}{A\dsep x\dsep 0}
  \DeltaRow{Funktionen}{\JoinOp\dsep\CanonicalProfileMap{\JoinOp,A,0}}
  \DeltaRow{Funktionensymbol}{\widehat c_{\JoinOp,A,0}}
  \DeltaRow{\textbf{Neues Symbol}}{}
  \DeltaPrem{Profilabbildung}{\CanonicalProfileMap{\JoinOp,A,0}}
}

\FormulaThmDeltaK[Die kanonische Profilabbildung ist eine Funktion]{%
  \ZeroJoinSemi{\JoinOp,A,0}\vdash
  \CanonicalProfileMap{\JoinOp,A,0}
    \colon A\to\mathcal U_{\JoinOp,A,0}
}{CanonicalProfileFunction}{%
  \DeltaRow{Mengen}{A\dsep x\dsep 0}
  \DeltaRow{Funktionen}{\JoinOp\dsep\CanonicalProfileMap{\JoinOp,A,0}}
  \DeltaRow{Funktionensymbol}{\widehat c_{\JoinOp,A,0}}
}
\begin{tabproofwide}
  \proofstepwidestar[1]{\ZeroJoinSemi{\JoinOp,A,0}}{\rA}
  \proofstepwidestar[1]{%
    \forall x\in A\,
      \widehat c_{\JoinOp,A,0}(x)\in\mathcal U_{\JoinOp,A,0}}{%
    \FormulaRefAuto{CanonicalProfileRuleTypedUniversal}[thm]{1}}
  \proofstepwidestar[1]{%
    \Graph(\widehat c_{\JoinOp,A,0})
      \colon A\to\mathcal U_{\JoinOp,A,0}}{%
    \FormulaRefAuto{\forall u\in A\,F(u)\in B
      \vdash \Graph(F)\colon A\to B}{2}}
  \proofstepwidestar[1]{%
    \CanonicalProfileMap{\JoinOp,A,0}
      =\Graph(\widehat c_{\JoinOp,A,0})}{%
    \FormulaRefAuto{CanonicalProfileMapDefinition}[def]{1}}
  \proofstepwidestar[1]{%
    \CanonicalProfileMap{\JoinOp,A,0}
      \colon A\to\mathcal U_{\JoinOp,A,0}}{%
    \rIE{4,3}}
\end{tabproofwide}

\FormulaThmDeltaK[Auswertung der kanonischen Profilabbildung]{%
  \begin{aligned}[t]
    &\ZeroJoinSemi{\JoinOp,A,0}\dsep x\in A\vdash{}\\
    &\CanonicalProfileMap{\JoinOp,A,0}(x)=
      \JoinIrrSet{\JoinOp,A,0}\setminus\IrrBelow{\JoinOp,A,0}{x}
  \end{aligned}
}{CanonicalProfileEvaluation}{%
  \DeltaRow{Mengen}{A\dsep x\dsep 0}
  \DeltaRow{Funktionen}{\JoinOp\dsep\CanonicalProfileMap{\JoinOp,A,0}}
  \DeltaRow{Funktionensymbol}{\widehat c_{\JoinOp,A,0}}
}
\begin{tabproofwide}
  \proofstepwidestar[1]{\ZeroJoinSemi{\JoinOp,A,0}}{\rA}
  \proofstepwidestar[2]{x\in A}{\rA}
  \proofstepwidestar[1]{%
    \CanonicalProfileMap{\JoinOp,A,0}
      =\Graph(\widehat c_{\JoinOp,A,0})}{%
    \FormulaRefAuto{CanonicalProfileMapDefinition}[def]{1}}
  \proofstepwidestar[1,2]{%
    \Graph(\widehat c_{\JoinOp,A,0})(x)
      =\widehat c_{\JoinOp,A,0}(x)}{%
    \FormulaRefAuto{\forall u\in A\,F(u)\in B
      \vdash \forall x\in A\,\Graph(F)(x)=F(x)}{%
      \FormulaRefAuto{CanonicalProfileRuleTypedUniversal}[thm]{1},2}}
  \proofstepwidestar[1,2]{%
    \CanonicalProfileMap{\JoinOp,A,0}(x)
      =\widehat c_{\JoinOp,A,0}(x)}{%
    \rIE{3,4}}
  \proofstepwidestar[1,2]{%
    \widehat c_{\JoinOp,A,0}(x)=
      \JoinIrrSet{\JoinOp,A,0}\setminus\IrrBelow{\JoinOp,A,0}{x}}{%
    \FormulaRefAuto{CanonicalProfileRule}[def]{1,2}}
  \proofstepwidestar[1,2]{%
    \CanonicalProfileMap{\JoinOp,A,0}(x)=
      \JoinIrrSet{\JoinOp,A,0}\setminus\IrrBelow{\JoinOp,A,0}{x}}{%
    \FormulaRefAuto{a=b,\,b=c\vdash a=c}{5,6}}
\end{tabproofwide}

\FormulaThmDeltaK[Die kanonische Profilabbildung ist surjektiv]{%
  \ZeroJoinSemi{\JoinOp,A,0}\vdash
  \CanonicalProfileMap{\JoinOp,A,0}
    \colon A\sur\mathcal U_{\JoinOp,A,0}
}{CanonicalProfileSurjective}{%
  \DeltaRow{Mengen}{A\dsep Y\dsep x\dsep 0}
  \DeltaRow{Funktionen}{\JoinOp\dsep\CanonicalProfileMap{\JoinOp,A,0}}
}
\begin{tabproofwide}
  \proofstepwidestar[1]{\ZeroJoinSemi{\JoinOp,A,0}}{\rA}
  \proofstepwidestar[1]{%
    \CanonicalProfileMap{\JoinOp,A,0}
      \colon A\to\mathcal U_{\JoinOp,A,0}}{%
    \FormulaRefAuto{CanonicalProfileFunction}[thm]{1}}
  \proofstepwidestar[3]{Y\in\mathcal U_{\JoinOp,A,0}}{\rA}
  \proofstepwidestar[3]{%
    \exists x\in A\,\bigl(
      Y=\JoinIrrSet{\JoinOp,A,0}\setminus
        \IrrBelow{\JoinOp,A,0}{x}\bigr)}{%
    \FormulaRefAuto{CanonicalProfileFamilyMembership}[thm]{3}}
  \proofstepwidestar[5]{%
    x\in A\land
    Y=\JoinIrrSet{\JoinOp,A,0}\setminus
      \IrrBelow{\JoinOp,A,0}{x}}{\rA}
  \proofstepwidestar[5]{x\in A}{\rAEa{5}}
  \proofstepwidestar[1,5]{%
    \CanonicalProfileMap{\JoinOp,A,0}(x)=
      \JoinIrrSet{\JoinOp,A,0}\setminus\IrrBelow{\JoinOp,A,0}{x}}{%
    \FormulaRefAuto{CanonicalProfileEvaluation}[thm]{1,6}}
  \proofstepwidestar[1,5]{%
    \CanonicalProfileMap{\JoinOp,A,0}(x)=Y}{%
    \FormulaRefAuto{a=b,\,c=b\vdash a=c}{7,\rAEb{5}}}
  \proofstepwidestar[1,3]{%
    \exists x\in A\,\bigl(
      \CanonicalProfileMap{\JoinOp,A,0}(x)=Y\bigr)}{%
    \rEE{4,5,\rEI{\rAI{6,8}}}}
  \proofstepwidestar[1]{%
    \forall Y\in\mathcal U_{\JoinOp,A,0}\,
      \exists x\in A\,\bigl(
        \CanonicalProfileMap{\JoinOp,A,0}(x)=Y\bigr)}{%
    \rUI{\rRI{3,9}}}
  \proofstepwidestar[1]{%
    \CanonicalProfileMap{\JoinOp,A,0}
      \colon A\sur\mathcal U_{\JoinOp,A,0}}{%
    \FormulaRefAuto{Surjektive Funktion}[def]{2,10}}
\end{tabproofwide}

\FormulaThmDeltaK[Irreduziblenprofile liegen im Irreduziblenträger]{%
  \ZeroJoinSemi{\JoinOp,A,0}\dsep x\in A\vdash
  \IrrBelow{\JoinOp,A,0}{x}\subseteq\JoinIrrSet{\JoinOp,A,0}
}{IrreducibleProfileSubset}{%
  \DeltaRow{Mengen}{A\dsep j\dsep x\dsep 0}
  \DeltaRow{Funktionen}{\JoinOp}
}
\begin{tabproofwide}
  \proofstepwidestar[1]{\ZeroJoinSemi{\JoinOp,A,0}}{\rA}
  \proofstepwidestar[2]{x\in A}{\rA}
  \proofstepwidestar[3]{j\in\IrrBelow{\JoinOp,A,0}{x}}{\rA}
  \proofstepwidestar[3]{%
    j\in\JoinIrrSet{\JoinOp,A,0}\land
      (j,x)\in\JoinOrd{\JoinOp}{A}}{%
    \FormulaRefAuto{IrreduciblesBelow}[def]{3}}
  \proofstepwidestar[3]{j\in\JoinIrrSet{\JoinOp,A,0}}{\rAEa{4}}
  \proofstepwidestar[]{%
    \IrrBelow{\JoinOp,A,0}{x}\subseteq\JoinIrrSet{\JoinOp,A,0}}{%
    \FormulaRefAuto{A \subseteq B \coloneq
      \forall x\,(x\in A\rightarrow x\in B)}{\rUI{\rRI{3,5}}}}
\end{tabproofwide}

\FormulaThmDeltaK[Profilinklusion reflektiert die Halbverbandsordnung]{%
  \begin{aligned}[t]
    &\Finite A\dsep\ZeroJoinSemi{\JoinOp,A,0}\dsep x,y\in A\dsep{}\\
    &\IrrBelow{\JoinOp,A,0}{x}\subseteq\IrrBelow{\JoinOp,A,0}{y}
      \vdash (x,y)\in\JoinOrd{\JoinOp}{A}
  \end{aligned}
}{IrreducibleProfileReflectsOrder}{%
  \DeltaRow{Mengen}{A\dsep j\dsep x\dsep y\dsep 0}
  \DeltaRow{Funktionen}{\JoinOp}
}
\begin{tabproofwide}
  \proofstepwidestar[1]{\Finite A}{\rA}
  \proofstepwidestar[2]{\ZeroJoinSemi{\JoinOp,A,0}}{\rA}
  \proofstepwidestar[3]{x\in A}{\rA}
  \proofstepwidestar[4]{y\in A}{\rA}
  \proofstepwidestar[5]{%
    \IrrBelow{\JoinOp,A,0}{x}\subseteq\IrrBelow{\JoinOp,A,0}{y}}{\rA}
  \proofstepwidestar[6]{(x,y)\notin\JoinOrd{\JoinOp}{A}}{\rA}
  \proofstepwidestar[1,2,3,4,6]{%
    \exists j\in\JoinIrrSet{\JoinOp,A,0}\,\bigl(
      (j,x)\in\JoinOrd{\JoinOp}{A}\land
      (j,y)\notin\JoinOrd{\JoinOp}{A}\bigr)}{%
    \FormulaRefAuto{JoinIrreduciblesSeparate}[thm]{1,2,3,4,6}}
  \proofstepwidestar[8]{%
    j\in\JoinIrrSet{\JoinOp,A,0}\land\bigl(
      (j,x)\in\JoinOrd{\JoinOp}{A}\land
      (j,y)\notin\JoinOrd{\JoinOp}{A}\bigr)}{\rA}
  \proofstepwidestar[8]{j\in\JoinIrrSet{\JoinOp,A,0}}{\rAEa{8}}
  \proofstepwidestar[8]{(j,x)\in\JoinOrd{\JoinOp}{A}}{\rAEa{\rAEb{8}}}
  \proofstepwidestar[8]{(j,y)\notin\JoinOrd{\JoinOp}{A}}{\rAEb{\rAEb{8}}}
  \proofstepwidestar[8]{j\in\IrrBelow{\JoinOp,A,0}{x}}{%
    \FormulaRefAuto{IrreduciblesBelow}[def]{\rAI{9,10}}}
  \proofstepwidestar[5,8]{j\in\IrrBelow{\JoinOp,A,0}{y}}{%
    \FormulaRefAuto{A\subseteq B,\,x\in A\vdash x\in B}{5,12}}
  \proofstepwidestar[5,8]{(j,y)\in\JoinOrd{\JoinOp}{A}}{%
    \rAEb{\FormulaRefAuto{IrreduciblesBelow}[def]{13}}}
  \proofstepwidestar[5,8]{\bot}{\rBI{11,14}}
  \proofstepwidestar[1,2,3,4,5,6]{\bot}{\rEE{7,8,15}}
  \proofstepwidestar[1,2,3,4,5]{(x,y)\in\JoinOrd{\JoinOp}{A}}{%
    \rCE{6,16}}
\end{tabproofwide}

\FormulaThmDeltaK[Die kanonische Profilabbildung ist injektiv]{%
  \Finite A\dsep\ZeroJoinSemi{\JoinOp,A,0}\vdash
  \CanonicalProfileMap{\JoinOp,A,0}
    \colon A\inj\mathcal U_{\JoinOp,A,0}
}{CanonicalProfileInjective}{%
  \DeltaRow{Mengen}{A\dsep x\dsep y\dsep 0}
  \DeltaRow{Funktionen}{\JoinOp\dsep\CanonicalProfileMap{\JoinOp,A,0}}
}
\begin{tabproofwide}
  \proofstepwidestar[1]{\Finite A}{\rA}
  \proofstepwidestar[2]{\ZeroJoinSemi{\JoinOp,A,0}}{\rA}
  \proofstepwidestar[2]{%
    \CanonicalProfileMap{\JoinOp,A,0}
      \colon A\to\mathcal U_{\JoinOp,A,0}}{%
    \FormulaRefAuto{CanonicalProfileFunction}[thm]{2}}
  \proofstepwidestar[4]{x\in A}{\rA}
  \proofstepwidestar[5]{y\in A}{\rA}
  \proofstepwidestar[6]{%
    \CanonicalProfileMap{\JoinOp,A,0}(x)
      =\CanonicalProfileMap{\JoinOp,A,0}(y)}{\rA}
  \proofstepwidestar[2,4]{%
    \CanonicalProfileMap{\JoinOp,A,0}(x)=
      \JoinIrrSet{\JoinOp,A,0}\setminus\IrrBelow{\JoinOp,A,0}{x}}{%
    \FormulaRefAuto{CanonicalProfileEvaluation}[thm]{2,4}}
  \proofstepwidestar[2,5]{%
    \CanonicalProfileMap{\JoinOp,A,0}(y)=
      \JoinIrrSet{\JoinOp,A,0}\setminus\IrrBelow{\JoinOp,A,0}{y}}{%
    \FormulaRefAuto{CanonicalProfileEvaluation}[thm]{2,5}}
  \proofstepwidestar[2,4,5,6]{%
    \JoinIrrSet{\JoinOp,A,0}\setminus\IrrBelow{\JoinOp,A,0}{x}
      =\CanonicalProfileMap{\JoinOp,A,0}(y)}{%
    \FormulaRefAuto{a=b,\,a=c\vdash b=c}{7,6}}
  \proofstepwidestar[2,4,5,6]{%
    \JoinIrrSet{\JoinOp,A,0}\setminus\IrrBelow{\JoinOp,A,0}{x}
      =\JoinIrrSet{\JoinOp,A,0}\setminus\IrrBelow{\JoinOp,A,0}{y}}{%
    \FormulaRefAuto{a=b,\,b=c\vdash a=c}{9,8}}
  \proofstepwidestar[2,4]{%
    \IrrBelow{\JoinOp,A,0}{x}\subseteq\JoinIrrSet{\JoinOp,A,0}}{%
    \FormulaRefAuto{IrreducibleProfileSubset}[thm]{2,4}}
  \proofstepwidestar[2,5]{%
    \IrrBelow{\JoinOp,A,0}{y}\subseteq\JoinIrrSet{\JoinOp,A,0}}{%
    \FormulaRefAuto{IrreducibleProfileSubset}[thm]{2,5}}
  \proofstepwidestar[2,4,5,6]{%
    \IrrBelow{\JoinOp,A,0}{x}=\IrrBelow{\JoinOp,A,0}{y}}{%
    \FormulaRefAuto{RelCompEqualityCancel}[thm]{12,11,10}}
  \proofstepwidestar[2,4,5,6]{%
    \IrrBelow{\JoinOp,A,0}{x}\subseteq\IrrBelow{\JoinOp,A,0}{y}}{%
    \FormulaRefAuto{A=B\vdash A\subseteq B}{13}}
  \proofstepwidestar[1,2,4,5,6]{(x,y)\in\JoinOrd{\JoinOp}{A}}{%
    \FormulaRefAuto{IrreducibleProfileReflectsOrder}[thm]{1,2,4,5,14}}
  \proofstepwidestar[2,4,5,6]{%
    \IrrBelow{\JoinOp,A,0}{y}\subseteq\IrrBelow{\JoinOp,A,0}{x}}{%
    \FormulaRefAuto{A=B\vdash B\subseteq A}{13}}
  \proofstepwidestar[1,2,4,5,6]{(y,x)\in\JoinOrd{\JoinOp}{A}}{%
    \FormulaRefAuto{IrreducibleProfileReflectsOrder}[thm]{1,2,5,4,16}}
  \proofstepwidestar[1,2,4,5,6]{x=y}{%
    \FormulaRefAuto{JoinOrderAntisymmetric}[thm]{%
      \FormulaRefAuto{ZeroJoinBaseSemilatticeAxiom}[ax]{2},15,17}}
  \proofstepwidestar[1,2]{%
    \CanonicalProfileMap{\JoinOp,A,0}
      \colon A\inj\mathcal U_{\JoinOp,A,0}}{%
    \FormulaRefAuto{Injektive Funktion}[def]{%
      3,\rUI{\rRI{4,\rRI{5,\rRI{6,18}}}}}}
\end{tabproofwide}

\FormulaThmDeltaK[Die kanonische Profilabbildung ist bijektiv]{%
  \Finite A\dsep\ZeroJoinSemi{\JoinOp,A,0}\vdash
  \CanonicalProfileMap{\JoinOp,A,0}
    \colon A\bij\mathcal U_{\JoinOp,A,0}
}{CanonicalProfileBijection}{%
  \DeltaRow{Mengen}{A\dsep 0}
  \DeltaRow{Funktionen}{\JoinOp\dsep\CanonicalProfileMap{\JoinOp,A,0}}
}
\begin{tabproofwide}
  \proofstepwidestar[1]{\Finite A}{\rA}
  \proofstepwidestar[2]{\ZeroJoinSemi{\JoinOp,A,0}}{\rA}
  \proofstepwidestar[1,2]{%
    \CanonicalProfileMap{\JoinOp,A,0}
      \colon A\inj\mathcal U_{\JoinOp,A,0}}{%
    \FormulaRefAuto{CanonicalProfileInjective}[thm]{1,2}}
  \proofstepwidestar[2]{%
    \CanonicalProfileMap{\JoinOp,A,0}
      \colon A\sur\mathcal U_{\JoinOp,A,0}}{%
    \FormulaRefAuto{CanonicalProfileSurjective}[thm]{2}}
  \proofstepwidestar[1,2]{%
    \CanonicalProfileMap{\JoinOp,A,0}
      \colon A\bij\mathcal U_{\JoinOp,A,0}}{%
    \FormulaRefAuto{F\colon A\inj B\dsep F\colon A\sur B
      \vdash F\colon A\bij B}{3,4}}
\end{tabproofwide}

\FormulaThmDeltaK[Repräsentant eines kanonischen Profils]{%
  \begin{aligned}[t]
    &\ZeroJoinSemi{\JoinOp,A,0}\dsep
      X\in\mathcal U_{\JoinOp,A,0}\vdash{}\\
    &\exists x\in A\,\bigl(
      X=\CanonicalProfileMap{\JoinOp,A,0}(x)\bigr)
  \end{aligned}
}{CanonicalProfileRepresentative}{%
  \DeltaRow{Mengen}{A\dsep X\dsep x\dsep 0}
  \DeltaRow{Funktionen}{\JoinOp\dsep\CanonicalProfileMap{\JoinOp,A,0}}
}
\begin{tabproofwide}
  \proofstepwidestar[1]{\ZeroJoinSemi{\JoinOp,A,0}}{\rA}
  \proofstepwidestar[2]{X\in\mathcal U_{\JoinOp,A,0}}{\rA}
  \proofstepwidestar[1]{%
    \CanonicalProfileMap{\JoinOp,A,0}
      \colon A\sur\mathcal U_{\JoinOp,A,0}}{%
    \FormulaRefAuto{CanonicalProfileSurjective}[thm]{1}}
  \proofstepwidestar[1,2]{%
    \exists x\in A\,\bigl(
      \CanonicalProfileMap{\JoinOp,A,0}(x)=X\bigr)}{%
    \FormulaRefAuto{Surjektivität}{3,2}}
  \proofstepwidestar[5]{%
    x\in A\land\CanonicalProfileMap{\JoinOp,A,0}(x)=X}{\rA}
  \proofstepwidestar[5]{x\in A}{\rAEa{5}}
  \proofstepwidestar[5]{%
    X=\CanonicalProfileMap{\JoinOp,A,0}(x)}{%
    \FormulaRefAuto{a=b\vdash b=a}{\rAEb{5}}}
  \proofstepwidestar[5]{%
    \exists x\in A\,\bigl(
      X=\CanonicalProfileMap{\JoinOp,A,0}(x)\bigr)}{%
    \rEI{\rAI{6,7}}}
  \proofstepwidestar[1,2]{%
    \exists x\in A\,\bigl(
      X=\CanonicalProfileMap{\JoinOp,A,0}(x)\bigr)}{%
    \rEE{4,5,8}}
\end{tabproofwide}

\FormulaThmDeltaK[Relatives de-Morgan-Gesetz]{%
  A\setminus(B\cap C)
  =
  (A\setminus B)\cup(A\setminus C)
}{RelativeComplementIntersectionDeMorgan}{%
  \DeltaRow{Mengen}{A\dsep B\dsep C\dsep x}
}
\begin{tabproofwide}
  \proofstepwide{x\in A\setminus(B\cap C)}{\leftrightarrow}{%
    x\in A\land x\notin B\cap C}{%
    \FormulaRefAuto{x \in A \setminus B \eqvdash
      x \in A \land x \notin B}}
  \proofstepwide{} {\leftrightarrow}{%
    x\in A\land\neg(x\in B\land x\in C)}{%
    \rLRS{\FormulaRefAuto{P \leftrightarrow Q \eqvdash
      \neg P \leftrightarrow \neg Q}{%
        \FormulaRefAuto{x \in A \cap B \eqvdash
          x \in A \land x \in B}},2}}
  \proofstepwide{} {\leftrightarrow}{%
    x\in A\land(x\notin B\lor x\notin C)}{%
    \rLRS{\FormulaRefAuto{\neg(P \land Q) \eqvdash
      \neg P \lor \neg Q},2}}
  \proofstepwide{} {\leftrightarrow}{%
    (x\in A\land x\notin B)\lor(x\in A\land x\notin C)}{%
    \FormulaRefAuto{P \land (Q\lor R) \eqvdash
      (P \land Q) \lor (P\land R)}}
  \proofstepwide{} {\leftrightarrow}{%
    x\in A\setminus B\lor(x\in A\land x\notin C)}{%
    \rLRS{\FormulaRefAuto{x \in A \setminus B \eqvdash
      x \in A \land x \notin B},1}}
  \proofstepwide{} {\leftrightarrow}{%
    x\in A\setminus B\lor x\in A\setminus C}{%
    \rLRS{\FormulaRefAuto{x \in A \setminus B \eqvdash
      x \in A \land x \notin B},1}}
  \proofstepwide{} {\leftrightarrow}{%
    x\in(A\setminus B)\cup(A\setminus C)}{%
    \FormulaRefAuto{x \in A \cup B \eqvdash
      x \in A \lor x \in B}}
  \proofstepwide{A\setminus(B\cap C)}{=}{%
    (A\setminus B)\cup(A\setminus C)}{%
    \FormulaRefAuto{\forall x\, (x \in A \leftrightarrow x \in B)
      \eqvdash A = B}{\rUI{\rChain{1,2,3,4,5,6,7}}}}
\end{tabproofwide}

\FormulaThmDeltaK[Komplementprofile eines Schnitts vereinigen sich]{%
  \begin{aligned}[t]
    &D=B\cap C\dsep X=U\setminus D\dsep
      Y=U\setminus B\dsep Z=U\setminus C\vdash{}\\
    &X=Y\cup Z
  \end{aligned}
}{RelativeComplementProfileUnionTransport}{%
  \DeltaRow{Mengen}{B\dsep C\dsep D\dsep U\dsep X\dsep Y\dsep Z}
}
\begin{tabproofwide}
  \proofstepwidestar[1]{D=B\cap C}{\rA}
  \proofstepwidestar[2]{X=U\setminus D}{\rA}
  \proofstepwidestar[3]{Y=U\setminus B}{\rA}
  \proofstepwidestar[4]{Z=U\setminus C}{\rA}
  \proofstepwidestar[]{%
    U\setminus(B\cap C)
      =(U\setminus B)\cup(U\setminus C)}{%
    \FormulaRefAuto{RelativeComplementIntersectionDeMorgan}[thm]}
  \proofstepwidestar[1]{%
    U\setminus D=U\setminus(B\cap C)}{\rIE{1,\rII}}
  \proofstepwidestar[1,2]{%
    X=U\setminus(B\cap C)}{%
    \FormulaRefAuto{a=b,\,b=c\vdash a=c}{2,6}}
  \proofstepwidestar[1,2]{%
    X=(U\setminus B)\cup(U\setminus C)}{%
    \FormulaRefAuto{a=b,\,b=c\vdash a=c}{7,5}}
  \proofstepwidestar[3,4]{%
    Y\cup Z=(U\setminus B)\cup(U\setminus C)}{%
    \FormulaRefAuto{a=b\dsep c=d \dsep P(b,d)\vdash P(a,c)}{%
      3,4,\rII}}
  \proofstepwidestar[1,2,3,4]{X=Y\cup Z}{%
    \FormulaRefAuto{a=b,\,c=b\vdash a=c}{8,9}}
\end{tabproofwide}

\FormulaThmDeltaK[Paarinfimum und Schnitt der Irreduziblenprofile]{%
  \begin{aligned}[t]
    &\ZeroJoinSemi{\JoinOp,A,0}\dsep
      \PairInf{\JoinOrd{\JoinOp}{A}}{A}{x}{y}{m}\vdash{}\\
    &\IrrBelow{\JoinOp,A,0}{m}
      =\IrrBelow{\JoinOp,A,0}{x}\cap
       \IrrBelow{\JoinOp,A,0}{y}
  \end{aligned}
}{CanonicalProfilePairInfimumIntersection}{%
  \DeltaRow{Mengen}{A\dsep j\dsep m\dsep x\dsep y\dsep 0}
  \DeltaRow{Funktionen}{\JoinOp}
}
\begin{tabproofsplitwide}
  \proofpartwideindR[Erste Teilmengenrichtung]{%
    \begin{aligned}[t]
      &\ZeroJoinSemi{\JoinOp,A,0}\dsep
        \PairInf{\JoinOrd{\JoinOp}{A}}{A}{x}{y}{m}\vdash{}\\[-2pt]
      &\IrrBelow{\JoinOp,A,0}{m}\subseteq
        \IrrBelow{\JoinOp,A,0}{x}\cap
        \IrrBelow{\JoinOp,A,0}{y}
    \end{aligned}
  }{\ZeroJoinSemi{\JoinOp,A,0}\dsep
    \PairInf{\JoinOrd{\JoinOp}{A}}{A}{x}{y}{m}\vdash
    \IrrBelow{\JoinOp,A,0}{m}\subseteq
    \IrrBelow{\JoinOp,A,0}{x}\cap
    \IrrBelow{\JoinOp,A,0}{y}}[CanonicalProfilePairInfimumLowerInclusion]
    \proofstepwidestar[1]{\ZeroJoinSemi{\JoinOp,A,0}}{\rA}
    \proofstepwidestar[2]{%
      \PairInf{\JoinOrd{\JoinOp}{A}}{A}{x}{y}{m}}{\rA}
    \proofstepwidestar[1]{\Semi{\JoinOp,A}}{%
      \FormulaRefAuto{ZeroJoinBaseSemilatticeAxiom}[ax]{1}}
    \proofstepwidestar[2]{(m,x)\in\JoinOrd{\JoinOp}{A}}{%
      \FormulaRefAuto{PairInfimum}[def]{2}}
    \proofstepwidestar[2]{(m,y)\in\JoinOrd{\JoinOp}{A}}{%
      \FormulaRefAuto{PairInfimum}[def]{2}}
    \proofstepwidestar[6]{j\in\IrrBelow{\JoinOp,A,0}{m}}{\rA}
    \proofstepwidestar[6]{%
      j\in\JoinIrrSet{\JoinOp,A,0}\land
      (j,m)\in\JoinOrd{\JoinOp}{A}}{%
      \FormulaRefAuto{IrreduciblesBelow}[def]{6}}
    \proofstepwidestar[6]{j\in\JoinIrrSet{\JoinOp,A,0}}{\rAEa{7}}
    \proofstepwidestar[6]{(j,m)\in\JoinOrd{\JoinOp}{A}}{\rAEb{7}}
    \proofstepwidestar[1,2,6]{(j,x)\in\JoinOrd{\JoinOp}{A}}{%
      \FormulaRefAuto{JoinOrderTransitive}[thm]{3,9,4}}
    \proofstepwidestar[1,2,6]{(j,y)\in\JoinOrd{\JoinOp}{A}}{%
      \FormulaRefAuto{JoinOrderTransitive}[thm]{3,9,5}}
    \proofstepwidestar[1,2,6]{j\in\IrrBelow{\JoinOp,A,0}{x}}{%
      \FormulaRefAuto{IrreduciblesBelow}[def]{\rAI{8,10}}}
    \proofstepwidestar[1,2,6]{j\in\IrrBelow{\JoinOp,A,0}{y}}{%
      \FormulaRefAuto{IrreduciblesBelow}[def]{\rAI{8,11}}}
    \proofstepwidestar[1,2,6]{%
      j\in\IrrBelow{\JoinOp,A,0}{x}\cap
        \IrrBelow{\JoinOp,A,0}{y}}{%
      \FormulaRefAuto{x \in A \cap B \eqvdash
        x \in A \land x \in B}{\rAI{12,13}}}
    \proofstepwidestar[1,2]{%
      \begin{aligned}
        \IrrBelow{\JoinOp,A,0}{m}&\subseteq
          \IrrBelow{\JoinOp,A,0}{x}\\[-2pt]
        &\quad\cap\IrrBelow{\JoinOp,A,0}{y}
      \end{aligned}}{%
      \FormulaRefAuto{A \subseteq B \coloneq
        \forall x\,(x\in A\rightarrow x\in B)}{\rUI{\rRI{6,14}}}}
  \closeproofpartwideindR

  \proofpartwideindR[Zweite Teilmengenrichtung]{%
    \begin{aligned}[t]
      &\ZeroJoinSemi{\JoinOp,A,0}\dsep
        \PairInf{\JoinOrd{\JoinOp}{A}}{A}{x}{y}{m}\vdash{}\\[-2pt]
      &\IrrBelow{\JoinOp,A,0}{x}\cap
        \IrrBelow{\JoinOp,A,0}{y}\subseteq
        \IrrBelow{\JoinOp,A,0}{m}
    \end{aligned}
  }{\ZeroJoinSemi{\JoinOp,A,0}\dsep
    \PairInf{\JoinOrd{\JoinOp}{A}}{A}{x}{y}{m}\vdash
    \IrrBelow{\JoinOp,A,0}{x}\cap
    \IrrBelow{\JoinOp,A,0}{y}\subseteq
    \IrrBelow{\JoinOp,A,0}{m}}[CanonicalProfilePairInfimumUpperInclusion]
    \proofstepwidestar[1]{\ZeroJoinSemi{\JoinOp,A,0}}{\rA}
    \proofstepwidestar[2]{%
      \PairInf{\JoinOrd{\JoinOp}{A}}{A}{x}{y}{m}}{\rA}
    \proofstepwidestar[3]{%
      j\in\IrrBelow{\JoinOp,A,0}{x}\cap
        \IrrBelow{\JoinOp,A,0}{y}}{\rA}
    \proofstepwidestar[3]{j\in\IrrBelow{\JoinOp,A,0}{x}}{%
      \FormulaRefAuto{x \in A \cap B \vdash x \in A}{3}}
    \proofstepwidestar[3]{j\in\IrrBelow{\JoinOp,A,0}{y}}{%
      \FormulaRefAuto{x \in A \cap B \vdash x \in B}{3}}
    \proofstepwidestar[3]{%
      j\in\JoinIrrSet{\JoinOp,A,0}\land
      (j,x)\in\JoinOrd{\JoinOp}{A}}{%
      \FormulaRefAuto{IrreduciblesBelow}[def]{4}}
    \proofstepwidestar[3]{%
      j\in\JoinIrrSet{\JoinOp,A,0}\land
      (j,y)\in\JoinOrd{\JoinOp}{A}}{%
      \FormulaRefAuto{IrreduciblesBelow}[def]{5}}
    \proofstepwidestar[3]{j\in\JoinIrrSet{\JoinOp,A,0}}{\rAEa{6}}
    \proofstepwidestar[3]{(j,x)\in\JoinOrd{\JoinOp}{A}}{\rAEb{6}}
    \proofstepwidestar[3]{(j,y)\in\JoinOrd{\JoinOp}{A}}{\rAEb{7}}
    \proofstepwidestar[3]{\JoinIrr{\JoinOp,A,0}{j}}{%
      \FormulaRefAuto{JoinIrreducibleSet}[def]{8}}
    \proofstepwidestar[3]{j\in A}{%
      \FormulaRefAuto{JoinIrreducible}[def]{11}}
    \proofstepwidestar[2,3]{(j,m)\in\JoinOrd{\JoinOp}{A}}{%
      \FormulaRefAuto{PairInfimum}[def]{2,12,9,10}}
    \proofstepwidestar[2,3]{j\in\IrrBelow{\JoinOp,A,0}{m}}{%
      \FormulaRefAuto{IrreduciblesBelow}[def]{\rAI{8,13}}}
    \proofstepwidestar[2]{%
      \begin{aligned}
        \IrrBelow{\JoinOp,A,0}{x}&\cap
          \IrrBelow{\JoinOp,A,0}{y}\\[-2pt]
        &\quad\subseteq\IrrBelow{\JoinOp,A,0}{m}
      \end{aligned}}{%
      \FormulaRefAuto{A \subseteq B \coloneq
        \forall x\,(x\in A\rightarrow x\in B)}{\rUI{\rRI{3,14}}}}
  \closeproofpartwideindR

  \proofpartwide{Schluss}
    \proofstepwidestar[1]{\ZeroJoinSemi{\JoinOp,A,0}}{\rA}
    \proofstepwidestar[2]{%
      \PairInf{\JoinOrd{\JoinOp}{A}}{A}{x}{y}{m}}{\rA}
    \proofstepwidestar[1,2]{%
      \IrrBelow{\JoinOp,A,0}{m}\subseteq
        \IrrBelow{\JoinOp,A,0}{x}\cap
        \IrrBelow{\JoinOp,A,0}{y}}{%
      \FormulaRefAuto{CanonicalProfilePairInfimumLowerInclusion}[thm]{1,2}}
    \proofstepwidestar[1,2]{%
      \IrrBelow{\JoinOp,A,0}{x}\cap
        \IrrBelow{\JoinOp,A,0}{y}\subseteq
        \IrrBelow{\JoinOp,A,0}{m}}{%
      \FormulaRefAuto{CanonicalProfilePairInfimumUpperInclusion}[thm]{1,2}}
    \proofstepwidestar[1,2]{%
      \begin{aligned}
        \IrrBelow{\JoinOp,A,0}{m}
          &=\IrrBelow{\JoinOp,A,0}{x}\\[-2pt]
          &\quad\cap\IrrBelow{\JoinOp,A,0}{y}
      \end{aligned}}{%
      \FormulaRefAuto{A \subseteq B, B \subseteq A \vdash A = B}{3,4}}
  \closeproofpartwide
\end{tabproofsplitwide}

\FormulaThmDeltaK[Profil eines Paarinfimums als Vereinigung]{%
  \begin{aligned}[t]
    &\ZeroJoinSemi{\JoinOp,A,0}\dsep
      \PairInf{\JoinOrd{\JoinOp}{A}}{A}{x}{y}{m}\vdash{}\\
    &\CanonicalProfileMap{\JoinOp,A,0}(m)
      =\CanonicalProfileMap{\JoinOp,A,0}(x)
       \cup\CanonicalProfileMap{\JoinOp,A,0}(y)
  \end{aligned}
}{CanonicalProfilePairInfimumUnion}{%
  \DeltaRow{Mengen}{A\dsep m\dsep x\dsep y\dsep 0}
  \DeltaRow{Funktionen}{\JoinOp\dsep\CanonicalProfileMap{\JoinOp,A,0}}
}
\begin{tabproofwide}
  \proofstepwidestar[1]{\ZeroJoinSemi{\JoinOp,A,0}}{\rA}
  \proofstepwidestar[2]{%
    \PairInf{\JoinOrd{\JoinOp}{A}}{A}{x}{y}{m}}{\rA}
  \proofstepwidestar[2]{x\in A}{\FormulaRefAuto{PairInfimum}[def]{2}}
  \proofstepwidestar[2]{y\in A}{\FormulaRefAuto{PairInfimum}[def]{2}}
  \proofstepwidestar[2]{m\in A}{\FormulaRefAuto{PairInfimum}[def]{2}}
  \proofstepwidestar[1,2]{%
    \IrrBelow{\JoinOp,A,0}{m}
      =\IrrBelow{\JoinOp,A,0}{x}\cap\IrrBelow{\JoinOp,A,0}{y}}{%
    \FormulaRefAuto{CanonicalProfilePairInfimumIntersection}[thm]{1,2}}
  \proofstepwidestar[1,2]{%
    \CanonicalProfileMap{\JoinOp,A,0}(m)=
    \JoinIrrSet{\JoinOp,A,0}\setminus\IrrBelow{\JoinOp,A,0}{m}}{%
    \FormulaRefAuto{CanonicalProfileEvaluation}[thm]{1,5}}
  \proofstepwidestar[1,2]{%
    \CanonicalProfileMap{\JoinOp,A,0}(x)=
    \JoinIrrSet{\JoinOp,A,0}\setminus\IrrBelow{\JoinOp,A,0}{x}}{%
    \FormulaRefAuto{CanonicalProfileEvaluation}[thm]{1,3}}
  \proofstepwidestar[1,2]{%
    \CanonicalProfileMap{\JoinOp,A,0}(y)=
    \JoinIrrSet{\JoinOp,A,0}\setminus\IrrBelow{\JoinOp,A,0}{y}}{%
    \FormulaRefAuto{CanonicalProfileEvaluation}[thm]{1,4}}
  \proofstepwidestar[1,2]{%
    \begin{aligned}
      \CanonicalProfileMap{\JoinOp,A,0}(m)
        &=\CanonicalProfileMap{\JoinOp,A,0}(x)\\[-2pt]
        &\quad\cup\CanonicalProfileMap{\JoinOp,A,0}(y)
    \end{aligned}}{%
    \FormulaRefAuto{RelativeComplementProfileUnionTransport}[thm]{%
      6,7,8,9}}
\end{tabproofwide}

\FormulaThmDeltaK[Abgeschlossenheit der kanonischen Profilfamilie]{%
  \Finite A\dsep\ZeroJoinSemi{\JoinOp,A,0}
  \vdash\UnionClosedFam(\mathcal U_{\JoinOp,A,0})
}{CanonicalProfileFamilyUnionClosed}{%
  \DeltaRow{Mengen}{A\dsep X\dsep Y\dsep m\dsep x\dsep y\dsep 0}
  \DeltaRow{Funktionen}{\JoinOp\dsep\CanonicalProfileMap{\JoinOp,A,0}}
}
\begin{tabproofsplitwide}
  \proofpartwideindR[Vereinigung zweier kanonischer Profile]{%
    \begin{aligned}[t]
      &\Finite A\dsep\ZeroJoinSemi{\JoinOp,A,0}\dsep
        X,Y\in\mathcal U_{\JoinOp,A,0}\vdash{}\\[-2pt]
      &X\cup Y\in\mathcal U_{\JoinOp,A,0}
    \end{aligned}
  }{\Finite A\dsep\ZeroJoinSemi{\JoinOp,A,0}\dsep
    X,Y\in\mathcal U_{\JoinOp,A,0}\vdash
    X\cup Y\in\mathcal U_{\JoinOp,A,0}}[CanonicalProfileUnionWitness]
    \proofstepwidestar[1]{\Finite A}{\rA}
    \proofstepwidestar[2]{\ZeroJoinSemi{\JoinOp,A,0}}{\rA}
    \proofstepwidestar[3]{X\in\mathcal U_{\JoinOp,A,0}}{\rA}
    \proofstepwidestar[4]{Y\in\mathcal U_{\JoinOp,A,0}}{\rA}
    \proofstepwidestar[2,3]{%
      \exists x\in A\,\bigl(
        X=\CanonicalProfileMap{\JoinOp,A,0}(x)\bigr)}{%
      \FormulaRefAuto{CanonicalProfileRepresentative}[thm]{2,3}}
    \proofstepwidestar[5]{%
      x\in A\land X=\CanonicalProfileMap{\JoinOp,A,0}(x)}{\rA}
    \proofstepwidestar[2,4]{%
      \exists y\in A\,\bigl(
        Y=\CanonicalProfileMap{\JoinOp,A,0}(y)\bigr)}{%
      \FormulaRefAuto{CanonicalProfileRepresentative}[thm]{2,4}}
    \proofstepwidestar[7]{%
      y\in A\land Y=\CanonicalProfileMap{\JoinOp,A,0}(y)}{\rA}
    \proofstepwidestar[1,2,5,7]{%
      \exists m\in A\,
        \PairInf{\JoinOrd{\JoinOp}{A}}{A}{x}{y}{m}}{%
      \FormulaRefAuto{FiniteZeroJoinSemilatticePairInfima}[thm]{%
        1,2,\rAEa{6},\rAEa{8}}}
    \proofstepwidestar[9]{%
      m\in A\land
        \PairInf{\JoinOrd{\JoinOp}{A}}{A}{x}{y}{m}}{\rA}
    \proofstepwidestar[2,9]{%
      \begin{aligned}
        \CanonicalProfileMap{\JoinOp,A,0}(m)
          &=\CanonicalProfileMap{\JoinOp,A,0}(x)\\[-2pt]
          &\quad\cup\CanonicalProfileMap{\JoinOp,A,0}(y)
      \end{aligned}}{%
      \FormulaRefAuto{CanonicalProfilePairInfimumUnion}[thm]{2,\rAEb{10}}}
    \proofstepwidestar[5,7]{%
      \begin{aligned}
        X\cup Y&=\CanonicalProfileMap{\JoinOp,A,0}(x)\\[-2pt]
          &\quad\cup\CanonicalProfileMap{\JoinOp,A,0}(y)
      \end{aligned}}{%
      \FormulaRefAuto{a=b\dsep c=d \dsep P(b,d)\vdash P(a,c)}{%
        \rAEb{6},\rAEb{8},\rII}}
    \proofstepwidestar[2,5,7,9]{%
      X\cup Y=\CanonicalProfileMap{\JoinOp,A,0}(m)}{%
      \FormulaRefAuto{a=b,\,c=b\vdash a=c}{12,11}}
    \proofstepwidestar[2]{%
      \CanonicalProfileMap{\JoinOp,A,0}
        \colon A\to\mathcal U_{\JoinOp,A,0}}{%
      \FormulaRefAuto{CanonicalProfileFunction}[thm]{2}}
    \proofstepwidestar[2,9]{%
      \CanonicalProfileMap{\JoinOp,A,0}(m)
        \in\mathcal U_{\JoinOp,A,0}}{%
      \FormulaRefAuto{F\colon A\to B\dsep x\in A\vdash F(x)\in B}{%
        14,\rAEa{10}}}
    \proofstepwidestar[2,5,7,9]{%
      X\cup Y\in\mathcal U_{\JoinOp,A,0}}{\rIE{13,15}}
    \proofstepwidestar[1,2,5,7]{%
      X\cup Y\in\mathcal U_{\JoinOp,A,0}}{\rEE{9,10,16}}
    \proofstepwidestar[1,2,5,4]{%
      X\cup Y\in\mathcal U_{\JoinOp,A,0}}{\rEE{7,8,17}}
    \proofstepwidestar[1,2,3,4]{%
      X\cup Y\in\mathcal U_{\JoinOp,A,0}}{\rEE{5,6,18}}
  \closeproofpartwideindR

  \proofpartwide{Schluss}
    \proofstepwidestar[1]{\Finite A}{\rA}
    \proofstepwidestar[2]{\ZeroJoinSemi{\JoinOp,A,0}}{\rA}
    \proofstepwidestar[3]{X\in\mathcal U_{\JoinOp,A,0}}{\rA}
    \proofstepwidestar[4]{Y\in\mathcal U_{\JoinOp,A,0}}{\rA}
    \proofstepwidestar[1,2,3,4]{%
      X\cup Y\in\mathcal U_{\JoinOp,A,0}}{%
      \FormulaRefAuto{CanonicalProfileUnionWitness}[thm]{1,2,3,4}}
    \proofstepwidestar[1,2]{\UnionClosedFam(\mathcal U_{\JoinOp,A,0})}{%
      \FormulaRefAuto{Vereinigungsabgeschlossene Familie}{%
        \rUI{\rRI{3,\rRI{4,5}}}}}
  \closeproofpartwide
\end{tabproofsplitwide}

\FormulaThmDeltaK[Das Nullprofil ist leer]{%
  \ZeroJoinSemi{\JoinOp,A,0}\vdash
  \IrrBelow{\JoinOp,A,0}{0}=\varnothing
}{CanonicalZeroProfileEmpty}{%
  \DeltaRow{Mengen}{A\dsep j\dsep 0}
  \DeltaRow{Funktionen}{\JoinOp}
}
\begin{tabproofwide}
  \proofstepwidestar[1]{\ZeroJoinSemi{\JoinOp,A,0}}{\rA}
  \proofstepwidestar[2]{j\in\IrrBelow{\JoinOp,A,0}{0}}{\rA}
  \proofstepwidestar[2]{%
    j\in\JoinIrrSet{\JoinOp,A,0}\land
      (j,0)\in\JoinOrd{\JoinOp}{A}}{%
    \FormulaRefAuto{IrreduciblesBelow}[def]{2}}
  \proofstepwidestar[2]{j\in\JoinIrrSet{\JoinOp,A,0}}{\rAEa{3}}
  \proofstepwidestar[2]{(j,0)\in\JoinOrd{\JoinOp}{A}}{\rAEb{3}}
  \proofstepwidestar[2]{\JoinIrr{\JoinOp,A,0}{j}}{%
    \FormulaRefAuto{JoinIrreducibleSet}[def]{4}}
  \proofstepwidestar[2]{j\in A}{\FormulaRefAuto{JoinIrreducible}[def]{6}}
  \proofstepwidestar[2]{j\neq0}{\FormulaRefAuto{JoinIrreducible}[def]{6}}
  \proofstepwidestar[1,2]{(0,j)\in\JoinOrd{\JoinOp}{A}}{%
    \FormulaRefAuto{ZeroJoinLeast}[thm]{1,7}}
  \proofstepwidestar[1]{\Semi{\JoinOp,A}}{%
    \FormulaRefAuto{ZeroJoinBaseSemilatticeAxiom}[ax]{1}}
  \proofstepwidestar[1,2]{j=0}{%
    \FormulaRefAuto{JoinOrderAntisymmetric}[thm]{10,5,9}}
  \proofstepwidestar[1,2]{\bot}{\rBI{11,8}}
  \proofstepwidestar[1]{j\notin\IrrBelow{\JoinOp,A,0}{0}}{\rCI{2,12}}
  \proofstepwidestar[1]{%
    \forall j\,j\notin\IrrBelow{\JoinOp,A,0}{0}}{\rUI{13}}
  \proofstepwidestar[1]{\IrrBelow{\JoinOp,A,0}{0}=\varnothing}{%
    \FormulaRefAuto{\forall x\, x\notin A\vdash A=\varnothing}{14}}
\end{tabproofwide}

\FormulaThmDeltaK[Die kanonische Profilfamilie ist nichtleer]{%
  \ZeroJoinSemi{\JoinOp,A,0}\vdash
  \mathcal U_{\JoinOp,A,0}\neq\varnothing
}{CanonicalProfileFamilyNonempty}{%
  \DeltaRow{Mengen}{A\dsep 0}
  \DeltaRow{Funktionen}{\JoinOp\dsep\CanonicalProfileMap{\JoinOp,A,0}}
}
\begin{tabproofwide}
  \proofstepwidestar[1]{\ZeroJoinSemi{\JoinOp,A,0}}{\rA}
  \proofstepwidestar[1]{0\in A}{\FormulaRefAuto{ZeroJoinCarrierAxiom}[ax]{1}}
  \proofstepwidestar[1]{%
    \CanonicalProfileMap{\JoinOp,A,0}\colon
      A\to\mathcal U_{\JoinOp,A,0}}{%
    \FormulaRefAuto{CanonicalProfileFunction}[thm]{1}}
  \proofstepwidestar[1]{%
    \CanonicalProfileMap{\JoinOp,A,0}(0)\in\mathcal U_{\JoinOp,A,0}}{%
    \FormulaRefAuto{F\colon A\to B\dsep x\in A\vdash F(x)\in B}{3,2}}
  \proofstepwidestar[1]{%
    \exists X\,X\in\mathcal U_{\JoinOp,A,0}}{\rEI{4}}
  \proofstepwidestar[1]{\mathcal U_{\JoinOp,A,0}\neq\varnothing}{%
    \FormulaRefAuto{A\neq\varnothing \eqvdash \exists x\,(x \in A)}{5}}
\end{tabproofwide}

\FormulaThmDeltaK[Ein nichttrivialer endlicher Halbverband besitzt Irreduzible]{%
  \begin{aligned}[t]
    &\Finite A\dsep\ZeroJoinSemi{\JoinOp,A,0}\dsep
      \exists x\in A\,(x\neq0)\vdash{}\\
    &\JoinIrrSet{\JoinOp,A,0}\neq\varnothing
  \end{aligned}
}{NontrivialFiniteZeroJoinHasIrreducible}{%
  \DeltaRow{Mengen}{A\dsep j\dsep x\dsep 0}
  \DeltaRow{Funktionen}{\JoinOp}
}
\begin{tabproofwide}
  \proofstepwidestar[1]{\Finite A}{\rA}
  \proofstepwidestar[2]{\ZeroJoinSemi{\JoinOp,A,0}}{\rA}
  \proofstepwidestar[3]{\exists x\in A\,(x\neq0)}{\rA}
  \proofstepwidestar[4]{x\in A\land x\neq0}{\rA}
  \proofstepwidestar[4]{x\in A}{\rAEa{4}}
  \proofstepwidestar[4]{x\neq0}{\rAEb{4}}
  \proofstepwidestar[2]{\Semi{\JoinOp,A}}{%
    \FormulaRefAuto{ZeroJoinBaseSemilatticeAxiom}[ax]{2}}
  \proofstepwidestar[8]{(x,0)\in\JoinOrd{\JoinOp}{A}}{\rA}
  \proofstepwidestar[2,4]{(0,x)\in\JoinOrd{\JoinOp}{A}}{%
    \FormulaRefAuto{ZeroJoinLeast}[thm]{2,5}}
  \proofstepwidestar[2,4,8]{x=0}{%
    \FormulaRefAuto{JoinOrderAntisymmetric}[thm]{7,8,9}}
  \proofstepwidestar[2,4,8]{\bot}{\rBI{10,6}}
  \proofstepwidestar[2,4]{(x,0)\notin\JoinOrd{\JoinOp}{A}}{\rCI{8,11}}
  \proofstepwidestar[1,2,4]{%
    \exists j\in\JoinIrrSet{\JoinOp,A,0}\,\bigl(
      (j,x)\in\JoinOrd{\JoinOp}{A}\land
      (j,0)\notin\JoinOrd{\JoinOp}{A}\bigr)}{%
    \FormulaRefAuto{JoinIrreduciblesSeparate}[thm]{%
      1,2,5,\FormulaRefAuto{ZeroJoinCarrierAxiom}[ax]{2},12}}
  \proofstepwidestar[14]{%
    j\in\JoinIrrSet{\JoinOp,A,0}\land\bigl(
      (j,x)\in\JoinOrd{\JoinOp}{A}\land
      (j,0)\notin\JoinOrd{\JoinOp}{A}\bigr)}{\rA}
  \proofstepwidestar[14]{j\in\JoinIrrSet{\JoinOp,A,0}}{\rAEa{14}}
  \proofstepwidestar[14]{\exists j\,j\in\JoinIrrSet{\JoinOp,A,0}}{\rEI{15}}
  \proofstepwidestar[1,2,4]{%
    \exists j\,j\in\JoinIrrSet{\JoinOp,A,0}}{\rEE{13,14,16}}
  \proofstepwidestar[1,2,4]{\JoinIrrSet{\JoinOp,A,0}\neq\varnothing}{%
    \FormulaRefAuto{A\neq\varnothing \eqvdash \exists x\,(x \in A)}{17}}
  \proofstepwidestar[1,2,3]{\JoinIrrSet{\JoinOp,A,0}\neq\varnothing}{%
    \rEE{3,4,18}}
\end{tabproofwide}

\FormulaThmDeltaK[Der Träger der kanonischen Profilfamilie ist nichtleer]{%
  \begin{aligned}[t]
    &\Finite A\dsep\ZeroJoinSemi{\JoinOp,A,0}\dsep
      \exists x\in A\,(x\neq0)\vdash{}\\
    &\bigcup\mathcal U_{\JoinOp,A,0}\neq\varnothing
  \end{aligned}
}{CanonicalProfileCarrierNonempty}{%
  \DeltaRow{Mengen}{A\dsep j\dsep x\dsep 0}
  \DeltaRow{Funktionen}{\JoinOp\dsep\CanonicalProfileMap{\JoinOp,A,0}}
}
\begin{tabproofwide}
  \proofstepwidestar[1]{\Finite A}{\rA}
  \proofstepwidestar[2]{\ZeroJoinSemi{\JoinOp,A,0}}{\rA}
  \proofstepwidestar[3]{\exists x\in A\,(x\neq0)}{\rA}
  \proofstepwidestar[1,2,3]{\JoinIrrSet{\JoinOp,A,0}\neq\varnothing}{%
    \FormulaRefAuto{NontrivialFiniteZeroJoinHasIrreducible}[thm]{1,2,3}}
  \proofstepwidestar[1,2,3]{%
    \exists j\,j\in\JoinIrrSet{\JoinOp,A,0}}{%
    \FormulaRefAuto{A\neq\varnothing \eqvdash \exists x\,(x \in A)}{4}}
  \proofstepwidestar[6]{j\in\JoinIrrSet{\JoinOp,A,0}}{\rA}
  \proofstepwidestar[2]{0\in A}{\FormulaRefAuto{ZeroJoinCarrierAxiom}[ax]{2}}
  \proofstepwidestar[2]{%
    \CanonicalProfileMap{\JoinOp,A,0}(0)=
      \JoinIrrSet{\JoinOp,A,0}\setminus\IrrBelow{\JoinOp,A,0}{0}}{%
    \FormulaRefAuto{CanonicalProfileEvaluation}[thm]{2,7}}
  \proofstepwidestar[2]{\IrrBelow{\JoinOp,A,0}{0}=\varnothing}{%
    \FormulaRefAuto{CanonicalZeroProfileEmpty}[thm]{2}}
  \proofstepwidestar[]{j\notin\varnothing}{\FormulaRefAuto{x \not\in \varnothing}}
  \proofstepwidestar[2,6]{j\notin\IrrBelow{\JoinOp,A,0}{0}}{%
    \rIE{\FormulaRefAuto{a=b\vdash b=a}{9},10}}
  \proofstepwidestar[2,6]{%
    j\in\JoinIrrSet{\JoinOp,A,0}\setminus\IrrBelow{\JoinOp,A,0}{0}}{%
    \FormulaRefAuto{x \in A\dsep x\notin B\vdash x \in A\setminus B}{6,11}}
  \proofstepwidestar[2,6]{%
    j\in\CanonicalProfileMap{\JoinOp,A,0}(0)}{%
    \rIE{\FormulaRefAuto{a=b\vdash b=a}{8},12}}
  \proofstepwidestar[2]{%
    \CanonicalProfileMap{\JoinOp,A,0}(0)\in\mathcal U_{\JoinOp,A,0}}{%
    \FormulaRefAuto{F\colon A\to B\dsep x\in A\vdash F(x)\in B}{%
      \FormulaRefAuto{CanonicalProfileFunction}[thm]{2},7}}
  \proofstepwidestar[2,6]{j\in\bigcup\mathcal U_{\JoinOp,A,0}}{%
    \FormulaRefAuto{B\in A\dsep x\in B\vdash x \in \bigcup A}{14,13}}
  \proofstepwidestar[2,6]{%
    \exists j\,j\in\bigcup\mathcal U_{\JoinOp,A,0}}{\rEI{15}}
  \proofstepwidestar[1,2,3]{%
    \exists j\,j\in\bigcup\mathcal U_{\JoinOp,A,0}}{\rEE{5,6,16}}
  \proofstepwidestar[1,2,3]{\bigcup\mathcal U_{\JoinOp,A,0}\neq\varnothing}{%
    \FormulaRefAuto{A\neq\varnothing \eqvdash \exists x\,(x \in A)}{17}}
\end{tabproofwide}

\FormulaThmDeltaK[Kanonische Familie ist eine Frankl-Familie]{%
  \begin{aligned}[t]
    &\Finite A\dsep\ZeroJoinSemi{\JoinOp,A,0}\dsep
      \exists x\in A\,(x\neq0)\vdash{}\\
    &\FranklFam{\mathcal U_{\JoinOp,A,0}}
  \end{aligned}
}{CanonicalUnionFamilyIsFranklFamily}{%
  \DeltaRow{Mengen}{A\dsep X\dsep Y\dsep x\dsep y\dsep m\dsep 0}
  \DeltaRow{Funktionen}{\JoinOp\dsep\CanonicalProfileMap{\JoinOp,A,0}}
}
\begin{tabproofwide}
  \proofstepwidestar[1]{\Finite A}{\rA}
  \proofstepwidestar[2]{\ZeroJoinSemi{\JoinOp,A,0}}{\rA}
  \proofstepwidestar[3]{\exists x\in A\,(x\neq0)}{\rA}
  \proofstepwidestar[1,2]{%
    \CanonicalProfileMap{\JoinOp,A,0}
      \colon A\bij\mathcal U_{\JoinOp,A,0}}{%
    \FormulaRefAuto{CanonicalProfileBijection}[thm]{1,2}}
  \proofstepwidestar[1,2]{\Finite{\mathcal U_{\JoinOp,A,0}}}{%
    \FormulaRefAuto{FiniteBijectionTransport}[thm]{1,4}}
  \proofstepwidestar[2]{\mathcal U_{\JoinOp,A,0}\neq\varnothing}{%
    \FormulaRefAuto{CanonicalProfileFamilyNonempty}[thm]{2}}
  \proofstepwidestar[1,2]{\UnionClosedFam(\mathcal U_{\JoinOp,A,0})}{%
    \FormulaRefAuto{CanonicalProfileFamilyUnionClosed}[thm]{1,2}}
  \proofstepwidestar[1,2,3]{%
    \bigcup\mathcal U_{\JoinOp,A,0}\neq\varnothing}{%
    \FormulaRefAuto{CanonicalProfileCarrierNonempty}[thm]{1,2,3}}
  \proofstepwidestar[1,2,3]{\FranklFam{\mathcal U_{\JoinOp,A,0}}}{%
    \FormulaRefAuto{Frankl-Familie}[def]{6,8,5,7}}
\end{tabproofwide}

\FormulaThmDeltaK[Mitglieder der kanonischen Profilfamilie liegen im Irreduziblenträger]{%
  X\in\mathcal U_{\JoinOp,A,0}\vdash
  X\subseteq\JoinIrrSet{\JoinOp,A,0}
}{CanonicalProfileMemberSubsetIrr}{%
  \DeltaRow{Mengen}{A\dsep X\dsep x\dsep 0}
  \DeltaRow{Funktionen}{\JoinOp}
}
\begin{tabproofwide}
  \proofstepwidestar[1]{X\in\mathcal U_{\JoinOp,A,0}}{\rA}
  \proofstepwidestar[1]{%
    \exists x\in A\,\bigl(
      X=\JoinIrrSet{\JoinOp,A,0}\setminus
        \IrrBelow{\JoinOp,A,0}{x}\bigr)}{%
    \FormulaRefAuto{CanonicalProfileFamilyMembership}[thm]{1}}
  \proofstepwidestar[3]{%
    x\in A\land
    X=\JoinIrrSet{\JoinOp,A,0}\setminus
      \IrrBelow{\JoinOp,A,0}{x}}{\rA}
  \proofstepwidestar[]{%
    \JoinIrrSet{\JoinOp,A,0}\setminus
      \IrrBelow{\JoinOp,A,0}{x}
      \subseteq\JoinIrrSet{\JoinOp,A,0}}{%
    \FormulaRefAuto{A\setminus B\subseteq A}}
  \proofstepwidestar[3]{X\subseteq\JoinIrrSet{\JoinOp,A,0}}{%
    \rIE{\rAEb{3},4}}
  \proofstepwidestar[1]{X\subseteq\JoinIrrSet{\JoinOp,A,0}}{%
    \rEE{2,3,5}}
\end{tabproofwide}

\FormulaThmDeltaK[Elemente des kanonischen Familienträgers sind irreduzibel]{%
  j\in\bigcup\mathcal U_{\JoinOp,A,0}\vdash
  j\in\JoinIrrSet{\JoinOp,A,0}
}{CanonicalCarrierElementIrreducible}{%
  \DeltaRow{Mengen}{A\dsep X\dsep j\dsep 0}
  \DeltaRow{Funktionen}{\JoinOp}
}
\begin{tabproofwide}
  \proofstepwidestar[1]{j\in\bigcup\mathcal U_{\JoinOp,A,0}}{\rA}
  \proofstepwidestar[1]{%
    \exists X\in\mathcal U_{\JoinOp,A,0}\,(j\in X)}{%
    \FormulaRefAuto{x \in \bigcup A \;\eqvdash\;
      \exists B\in A\,(x \in B)}{1}}
  \proofstepwidestar[3]{%
    X\in\mathcal U_{\JoinOp,A,0}\land j\in X}{\rA}
  \proofstepwidestar[3]{X\subseteq\JoinIrrSet{\JoinOp,A,0}}{%
    \FormulaRefAuto{CanonicalProfileMemberSubsetIrr}[thm]{\rAEa{3}}}
  \proofstepwidestar[3]{j\in\JoinIrrSet{\JoinOp,A,0}}{%
    \FormulaRefAuto{A\subseteq B,\,x\in A\vdash x\in B}{4,\rAEb{3}}}
  \proofstepwidestar[1]{j\in\JoinIrrSet{\JoinOp,A,0}}{%
    \rEE{2,3,5}}
\end{tabproofwide}

\FormulaThmDeltaK[Hauptfilter liegen im Halbverbandsträger]{%
  \PrincipalFilter{R,A}{j}\subseteq A
}{PrincipalFilterSubsetCarrier}{%
  \DeltaRow{Mengen}{A\dsep j\dsep x}
  \DeltaRow{Relationen}{R}
}
\begin{tabproof}
  \proofstep{1}{x\in\PrincipalFilter{R,A}{j}}{\rA}
  \proofstep{1}{x\in A\land(j,x)\in R}{%
    \FormulaRefAuto{PrincipalFilter}[def]{1}}
  \proofstep{1}{x\in A}{\rAEa{2}}
  \proofstep{}{\PrincipalFilter{R,A}{j}\subseteq A}{%
    \FormulaRefAuto{A \subseteq B \coloneq
      \forall x\,(x\in A\rightarrow x\in B)}{\rUI{\rRI{1,3}}}}
\end{tabproof}

\FormulaThmDeltaK[Vermeidende Profile entsprechen dem Hauptfilter]{%
  \begin{aligned}[t]
    &\ZeroJoinSemi{\JoinOp,A,0}\dsep
      j\in\JoinIrrSet{\JoinOp,A,0}\dsep x\in A\vdash{}\\
    &\CanonicalProfileMap{\JoinOp,A,0}(x)
      \in\FamAvoiding{\mathcal U_{\JoinOp,A,0}}{j}
      \leftrightarrow
      x\in\PrincipalFilter{\JoinOrd{\JoinOp}{A},A}{j}
  \end{aligned}
}{CanonicalAvoidingMembership}{%
  \DeltaRow{Mengen}{A\dsep j\dsep x\dsep 0}
  \DeltaRow{Funktionen}{\JoinOp\dsep\CanonicalProfileMap{\JoinOp,A,0}}
}
\begin{tabproofsplitwide}
  \proofpartwide{\(\vdash\)}
    \proofstepwidestar[1]{\ZeroJoinSemi{\JoinOp,A,0}}{\rA}
    \proofstepwidestar[2]{j\in\JoinIrrSet{\JoinOp,A,0}}{\rA}
    \proofstepwidestar[3]{x\in A}{\rA}
    \proofstepwidestar[4]{%
      \CanonicalProfileMap{\JoinOp,A,0}(x)
        \in\FamAvoiding{\mathcal U_{\JoinOp,A,0}}{j}}{\rA}
    \proofstepwidestar[4]{%
      \CanonicalProfileMap{\JoinOp,A,0}(x)\in\mathcal U_{\JoinOp,A,0}
      \land j\notin\CanonicalProfileMap{\JoinOp,A,0}(x)}{%
      \FormulaRefAuto{Vermeidende Teilfamilie}[def]{4}}
    \proofstepwidestar[4]{%
      j\notin\CanonicalProfileMap{\JoinOp,A,0}(x)}{\rAEb{5}}
    \proofstepwidestar[1,3]{%
      \CanonicalProfileMap{\JoinOp,A,0}(x)=
      \JoinIrrSet{\JoinOp,A,0}\setminus\IrrBelow{\JoinOp,A,0}{x}}{%
      \FormulaRefAuto{CanonicalProfileEvaluation}[thm]{1,3}}
    \proofstepwidestar[1,3,4]{%
      j\notin\JoinIrrSet{\JoinOp,A,0}\setminus
        \IrrBelow{\JoinOp,A,0}{x}}{\rIE{7,6}}
    \proofstepwidestar[2]{%
      j\notin\JoinIrrSet{\JoinOp,A,0}\setminus
        \IrrBelow{\JoinOp,A,0}{x}
      \leftrightarrow j\in\IrrBelow{\JoinOp,A,0}{x}}{%
      \FormulaRefAuto{RelCompNotMembership}[thm]{2}}
    \proofstepwidestar[1,2,3,4]{j\in\IrrBelow{\JoinOp,A,0}{x}}{%
      \FormulaRefAuto{P\leftrightarrow Q, P\vdash Q}{9,8}}
    \proofstepwidestar[1,2,3,4]{(j,x)\in\JoinOrd{\JoinOp}{A}}{%
      \rAEb{\FormulaRefAuto{IrreduciblesBelow}[def]{10}}}
    \proofstepwidestar[1,2,3,4]{%
      x\in\PrincipalFilter{\JoinOrd{\JoinOp}{A},A}{j}}{%
      \FormulaRefAuto{PrincipalFilter}[def]{\rAI{3,11}}}
  \closeproofpartwide

  \proofpartwide{\(\dashv\)}
    \proofstepwidestar[1]{\ZeroJoinSemi{\JoinOp,A,0}}{\rA}
    \proofstepwidestar[2]{j\in\JoinIrrSet{\JoinOp,A,0}}{\rA}
    \proofstepwidestar[3]{x\in A}{\rA}
    \proofstepwidestar[4]{%
      x\in\PrincipalFilter{\JoinOrd{\JoinOp}{A},A}{j}}{\rA}
    \proofstepwidestar[4]{(j,x)\in\JoinOrd{\JoinOp}{A}}{%
      \rAEb{\FormulaRefAuto{PrincipalFilter}[def]{4}}}
    \proofstepwidestar[2,4]{j\in\IrrBelow{\JoinOp,A,0}{x}}{%
      \FormulaRefAuto{IrreduciblesBelow}[def]{\rAI{2,5}}}
    \proofstepwidestar[2]{%
      j\notin\JoinIrrSet{\JoinOp,A,0}\setminus
        \IrrBelow{\JoinOp,A,0}{x}
      \leftrightarrow j\in\IrrBelow{\JoinOp,A,0}{x}}{%
      \FormulaRefAuto{RelCompNotMembership}[thm]{2}}
    \proofstepwidestar[2,4]{%
      j\notin\JoinIrrSet{\JoinOp,A,0}\setminus
        \IrrBelow{\JoinOp,A,0}{x}}{%
      \FormulaRefAuto{P\leftrightarrow Q, Q\vdash P}{7,6}}
    \proofstepwidestar[1,3]{%
      \CanonicalProfileMap{\JoinOp,A,0}(x)=
      \JoinIrrSet{\JoinOp,A,0}\setminus\IrrBelow{\JoinOp,A,0}{x}}{%
      \FormulaRefAuto{CanonicalProfileEvaluation}[thm]{1,3}}
    \proofstepwidestar[1,2,3,4]{%
      j\notin\CanonicalProfileMap{\JoinOp,A,0}(x)}{%
      \rIE{\FormulaRefAuto{a=b\vdash b=a}{9},8}}
    \proofstepwidestar[1]{%
      \CanonicalProfileMap{\JoinOp,A,0}\colon
        A\to\mathcal U_{\JoinOp,A,0}}{%
      \FormulaRefAuto{CanonicalProfileFunction}[thm]{1}}
    \proofstepwidestar[1,3]{%
      \CanonicalProfileMap{\JoinOp,A,0}(x)\in\mathcal U_{\JoinOp,A,0}}{%
      \FormulaRefAuto{F\colon A\to B\dsep x\in A\vdash F(x)\in B}{11,3}}
    \proofstepwidestar[1,2,3,4]{%
      \CanonicalProfileMap{\JoinOp,A,0}(x)
        \in\FamAvoiding{\mathcal U_{\JoinOp,A,0}}{j}}{%
      \FormulaRefAuto{Vermeidende Teilfamilie}[def]{\rAI{12,10}}}
  \closeproofpartwide
\end{tabproofsplitwide}

\FormulaThmDeltaK[Enthaltende Profile entsprechen dem Filterkomplement]{%
  \begin{aligned}[t]
    &\ZeroJoinSemi{\JoinOp,A,0}\dsep
      j\in\JoinIrrSet{\JoinOp,A,0}\dsep x\in A\vdash{}\\
    &\CanonicalProfileMap{\JoinOp,A,0}(x)
      \in\FamContaining{\mathcal U_{\JoinOp,A,0}}{j}
      \leftrightarrow
      x\in A\setminus\PrincipalFilter{\JoinOrd{\JoinOp}{A},A}{j}
  \end{aligned}
}{CanonicalContainingMembership}{%
  \DeltaRow{Mengen}{A\dsep j\dsep x\dsep 0}
  \DeltaRow{Funktionen}{\JoinOp\dsep\CanonicalProfileMap{\JoinOp,A,0}}
}
\begin{tabproofsplitwide}
  \proofpartwide{\(\vdash\)}
    \proofstepwidestar[1]{\ZeroJoinSemi{\JoinOp,A,0}}{\rA}
    \proofstepwidestar[2]{j\in\JoinIrrSet{\JoinOp,A,0}}{\rA}
    \proofstepwidestar[3]{x\in A}{\rA}
    \proofstepwidestar[4]{%
      \CanonicalProfileMap{\JoinOp,A,0}(x)
        \in\FamContaining{\mathcal U_{\JoinOp,A,0}}{j}}{\rA}
    \proofstepwidestar[4]{%
      j\in\CanonicalProfileMap{\JoinOp,A,0}(x)}{%
      \rAEb{\FormulaRefAuto{Enthaltende Teilfamilie}[def]{4}}}
    \proofstepwidestar[6]{%
      x\in\PrincipalFilter{\JoinOrd{\JoinOp}{A},A}{j}}{\rA}
    \proofstepwidestar[1,2,3,6]{%
      \CanonicalProfileMap{\JoinOp,A,0}(x)
        \in\FamAvoiding{\mathcal U_{\JoinOp,A,0}}{j}}{%
      \FormulaRefAuto{CanonicalAvoidingMembership}[thm]{1,2,3,6}}
    \proofstepwidestar[7]{%
      j\notin\CanonicalProfileMap{\JoinOp,A,0}(x)}{%
      \rAEb{\FormulaRefAuto{Vermeidende Teilfamilie}[def]{7}}}
    \proofstepwidestar[4,6]{\bot}{\rBI{5,8}}
    \proofstepwidestar[4]{%
      x\notin\PrincipalFilter{\JoinOrd{\JoinOp}{A},A}{j}}{%
      \rCI{6,9}}
    \proofstepwidestar[3,4]{%
      x\in A\setminus\PrincipalFilter{\JoinOrd{\JoinOp}{A},A}{j}}{%
      \FormulaRefAuto{x \in A\dsep x\notin B\vdash x \in A\setminus B}{3,10}}
  \closeproofpartwide

  \proofpartwide{\(\dashv\)}
    \proofstepwidestar[1]{\ZeroJoinSemi{\JoinOp,A,0}}{\rA}
    \proofstepwidestar[2]{j\in\JoinIrrSet{\JoinOp,A,0}}{\rA}
    \proofstepwidestar[3]{x\in A}{\rA}
    \proofstepwidestar[4]{%
      x\in A\setminus\PrincipalFilter{\JoinOrd{\JoinOp}{A},A}{j}}{\rA}
    \proofstepwidestar[4]{%
      x\notin\PrincipalFilter{\JoinOrd{\JoinOp}{A},A}{j}}{%
      \FormulaRefAuto{x\in A\setminus B\vdash x\notin B}{4}}
    \proofstepwidestar[1]{%
      \CanonicalProfileMap{\JoinOp,A,0}\colon
        A\to\mathcal U_{\JoinOp,A,0}}{%
      \FormulaRefAuto{CanonicalProfileFunction}[thm]{1}}
    \proofstepwidestar[1,3]{%
      \CanonicalProfileMap{\JoinOp,A,0}(x)\in\mathcal U_{\JoinOp,A,0}}{%
      \FormulaRefAuto{F\colon A\to B\dsep x\in A\vdash F(x)\in B}{6,3}}
    \proofstepwidestar[8]{%
      j\notin\CanonicalProfileMap{\JoinOp,A,0}(x)}{\rA}
    \proofstepwidestar[1,3,8]{%
      \CanonicalProfileMap{\JoinOp,A,0}(x)
        \in\FamAvoiding{\mathcal U_{\JoinOp,A,0}}{j}}{%
      \FormulaRefAuto{Vermeidende Teilfamilie}[def]{\rAI{7,8}}}
    \proofstepwidestar[1,2,3,8]{%
      x\in\PrincipalFilter{\JoinOrd{\JoinOp}{A},A}{j}}{%
      \FormulaRefAuto{CanonicalAvoidingMembership}[thm]{1,2,3,9}}
    \proofstepwidestar[4,8]{\bot}{\rBI{10,5}}
    \proofstepwidestar[4]{%
      \neg\bigl(j\notin\CanonicalProfileMap{\JoinOp,A,0}(x)\bigr)}{%
      \rCI{8,11}}
    \proofstepwidestar[4]{%
      j\in\CanonicalProfileMap{\JoinOp,A,0}(x)}{%
      \FormulaRefAuto{rDN}{12}}
    \proofstepwidestar[1,3,4]{%
      \CanonicalProfileMap{\JoinOp,A,0}(x)
        \in\FamContaining{\mathcal U_{\JoinOp,A,0}}{j}}{%
      \FormulaRefAuto{Enthaltende Teilfamilie}[def]{\rAI{7,13}}}
  \closeproofpartwide
\end{tabproofsplitwide}

\FormulaThmDeltaK[Urbild der vermeidenden kanonischen Profile]{%
  \begin{aligned}[t]
    &\ZeroJoinSemi{\JoinOp,A,0}\dsep
      j\in\JoinIrrSet{\JoinOp,A,0}\vdash{}\\
    &\bigl(\CanonicalProfileMap{\JoinOp,A,0}\bigr)^{-1}
      [\FamAvoiding{\mathcal U_{\JoinOp,A,0}}{j}]
      =\PrincipalFilter{\JoinOrd{\JoinOp}{A},A}{j}
  \end{aligned}
}{CanonicalAvoidingPreimage}{%
  \DeltaRow{Mengen}{A\dsep j\dsep x\dsep 0}
  \DeltaRow{Funktionen}{\JoinOp\dsep\CanonicalProfileMap{\JoinOp,A,0}}
}
\begin{tabproofwide}
  \proofstepwidestar[1]{\ZeroJoinSemi{\JoinOp,A,0}}{\rA}
  \proofstepwidestar[2]{j\in\JoinIrrSet{\JoinOp,A,0}}{\rA}
  \proofstepwidestar[1]{%
    \CanonicalProfileMap{\JoinOp,A,0}\colon
      A\to\mathcal U_{\JoinOp,A,0}}{%
    \FormulaRefAuto{CanonicalProfileFunction}[thm]{1}}
  \proofstepwidestar[]{%
    \FamAvoiding{\mathcal U_{\JoinOp,A,0}}{j}
      \subseteq\mathcal U_{\JoinOp,A,0}}{%
    \FormulaRefAuto{\FamAvoiding{M}{a}\subseteq M}}
  \proofstepwidestar[]{%
    \PrincipalFilter{\JoinOrd{\JoinOp}{A},A}{j}\subseteq A}{%
    \FormulaRefAuto{PrincipalFilterSubsetCarrier}[thm]}
  \proofstepwidestar[3]{x\in A}{\rA}
  \proofstepwidestar[1,2,3]{%
    \CanonicalProfileMap{\JoinOp,A,0}(x)
      \in\FamAvoiding{\mathcal U_{\JoinOp,A,0}}{j}
      \leftrightarrow
      x\in\PrincipalFilter{\JoinOrd{\JoinOp}{A},A}{j}}{%
    \FormulaRefAuto{CanonicalAvoidingMembership}[thm]{1,2,6}}
  \proofstepwidestar[1,2]{%
    \forall x\in A\,\bigl(
      \CanonicalProfileMap{\JoinOp,A,0}(x)
        \in\FamAvoiding{\mathcal U_{\JoinOp,A,0}}{j}
      \leftrightarrow
      x\in\PrincipalFilter{\JoinOrd{\JoinOp}{A},A}{j}\bigr)}{%
    \rUI{\rRI{6,7}}}
  \proofstepwidestar[1,2]{%
    \bigl(\CanonicalProfileMap{\JoinOp,A,0}\bigr)^{-1}
      [\FamAvoiding{\mathcal U_{\JoinOp,A,0}}{j}]
      =\PrincipalFilter{\JoinOrd{\JoinOp}{A},A}{j}}{%
    \FormulaRefAuto{PreimageEqFromPointwiseIff}[thm]{3,4,5,8}}
\end{tabproofwide}

\FormulaThmDeltaK[Urbild der enthaltenden kanonischen Profile]{%
  \begin{aligned}[t]
    &\ZeroJoinSemi{\JoinOp,A,0}\dsep
      j\in\JoinIrrSet{\JoinOp,A,0}\vdash{}\\
    &\bigl(\CanonicalProfileMap{\JoinOp,A,0}\bigr)^{-1}
      [\FamContaining{\mathcal U_{\JoinOp,A,0}}{j}]
      =A\setminus\PrincipalFilter{\JoinOrd{\JoinOp}{A},A}{j}
  \end{aligned}
}{CanonicalContainingPreimage}{%
  \DeltaRow{Mengen}{A\dsep j\dsep x\dsep 0}
  \DeltaRow{Funktionen}{\JoinOp\dsep\CanonicalProfileMap{\JoinOp,A,0}}
}
\begin{tabproofwide}
  \proofstepwidestar[1]{\ZeroJoinSemi{\JoinOp,A,0}}{\rA}
  \proofstepwidestar[2]{j\in\JoinIrrSet{\JoinOp,A,0}}{\rA}
  \proofstepwidestar[1]{%
    \CanonicalProfileMap{\JoinOp,A,0}\colon
      A\to\mathcal U_{\JoinOp,A,0}}{%
    \FormulaRefAuto{CanonicalProfileFunction}[thm]{1}}
  \proofstepwidestar[]{%
    \FamContaining{\mathcal U_{\JoinOp,A,0}}{j}
      \subseteq\mathcal U_{\JoinOp,A,0}}{%
    \FormulaRefAuto{\FamContaining{M}{a}\subseteq M}}
  \proofstepwidestar[]{%
    A\setminus\PrincipalFilter{\JoinOrd{\JoinOp}{A},A}{j}\subseteq A}{%
    \FormulaRefAuto{A\setminus B\subseteq A}}
  \proofstepwidestar[3]{x\in A}{\rA}
  \proofstepwidestar[1,2,3]{%
    \CanonicalProfileMap{\JoinOp,A,0}(x)
      \in\FamContaining{\mathcal U_{\JoinOp,A,0}}{j}
      \leftrightarrow
      x\in A\setminus\PrincipalFilter{\JoinOrd{\JoinOp}{A},A}{j}}{%
    \FormulaRefAuto{CanonicalContainingMembership}[thm]{1,2,6}}
  \proofstepwidestar[1,2]{%
    \forall x\in A\,\bigl(
      \CanonicalProfileMap{\JoinOp,A,0}(x)
        \in\FamContaining{\mathcal U_{\JoinOp,A,0}}{j}
      \leftrightarrow
      x\in A\setminus\PrincipalFilter{\JoinOrd{\JoinOp}{A},A}{j}\bigr)}{%
    \rUI{\rRI{6,7}}}
  \proofstepwidestar[1,2]{%
    \bigl(\CanonicalProfileMap{\JoinOp,A,0}\bigr)^{-1}
      [\FamContaining{\mathcal U_{\JoinOp,A,0}}{j}]
      =A\setminus\PrincipalFilter{\JoinOrd{\JoinOp}{A},A}{j}}{%
    \FormulaRefAuto{PreimageEqFromPointwiseIff}[thm]{3,4,5,8}}
\end{tabproofwide}

\FormulaThmDeltaK[Injektionstransport entlang einer Bijektion]{%
  c\colon A\bij B\dsep H\subseteq B\dsep\AtMostHalf{H}{B}
  \vdash\exists F\colon c^{-1}[H]\inj c^{-1}[B\setminus H]
}{BijectionAtMostHalfTransport}{%
  \DeltaRow{Mengen}{A\dsep B\dsep H}
  \DeltaRow{Funktionen}{c\dsep F\dsep G}
}
\begin{tabproofwide}
  \proofstepwidestar[1]{c\colon A\bij B}{\rA}
  \proofstepwidestar[2]{H\subseteq B}{\rA}
  \proofstepwidestar[3]{\AtMostHalf{H}{B}}{\rA}
  \proofstepwidestar[1,2]{%
    c\restriction_{c^{-1}[H]}\colon c^{-1}[H]\bij H}{%
    \FormulaRefAuto{F\colon A\bij B \dsep T\subseteq B \vdash
      F\restriction_{F^{-1}[T]}\colon F^{-1}[T]\bij T}{1,2}}
  \proofstepwidestar[]{B\setminus H\subseteq B}{%
    \FormulaRefAuto{A\setminus B\subseteq A}}
  \proofstepwidestar[1]{%
    c\restriction_{c^{-1}[B\setminus H]}\colon
    c^{-1}[B\setminus H]\bij B\setminus H}{%
    \FormulaRefAuto{F\colon A\bij B \dsep T\subseteq B \vdash
      F\restriction_{F^{-1}[T]}\colon F^{-1}[T]\bij T}{1,5}}
  \proofstepwidestar[1]{%
    \left(c\restriction_{c^{-1}[B\setminus H]}\right)^{-1}
    \colon B\setminus H\bij c^{-1}[B\setminus H]}{%
    \FormulaRefAuto{F\colon A\bij B\vdash F^{-1}\colon B\bij A}{6}}
  \proofstepwidestar[2,3]{\exists G\colon H\inj B\setminus H}{%
    \FormulaRefAuto{AtMostHalf}[def]{2,3}}
  \proofstepwidestar[9]{G\colon H\inj B\setminus H}{\rA}
  \proofstepwidestar[1,2]{%
    c\restriction_{c^{-1}[H]}\colon c^{-1}[H]\inj H}{%
    \FormulaRefAuto{F\colon A\bij B\vdash F\colon A\inj B}{4}}
  \proofstepwidestar[1]{%
    \left(c\restriction_{c^{-1}[B\setminus H]}\right)^{-1}
      \colon B\setminus H\inj c^{-1}[B\setminus H]}{%
    \FormulaRefAuto{F\colon A\bij B\vdash F\colon A\inj B}{7}}
  \proofstepwidestar[1,2,9]{%
    G\circ(c\restriction_{c^{-1}[H]})
      \colon c^{-1}[H]\inj B\setminus H}{%
    \FormulaRefAuto{F\colon A\inj B \dsep G\colon B\inj C
      \vdash G\circ F\colon A\inj C}{10,9}}
  \proofstepwidestar[1,2,9]{%
    \left(c\restriction_{c^{-1}[B\setminus H]}\right)^{-1}
      \circ G\circ(c\restriction_{c^{-1}[H]})
      \colon c^{-1}[H]\inj c^{-1}[B\setminus H]}{%
    \FormulaRefAuto{F\colon A\inj B \dsep G\colon B\inj C
      \vdash G\circ F\colon A\inj C}{12,11}}
  \proofstepwidestar[1,2,9]{%
    \exists F\colon c^{-1}[H]\inj c^{-1}[B\setminus H]}{\rEI{13}}
  \proofstepwidestar[1,2,3]{%
    \exists F\colon c^{-1}[H]\inj c^{-1}[B\setminus H]}{%
    \rEE{8,9,14}}
\end{tabproofwide}

\FormulaThmDeltaK[Transport eines halbhäufigen Profilelements]{%
  \begin{aligned}[t]
    &\Finite A\dsep\ZeroJoinSemi{\JoinOp,A,0}\dsep
      \exists x\in A\,(x\neq0)\dsep
      \HalfOcc{j}{\mathcal U_{\JoinOp,A,0}}\vdash{}\\
    &\FranklSemi{\JoinOp,A,0}
  \end{aligned}
}{CanonicalHalfOccurrenceTransport}{%
  \DeltaRow{Mengen}{A\dsep j\dsep x\dsep 0}
  \DeltaRow{Funktionen}{%
    \JoinOp\dsep\CanonicalProfileMap{\JoinOp,A,0}\dsep F\dsep G}
}

Für die kanonische Bijektion
\(\CanonicalProfileMap{\JoinOp,A,0}\) gilt für jedes \(x\in A\)
\[
  \begin{aligned}
  j\notin \CanonicalProfileMap{\JoinOp,A,0}(x)
  &\ \Longleftrightarrow\
    j\in\IrrBelow{\JoinOp,A,0}{x}\\
  &\ \Longleftrightarrow\
    x\in\PrincipalFilter{\JoinOrd{\JoinOp}{A},A}{j}.
  \end{aligned}
\]
Da jedes Mitglied von \(\mathcal U_{\JoinOp,A,0}\) eine Teilmenge von
\(\JoinIrrSet{\JoinOp,A,0}\) ist, liefert
\(j\in\bigcup\mathcal U_{\JoinOp,A,0}\) zugleich die benötigte
\(\JoinOp\)-Irreduzibilität. Die in \(\HalfOcc{j}{\mathcal U_{\JoinOp,A,0}}\) enthaltene
Injektion wird nun entlang der kanonischen Profilabbildung transportiert.

\begin{tabproofsplitwide}
  \proofpartwideindR[Halbhäufigkeit der vermeidenden Teilfamilie]{%
    \begin{aligned}[t]
      &\HalfOcc{j}{\mathcal U_{\JoinOp,A,0}}\vdash{}\\[-2pt]
      &\AtMostHalf{\FamAvoiding{\mathcal U_{\JoinOp,A,0}}{j}}{
        \mathcal U_{\JoinOp,A,0}}
    \end{aligned}
  }{\HalfOcc{j}{\mathcal U_{\JoinOp,A,0}}\vdash
    \AtMostHalf{\FamAvoiding{\mathcal U_{\JoinOp,A,0}}{j}}{
      \mathcal U_{\JoinOp,A,0}}}[CanonicalAvoidingAtMostHalf]
    \proofstepwidestar[1]{\HalfOcc{j}{\mathcal U_{\JoinOp,A,0}}}{\rA}
    \proofstepwidestar[]{%
      \FamAvoiding{\mathcal U_{\JoinOp,A,0}}{j}
        \subseteq\mathcal U_{\JoinOp,A,0}}{%
      \FormulaRefAuto{\FamAvoiding{M}{a}\subseteq M}}
    \proofstepwidestar[1]{%
      \exists F\colon
        \FamAvoiding{\mathcal U_{\JoinOp,A,0}}{j}
        \inj\FamContaining{\mathcal U_{\JoinOp,A,0}}{j}}{%
      \FormulaRefAuto{Halbhäufiges Element}[def]{1}}
    \proofstepwidestar[]{%
      \mathcal U_{\JoinOp,A,0}\setminus
        \FamAvoiding{\mathcal U_{\JoinOp,A,0}}{j}
        =\FamContaining{\mathcal U_{\JoinOp,A,0}}{j}}{%
      \FormulaRefAuto{FamContainingAsAvoidingRelativeComplement}[thm]}
    \proofstepwidestar[3]{%
      F\colon\FamAvoiding{\mathcal U_{\JoinOp,A,0}}{j}
        \inj\FamContaining{\mathcal U_{\JoinOp,A,0}}{j}}{\rA}
    \proofstepwidestar[]{%
      \FamContaining{\mathcal U_{\JoinOp,A,0}}{j}
        =\mathcal U_{\JoinOp,A,0}\setminus
          \FamAvoiding{\mathcal U_{\JoinOp,A,0}}{j}}{%
      \FormulaRefAuto{a=b\vdash b=a}{4}}
    \proofstepwidestar[3]{%
      \begin{aligned}
        F\colon\FamAvoiding{\mathcal U_{\JoinOp,A,0}}{j}
          \inj{}&\mathcal U_{\JoinOp,A,0}\setminus\\[-2pt]
          &\FamAvoiding{\mathcal U_{\JoinOp,A,0}}{j}
      \end{aligned}}{\rIE{6,5}}
    \proofstepwidestar[3]{%
      \exists F\colon
        \FamAvoiding{\mathcal U_{\JoinOp,A,0}}{j}
        \inj\mathcal U_{\JoinOp,A,0}\setminus
          \FamAvoiding{\mathcal U_{\JoinOp,A,0}}{j}}{\rEI{7}}
    \proofstepwidestar[1]{%
      \exists F\colon
        \FamAvoiding{\mathcal U_{\JoinOp,A,0}}{j}
        \inj\mathcal U_{\JoinOp,A,0}\setminus
          \FamAvoiding{\mathcal U_{\JoinOp,A,0}}{j}}{\rEE{3,5,8}}
    \proofstepwidestar[1]{%
      \AtMostHalf{\FamAvoiding{\mathcal U_{\JoinOp,A,0}}{j}}{
        \mathcal U_{\JoinOp,A,0}}}{%
      \FormulaRefAuto{AtMostHalf}[def]{2,9}}
  \closeproofpartwideindR

  \proofpartwideindR[Transport in den Hauptfilter]{%
    \begin{aligned}[t]
      &\Finite A\dsep\ZeroJoinSemi{\JoinOp,A,0}\dsep
        j\in\JoinIrrSet{\JoinOp,A,0}\dsep
        \HalfOcc{j}{\mathcal U_{\JoinOp,A,0}}\vdash{}\\[-2pt]
      &\exists G\colon
        \PrincipalFilter{\JoinOrd{\JoinOp}{A},A}{j}
        \inj A\setminus\PrincipalFilter{\JoinOrd{\JoinOp}{A},A}{j}
    \end{aligned}
  }{\Finite A\dsep\ZeroJoinSemi{\JoinOp,A,0}\dsep
    j\in\JoinIrrSet{\JoinOp,A,0}\dsep
    \HalfOcc{j}{\mathcal U_{\JoinOp,A,0}}\vdash
    \exists G\colon\PrincipalFilter{\JoinOrd{\JoinOp}{A},A}{j}
    \inj A\setminus\PrincipalFilter{\JoinOrd{\JoinOp}{A},A}{j}}%
  [CanonicalHalfOccurrenceFilterInjection]
    \proofstepwidestar[1]{\Finite A}{\rA}
    \proofstepwidestar[2]{\ZeroJoinSemi{\JoinOp,A,0}}{\rA}
    \proofstepwidestar[3]{j\in\JoinIrrSet{\JoinOp,A,0}}{\rA}
    \proofstepwidestar[4]{\HalfOcc{j}{\mathcal U_{\JoinOp,A,0}}}{\rA}
    \proofstepwidestar[1,2]{%
      \CanonicalProfileMap{\JoinOp,A,0}
        \colon A\bij\mathcal U_{\JoinOp,A,0}}{%
      \FormulaRefAuto{CanonicalProfileBijection}[thm]{1,2}}
    \proofstepwidestar[]{%
      \FamAvoiding{\mathcal U_{\JoinOp,A,0}}{j}
        \subseteq\mathcal U_{\JoinOp,A,0}}{%
      \FormulaRefAuto{\FamAvoiding{M}{a}\subseteq M}}
    \proofstepwidestar[4]{%
      \AtMostHalf{\FamAvoiding{\mathcal U_{\JoinOp,A,0}}{j}}{
        \mathcal U_{\JoinOp,A,0}}}{%
      \FormulaRefAuto{CanonicalAvoidingAtMostHalf}[thm]{4}}
    \proofstepwidestar[1,2,4]{%
      \begin{aligned}
        &\exists G\colon
        \bigl(\CanonicalProfileMap{\JoinOp,A,0}\bigr)^{-1}
          [\FamAvoiding{\mathcal U_{\JoinOp,A,0}}{j}]\inj{}\\[-2pt]
        &\quad\bigl(\CanonicalProfileMap{\JoinOp,A,0}\bigr)^{-1}
          [\mathcal U_{\JoinOp,A,0}\setminus
            \FamAvoiding{\mathcal U_{\JoinOp,A,0}}{j}]
      \end{aligned}}{%
      \FormulaRefAuto{BijectionAtMostHalfTransport}[thm]{5,6,7}}
    \proofstepwidestar[]{%
      \mathcal U_{\JoinOp,A,0}\setminus
        \FamAvoiding{\mathcal U_{\JoinOp,A,0}}{j}
        =\FamContaining{\mathcal U_{\JoinOp,A,0}}{j}}{%
      \FormulaRefAuto{FamContainingAsAvoidingRelativeComplement}[thm]}
    \proofstepwidestar[1,2,4]{%
      \begin{aligned}
        &\exists G\colon
        \bigl(\CanonicalProfileMap{\JoinOp,A,0}\bigr)^{-1}
          [\FamAvoiding{\mathcal U_{\JoinOp,A,0}}{j}]\inj{}\\[-2pt]
        &\quad\bigl(\CanonicalProfileMap{\JoinOp,A,0}\bigr)^{-1}
          [\FamContaining{\mathcal U_{\JoinOp,A,0}}{j}]
      \end{aligned}}{\rIE{9,8}}
    \proofstepwidestar[2,3]{%
      \bigl(\CanonicalProfileMap{\JoinOp,A,0}\bigr)^{-1}
        [\FamAvoiding{\mathcal U_{\JoinOp,A,0}}{j}]
        =\PrincipalFilter{\JoinOrd{\JoinOp}{A},A}{j}}{%
      \FormulaRefAuto{CanonicalAvoidingPreimage}[thm]{2,3}}
    \proofstepwidestar[1,2,3,4]{%
      \exists G\colon
        \PrincipalFilter{\JoinOrd{\JoinOp}{A},A}{j}\inj
        \bigl(\CanonicalProfileMap{\JoinOp,A,0}\bigr)^{-1}
          [\FamContaining{\mathcal U_{\JoinOp,A,0}}{j}]}{\rIE{11,10}}
    \proofstepwidestar[2,3]{%
      \bigl(\CanonicalProfileMap{\JoinOp,A,0}\bigr)^{-1}
        [\FamContaining{\mathcal U_{\JoinOp,A,0}}{j}]
        =A\setminus\PrincipalFilter{\JoinOrd{\JoinOp}{A},A}{j}}{%
      \FormulaRefAuto{CanonicalContainingPreimage}[thm]{2,3}}
    \proofstepwidestar[1,2,3,4]{%
      \exists G\colon
        \PrincipalFilter{\JoinOrd{\JoinOp}{A},A}{j}
        \inj A\setminus\PrincipalFilter{\JoinOrd{\JoinOp}{A},A}{j}}{%
      \rIE{13,12}}
  \closeproofpartwideindR

  \proofpartwide{Schluss}
    \proofstepwidestar[1]{\Finite A}{\rA}
    \proofstepwidestar[2]{\ZeroJoinSemi{\JoinOp,A,0}}{\rA}
    \proofstepwidestar[3]{\exists x\in A\,(x\neq0)}{\rA}
    \proofstepwidestar[4]{\HalfOcc{j}{\mathcal U_{\JoinOp,A,0}}}{\rA}
    \proofstepwidestar[4]{j\in\bigcup\mathcal U_{\JoinOp,A,0}}{%
      \FormulaRefAuto{Halbhäufiges Element}[def]{4}}
    \proofstepwidestar[5]{j\in\JoinIrrSet{\JoinOp,A,0}}{%
      \FormulaRefAuto{CanonicalCarrierElementIrreducible}[thm]{5}}
    \proofstepwidestar[2,5]{\JoinIrr{\JoinOp,A,0}{j}}{%
      \FormulaRefAuto{JoinIrreducibleSet}[def]{6}}
    \proofstepwidestar[1,2,4,5]{%
      \exists G\colon
        \PrincipalFilter{\JoinOrd{\JoinOp}{A},A}{j}
        \inj A\setminus\PrincipalFilter{\JoinOrd{\JoinOp}{A},A}{j}}{%
      \FormulaRefAuto{CanonicalHalfOccurrenceFilterInjection}[thm]{1,2,6,4}}
    \proofstepwidestar[]{%
      \PrincipalFilter{\JoinOrd{\JoinOp}{A},A}{j}\subseteq A}{%
      \FormulaRefAuto{PrincipalFilterSubsetCarrier}[thm]}
    \proofstepwidestar[1,2,4,5]{%
      \AtMostHalf{\PrincipalFilter{\JoinOrd{\JoinOp}{A},A}{j}}{A}}{%
      \FormulaRefAuto{AtMostHalf}[def]{9,8}}
    \proofstepwidestar[1,2,3,4]{\FranklSemi{\JoinOp,A,0}}{%
      \FormulaRefAuto{FranklSemilatticeProperty}[def]{%
        \rEI{\rAI{7,10}}}}
  \closeproofpartwide
\end{tabproofsplitwide}

\FormulaThmDeltaK[Mengenform impliziert Halbverbandsform]{%
  \begin{aligned}[t]
    &\forall M\,(\FranklFam M\to\Frankl M)\dsep
      \Finite A\dsep{}\\
    &\ZeroJoinSemi{\JoinOp,A,0}\dsep
      \exists x\in A\,(x\neq0)\vdash{}\\
    &\FranklSemi{\JoinOp,A,0}
  \end{aligned}
}{FranklSetImpliesSemilattice}{%
  \DeltaRow{Mengen}{A\dsep M\dsep x\dsep 0}
  \DeltaRow{Funktionen}{\JoinOp}
}
\begin{tabproofwide}
  \proofstepwidestar[1]{\forall M\,(\FranklFam M\to\Frankl M)}{\rA}
  \proofstepwidestar[2]{\Finite A}{\rA}
  \proofstepwidestar[3]{\ZeroJoinSemi{\JoinOp,A,0}}{\rA}
  \proofstepwidestar[4]{\exists x\in A\,(x\neq0)}{\rA}
  \proofstepwidestar[2,3,4]{\FranklFam{\mathcal U_{\JoinOp,A,0}}}{%
    \FormulaRefAuto{CanonicalUnionFamilyIsFranklFamily}[thm]{2,3,4}}
  \proofstepwidestar[1,2,3,4]{\Frankl{\mathcal U_{\JoinOp,A,0}}}{%
    \rRE{\rUE{1},5}}
  \proofstepwidestar[1,2,3,4]{%
    \exists j\,\HalfOcc{j}{\mathcal U_{\JoinOp,A,0}}}{%
    \FormulaRefAuto{Frankl-Eigenschaft}[def]{6}}
  \proofstepwidestar[8]{\HalfOcc{j}{\mathcal U_{\JoinOp,A,0}}}{\rA}
  \proofstepwidestar[2,3,4,8]{\FranklSemi{\JoinOp,A,0}}{%
    \FormulaRefAuto{CanonicalHalfOccurrenceTransport}[thm]{2,3,4,8}}
  \proofstepwidestar[1,2,3,4]{\FranklSemi{\JoinOp,A,0}}{\rEE{7,8,9}}
\end{tabproofwide}

\section{Von der Vereinigungsfamilie zum Halbverband}

\FormulaDefDeltaK[Schnittabgeschlossene Familie]{%
  \InterClosedFam L
  \coloneqq
  \forall X\in L\,\forall Y\in L\,(X\cap Y\in L)
}{IntersectionClosedFamily}{%
  \DeltaRow{Mengen}{L\dsep X\dsep Y}
  \DeltaRow{\textbf{Neues Symbol}}{}
  \DeltaPrem{Schnittabgeschlossene Familien}{\InterClosedFam L}
}

\par\medskip
Sei nun \(M\) eine Frankl-Familie. Wir verwenden die folgenden vier
kanonischen Objekte:
\[
  \begin{aligned}
    U_M&\coloneqq\bigcup M,&
    M^\circ&\coloneqq M\cup\{\varnothing\},\\
    \mathcal L_M
      &\coloneqq\{U_M\setminus X\mid X\in M^\circ\},&
    \mathcal R_M
      &\coloneqq
        \{(X,Y)\in\mathcal L_M\times\mathcal L_M\mid X\subseteq Y\}.
  \end{aligned}
\]
\par\noindent
Für \(X,Y\in\mathcal L_M\) setzen wir außerdem
\[
  X\FamilyJoinOp{M}{}Y
  \coloneqq
  \bigcap\{Z\in\mathcal L_M\mid X\subseteq Z\land Y\subseteq Z\}.
\]
Die Menge unter dem Durchschnitt ist nichtleer, weil sie \(U_M\) enthält.
Die explizite Relation \(\mathcal R_M\) verhindert eine Mehrdeutigkeit:
Im folgenden Halbverband ist die induzierte Ordnung tatsächlich die
Inklusion.

\FormulaDefDeltaK[Komplementduale Konstruktion]{%
  \begin{aligned}[t]
    U_M&\coloneqq\bigcup M,&
    M^\circ&\coloneqq M\cup\{\varnothing\},\\
    \mathcal L_M
      &\coloneqq\{U_M\setminus X\mid X\in M^\circ\},&
    \mathcal R_M
      &\coloneqq
        \{(X,Y)\in\mathcal L_M^2\mid X\subseteq Y\},\\
    X\FamilyJoinOp{M}{}Y
      &\coloneqq
        \bigcap\{Z\in\mathcal L_M\mid X\cup Y\subseteq Z\}.
  \end{aligned}
}{FranklComplementConstruction}{%
  \DeltaRow{Mengen}{M\dsep X\dsep Y\dsep Z}
  \DeltaRow{Relationen}{\mathcal R_M}
  \DeltaRow{Funktionen}{\FamilyJoinOp{M}}
}

\FormulaThmDeltaK[Beschränkter Existenzquantor über einer Adjunktion]{%
  \exists Y\in C\cup\{X\}\,P(Y)
  \leftrightarrow
  \bigl(\exists Y\in C\,P(Y)\bigr)\lor P(X)
}{BoundedExistentialAdjunction}{%
  \DeltaRow{Mengen}{C\dsep X\dsep Y}
  \DeltaRow{Prädikatensymbole}{P}
}
\begin{tabproofwide}
  \proofstepwidestar[1]{\exists Y\in C\cup\{X\}\,P(Y)}{\rA}
  \proofstepwidestar[2]{Y\in C\cup\{X\}\land P(Y)}{\rA}
  \proofstepwidestar[2]{Y\in C\cup\{X\}}{\rAEa{2}}
  \proofstepwidestar[2]{P(Y)}{\rAEb{2}}
  \proofstepwidestar[2]{Y\in C\lor Y=X}{%
    \FormulaRefAuto{x\in A\cup\{a\}\vdash x\in A\lor x=a}{3}}
  \proofstepwidestar[6]{Y\in C}{\rA}
  \proofstepwidestar[2,6]{\exists Y\in C\,P(Y)}{\rEI{\rAI{6,4}}}
  \proofstepwidestar[2,6]{\bigl(\exists Y\in C\,P(Y)\bigr)\lor P(X)}{%
    \rOIa{7}}
  \proofstepwidestar[9]{Y=X}{\rA}
  \proofstepwidestar[2,9]{P(X)}{\rIE{9,4}}
  \proofstepwidestar[2,9]{\bigl(\exists Y\in C\,P(Y)\bigr)\lor P(X)}{%
    \rOIb{10}}
  \proofstepwidestar[2]{\bigl(\exists Y\in C\,P(Y)\bigr)\lor P(X)}{%
    \rOE{5,6,8,9,11}}
  \proofstepwidestar[1]{\bigl(\exists Y\in C\,P(Y)\bigr)\lor P(X)}{%
    \rEE{1,2,12}}
  \proofstepwidestar[]{\exists Y\in C\cup\{X\}\,P(Y)\to
    \bigl((\exists Y\in C\,P(Y))\lor P(X)\bigr)}{\rRI{1,13}}
  \proofstepwidestar[15]{\bigl(\exists Y\in C\,P(Y)\bigr)\lor P(X)}{\rA}
  \proofstepwidestar[16]{\exists Y\in C\,P(Y)}{\rA}
  \proofstepwidestar[17]{Y\in C\land P(Y)}{\rA}
  \proofstepwidestar[17]{Y\in C\cup\{X\}}{%
    \FormulaRefAuto{z\in A\vdash z\in A\cup B}{\rAEa{17}}}
  \proofstepwidestar[17]{\exists Y\in C\cup\{X\}\,P(Y)}{%
    \rEI{\rAI{18,\rAEb{17}}}}
  \proofstepwidestar[16]{\exists Y\in C\cup\{X\}\,P(Y)}{%
    \rEE{16,17,19}}
  \proofstepwidestar[21]{P(X)}{\rA}
  \proofstepwidestar[]{X\in C\cup\{X\}}{\FormulaRefAuto{a\in A\cup\{a\}}}
  \proofstepwidestar[21]{\exists Y\in C\cup\{X\}\,P(Y)}{%
    \rEI{\rAI{22,21}}}
  \proofstepwidestar[15]{\exists Y\in C\cup\{X\}\,P(Y)}{%
    \rOE{15,16,20,21,23}}
  \proofstepwidestar[]{\bigl((\exists Y\in C\,P(Y))\lor P(X)\bigr)
    \to\exists Y\in C\cup\{X\}\,P(Y)}{\rRI{15,24}}
  \proofstepwidestar[]{\exists Y\in C\cup\{X\}\,P(Y)
    \leftrightarrow\bigl((\exists Y\in C\,P(Y))\lor P(X)\bigr)}{%
    \FormulaRefAuto{P\leftrightarrow Q\eqvdash
      (P\to Q)\land(Q\to P)}{\rAI{14,25}}}
\end{tabproofwide}

\FormulaThmDeltaK[Beschränkter Allquantor über einer Einermenge]{%
  \forall Y\in\{X\}\,P(Y)\leftrightarrow P(X)
}{BoundedUniversalSingleton}{%
  \DeltaRow{Mengen}{X\dsep Y}
  \DeltaRow{Prädikatensymbole}{P}
}
\begin{tabproof}
  \proofstep{1}{\forall Y\in\{X\}\,P(Y)}{\rA}
  \proofstep{}{X\in\{X\}}{\FormulaRefAuto{a\in\{a\}}}
  \proofstep{1}{P(X)}{\rRE{2,\rUE{1}}}
  \proofstep{}{\forall Y\in\{X\}\,P(Y)\to P(X)}{\rRI{1,3}}
  \proofstep{5}{P(X)}{\rA}
  \proofstep{6}{Y\in\{X\}}{\rA}
  \proofstep{6}{Y=X}{\FormulaRefAuto{x\in\{a\}\eqvdash x=a}{6}}
  \proofstep{5,6}{P(Y)}{\rIE{\FormulaRefAuto{a=b\vdash b=a}{7},5}}
  \proofstep{5}{\forall Y\in\{X\}\,P(Y)}{\rUI{\rRI{6,8}}}
  \proofstep{}{P(X)\to\forall Y\in\{X\}\,P(Y)}{\rRI{5,9}}
  \proofstep{}{\forall Y\in\{X\}\,P(Y)\leftrightarrow P(X)}{%
    \FormulaRefAuto{P\leftrightarrow Q\eqvdash
      (P\to Q)\land(Q\to P)}{\rAI{4,10}}}
\end{tabproof}

\FormulaThmDeltaK[Beschränkter Allquantor über einer Adjunktion]{%
  \forall Y\in C\cup\{X\}\,P(Y)
  \leftrightarrow
  \bigl(\forall Y\in C\,P(Y)\bigr)\land P(X)
}{BoundedUniversalAdjunction}{%
  \DeltaRow{Mengen}{C\dsep X\dsep Y}
  \DeltaRow{Prädikatensymbole}{P}
}
\begin{tabproofwide}
  \proofstepwidestar[1]{\forall Y\in C\cup\{X\}\,P(Y)}{\rA}
  \proofstepwidestar[2]{Y\in C}{\rA}
  \proofstepwidestar[2]{Y\in C\cup\{X\}}{%
    \FormulaRefAuto{z\in A\vdash z\in A\cup B}{2}}
  \proofstepwidestar[1,2]{P(Y)}{\rRE{3,\rUE{1}}}
  \proofstepwidestar[1]{\forall Y\in C\,P(Y)}{\rUI{\rRI{2,4}}}
  \proofstepwidestar[]{X\in C\cup\{X\}}{\FormulaRefAuto{a\in A\cup\{a\}}}
  \proofstepwidestar[1]{P(X)}{\rRE{6,\rUE{1}}}
  \proofstepwidestar[1]{\bigl(\forall Y\in C\,P(Y)\bigr)\land P(X)}{%
    \rAI{5,7}}
  \proofstepwidestar[]{\forall Y\in C\cup\{X\}\,P(Y)\to
    \bigl((\forall Y\in C\,P(Y))\land P(X)\bigr)}{\rRI{1,8}}
  \proofstepwidestar[10]{\bigl(\forall Y\in C\,P(Y)\bigr)\land P(X)}{\rA}
  \proofstepwidestar[11]{Y\in C\cup\{X\}}{\rA}
  \proofstepwidestar[11]{Y\in C\lor Y=X}{%
    \FormulaRefAuto{x\in A\cup\{a\}\vdash x\in A\lor x=a}{11}}
  \proofstepwidestar[13]{Y\in C}{\rA}
  \proofstepwidestar[10,13]{P(Y)}{\rRE{13,\rUE{\rAEa{10}}}}
  \proofstepwidestar[15]{Y=X}{\rA}
  \proofstepwidestar[10,15]{P(Y)}{%
    \rIE{\FormulaRefAuto{a=b\vdash b=a}{15},\rAEb{10}}}
  \proofstepwidestar[10,11]{P(Y)}{\rOE{12,13,14,15,16}}
  \proofstepwidestar[10]{\forall Y\in C\cup\{X\}\,P(Y)}{%
    \rUI{\rRI{11,17}}}
  \proofstepwidestar[]{\bigl((\forall Y\in C\,P(Y))\land P(X)\bigr)
    \to\forall Y\in C\cup\{X\}\,P(Y)}{\rRI{10,18}}
  \proofstepwidestar[]{\forall Y\in C\cup\{X\}\,P(Y)
    \leftrightarrow\bigl((\forall Y\in C\,P(Y))\land P(X)\bigr)}{%
    \FormulaRefAuto{P\leftrightarrow Q\eqvdash
      (P\to Q)\land(Q\to P)}{\rAI{9,19}}}
\end{tabproofwide}

\FormulaThmDeltaK[Vereinigung einer adjungierten Mengenfamilie]{%
  \bigcup(C\cup\{X\})=(\bigcup C)\cup X
}{UnionFamilyAdjunction}{%
  \DeltaRow{Mengen}{C\dsep X\dsep Y\dsep a}
}
\begin{tabproofwide}
  \proofstepwide{a\in\bigcup(C\cup\{X\})}{\leftrightarrow}{%
    \exists Y\in C\cup\{X\}\,(a\in Y)}{%
    \FormulaRefAuto{x \in \bigcup A \;\eqvdash\;\exists B\in A\,(x \in B)}}
  \proofstepwide{} {\leftrightarrow}{%
    \bigl(\exists Y\in C\,(a\in Y)\bigr)\lor a\in X}{%
    \FormulaRefAuto{BoundedExistentialAdjunction}[thm]}
  \proofstepwide{} {\leftrightarrow}{a\in\bigcup C\lor a\in X}{%
    \rLRS{\FormulaRefAuto{x \in \bigcup A \;\eqvdash\;\exists B\in A\,(x \in B)},2}}
  \proofstepwide{} {\leftrightarrow}{a\in(\bigcup C)\cup X}{%
    \FormulaRefAuto{z \in A \cup B \eqvdash z \in A \lor z \in B}{3}}
  \proofstepwidestar{\bigcup(C\cup\{X\})=(\bigcup C)\cup X}{%
    \FormulaRefAuto{\forall x\,(x \in A \leftrightarrow x \in B) \eqvdash A = B}{\rUI{4}}}
\end{tabproofwide}

\FormulaThmDeltaK[Durchschnitt einer einelementigen Mengenfamilie]{%
  \bigcap\{X\}=X
}{IntersectionSingletonFamily}{%
  \DeltaRow{Mengen}{X\dsep a}
}
\begin{tabproofwide}
  \proofstepwide{a\in\bigcap\{X\}}{\leftrightarrow}{%
    \forall Y\in\{X\}\,(a\in Y)}{%
    \FormulaRefAuto{M\neq\varnothing \vdash x\in\bigcap M \leftrightarrow
      \forall X\in M\,(x\in X)}{\FormulaRefAuto{a\in\{a\}}}}
  \proofstepwide{} {\leftrightarrow}{a\in X}{%
    \FormulaRefAuto{BoundedUniversalSingleton}[thm]}
  \proofstepwidestar{\bigcap\{X\}=X}{%
    \FormulaRefAuto{\forall x\,(x \in A \leftrightarrow x \in B) \eqvdash A = B}{\rUI{2}}}
\end{tabproofwide}

\FormulaThmDeltaK[Durchschnitt einer adjungierten nichtleeren Mengenfamilie]{%
  C\neq\varnothing\vdash
  \bigcap(C\cup\{X\})=(\bigcap C)\cap X
}{IntersectionFamilyAdjunction}{%
  \DeltaRow{Mengen}{C\dsep X\dsep Y\dsep a}
}
\begin{tabproofwide}
  \proofstepwidestar[1]{C\neq\varnothing}{\rA}
  \proofstepwide[1]{a\in\bigcap(C\cup\{X\})}{\leftrightarrow}{%
    \forall Y\in C\cup\{X\}\,(a\in Y)}{%
    \FormulaRefAuto{M\neq\varnothing \vdash x\in\bigcap M \leftrightarrow
      \forall X\in M\,(x\in X)}{%
      \FormulaRefAuto{a\in A\vdash A\neq\varnothing}[thm]{%
        \FormulaRefAuto{a\in A\cup\{a\}}[thm]}}}
  \proofstepwide[1]{} {\leftrightarrow}{%
    \bigl(\forall Y\in C\,(a\in Y)\bigr)\land a\in X}{%
    \FormulaRefAuto{BoundedUniversalAdjunction}[thm]}
  \proofstepwide[1]{} {\leftrightarrow}{a\in\bigcap C\land a\in X}{%
    \rLRS{\FormulaRefAuto{M\neq\varnothing \vdash x\in\bigcap M \leftrightarrow
      \forall X\in M\,(x\in X)}{1},3}}
  \proofstepwide[1]{} {\leftrightarrow}{a\in(\bigcap C)\cap X}{%
    \FormulaRefAuto{x\in A\cap B\eqvdash x\in A\land x\in B}{4}}
  \proofstepwidestar[1]{\bigcap(C\cup\{X\})=(\bigcap C)\cap X}{%
    \FormulaRefAuto{\forall x\,(x \in A \leftrightarrow x \in B) \eqvdash A = B}{\rUI{5}}}
\end{tabproofwide}

\FormulaThmDeltaK[Der Familiendurchschnitt liegt in jedem Familienglied]{%
  M\neq\varnothing\dsep X\in M\vdash\bigcap M\subseteq X
}{FamilyIntersectionSubsetMember}{%
  \DeltaRow{Mengen}{M\dsep X\dsep a}
}
\begin{tabproof}
  \proofstep{1}{M\neq\varnothing}{\rA}
  \proofstep{2}{X\in M}{\rA}
  \proofstep{3}{a\in\bigcap M}{\rA}
  \proofstep{1,2,3}{a\in X}{%
    \FormulaRefAuto{M\neq\varnothing\dsep a\in\bigcap M\dsep
      X\in M\vdash a\in X}[thm]{1,3,2}}
  \proofstep{1,2}{\bigcap M\subseteq X}{%
    \FormulaRefAuto{A\subseteq B\coloneq
      \forall x\,(x\in A\rightarrow x\in B)}[def]{\rUI{\rRI{3,4}}}}
\end{tabproof}

\FormulaThmDeltaK[Endlicher nichtleerer Schnitt in einer schnittabgeschlossenen Familie]{%
  \Finite C\dsep C\neq\varnothing\dsep C\subseteq L\dsep
  \InterClosedFam L\vdash \bigcap C\in L
}{FiniteIntersectionClosedFamilyIntersection}{%
  \DeltaRow{Mengen}{C\dsep L\dsep X\dsep Y}
}
Setze für die Endlichkeitsinduktion
\[
  P(C)\quad:\Longleftrightarrow\quad
  \bigl(C\subseteq L\land C\neq\varnothing\to\bigcap C\in L\bigr).
\]
\begin{tabproofsplitwide}
  \proofpartwideindR[Induktionsanfang]{P(\varnothing)}{P(\varnothing)}%
    [FiniteIntersectionClosedBase]
    \proofstepwidestar[1]{\varnothing\subseteq L\land
      \varnothing\neq\varnothing}{\rA}
    \proofstepwidestar[]{\varnothing=\varnothing}{\rII}
    \proofstepwidestar[1]{\bot}{\rCE{2,\rAEb{1}}}
    \proofstepwidestar[]{P(\varnothing)}{\rRI{1,\FormulaRefAuto{\bot\vdash P}{3}}}
  \closeproofpartwideindR

  \proofpartwideindR[Induktionsschritt]{%
    \Finite C\dsep P(C)\dsep X\notin C\dsep\InterClosedFam L
    \vdash P(C\cup\{X\})
  }{%
    \Finite C\dsep P(C)\dsep X\notin C\dsep\InterClosedFam L
    \vdash P(C\cup\{X\})
  }[FiniteIntersectionClosedStep]
    \proofstepwidestar[1]{\Finite C}{\rA}
    \proofstepwidestar[2]{P(C)}{\rA}
    \proofstepwidestar[3]{X\notin C}{\rA}
    \proofstepwidestar[4]{\InterClosedFam L}{\rA}
    \proofstepwidestar[5]{%
      C\cup\{X\}\subseteq L\land C\cup\{X\}\neq\varnothing}{\rA}
    \proofstepwidestar[5]{X\in L}{%
      \FormulaRefAuto{A\subseteq B,\,x\in A\vdash x\in B}{%
        \rAEa{5},\FormulaRefAuto{a\in A\cup\{a\}}}}
    \proofstepwidestar[7]{C=\varnothing}{\rA}
    \proofstepwidestar[5,7]{C\cup\{X\}=\{X\}}{\rIE{7,\rII}}
    \proofstepwidestar[5,7]{\bigcap(C\cup\{X\})=X}{%
      \rIE{8,\FormulaRefAuto{IntersectionSingletonFamily}[thm]}}
    \proofstepwidestar[5,7]{\bigcap(C\cup\{X\})\in L}{\rIE{9,6}}
    \proofstepwidestar[11]{C\neq\varnothing}{\rA}
    \proofstepwidestar[5]{C\subseteq L}{%
      \FormulaRefAuto{A\subseteq B\dsep B\subseteq C\vdash A\subseteq C}{%
        \FormulaRefAuto{A\subseteq A\cup B},\rAEa{5}}}
    \proofstepwidestar[2,5,11]{\bigcap C\in L}{\rRE{2,\rAI{12,11}}}
    \proofstepwidestar[4,6,13]{(\bigcap C)\cap X\in L}{%
      \FormulaRefAuto{IntersectionClosedFamily}[def]{4,13,6}}
    \proofstepwidestar[11]{\bigcap(C\cup\{X\})=(\bigcap C)\cap X}{%
      \FormulaRefAuto{IntersectionFamilyAdjunction}[thm]{11}}
    \proofstepwidestar[2,4,5,11]{\bigcap(C\cup\{X\})\in L}{\rIE{15,14}}
    \proofstepwidestar[]{C=\varnothing\lor C\neq\varnothing}{%
      \FormulaRefAuto{P\lor\neg P}}
    \proofstepwidestar[2,4,5]{\bigcap(C\cup\{X\})\in L}{%
      \rOE{17,7,10,11,16}}
    \proofstepwidestar[1,2,3,4]{P(C\cup\{X\})}{\rRI{5,18}}
  \closeproofpartwideindR

  \proofpartwide{Induktionsschluss}
    \proofstepwidestar[1]{\Finite C}{\rA}
    \proofstepwidestar[2]{C\neq\varnothing}{\rA}
    \proofstepwidestar[3]{C\subseteq L}{\rA}
    \proofstepwidestar[4]{\InterClosedFam L}{\rA}
    \proofstepwidestar[1,4]{P(C)}{%
      \FormulaRefAuto{FiniteInductionRule}[ax]{IA,IS,1}}
    \proofstepwidestar[1,2,3,4]{\bigcap C\in L}{\rRE{5,\rAI{3,2}}}
  \closeproofpartwide
\end{tabproofsplitwide}

\FormulaThmDeltaK[Gesamtvereinigung einer endlichen vereinigungsabgeschlossenen Familie]{%
  \Finite M\dsep M\neq\varnothing\dsep\UnionClosedFam(M)
  \vdash\bigcup M\in M
}{FiniteUnionClosedFamilyTotalUnion}{%
  \DeltaRow{Mengen}{M\dsep C\dsep X}
}
Setze
\[
  Q(C)\quad:\Longleftrightarrow\quad
  \bigl(C\subseteq M\land C\neq\varnothing\to\bigcup C\in M\bigr).
\]
\begin{tabproofsplitwide}
  \proofpartwideindR[Induktionsanfang]{Q(\varnothing)}{Q(\varnothing)}%
    [FiniteUnionClosedBase]
    \proofstepwidestar[1]{\varnothing\subseteq M\land
      \varnothing\neq\varnothing}{\rA}
    \proofstepwidestar[]{\varnothing=\varnothing}{\rII}
    \proofstepwidestar[1]{\bot}{\rCE{2,\rAEb{1}}}
    \proofstepwidestar[]{Q(\varnothing)}{\rRI{1,\FormulaRefAuto{\bot\vdash P}{3}}}
  \closeproofpartwideindR

  \proofpartwideindR[Induktionsschritt]{%
    \Finite C\dsep Q(C)\dsep X\notin C\dsep\UnionClosedFam(M)
    \vdash Q(C\cup\{X\})
  }{%
    \Finite C\dsep Q(C)\dsep X\notin C\dsep\UnionClosedFam(M)
    \vdash Q(C\cup\{X\})
  }[FiniteUnionClosedStep]
    \proofstepwidestar[1]{\Finite C}{\rA}
    \proofstepwidestar[2]{Q(C)}{\rA}
    \proofstepwidestar[3]{X\notin C}{\rA}
    \proofstepwidestar[4]{\UnionClosedFam(M)}{\rA}
    \proofstepwidestar[5]{%
      C\cup\{X\}\subseteq M\land C\cup\{X\}\neq\varnothing}{\rA}
    \proofstepwidestar[5]{X\in M}{%
      \FormulaRefAuto{A\subseteq B,\,x\in A\vdash x\in B}{%
        \rAEa{5},\FormulaRefAuto{a\in A\cup\{a\}}}}
    \proofstepwidestar[7]{C=\varnothing}{\rA}
    \proofstepwidestar[5,7]{\bigcup(C\cup\{X\})=X}{%
      \rIE{7,\FormulaRefAuto{UnionFamilyAdjunction}[thm]}}
    \proofstepwidestar[5,7]{\bigcup(C\cup\{X\})\in M}{\rIE{8,6}}
    \proofstepwidestar[10]{C\neq\varnothing}{\rA}
    \proofstepwidestar[5]{C\subseteq M}{%
      \FormulaRefAuto{A\subseteq B\dsep B\subseteq C\vdash A\subseteq C}{%
        \FormulaRefAuto{A\subseteq A\cup B},\rAEa{5}}}
    \proofstepwidestar[2,5,10]{\bigcup C\in M}{\rRE{2,\rAI{11,10}}}
    \proofstepwidestar[4,6,12]{(\bigcup C)\cup X\in M}{%
      \FormulaRefAuto{X\in M\dsep Y\in M\vdash X\cup Y\in M}[ax]{12,6}}
    \proofstepwidestar[]{\bigcup(C\cup\{X\})=(\bigcup C)\cup X}{%
      \FormulaRefAuto{UnionFamilyAdjunction}[thm]}
    \proofstepwidestar[2,4,5,10]{\bigcup(C\cup\{X\})\in M}{\rIE{14,13}}
    \proofstepwidestar[]{C=\varnothing\lor C\neq\varnothing}{%
      \FormulaRefAuto{P\lor\neg P}}
    \proofstepwidestar[2,4,5]{\bigcup(C\cup\{X\})\in M}{%
      \rOE{16,7,9,10,15}}
    \proofstepwidestar[1,2,3,4]{Q(C\cup\{X\})}{\rRI{5,17}}
  \closeproofpartwideindR

  \proofpartwide{Induktionsschluss}
    \proofstepwidestar[1]{\Finite M}{\rA}
    \proofstepwidestar[2]{M\neq\varnothing}{\rA}
    \proofstepwidestar[3]{\UnionClosedFam(M)}{\rA}
    \proofstepwidestar[1,3]{Q(M)}{%
      \FormulaRefAuto{FiniteInductionRule}[ax]{IA,IS,1}}
    \proofstepwidestar[]{M\subseteq M}{\FormulaRefAuto{A\subseteq A}}
    \proofstepwidestar[1,2,3]{\bigcup M\in M}{\rRE{4,\rAI{5,2}}}
  \closeproofpartwide
\end{tabproofsplitwide}

\FormulaThmDeltaK[de-Morgan-Gesetz für relative Komplemente]{%
  (U\setminus X)\cap(U\setminus Y)=U\setminus(X\cup Y)
}{RelativeComplementUnionDeMorgan}{%
  \DeltaRow{Mengen}{U\dsep X\dsep Y\dsep a}
}
\begin{tabproofwide}
  \proofstepwide{a\in(U\setminus X)\cap(U\setminus Y)}{\leftrightarrow}{%
    a\in U\land a\notin X\land a\notin Y}{%
    \FormulaRefAuto{x\in A\cap B\eqvdash x\in A\land x\in B};
    \FormulaRefAuto{x \in A \setminus B \eqvdash x \in A \land x \notin B}}
  \proofstepwide{} {\leftrightarrow}{a\in U\land a\notin X\cup Y}{%
    \FormulaRefAuto{z \in A \cup B \eqvdash z \in A \lor z \in B};
    \FormulaRefAuto{\neg(P\lor Q)\eqvdash\neg P\land\neg Q}}
  \proofstepwide{} {\leftrightarrow}{a\in U\setminus(X\cup Y)}{%
    \FormulaRefAuto{x \in A \setminus B \eqvdash x \in A \land x \notin B}{2}}
  \proofstepwidestar{(U\setminus X)\cap(U\setminus Y)=U\setminus(X\cup Y)}{%
    \FormulaRefAuto{\forall x\,(x \in A \leftrightarrow x \in B) \eqvdash A = B}{\rUI{3}}}
\end{tabproofwide}

\FormulaThmDeltaK[Differenz mit der leeren Menge]{%
  A\setminus\varnothing=A
}{RelativeComplementEmpty}{%
  \DeltaRow{Mengen}{A\dsep x}
}
\begin{tabproofwide}
  \proofstepwidestar[1]{x\in A\setminus\varnothing}{\rA}
  \proofstepwidestar[1]{x\in A}{%
    \FormulaRefAuto{x\in A\setminus B\vdash x\in A}[thm]{1}}
  \proofstepwidestar[]{x\in A\setminus\varnothing\to x\in A}{\rRI{1,2}}
  \proofstepwidestar[4]{x\in A}{\rA}
  \proofstepwidestar[]{x\notin\varnothing}{%
    \FormulaRefAuto{\varnothing \coloneq
      \iota O\bigl(\forall x\,(x \not\in O)\bigr)}[def]}
  \proofstepwidestar[4]{x\in A\setminus\varnothing}{%
    \FormulaRefAuto{x\in A\dsep x\notin B\vdash x\in A\setminus B}[thm]{4,5}}
  \proofstepwidestar[]{x\in A\to x\in A\setminus\varnothing}{\rRI{4,6}}
  \proofstepwidestar[]{x\in A\setminus\varnothing\leftrightarrow x\in A}{%
    \rLRI{3,7}}
  \proofstepwidestar[]{A\setminus\varnothing=A}{%
    \FormulaRefAuto{\forall x\,(x\in A\leftrightarrow x\in B)
      \eqvdash A=B}[thm]{\rUI{8}}}
\end{tabproofwide}

\FormulaThmDeltaK[Differenz einer Menge mit sich selbst]{%
  A\setminus A=\varnothing
}{RelativeComplementSelf}{%
  \DeltaRow{Mengen}{A\dsep x}
}
\begin{tabproofwide}
  \proofstepwidestar[1]{x\in A\setminus A}{\rA}
  \proofstepwidestar[1]{x\in A\land x\notin A}{%
    \FormulaRefAuto{x\in A\setminus B\eqvdash
      x\in A\land x\notin B}[thm]{1}}
  \proofstepwidestar[1]{\bot}{\rBI{\rAEa{2},\rAEb{2}}}
  \proofstepwidestar[1]{x\in\varnothing}{%
    \FormulaRefAuto{\bot\vdash P}[thm]{3}}
  \proofstepwidestar[]{x\in A\setminus A\to x\in\varnothing}{\rRI{1,4}}
  \proofstepwidestar[6]{x\in\varnothing}{\rA}
  \proofstepwidestar[]{x\notin\varnothing}{%
    \FormulaRefAuto{\varnothing \coloneq
      \iota O\bigl(\forall x\,(x \not\in O)\bigr)}[def]}
  \proofstepwidestar[6]{\bot}{\rBI{6,7}}
  \proofstepwidestar[6]{x\in A\setminus A}{%
    \FormulaRefAuto{\bot\vdash P}[thm]{8}}
  \proofstepwidestar[]{x\in\varnothing\to x\in A\setminus A}{\rRI{6,9}}
  \proofstepwidestar[]{x\in A\setminus A\leftrightarrow x\in\varnothing}{%
    \rLRI{5,10}}
  \proofstepwidestar[]{A\setminus A=\varnothing}{%
    \FormulaRefAuto{\forall x\,(x\in A\leftrightarrow x\in B)
      \eqvdash A=B}[thm]{\rUI{11}}}
\end{tabproofwide}

\FormulaDefDeltaK[Komplementabbildung der dualen Konstruktion]{%
  X\in M^\circ\vdash \kappa_M(X)\coloneqq U_M\setminus X
}{FranklComplementMap}{%
  \DeltaRow{Mengen}{M\dsep X}
  \DeltaRow{Funktionen}{\kappa_M}
  \DeltaRow{\textbf{Neues Symbol}}{}
  \DeltaPrem{Komplementabbildung}{\kappa_M}
}

\FormulaThmDeltaK[Die Komplementabbildung ist wohldefiniert]{%
  \FranklFam M\vdash \kappa_M\colon M^\circ\to\mathcal L_M
}{FranklComplementMapFunction}{%
  \DeltaRow{Mengen}{M\dsep X\dsep Y\dsep Z}
  \DeltaRow{Funktionen}{\kappa_M}
}
\begin{tabproofwide}
  \proofstepwidestar[1]{\FranklFam M}{\rA}
  \proofstepwidestar[2]{X\in M^\circ}{\rA}
  \proofstepwidestar[2]{X\subseteq U_M}{%
    \FormulaRefAuto{FranklAugmentedMemberSubsetCarrier}[thm]{1,2}}
  \proofstepwidestar[2]{\kappa_M(X)=U_M\setminus X}{%
    \FormulaRefAuto{FranklComplementMap}[def]{2}}
  \proofstepwidestar[2]{\kappa_M(X)\in\mathcal L_M}{%
    \FormulaRefAuto{FranklComplementConstruction}[def]{2,4}}
  \proofstepwidestar[]{\kappa_M\colon M^\circ\to\mathcal L_M}{%
    \FormulaRefAuto{Funktion}[def]{2,5}}
\end{tabproofwide}

\FormulaThmDeltaK[Die Komplementabbildung ist injektiv]{%
  \FranklFam M\vdash
  \forall X,Y\in M^\circ\,
    (\kappa_M(X)=\kappa_M(Y)\to X=Y)
}{FranklComplementMapInjective}{%
  \DeltaRow{Mengen}{M\dsep X\dsep Y}
  \DeltaRow{Funktionen}{\kappa_M}
}
\begin{tabproofwide}
  \proofstepwidestar[1]{\FranklFam M}{\rA}
  \proofstepwidestar[2]{X,Y\in M^\circ\land\kappa_M(X)=\kappa_M(Y)}{\rA}
  \proofstepwidestar[1,2]{X\subseteq U_M}{%
    \FormulaRefAuto{FranklAugmentedMemberSubsetCarrier}[thm]{1,\rAEa{2}}}
  \proofstepwidestar[1,2]{Y\subseteq U_M}{%
    \FormulaRefAuto{FranklAugmentedMemberSubsetCarrier}[thm]{1,\rAEa{\rAEb{2}}}}
  \proofstepwidestar[2]{\kappa_M(X)=U_M\setminus X}{%
    \FormulaRefAuto{FranklComplementMap}[def]{\rAEa{2}}}
  \proofstepwidestar[2]{\kappa_M(Y)=U_M\setminus Y}{%
    \FormulaRefAuto{FranklComplementMap}[def]{\rAEa{\rAEb{2}}}}
  \proofstepwidestar[2]{\kappa_M(X)=U_M\setminus Y}{%
    \rIE{6,\rAEb{\rAEb{2}}}}
  \proofstepwidestar[2]{U_M\setminus X=U_M\setminus Y}{\rIE{5,7}}
  \proofstepwidestar[2]{U_M\setminus Y=U_M\setminus X}{%
    \FormulaRefAuto{a=b\vdash b=a}{8}}
  \proofstepwidestar[1,2]{X=Y}{%
    \FormulaRefAuto{RelCompEqualityCancel}[thm]{3,4,9}}
  \proofstepwidestar[1]{\forall X,Y\in M^\circ
    (\kappa_M(X)=\kappa_M(Y)\to X=Y)}{%
    \rUI{\rUI{\rRI{2,10}}}}
\end{tabproofwide}

\FormulaThmDeltaK[Die Komplementabbildung ist surjektiv]{%
  \FranklFam M\vdash
  \forall Z\in\mathcal L_M\,\exists X\in M^\circ\,(\kappa_M(X)=Z)
}{FranklComplementMapSurjective}{%
  \DeltaRow{Mengen}{M\dsep X\dsep Z}
  \DeltaRow{Funktionen}{\kappa_M}
}
\begin{tabproofwide}
  \proofstepwidestar[1]{\FranklFam M}{\rA}
  \proofstepwidestar[2]{Z\in\mathcal L_M}{\rA}
  \proofstepwidestar[2]{\exists X\in M^\circ\,(Z=U_M\setminus X)}{%
    \FormulaRefAuto{FranklComplementConstruction}[def]{2}}
  \proofstepwidestar[4]{X\in M^\circ\land Z=U_M\setminus X}{\rA}
  \proofstepwidestar[4]{\kappa_M(X)=U_M\setminus X}{%
    \FormulaRefAuto{FranklComplementMap}[def]{\rAEa{4}}}
  \proofstepwidestar[4]{U_M\setminus X=Z}{%
    \FormulaRefAuto{a=b\vdash b=a}{\rAEb{4}}}
  \proofstepwidestar[4]{\kappa_M(X)=Z}{%
    \FormulaRefAuto{a=b,\,b=c\vdash a=c}{5,6}}
  \proofstepwidestar[2]{\exists X\in M^\circ\,(\kappa_M(X)=Z)}{%
    \rEE{3,4,\rEI{7}}}
  \proofstepwidestar[1]{\forall Z\in\mathcal L_M\exists X\in M^\circ
    (\kappa_M(X)=Z)}{\rUI{\rRI{2,8}}}
\end{tabproofwide}

\FormulaThmDeltaK[Die Komplementabbildung ist bijektiv]{%
  \FranklFam M\vdash \kappa_M\colon M^\circ\bij\mathcal L_M
}{FranklComplementMapBijective}{%
  \DeltaRow{Mengen}{M\dsep X\dsep Y\dsep Z}
  \DeltaRow{Funktionen}{\kappa_M}
}
\begin{tabproofwide}
  \proofstepwidestar[1]{\FranklFam M}{\rA}
  \proofstepwidestar[1]{\kappa_M\colon M^\circ\to\mathcal L_M}{%
    \FormulaRefAuto{FranklComplementMapFunction}[thm]{1}}
  \proofstepwidestar[1]{\forall X,Y\in M^\circ
    (\kappa_M(X)=\kappa_M(Y)\to X=Y)}{%
    \FormulaRefAuto{FranklComplementMapInjective}[thm]{1}}
  \proofstepwidestar[1]{\forall Z\in\mathcal L_M\exists X\in M^\circ
    (\kappa_M(X)=Z)}{%
    \FormulaRefAuto{FranklComplementMapSurjective}[thm]{1}}
  \proofstepwidestar[1]{\kappa_M\colon M^\circ\bij\mathcal L_M}{%
    \FormulaRefAuto{Bijektive Funktion}[def]{2,3,4}}
\end{tabproofwide}

\FormulaThmDeltaK[Mitglieder der erweiterten Familie liegen im Träger]{%
  \FranklFam M\dsep X\in M^\circ\vdash X\subseteq U_M
}{FranklAugmentedMemberSubsetCarrier}{%
  \DeltaRow{Mengen}{M\dsep X}
}
\begin{tabproof}
  \proofstep{1}{\FranklFam M}{\rA}
  \proofstep{2}{X\in M^\circ}{\rA}
  \proofstep{1}{M^\circ=M\cup\{\varnothing\}}{%
    \FormulaRefAuto{FranklComplementConstruction}[def]{1}}
  \proofstep{2}{X\in M\lor X=\varnothing}{%
    \FormulaRefAuto{x\in A\cup\{a\} \vdash x\in A\lor x=a}{\rIE{3,2}}}
  \proofstep{5}{X\in M}{\rA}
  \proofstep{5}{X\subseteq U_M}{\FormulaRefAuto{FamilyMemberSubsetUnion}[thm]{5}}
  \proofstep{7}{X=\varnothing}{\rA}
  \proofstep{7}{X\subseteq U_M}{%
    \rIE{7,\FormulaRefAuto{\varnothing\subseteq A}}}
  \proofstep{1,2}{X\subseteq U_M}{\rOE{4,5,6,7,8}}
\end{tabproof}

\FormulaThmDeltaK[Mitglieder der Komplementfamilie liegen im Träger]{%
  \FranklFam M\dsep X\in\mathcal L_M\vdash X\subseteq U_M
}{FranklComplementMemberSubsetCarrier}{%
  \DeltaRow{Mengen}{M\dsep X\dsep Y}
}
\begin{tabproof}
  \proofstep{1}{\FranklFam M}{\rA}
  \proofstep{2}{X\in\mathcal L_M}{\rA}
  \proofstep{2}{\exists Y\in M^\circ\,(X=U_M\setminus Y)}{%
    \FormulaRefAuto{FranklComplementConstruction}[def]{2}}
  \proofstep{4}{Y\in M^\circ\land X=U_M\setminus Y}{\rA}
  \proofstep{4}{X\subseteq U_M}{%
    \rIE{\rAEb{4},\FormulaRefAuto{A\setminus B\subseteq A}}}
  \proofstep{2}{X\subseteq U_M}{\rEE{3,4,5}}
\end{tabproof}

\FormulaThmDeltaK[Die erweiterte Familie ist endlich]{%
  \FranklFam M\vdash\Finite{M^\circ}
}{FranklAugmentedFamilyFinite}{%
  \DeltaRow{Mengen}{M}
}
\begin{tabproof}
  \proofstep{1}{\FranklFam M}{\rA}
  \proofstep{1}{\Finite M}{\FormulaRefAuto{FranklFamFinite}[thm]{1}}
  \proofstep{1}{\Finite{M\cup\{\varnothing\}}}{%
    \FormulaRefAuto{FiniteInsert}[thm]{2}}
  \proofstep{1}{M^\circ=M\cup\{\varnothing\}}{%
    \FormulaRefAuto{FranklComplementConstruction}[def]{1}}
  \proofstep{1}{\Finite{M^\circ}}{\rIE{4,3}}
\end{tabproof}

\FormulaThmDeltaK[Monotonie der vermeidenden Teilfamilie]{%
  A\subseteq B\vdash
  \FamAvoiding{A}{a}\subseteq\FamAvoiding{B}{a}
}{FamAvoidingMonotone}{%
  \DeltaRow{Mengen}{A\dsep B\dsep X\dsep a}
}
\begin{tabproof}
  \proofstep{1}{A\subseteq B}{\rA}
  \proofstep{2}{X\in\FamAvoiding{A}{a}}{\rA}
  \proofstep{2}{X\in A\land a\notin X}{%
    \FormulaRefAuto{\FamAvoiding{M}{a}\coloneqq
      \{X\in M\mid a\notin X\}}[def]{2}}
  \proofstep{1,2}{X\in B}{%
    \FormulaRefAuto{A\subseteq B\dsep x\in A\vdash x\in B}[thm]{%
      1,\rAEa{3}}}
  \proofstep{1,2}{X\in\FamAvoiding{B}{a}}{%
    \FormulaRefAuto{\FamAvoiding{M}{a}\coloneqq
      \{X\in M\mid a\notin X\}}[def]{\rAI{4,\rAEb{3}}}}
  \proofstep{1}{\FamAvoiding{A}{a}\subseteq\FamAvoiding{B}{a}}{%
    \FormulaRefAuto{A\subseteq B\coloneq
      \forall x\,(x\in A\rightarrow x\in B)}[def]{\rUI{\rRI{2,5}}}}
\end{tabproof}

\FormulaThmDeltaK[Adjunktion der leeren Menge erhält die enthaltende Teilfamilie]{%
  \FranklFam M\vdash
  \FamContaining{M^\circ}{a}=\FamContaining{M}{a}
}{FranklAugmentedContainingUnchanged}{%
  \DeltaRow{Mengen}{M\dsep X\dsep a}
}
\begin{tabproofsplitwide}
  \proofpartwideindR[Teilmengenrichtung zur ursprünglichen Familie]{%
    \FranklFam M\vdash
    \FamContaining{M^\circ}{a}\subseteq\FamContaining{M}{a}
  }{%
    \FranklFam M\vdash
    \FamContaining{M^\circ}{a}\subseteq\FamContaining{M}{a}
  }[FranklAugmentedContainingSubsetOriginal]
    \proofstepwidestar[1]{\FranklFam M}{\rA}
    \proofstepwidestar[2]{X\in\FamContaining{M^\circ}{a}}{\rA}
    \proofstepwidestar[2]{X\in M^\circ\land a\in X}{%
      \FormulaRefAuto{\FamContaining{M}{a}\coloneqq
        \{X\in M\mid a\in X\}}[def]{2}}
    \proofstepwidestar[1]{M^\circ=M\cup\{\varnothing\}}{%
      \FormulaRefAuto{FranklComplementConstruction}[def]{1}}
    \proofstepwidestar[1,2]{X\in M\cup\{\varnothing\}}{%
      \rIE{4,\rAEa{3}}}
    \proofstepwidestar[1,2]{X\in M\lor X=\varnothing}{%
      \FormulaRefAuto{x\in A\cup\{a\}\vdash x\in A\lor x=a}[thm]{5}}
    \proofstepwidestar[7]{X\in M}{\rA}
    \proofstepwidestar[2,7]{X\in\FamContaining{M}{a}}{%
      \FormulaRefAuto{\FamContaining{M}{a}\coloneqq
        \{X\in M\mid a\in X\}}[def]{\rAI{7,\rAEb{3}}}}
    \proofstepwidestar[9]{X=\varnothing}{\rA}
    \proofstepwidestar[2,9]{a\in\varnothing}{\rIE{9,\rAEb{3}}}
    \proofstepwidestar[]{a\notin\varnothing}{%
      \FormulaRefAuto{\varnothing \coloneq
        \iota O\bigl(\forall x\,(x \not\in O)\bigr)}[def]}
    \proofstepwidestar[2,9]{\bot}{\rBI{10,11}}
    \proofstepwidestar[2,9]{X\in\FamContaining{M}{a}}{%
      \FormulaRefAuto{\bot\vdash P}[thm]{12}}
    \proofstepwidestar[1,2]{X\in\FamContaining{M}{a}}{%
      \rOE{6,7,8,9,13}}
    \proofstepwidestar[1]{%
      \FamContaining{M^\circ}{a}\subseteq\FamContaining{M}{a}}{%
      \FormulaRefAuto{A\subseteq B\coloneq
        \forall x\,(x\in A\rightarrow x\in B)}[def]{\rUI{\rRI{2,14}}}}
  \closeproofpartwideindR

  \proofpartwideindR[Teilmengenrichtung zur erweiterten Familie]{%
    \FranklFam M\vdash
    \FamContaining{M}{a}\subseteq\FamContaining{M^\circ}{a}
  }{%
    \FranklFam M\vdash
    \FamContaining{M}{a}\subseteq\FamContaining{M^\circ}{a}
  }[FranklAugmentedContainingSubsetAugmented]
    \proofstepwidestar[1]{\FranklFam M}{\rA}
    \proofstepwidestar[2]{X\in\FamContaining{M}{a}}{\rA}
    \proofstepwidestar[2]{X\in M\land a\in X}{%
      \FormulaRefAuto{\FamContaining{M}{a}\coloneqq
        \{X\in M\mid a\in X\}}[def]{2}}
    \proofstepwidestar[1]{M^\circ=M\cup\{\varnothing\}}{%
      \FormulaRefAuto{FranklComplementConstruction}[def]{1}}
    \proofstepwidestar[2]{X\in M\cup\{\varnothing\}}{%
      \FormulaRefAuto{z\in A\vdash z\in A\cup B}[thm]{\rAEa{3}}}
    \proofstepwidestar[1,2]{X\in M^\circ}{\rIE{4,5}}
    \proofstepwidestar[1,2]{X\in\FamContaining{M^\circ}{a}}{%
      \FormulaRefAuto{\FamContaining{M}{a}\coloneqq
        \{X\in M\mid a\in X\}}[def]{\rAI{6,\rAEb{3}}}}
    \proofstepwidestar[1]{%
      \FamContaining{M}{a}\subseteq\FamContaining{M^\circ}{a}}{%
      \FormulaRefAuto{A\subseteq B\coloneq
        \forall x\,(x\in A\rightarrow x\in B)}[def]{\rUI{\rRI{2,7}}}}
  \closeproofpartwideindR

  \proofpartwide{Schluss}
    \proofstepwidestar[1]{%
      \FamContaining{M^\circ}{a}\subseteq\FamContaining{M}{a}}{%
      \FormulaRefAuto{FranklAugmentedContainingSubsetOriginal}[thm]{1}}
    \proofstepwidestar[1]{%
      \FamContaining{M}{a}\subseteq\FamContaining{M^\circ}{a}}{%
      \FormulaRefAuto{FranklAugmentedContainingSubsetAugmented}[thm]{1}}
    \proofstepwidestar[1]{%
      \FamContaining{M^\circ}{a}=\FamContaining{M}{a}}{%
      \FormulaRefAuto{A\subseteq B\dsep B\subseteq A\vdash A=B}[thm]{1,2}}
  \closeproofpartwide
\end{tabproofsplitwide}

\FormulaThmDeltaK[Die erweiterte Familie ist vereinigungsabgeschlossen]{%
  \FranklFam M\vdash\UnionClosedFam(M^\circ)
}{FranklAugmentedFamilyUnionClosed}{%
  \DeltaRow{Mengen}{M\dsep X\dsep Y}
}
\begin{tabproofwide}
  \proofstepwidestar[1]{\FranklFam M}{\rA}
  \proofstepwidestar[1]{\UnionClosedFam(M)}{%
    \FormulaRefAuto{FranklFamUnionClosed}[thm]{1}}
  \proofstepwidestar[]{M^\circ=M\cup\{\varnothing\}}{%
    \FormulaRefAuto{FranklComplementConstruction}[def]}
  \proofstepwidestar[4]{X,Y\in M^\circ}{\rA}
  \proofstepwidestar[4]{X\in M\lor X=\varnothing}{%
    \FormulaRefAuto{x\in A\cup\{a\} \vdash x\in A\lor x=a}{%
      \rIE{3,\rAEa{4}}}}
  \proofstepwidestar[4]{Y\in M\lor Y=\varnothing}{%
    \FormulaRefAuto{x\in A\cup\{a\} \vdash x\in A\lor x=a}{%
      \rIE{3,\rAEb{4}}}}

  \proofstepwidestar[7]{X\in M}{\rA}
  \proofstepwidestar[8]{Y\in M}{\rA}
  \proofstepwidestar[2,7,8]{X\cup Y\in M}{%
    \FormulaRefAuto{X\in M\dsep Y\in M\vdash X\cup Y\in M}[ax]{7,8}}
  \proofstepwidestar[1,7,8]{X\cup Y\in M\cup\{\varnothing\}}{%
    \FormulaRefAuto{z\in A\vdash z\in A\cup B}[thm]{9}}
  \proofstepwidestar[1,7,8]{X\cup Y\in M^\circ}{\rIE{3,10}}
  \proofstepwidestar[12]{Y=\varnothing}{\rA}
  \proofstepwidestar[12]{X\cup Y=X}{%
    \rIE{12,\FormulaRefAuto{A\cup\varnothing=A}}}
  \proofstepwidestar[7]{X\in M\cup\{\varnothing\}}{%
    \FormulaRefAuto{z\in A\vdash z\in A\cup B}[thm]{7}}
  \proofstepwidestar[7]{X\in M^\circ}{\rIE{3,14}}
  \proofstepwidestar[7,12]{X\cup Y\in M^\circ}{\rIE{13,15}}
  \proofstepwidestar[1,4,6,7]{X\cup Y\in M^\circ}{%
    \rOE{6,8,11,12,16}}

  \proofstepwidestar[18]{X=\varnothing}{\rA}
  \proofstepwidestar[19]{Y\in M}{\rA}
  \proofstepwidestar[18]{X\cup Y=Y}{%
    \rIE{18,\FormulaRefAuto{\varnothing\cup A=A}}}
  \proofstepwidestar[19]{Y\in M\cup\{\varnothing\}}{%
    \FormulaRefAuto{z\in A\vdash z\in A\cup B}[thm]{19}}
  \proofstepwidestar[19]{Y\in M^\circ}{\rIE{3,21}}
  \proofstepwidestar[18,19]{X\cup Y\in M^\circ}{\rIE{20,22}}
  \proofstepwidestar[24]{Y=\varnothing}{\rA}
  \proofstepwidestar[18,24]{X\cup Y=\varnothing}{%
    \rIE{18,\rIE{24,\FormulaRefAuto{\varnothing\cup\varnothing=\varnothing}}}}
  \proofstepwidestar[]{\varnothing\in M\cup\{\varnothing\}}{%
    \FormulaRefAuto{a\in A\cup\{a\}}[thm]}
  \proofstepwidestar[]{\varnothing\in M^\circ}{\rIE{3,26}}
  \proofstepwidestar[18,24]{X\cup Y\in M^\circ}{\rIE{25,27}}
  \proofstepwidestar[4,6,18]{X\cup Y\in M^\circ}{%
    \rOE{6,19,23,24,28}}
  \proofstepwidestar[1,4]{X\cup Y\in M^\circ}{%
    \rOE{5,7,17,18,29}}
  \proofstepwidestar[1]{\UnionClosedFam(M^\circ)}{%
    \FormulaRefAuto{Vereinigungsabgeschlossene Familie}[def]{%
      \rUI{\rUI{\rRI{4,30}}}}}
\end{tabproofwide}

\FormulaThmDeltaK[Die komplementduale Familie ist endlich]{%
  \FranklFam M\vdash\Finite{\mathcal L_M}
}{FranklComplementFamilyFinite}{%
  \DeltaRow{Mengen}{M}
  \DeltaRow{Funktionen}{\kappa_M}
}
\begin{tabproof}
  \proofstep{1}{\FranklFam M}{\rA}
  \proofstep{1}{\Finite{M^\circ}}{%
    \FormulaRefAuto{FranklAugmentedFamilyFinite}[thm]{1}}
  \proofstep{1}{\kappa_M\colon M^\circ\bij\mathcal L_M}{%
    \FormulaRefAuto{FranklComplementMapBijective}[thm]{1}}
  \proofstep{1}{\Finite{\mathcal L_M}}{%
    \FormulaRefAuto{FiniteBijectionTransport}[thm]{2,3}}
\end{tabproof}

\FormulaThmDeltaK[Die komplementduale Familie ist schnittabgeschlossen]{%
  \FranklFam M\vdash\InterClosedFam{\mathcal L_M}
}{FranklComplementFamilyIntersectionClosed}{%
  \DeltaRow{Mengen}{M\dsep X\dsep Y\dsep A\dsep B}
}
\begin{tabproofwide}
  \proofstepwidestar[1]{\FranklFam M}{\rA}
  \proofstepwidestar[2]{X,Y\in\mathcal L_M}{\rA}
  \proofstepwidestar[1]{\UnionClosedFam(M^\circ)}{%
    \FormulaRefAuto{FranklAugmentedFamilyUnionClosed}[thm]{1}}
  \proofstepwidestar[2]{\exists A\in M^\circ\,(X=U_M\setminus A)}{%
    \FormulaRefAuto{FranklComplementConstruction}[def]{\rAEa{2}}}
  \proofstepwidestar[2]{\exists B\in M^\circ\,(Y=U_M\setminus B)}{%
    \FormulaRefAuto{FranklComplementConstruction}[def]{\rAEb{2}}}
  \proofstepwidestar[6]{A\in M^\circ\land X=U_M\setminus A}{\rA}
  \proofstepwidestar[7]{B\in M^\circ\land Y=U_M\setminus B}{\rA}
  \proofstepwidestar[1,6,7]{A\cup B\in M^\circ}{%
    \FormulaRefAuto{X\in M\dsep Y\in M\vdash X\cup Y\in M}[ax]{%
      \rAEa{6},\rAEa{7}}}
  \proofstepwidestar[6,7]{X\cap Y=U_M\setminus(A\cup B)}{%
    \rIE{\rAEb{6},\rIE{\rAEb{7},
      \FormulaRefAuto{RelativeComplementUnionDeMorgan}[thm]}}}
  \proofstepwidestar[1,2,6,7]{X\cap Y\in\mathcal L_M}{%
    \FormulaRefAuto{FranklComplementConstruction}[def]{8,9}}
  \proofstepwidestar[1,2]{X\cap Y\in\mathcal L_M}{%
    \rEE{4,6,\rEE{5,7,10}}}
  \proofstepwidestar[1]{\InterClosedFam{\mathcal L_M}}{%
    \FormulaRefAuto{IntersectionClosedFamily}[def]{\rUI{\rUI{\rRI{2,11}}}}}
\end{tabproofwide}

\FormulaThmDeltaK[Die explizite Komplementordnung ist partiell]{%
  \FranklFam M\vdash\PartOrd{\mathcal R_M,\mathcal L_M}
}{FranklComplementOrderPartial}{%
  \DeltaRow{Mengen}{M\dsep X\dsep Y\dsep Z}
  \DeltaRow{Relationen}{\mathcal R_M}
}
\begin{tabproofwide}
  \proofstepwidestar[1]{\FranklFam M}{\rA}
  \proofstepwidestar[]{\mathcal R_M\subseteq
    \mathcal L_M\times\mathcal L_M}{%
    \FormulaRefAuto{FranklComplementConstruction}[def]}
  \proofstepwidestar[3]{X\in\mathcal L_M}{\rA}
  \proofstepwidestar[3]{(X,X)\in\mathcal R_M}{%
    \FormulaRefAuto{FranklComplementConstruction}[def]{3,
      \FormulaRefAuto{A\subseteq A}}}
  \proofstepwidestar[]{X\in\mathcal L_M\vdash(X,X)\in\mathcal R_M}{%
    \rRI{3,4}}
  \proofstepwidestar[6]{(X,Y)\in\mathcal R_M}{\rA}
  \proofstepwidestar[7]{(Y,Z)\in\mathcal R_M}{\rA}
  \proofstepwidestar[6]{X\subseteq Y}{%
    \FormulaRefAuto{FranklComplementConstruction}[def]{6}}
  \proofstepwidestar[7]{Y\subseteq Z}{%
    \FormulaRefAuto{FranklComplementConstruction}[def]{7}}
  \proofstepwidestar[6,7]{X\subseteq Z}{%
    \FormulaRefAuto{A\subseteq B\dsep B\subseteq C\vdash A\subseteq C}{8,9}}
  \proofstepwidestar[6,7]{(X,Z)\in\mathcal R_M}{%
    \FormulaRefAuto{FranklComplementConstruction}[def]{10}}
  \proofstepwidestar[]{(X,Y)\in\mathcal R_M\dsep(Y,Z)\in\mathcal R_M
    \vdash(X,Z)\in\mathcal R_M}{\rRI{6,\rRI{7,11}}}
  \proofstepwidestar[13]{(X,Y)\in\mathcal R_M}{\rA}
  \proofstepwidestar[14]{(Y,X)\in\mathcal R_M}{\rA}
  \proofstepwidestar[13]{X\subseteq Y}{%
    \FormulaRefAuto{FranklComplementConstruction}[def]{13}}
  \proofstepwidestar[14]{Y\subseteq X}{%
    \FormulaRefAuto{FranklComplementConstruction}[def]{14}}
  \proofstepwidestar[13,14]{X=Y}{%
    \FormulaRefAuto{A \subseteq B \land B \subseteq A \eqvdash A = B}{15,16}}
  \proofstepwidestar[]{(X,Y)\in\mathcal R_M\dsep(Y,X)\in\mathcal R_M
    \vdash X=Y}{\rRI{13,\rRI{14,17}}}
  \proofstepwidestar[1]{\PartOrd{\mathcal R_M,\mathcal L_M}}{%
    \FormulaRefAuto{Partielle Ordnung}[def]{2,5,12,18}}
\end{tabproofwide}

\FormulaDefDeltaK[Gemeinsame obere Elemente der Komplementordnung]{%
  \mathcal C_M(X,Y)\coloneqq
  \{Z\in\mathcal L_M\mid X\cup Y\subseteq Z\}
}{FranklComplementCommonUpperFamily}{%
  \DeltaRow{Mengen}{M\dsep X\dsep Y\dsep Z}
  \DeltaRow{\textbf{Neues Symbol}}{}
  \DeltaPrem{Gemeinsame obere Elemente}{\mathcal C_M(X,Y)}
}

\FormulaThmDeltaK[Die Familie gemeinsamer oberer Elemente ist endlich und nichtleer]{%
  \FranklFam M\dsep X,Y\in\mathcal L_M
  \vdash\Finite{\mathcal C_M(X,Y)}\land\mathcal C_M(X,Y)\neq\varnothing
}{FranklComplementCommonUpperFiniteNonempty}{%
  \DeltaRow{Mengen}{M\dsep X\dsep Y\dsep Z}
}
\begin{tabproofwide}
  \proofstepwidestar[1]{\FranklFam M}{\rA}
  \proofstepwidestar[2]{X,Y\in\mathcal L_M}{\rA}
  \proofstepwidestar[1]{\Finite{\mathcal L_M}}{%
    \FormulaRefAuto{FranklComplementFamilyFinite}[thm]{1}}
  \proofstepwidestar[]{\mathcal C_M(X,Y)\subseteq\mathcal L_M}{%
    \FormulaRefAuto{FranklComplementCommonUpperFamily}[def]}
  \proofstepwidestar[1]{U_M\in M}{%
    \FormulaRefAuto{FiniteUnionClosedFamilyTotalUnion}[thm]{%
      \FormulaRefAuto{FranklFamFinite}[thm]{1},
      \FormulaRefAuto{FranklFamNonempty}[thm]{1},
      \FormulaRefAuto{FranklFamUnionClosed}[thm]{1}}}
  \proofstepwidestar[1]{U_M\in\mathcal L_M}{%
    \FormulaRefAuto{FranklComplementConstruction}[def]{%
      \FormulaRefAuto{a\in A\cup\{a\}}[thm],
      \FormulaRefAuto{RelativeComplementEmpty}[thm]}}
  \proofstepwidestar[1,2]{X\cup Y\subseteq U_M}{%
    \FormulaRefAuto{A \subseteq C,\, B \subseteq C \vdash A \cup B \subseteq C}{%
      \FormulaRefAuto{FranklComplementMemberSubsetCarrier}[thm]{1,\rAEa{2}},
      \FormulaRefAuto{FranklComplementMemberSubsetCarrier}[thm]{1,\rAEb{2}}}}
  \proofstepwidestar[1,2]{U_M\in\mathcal C_M(X,Y)}{%
    \FormulaRefAuto{FranklComplementCommonUpperFamily}[def]{6,7}}
  \proofstepwidestar[1,2]{\mathcal C_M(X,Y)\neq\varnothing}{%
    \FormulaRefAuto{a\in S\vdash S\neq\varnothing}{8}}
  \proofstepwidestar[1]{\Finite{\mathcal C_M(X,Y)}}{%
    \FormulaRefAuto{FiniteSubset}[thm]{4,3}}
  \proofstepwidestar[1,2]{\Finite{\mathcal C_M(X,Y)}\land
    \mathcal C_M(X,Y)\neq\varnothing}{\rAI{10,9}}
\end{tabproofwide}

\FormulaThmDeltaK[Die konstruierte Operation ist das Paarsupremum]{%
  \FranklFam M\dsep X,Y\in\mathcal L_M\vdash
  \PairSup{\mathcal R_M}{\mathcal L_M}{X}{Y}{X\FamilyJoinOp{M}{}Y}
}{FranklComplementPairSupremum}{%
  \DeltaRow{Mengen}{M\dsep X\dsep Y\dsep Z\dsep W}
  \DeltaRow{Relationen}{\mathcal R_M}
  \DeltaRow{Funktionen}{\FamilyJoinOp{M}}
}
\begin{tabproofsplitwide}
  \proofpartwideindR[Beide Operanden liegen im Schnitt der gemeinsamen oberen Elemente]{%
    \begin{aligned}[t]
      &\FranklFam M\dsep X,Y\in\mathcal L_M\vdash{}\\
      &X\subseteq\bigcap\mathcal C_M(X,Y)\land
       Y\subseteq\bigcap\mathcal C_M(X,Y)
    \end{aligned}
  }{%
    \FranklFam M\dsep X,Y\in\mathcal L_M\vdash
    X\subseteq\bigcap\mathcal C_M(X,Y)\land
    Y\subseteq\bigcap\mathcal C_M(X,Y)
  }[FranklComplementCommonUpperIntersectionBounds]
    \proofstepwidestar[1]{\FranklFam M}{\rA}
    \proofstepwidestar[2]{X,Y\in\mathcal L_M}{\rA}
    \proofstepwidestar[1,2]{\mathcal C_M(X,Y)\neq\varnothing}{%
      \rAEb{\FormulaRefAuto{FranklComplementCommonUpperFiniteNonempty}[thm]{1,2}}}
    \proofstepwidestar[4]{a\in X}{\rA}
    \proofstepwidestar[5]{Z\in\mathcal C_M(X,Y)}{\rA}
    \proofstepwidestar[5]{X\cup Y\subseteq Z}{%
      \FormulaRefAuto{FranklComplementCommonUpperFamily}[def]{5}}
    \proofstepwidestar[]{X\subseteq X\cup Y}{%
      \FormulaRefAuto{A\subseteq A\cup B}}
    \proofstepwidestar[5]{X\subseteq Z}{%
      \FormulaRefAuto{A\subseteq B\dsep B\subseteq C\vdash A\subseteq C}{7,6}}
    \proofstepwidestar[4,5]{a\in Z}{%
      \FormulaRefAuto{A\subseteq B,\,x\in A\vdash x\in B}{8,4}}
    \proofstepwidestar[4]{\forall Z\in\mathcal C_M(X,Y)\,(a\in Z)}{%
      \rUI{\rRI{5,9}}}
    \proofstepwidestar[1,2]{%
      a\in\bigcap\mathcal C_M(X,Y)\leftrightarrow
      \forall Z\in\mathcal C_M(X,Y)\,(a\in Z)}{%
      \FormulaRefAuto{M\neq\varnothing \vdash x\in\bigcap M \leftrightarrow
        \forall X\in M\,(x\in X)}{3}}
    \proofstepwidestar[1,2,4]{a\in\bigcap\mathcal C_M(X,Y)}{%
      \FormulaRefAuto{P\leftrightarrow Q,Q\vdash P}{11,10}}
    \proofstepwidestar[1,2]{X\subseteq\bigcap\mathcal C_M(X,Y)}{%
      \FormulaRefAuto{A \subseteq B \coloneq
        \forall x\,(x\in A \rightarrow x\in B)}{\rUI{\rRI{4,12}}}}
    \proofstepwidestar[14]{a\in Y}{\rA}
    \proofstepwidestar[15]{Z\in\mathcal C_M(X,Y)}{\rA}
    \proofstepwidestar[15]{X\cup Y\subseteq Z}{%
      \FormulaRefAuto{FranklComplementCommonUpperFamily}[def]{15}}
    \proofstepwidestar[]{Y\subseteq X\cup Y}{%
      \FormulaRefAuto{B\subseteq A\cup B}}
    \proofstepwidestar[15]{Y\subseteq Z}{%
      \FormulaRefAuto{A\subseteq B\dsep B\subseteq C\vdash A\subseteq C}{17,16}}
    \proofstepwidestar[14,15]{a\in Z}{%
      \FormulaRefAuto{A\subseteq B,\,x\in A\vdash x\in B}{18,14}}
    \proofstepwidestar[14]{\forall Z\in\mathcal C_M(X,Y)\,(a\in Z)}{%
      \rUI{\rRI{15,19}}}
    \proofstepwidestar[1,2]{%
      a\in\bigcap\mathcal C_M(X,Y)\leftrightarrow
      \forall Z\in\mathcal C_M(X,Y)\,(a\in Z)}{%
      \FormulaRefAuto{M\neq\varnothing \vdash x\in\bigcap M \leftrightarrow
        \forall X\in M\,(x\in X)}{3}}
    \proofstepwidestar[1,2,14]{a\in\bigcap\mathcal C_M(X,Y)}{%
      \FormulaRefAuto{P\leftrightarrow Q,Q\vdash P}{21,20}}
    \proofstepwidestar[1,2]{Y\subseteq\bigcap\mathcal C_M(X,Y)}{%
      \FormulaRefAuto{A \subseteq B \coloneq
        \forall x\,(x\in A \rightarrow x\in B)}{\rUI{\rRI{14,22}}}}
    \proofstepwidestar[1,2]{%
      X\subseteq\bigcap\mathcal C_M(X,Y)\land
      Y\subseteq\bigcap\mathcal C_M(X,Y)}{\rAI{13,23}}
  \closeproofpartwideindR

  \proofpartwideindR[Träger und obere Schranken]{%
    \begin{aligned}[t]
      &\FranklFam M\dsep X,Y\in\mathcal L_M\vdash{}\\
      &X\FamilyJoinOp{M}{}Y\in\mathcal L_M\land
       (X,X\FamilyJoinOp{M}{}Y)\in\mathcal R_M\land
       (Y,X\FamilyJoinOp{M}{}Y)\in\mathcal R_M
    \end{aligned}
  }{%
    \FranklFam M\dsep X,Y\in\mathcal L_M\vdash
    X\FamilyJoinOp{M}{}Y\in\mathcal L_M\land
    (X,X\FamilyJoinOp{M}{}Y)\in\mathcal R_M\land
    (Y,X\FamilyJoinOp{M}{}Y)\in\mathcal R_M
  }[FranklComplementPairSupremumBounds]
    \proofstepwidestar[1]{\FranklFam M}{\rA}
    \proofstepwidestar[2]{X,Y\in\mathcal L_M}{\rA}
    \proofstepwidestar[1]{\InterClosedFam{\mathcal L_M}}{%
      \FormulaRefAuto{FranklComplementFamilyIntersectionClosed}[thm]{1}}
    \proofstepwidestar[1,2]{\Finite{\mathcal C_M(X,Y)}\land
      \mathcal C_M(X,Y)\neq\varnothing}{%
      \FormulaRefAuto{FranklComplementCommonUpperFiniteNonempty}[thm]{1,2}}
    \proofstepwidestar[1,2]{\Finite{\mathcal C_M(X,Y)}}{\rAEa{4}}
    \proofstepwidestar[1,2]{\mathcal C_M(X,Y)\neq\varnothing}{\rAEb{4}}
    \proofstepwidestar[]{\mathcal C_M(X,Y)\subseteq\mathcal L_M}{%
      \FormulaRefAuto{FranklComplementCommonUpperFamily}[def]}
    \proofstepwidestar[1,2]{\bigcap\mathcal C_M(X,Y)\in\mathcal L_M}{%
      \FormulaRefAuto{FiniteIntersectionClosedFamilyIntersection}[thm]{%
        5,6,7,3}}
    \proofstepwidestar[1,2]{X\FamilyJoinOp{M}{}Y=\bigcap\mathcal C_M(X,Y)}{%
      \FormulaRefAuto{FranklComplementConstruction}[def]{1,2}}
    \proofstepwidestar[1,2]{X\FamilyJoinOp{M}{}Y\in\mathcal L_M}{\rIE{9,8}}
    \proofstepwidestar[1,2]{%
      X\subseteq\bigcap\mathcal C_M(X,Y)\land
      Y\subseteq\bigcap\mathcal C_M(X,Y)}{%
      \FormulaRefAuto{FranklComplementCommonUpperIntersectionBounds}[thm]{1,2}}
    \proofstepwidestar[1,2]{X\subseteq\bigcap\mathcal C_M(X,Y)}{\rAEa{11}}
    \proofstepwidestar[1,2]{Y\subseteq\bigcap\mathcal C_M(X,Y)}{\rAEb{11}}
    \proofstepwidestar[1,2]{X\subseteq X\FamilyJoinOp{M}{}Y}{%
      \rIE{\FormulaRefAuto{a=b\vdash b=a}{9},12}}
    \proofstepwidestar[1,2]{Y\subseteq X\FamilyJoinOp{M}{}Y}{%
      \rIE{\FormulaRefAuto{a=b\vdash b=a}{9},13}}
    \proofstepwidestar[1,2]{(X,X\FamilyJoinOp{M}{}Y)\in\mathcal R_M}{%
      \FormulaRefAuto{FranklComplementConstruction}[def]{\rAEa{2},10,14}}
    \proofstepwidestar[1,2]{(Y,X\FamilyJoinOp{M}{}Y)\in\mathcal R_M}{%
      \FormulaRefAuto{FranklComplementConstruction}[def]{\rAEb{2},10,15}}
    \proofstepwidestar[1,2]{%
      X\FamilyJoinOp{M}{}Y\in\mathcal L_M\land
      (X,X\FamilyJoinOp{M}{}Y)\in\mathcal R_M\land
      (Y,X\FamilyJoinOp{M}{}Y)\in\mathcal R_M}{%
      \rAI{10,\rAI{16,17}}}
  \closeproofpartwideindR
  \proofpartwideindR[Kleinste obere Schranke]{%
    \begin{aligned}[t]
      &\FranklFam M\dsep X,Y\in\mathcal L_M\dsep
       W\in\mathcal L_M\dsep(X,W)\in\mathcal R_M\dsep
       (Y,W)\in\mathcal R_M\vdash{}\\
      &(X\FamilyJoinOp{M}{}Y,W)\in\mathcal R_M
    \end{aligned}
  }{%
    \FranklFam M\dsep X,Y\in\mathcal L_M\dsep
    W\in\mathcal L_M\dsep(X,W)\in\mathcal R_M\dsep
    (Y,W)\in\mathcal R_M\vdash
    (X\FamilyJoinOp{M}{}Y,W)\in\mathcal R_M
  }[FranklComplementPairSupremumLeast]
    \proofstepwidestar[1]{\FranklFam M}{\rA}
    \proofstepwidestar[2]{X,Y\in\mathcal L_M}{\rA}
    \proofstepwidestar[3]{W\in\mathcal L_M}{\rA}
    \proofstepwidestar[4]{(X,W)\in\mathcal R_M}{\rA}
    \proofstepwidestar[5]{(Y,W)\in\mathcal R_M}{\rA}
    \proofstepwidestar[4]{X\subseteq W}{%
      \FormulaRefAuto{FranklComplementConstruction}[def]{4}}
    \proofstepwidestar[5]{Y\subseteq W}{%
      \FormulaRefAuto{FranklComplementConstruction}[def]{5}}
    \proofstepwidestar[4,5]{X\cup Y\subseteq W}{%
      \FormulaRefAuto{A \subseteq C,\, B \subseteq C \vdash A \cup B \subseteq C}{6,7}}
    \proofstepwidestar[3,4,5]{W\in\mathcal C_M(X,Y)}{%
      \FormulaRefAuto{FranklComplementCommonUpperFamily}[def]{3,8}}
    \proofstepwidestar[1,2]{\mathcal C_M(X,Y)\neq\varnothing}{%
      \rAEb{\FormulaRefAuto{FranklComplementCommonUpperFiniteNonempty}[thm]{1,2}}}
    \proofstepwidestar[1,2,3,4,5]{\bigcap\mathcal C_M(X,Y)\subseteq W}{%
      \FormulaRefAuto{FamilyIntersectionSubsetMember}[thm]{%
        10,9}}
    \proofstepwidestar[1,2]{X\FamilyJoinOp{M}{}Y=\bigcap\mathcal C_M(X,Y)}{%
      \FormulaRefAuto{FranklComplementConstruction}[def]{1,2}}
    \proofstepwidestar[1,2,3,4,5]{X\FamilyJoinOp{M}{}Y\subseteq W}{%
      \rIE{\FormulaRefAuto{a=b\vdash b=a}{12},11}}
    \proofstepwidestar[1,2]{X\FamilyJoinOp{M}{}Y\in\mathcal L_M}{%
      \rAEa{\FormulaRefAuto{FranklComplementPairSupremumBounds}[thm]{1,2}}}
    \proofstepwidestar[1,2,3,4,5]{%
      (X\FamilyJoinOp{M}{}Y,W)\in\mathcal R_M}{%
      \FormulaRefAuto{FranklComplementConstruction}[def]{14,3,13}}
  \closeproofpartwideindR
  \proofpartwideindR[Zusammenführung der Paarsupremumseigenschaften]{%
    \FranklFam M\dsep X,Y\in\mathcal L_M\vdash
    \PairSup{\mathcal R_M}{\mathcal L_M}{X}{Y}{X\FamilyJoinOp{M}{}Y}
  }{%
    \FranklFam M\dsep X,Y\in\mathcal L_M\vdash
    \PairSup{\mathcal R_M}{\mathcal L_M}{X}{Y}{X\FamilyJoinOp{M}{}Y}
  }[FranklComplementPairSupremumAssembly]
    \proofstepwidestar[1]{\FranklFam M}{\rA}
    \proofstepwidestar[2]{X,Y\in\mathcal L_M}{\rA}
    \proofstepwidestar[1,2]{X\FamilyJoinOp{M}{}Y\in\mathcal L_M}{%
      \rAEa{\FormulaRefAuto{FranklComplementPairSupremumBounds}[thm]{1,2}}}
    \proofstepwidestar[1,2]{(X,X\FamilyJoinOp{M}{}Y)\in\mathcal R_M}{%
      \rAEa{\rAEb{\FormulaRefAuto{FranklComplementPairSupremumBounds}[thm]{1,2}}}}
    \proofstepwidestar[1,2]{(Y,X\FamilyJoinOp{M}{}Y)\in\mathcal R_M}{%
      \rAEb{\rAEb{\FormulaRefAuto{FranklComplementPairSupremumBounds}[thm]{1,2}}}}
    \proofstepwidestar[1,2]{W\in\mathcal L_M\dsep(X,W)\in\mathcal R_M
      \dsep(Y,W)\in\mathcal R_M\vdash
      (X\FamilyJoinOp{M}{}Y,W)\in\mathcal R_M}{%
      \FormulaRefAuto{FranklComplementPairSupremumLeast}[thm]{1,2}}
    \proofstepwidestar[2]{X\in\mathcal L_M}{\rAEa{2}}
    \proofstepwidestar[2]{Y\in\mathcal L_M}{\rAEb{2}}
    \proofstepwidestar[2]{X\in\mathcal L_M\land Y\in\mathcal L_M}{\rAI{7,8}}
    \proofstepwidestar[1,2]{%
      (X\in\mathcal L_M\land Y\in\mathcal L_M)\land
      X\FamilyJoinOp{M}{}Y\in\mathcal L_M}{\rAI{9,3}}
    \proofstepwidestar[1,2]{%
      \PairSup{\mathcal R_M}{\mathcal L_M}{X}{Y}{X\FamilyJoinOp{M}{}Y}}{%
      \FormulaRefAuto{PairSupremum}[def]{10,4,5,6}}
  \closeproofpartwideindR
\end{tabproofsplitwide}

\FormulaThmDeltaK[Die konstruierte Operation ist eine Paarsupremumsoperation]{%
  \FranklFam M\vdash
  \PairJoinOp{\FamilyJoinOp{M},\mathcal R_M,\mathcal L_M}
}{FranklComplementPairJoinOperation}{%
  \DeltaRow{Mengen}{M\dsep X\dsep Y}
  \DeltaRow{Relationen}{\mathcal R_M}
  \DeltaRow{Funktionen}{\FamilyJoinOp{M}}
}
\begin{tabproof}
  \proofstep{1}{\FranklFam M}{\rA}
  \proofstep{1}{X,Y\in\mathcal L_M\vdash
    \PairSup{\mathcal R_M}{\mathcal L_M}{X}{Y}{X\FamilyJoinOp{M}{}Y}}{%
    \FormulaRefAuto{FranklComplementPairSupremum}[thm]{1}}
  \proofstep{1}{\PairJoinOp{\FamilyJoinOp{M},\mathcal R_M,\mathcal L_M}}{%
    \FormulaRefAuto{PairJoinOperation}[def]{2}}
\end{tabproof}

\FormulaThmDeltaK[Die konstruierte Operation bildet einen Halbverband]{%
  \FranklFam M\vdash\Semi{\FamilyJoinOp{M},\mathcal L_M}
}{FranklComplementSemilattice}{%
  \DeltaRow{Mengen}{M}
  \DeltaRow{Relationen}{\mathcal R_M}
  \DeltaRow{Funktionen}{\FamilyJoinOp{M}}
}
\begin{tabproof}
  \proofstep{1}{\FranklFam M}{\rA}
  \proofstep{1}{\PartOrd{\mathcal R_M,\mathcal L_M}}{%
    \FormulaRefAuto{FranklComplementOrderPartial}[thm]{1}}
  \proofstep{1}{\PairJoinOp{\FamilyJoinOp{M},\mathcal R_M,\mathcal L_M}}{%
    \FormulaRefAuto{FranklComplementPairJoinOperation}[thm]{1}}
  \proofstep{1}{\Semi{\FamilyJoinOp{M},\mathcal L_M}}{%
    \FormulaRefAuto{PairJoinCharacterization}[thm]{2,3}}
\end{tabproof}

\FormulaThmDeltaK[Die induzierte Ordnung ist die Inklusionsordnung]{%
  \FranklFam M\vdash
  \JoinOrd{\FamilyJoinOp{M}}{\mathcal L_M}=\mathcal R_M
}{FranklComplementOrderRecovered}{%
  \DeltaRow{Mengen}{M\dsep X\dsep Y\dsep q}
  \DeltaRow{Relationen}{\mathcal R_M}
  \DeltaRow{Funktionen}{\FamilyJoinOp{M}}
}
\begin{tabproofsplitwide}
  \proofpartwideindR[Von der expliziten zur induzierten Relation]{%
    \FranklFam M\vdash
    \mathcal R_M\subseteq\JoinOrd{\FamilyJoinOp{M}}{\mathcal L_M}
  }{%
    \FranklFam M\vdash
    \mathcal R_M\subseteq\JoinOrd{\FamilyJoinOp{M}}{\mathcal L_M}
  }[FranklComplementOrderRecoveredForward]
    \proofstepwidestar[1]{\FranklFam M}{\rA}
    \proofstepwidestar[1]{\PartOrd{\mathcal R_M,\mathcal L_M}}{%
      \FormulaRefAuto{FranklComplementOrderPartial}[thm]{1}}
    \proofstepwidestar[1]{%
      \PairJoinOp{\FamilyJoinOp{M},\mathcal R_M,\mathcal L_M}}{%
      \FormulaRefAuto{FranklComplementPairJoinOperation}[thm]{1}}
    \proofstepwidestar[4]{q\in\mathcal R_M}{\rA}
    \proofstepwidestar[]{\mathcal R_M\subseteq\mathcal L_M\times\mathcal L_M}{%
      \FormulaRefAuto{FranklComplementConstruction}[def]}
    \proofstepwidestar[4]{q\in\mathcal L_M\times\mathcal L_M}{%
      \FormulaRefAuto{A\subseteq B,\,x\in A\vdash x\in B}{5,4}}
    \proofstepwidestar[4]{%
      \exists X\in\mathcal L_M\,\exists Y\in\mathcal L_M\,q=(X,Y)}{%
      \FormulaRefAuto{x\in A\times B\eqvdash
        \exists a \in A\, \exists b \in B\,x=(a,b)}{6}}
    \proofstepwidestar[8]{%
      X\in\mathcal L_M\land Y\in\mathcal L_M\land q=(X,Y)}{\rA}
    \proofstepwidestar[8]{X\in\mathcal L_M}{%
      \FormulaRefAuto{P\land Q\land R\vdash P}{8}}
    \proofstepwidestar[8]{Y\in\mathcal L_M}{%
      \FormulaRefAuto{P\land Q\land R\vdash Q}{8}}
    \proofstepwidestar[8]{q=(X,Y)}{%
      \FormulaRefAuto{P\land Q\land R\vdash R}{8}}
    \proofstepwidestar[4,8]{(X,Y)\in\mathcal R_M}{\rIE{11,4}}
    \proofstepwidestar[1,8]{%
      (X,Y)\in\mathcal R_M\leftrightarrow
      (X,Y)\in\JoinOrd{\FamilyJoinOp{M}}{\mathcal L_M}}{%
      \FormulaRefAuto{PairJoinRecoversOrder}[thm]{2,3,\rAI{9,10}}}
    \proofstepwidestar[1,4,8]{%
      (X,Y)\in\JoinOrd{\FamilyJoinOp{M}}{\mathcal L_M}}{%
      \FormulaRefAuto{P\leftrightarrow Q,P\vdash Q}{13,12}}
    \proofstepwidestar[1,4,8]{%
      q\in\JoinOrd{\FamilyJoinOp{M}}{\mathcal L_M}}{%
      \rIE{\FormulaRefAuto{a=b\vdash b=a}{11},14}}
    \proofstepwidestar[1,4]{%
      q\in\JoinOrd{\FamilyJoinOp{M}}{\mathcal L_M}}{\rEE{7,8,15}}
    \proofstepwidestar[1]{%
      \mathcal R_M\subseteq\JoinOrd{\FamilyJoinOp{M}}{\mathcal L_M}}{%
      \FormulaRefAuto{A \subseteq B \coloneq
        \forall x\,(x\in A \rightarrow x\in B)}{\rUI{\rRI{4,16}}}}
  \closeproofpartwideindR

  \proofpartwideindR[Von der induzierten zur expliziten Relation]{%
    \FranklFam M\vdash
    \JoinOrd{\FamilyJoinOp{M}}{\mathcal L_M}\subseteq\mathcal R_M
  }{%
    \FranklFam M\vdash
    \JoinOrd{\FamilyJoinOp{M}}{\mathcal L_M}\subseteq\mathcal R_M
  }[FranklComplementOrderRecoveredBackward]
    \proofstepwidestar[1]{\FranklFam M}{\rA}
    \proofstepwidestar[1]{\PartOrd{\mathcal R_M,\mathcal L_M}}{%
      \FormulaRefAuto{FranklComplementOrderPartial}[thm]{1}}
    \proofstepwidestar[1]{%
      \PairJoinOp{\FamilyJoinOp{M},\mathcal R_M,\mathcal L_M}}{%
      \FormulaRefAuto{FranklComplementPairJoinOperation}[thm]{1}}
    \proofstepwidestar[4]{%
      q\in\JoinOrd{\FamilyJoinOp{M}}{\mathcal L_M}}{\rA}
    \proofstepwidestar[]{%
      \JoinOrd{\FamilyJoinOp{M}}{\mathcal L_M}
      \subseteq\mathcal L_M\times\mathcal L_M}{%
      \FormulaRefAuto{JoinOrderCarrierPart}[thm]}
    \proofstepwidestar[4]{q\in\mathcal L_M\times\mathcal L_M}{%
      \FormulaRefAuto{A\subseteq B,\,x\in A\vdash x\in B}{5,4}}
    \proofstepwidestar[4]{%
      \exists X\in\mathcal L_M\,\exists Y\in\mathcal L_M\,q=(X,Y)}{%
      \FormulaRefAuto{x\in A\times B\eqvdash
        \exists a \in A\, \exists b \in B\,x=(a,b)}{6}}
    \proofstepwidestar[8]{%
      X\in\mathcal L_M\land Y\in\mathcal L_M\land q=(X,Y)}{\rA}
    \proofstepwidestar[8]{X\in\mathcal L_M}{%
      \FormulaRefAuto{P\land Q\land R\vdash P}{8}}
    \proofstepwidestar[8]{Y\in\mathcal L_M}{%
      \FormulaRefAuto{P\land Q\land R\vdash Q}{8}}
    \proofstepwidestar[8]{q=(X,Y)}{%
      \FormulaRefAuto{P\land Q\land R\vdash R}{8}}
    \proofstepwidestar[4,8]{%
      (X,Y)\in\JoinOrd{\FamilyJoinOp{M}}{\mathcal L_M}}{\rIE{11,4}}
    \proofstepwidestar[1,8]{%
      (X,Y)\in\mathcal R_M\leftrightarrow
      (X,Y)\in\JoinOrd{\FamilyJoinOp{M}}{\mathcal L_M}}{%
      \FormulaRefAuto{PairJoinRecoversOrder}[thm]{2,3,\rAI{9,10}}}
    \proofstepwidestar[1,4,8]{(X,Y)\in\mathcal R_M}{%
      \FormulaRefAuto{P\leftrightarrow Q,Q\vdash P}{13,12}}
    \proofstepwidestar[1,4,8]{q\in\mathcal R_M}{%
      \rIE{\FormulaRefAuto{a=b\vdash b=a}{11},14}}
    \proofstepwidestar[1,4]{q\in\mathcal R_M}{\rEE{7,8,15}}
    \proofstepwidestar[1]{%
      \JoinOrd{\FamilyJoinOp{M}}{\mathcal L_M}\subseteq\mathcal R_M}{%
      \FormulaRefAuto{A \subseteq B \coloneq
        \forall x\,(x\in A \rightarrow x\in B)}{\rUI{\rRI{4,16}}}}
  \closeproofpartwideindR

  \proofpartwide{Schluss}
    \proofstepwidestar[1]{\FranklFam M}{\rA}
    \proofstepwidestar[1]{%
      \mathcal R_M\subseteq\JoinOrd{\FamilyJoinOp{M}}{\mathcal L_M}}{%
      \FormulaRefAuto{FranklComplementOrderRecoveredForward}[thm]{1}}
    \proofstepwidestar[1]{%
      \JoinOrd{\FamilyJoinOp{M}}{\mathcal L_M}\subseteq\mathcal R_M}{%
      \FormulaRefAuto{FranklComplementOrderRecoveredBackward}[thm]{1}}
    \proofstepwidestar[1]{%
      \JoinOrd{\FamilyJoinOp{M}}{\mathcal L_M}=\mathcal R_M}{%
      \FormulaRefAuto{A\subseteq B\dsep B\subseteq A\vdash A=B}{3,2}}
  \closeproofpartwide
\end{tabproofsplitwide}

\FormulaThmDeltaK[Komplementdual einer Frankl-Familie]{%
  \FranklFam M\vdash
  \ZeroJoinSemi{\FamilyJoinOp{M},\mathcal L_M,\varnothing}
}{FranklComplementFiniteLattice}{%
  \DeltaRow{Mengen}{M\dsep X}
  \DeltaRow{Relationen}{\mathcal R_M}
  \DeltaRow{Funktionen}{\FamilyJoinOp{M}}
}
\begin{tabproofsplitwide}
  \proofpartwideindR[Träger des Nullelements]{%
    \FranklFam M\vdash\varnothing\in\mathcal L_M
  }{%
    \FranklFam M\vdash\varnothing\in\mathcal L_M
  }[FranklComplementZeroCarrier]
    \proofstepwidestar[1]{\FranklFam M}{\rA}
    \proofstepwidestar[1]{U_M\in M}{%
      \FormulaRefAuto{FiniteUnionClosedFamilyTotalUnion}[thm]{%
        \FormulaRefAuto{FranklFamFinite}[thm]{1},
        \FormulaRefAuto{FranklFamNonempty}[thm]{1},
        \FormulaRefAuto{FranklFamUnionClosed}[thm]{1}}}
    \proofstepwidestar[1]{\varnothing\in\mathcal L_M}{%
      \FormulaRefAuto{FranklComplementConstruction}[def]{2,
        \FormulaRefAuto{RelativeComplementSelf}[thm]}}
  \closeproofpartwideindR
  \proofpartwideindR[Ein Trägerelement ist das Supremum von Null und sich selbst]{%
    \FranklFam M\dsep X\in\mathcal L_M\vdash
    \PairSup{\mathcal R_M}{\mathcal L_M}{\varnothing}{X}{X}
  }{%
    \FranklFam M\dsep X\in\mathcal L_M\vdash
    \PairSup{\mathcal R_M}{\mathcal L_M}{\varnothing}{X}{X}
  }[FranklComplementZeroSelfSupremum]
    \proofstepwidestar[1]{\FranklFam M}{\rA}
    \proofstepwidestar[2]{X\in\mathcal L_M}{\rA}
    \proofstepwidestar[1]{\varnothing\in\mathcal L_M}{%
      \FormulaRefAuto{FranklComplementZeroCarrier}[thm]{1}}
    \proofstepwidestar[1,2]{(\varnothing,X)\in\mathcal R_M}{%
      \FormulaRefAuto{FranklComplementConstruction}[def]{%
        3,2,\FormulaRefAuto{\varnothing\subseteq A}}}
    \proofstepwidestar[2]{(X,X)\in\mathcal R_M}{%
      \FormulaRefAuto{FranklComplementConstruction}[def]{%
        2,2,\FormulaRefAuto{A\subseteq A}}}
    \proofstepwidestar[]{%
      W\in\mathcal L_M\dsep(\varnothing,W)\in\mathcal R_M\dsep
      (X,W)\in\mathcal R_M\vdash(X,W)\in\mathcal R_M}{\rA}
    \proofstepwidestar[1,2]{%
      \varnothing\in\mathcal L_M\land X\in\mathcal L_M}{\rAI{3,2}}
    \proofstepwidestar[1,2]{%
      (\varnothing\in\mathcal L_M\land X\in\mathcal L_M)\land
      X\in\mathcal L_M}{\rAI{7,2}}
    \proofstepwidestar[1,2]{%
      \PairSup{\mathcal R_M}{\mathcal L_M}{\varnothing}{X}{X}}{%
      \FormulaRefAuto{PairSupremum}[def]{8,4,5,6}}
  \closeproofpartwideindR
  \proofpartwideindR[Neutralität des Nullelements]{%
    \FranklFam M\dsep X\in\mathcal L_M\vdash
    \varnothing\FamilyJoinOp{M}{}X=X
  }{%
    \FranklFam M\dsep X\in\mathcal L_M\vdash
    \varnothing\FamilyJoinOp{M}{}X=X
  }[FranklComplementZeroIdentity]
    \proofstepwidestar[1]{\FranklFam M}{\rA}
    \proofstepwidestar[2]{X\in\mathcal L_M}{\rA}
    \proofstepwidestar[1]{\varnothing\in\mathcal L_M}{%
      \FormulaRefAuto{FranklComplementZeroCarrier}[thm]{1}}
    \proofstepwidestar[1,2]{%
      \PairSup{\mathcal R_M}{\mathcal L_M}{\varnothing}{X}{%
        \varnothing\FamilyJoinOp{M}{}X}}{%
      \FormulaRefAuto{FranklComplementPairSupremum}[thm]{1,3,2}}
    \proofstepwidestar[1,2]{%
      \PairSup{\mathcal R_M}{\mathcal L_M}{\varnothing}{X}{X}}{%
      \FormulaRefAuto{FranklComplementZeroSelfSupremum}[thm]{1,2}}
    \proofstepwidestar[1,2]{\varnothing\FamilyJoinOp{M}{}X=X}{%
      \FormulaRefAuto{PairSupremumUnique}[thm]{%
        \FormulaRefAuto{FranklComplementOrderPartial}[thm]{1},4,5}}
  \closeproofpartwideindR
  \proofpartwide{Schluss}
    \proofstepwidestar[1]{\FranklFam M}{\rA}
    \proofstepwidestar[1]{\Semi{\FamilyJoinOp{M},\mathcal L_M}}{%
      \FormulaRefAuto{FranklComplementSemilattice}[thm]{1}}
    \proofstepwidestar[1]{\varnothing\in\mathcal L_M}{%
      \FormulaRefAuto{FranklComplementZeroCarrier}[thm]{1}}
    \proofstepwidestar[1]{X\in\mathcal L_M\vdash
      \varnothing\FamilyJoinOp{M}{}X=X}{%
      \FormulaRefAuto{FranklComplementZeroIdentity}[thm]{1}}
    \proofstepwidestar[1]{%
      \ZeroJoinSemi{\FamilyJoinOp{M},\mathcal L_M,\varnothing}}{%
      \FormulaRefAuto{ZeroJoinSemilattice}[def]{2,3,4}}
  \closeproofpartwide
\end{tabproofsplitwide}

\FormulaThmDeltaK[Der Komplementdual ist nichttrivial]{%
  \FranklFam M\vdash\exists X\in\mathcal L_M\,(X\neq\varnothing)
}{FranklComplementNontrivial}{%
  \DeltaRow{Mengen}{M\dsep X}
}
\begin{tabproof}
  \proofstep{1}{\FranklFam M}{\rA}
  \proofstep{1}{U_M\neq\varnothing}{%
    \FormulaRefAuto{FranklFamCarrierNonempty}[thm]{1}}
  \proofstep{1}{U_M\in\mathcal L_M}{%
    \FormulaRefAuto{FranklComplementConstruction}[def]{%
      \FormulaRefAuto{a\in A\cup\{a\}}[thm],
      \FormulaRefAuto{RelativeComplementEmpty}[thm]}}
  \proofstep{1}{\exists X\in\mathcal L_M\,(X\neq\varnothing)}{%
    \rEI{\rAI{3,2}}}
\end{tabproof}

\FormulaDefDeltaK[Echte untere Elemente eines Irreduziblen]{%
  D_L(P)\coloneqq
  \{Q\in L\mid (Q,P)\in\JoinOrd{\JoinOp}{L}\land Q\neq P\}
}{JoinIrreducibleProperLowerSet}{%
  \DeltaRow{Mengen}{L\dsep P\dsep Q}
  \DeltaRow{Funktionen}{\JoinOp}
  \DeltaRow{\textbf{Neues Symbol}}{}
  \DeltaPrem{Echte untere Elemente}{D_L(P)}
}

\FormulaThmDeltaK[Echte untere Elemente liegen im Träger]{%
  Q\in D_L(P)\vdash Q\in L
}{JoinIrreducibleProperLowerCarrier}{%
  \DeltaRow{Mengen}{L\dsep P\dsep Q}
  \DeltaRow{Funktionen}{\JoinOp}
}
\begin{tabproof}
  \proofstep{1}{Q\in D_L(P)}{\rA}
  \proofstep{1}{Q\in L\land
    (Q,P)\in\JoinOrd{\JoinOp}{L}\land Q\neq P}{%
    \FormulaRefAuto{JoinIrreducibleProperLowerSet}[def]{1}}
  \proofstep{1}{Q\in L}{\rAEa{2}}
\end{tabproof}

\FormulaThmDeltaK[Die echten unteren Elemente bilden eine endliche nichtleere Teilmenge]{%
  \Finite L\dsep\ZeroJoinSemi{\JoinOp,L,\varnothing}\dsep
  \JoinIrr{\JoinOp,L,\varnothing}{P}\vdash
  \Finite{D_L(P)}\land D_L(P)\neq\varnothing\land D_L(P)\subseteq L
}{JoinIrreducibleProperLowerFinite}{%
  \DeltaRow{Mengen}{L\dsep P\dsep Q}
  \DeltaRow{Funktionen}{\JoinOp}
}
\begin{tabproofwide}
  \proofstepwidestar[1]{\Finite L}{\rA}
  \proofstepwidestar[2]{\ZeroJoinSemi{\JoinOp,L,\varnothing}}{\rA}
  \proofstepwidestar[3]{\JoinIrr{\JoinOp,L,\varnothing}{P}}{\rA}
  \proofstepwidestar[3]{P\in L}{\FormulaRefAuto{JoinIrreducible}[def]{3}}
  \proofstepwidestar[3]{P\neq\varnothing}{\FormulaRefAuto{JoinIrreducible}[def]{3}}
  \proofstepwidestar[2]{\varnothing\in L}{%
    \FormulaRefAuto{ZeroJoinCarrierAxiom}[ax]{2}}
  \proofstepwidestar[2,4]{(\varnothing,P)\in\JoinOrd{\JoinOp}{L}}{%
    \FormulaRefAuto{ZeroJoinLeast}[thm]{2,4}}
  \proofstepwidestar[2,3]{\varnothing\in D_L(P)}{%
    \FormulaRefAuto{JoinIrreducibleProperLowerSet}[def]{6,7,5}}
  \proofstepwidestar[2,3]{D_L(P)\neq\varnothing}{%
    \FormulaRefAuto{a\in S\vdash S\neq\varnothing}{8}}
  \proofstepwidestar[]{D_L(P)\subseteq L}{%
    \FormulaRefAuto{JoinIrreducibleProperLowerSet}[def]}
  \proofstepwidestar[1]{\Finite{D_L(P)}}{%
    \FormulaRefAuto{FiniteSubset}[thm]{10,1}}
  \proofstepwidestar[1,2,3]{\Finite{D_L(P)}\land D_L(P)\neq\varnothing
    \land D_L(P)\subseteq L}{\rAI{11,\rAI{9,10}}}
\end{tabproofwide}

\FormulaThmDeltaK[Die echten unteren Elemente sind unter dem Supremum abgeschlossen]{%
  \ZeroJoinSemi{\JoinOp,L,\varnothing}\dsep
  \JoinIrr{\JoinOp,L,\varnothing}{P}\dsep Q,S\in D_L(P)
  \vdash Q\JoinOp{}S\in D_L(P)
}{JoinIrreducibleProperLowerJoinClosed}{%
  \DeltaRow{Mengen}{L\dsep P\dsep Q\dsep S}
  \DeltaRow{Funktionen}{\JoinOp}
}
\begin{tabproofsplitwide}
  \proofpartwideindR[Das Supremum liegt im Träger und unter \(P\)]{%
    \begin{aligned}[t]
      &\ZeroJoinSemi{\JoinOp,L,\varnothing}\dsep
        \JoinIrr{\JoinOp,L,\varnothing}{P}\dsep
        Q,S\in D_L(P)\vdash{}\\
      &Q\JoinOp{}S\in L\land
        (Q\JoinOp{}S,P)\in\JoinOrd{\JoinOp}{L}
    \end{aligned}
  }{%
    \ZeroJoinSemi{\JoinOp,L,\varnothing}\dsep
    \JoinIrr{\JoinOp,L,\varnothing}{P}\dsep Q,S\in D_L(P)\vdash
    Q\JoinOp{}S\in L\land
    (Q\JoinOp{}S,P)\in\JoinOrd{\JoinOp}{L}
  }[JoinIrreducibleProperLowerJoinBound]
    \proofstepwidestar[1]{\ZeroJoinSemi{\JoinOp,L,\varnothing}}{\rA}
    \proofstepwidestar[2]{\JoinIrr{\JoinOp,L,\varnothing}{P}}{\rA}
    \proofstepwidestar[3]{Q,S\in D_L(P)}{\rA}
    \proofstepwidestar[3]{Q\in D_L(P)}{\rAEa{3}}
    \proofstepwidestar[3]{S\in D_L(P)}{\rAEb{3}}
    \proofstepwidestar[3]{Q\in L}{%
      \FormulaRefAuto{JoinIrreducibleProperLowerCarrier}[thm]{4}}
    \proofstepwidestar[3]{S\in L}{%
      \FormulaRefAuto{JoinIrreducibleProperLowerCarrier}[thm]{5}}
    \proofstepwidestar[2]{P\in L}{%
      \FormulaRefAuto{JoinIrreducible}[def]{2}}
    \proofstepwidestar[3]{%
      Q\in L\land(Q,P)\in\JoinOrd{\JoinOp}{L}\land Q\neq P}{%
      \FormulaRefAuto{JoinIrreducibleProperLowerSet}[def]{4}}
    \proofstepwidestar[3]{(Q,P)\in\JoinOrd{\JoinOp}{L}}{\rAEa{\rAEb{9}}}
    \proofstepwidestar[3]{%
      S\in L\land(S,P)\in\JoinOrd{\JoinOp}{L}\land S\neq P}{%
      \FormulaRefAuto{JoinIrreducibleProperLowerSet}[def]{5}}
    \proofstepwidestar[3]{(S,P)\in\JoinOrd{\JoinOp}{L}}{\rAEa{\rAEb{11}}}
    \proofstepwidestar[1]{\Semi{\JoinOp,L}}{%
      \FormulaRefAuto{ZeroJoinBaseSemilatticeAxiom}[ax]{1}}
    \proofstepwidestar[1,3]{Q\JoinOp{}S\in L}{%
      \FormulaRefAuto{SemilatticeClosureAxiom}[ax]{13,6,7}}
    \proofstepwidestar[1,2,3]{(Q\JoinOp{}S,P)\in\JoinOrd{\JoinOp}{L}}{%
      \FormulaRefAuto{JoinLeastUpper}[thm]{13,6,7,8,10,12}}
    \proofstepwidestar[1,2,3]{%
      Q\JoinOp{}S\in L\land
      (Q\JoinOp{}S,P)\in\JoinOrd{\JoinOp}{L}}{\rAI{14,15}}
  \closeproofpartwideindR

  \proofpartwideindR[Das Supremum bleibt von \(P\) verschieden]{%
    \ZeroJoinSemi{\JoinOp,L,\varnothing}\dsep
    \JoinIrr{\JoinOp,L,\varnothing}{P}\dsep Q,S\in D_L(P)
    \vdash Q\JoinOp{}S\neq P
  }{%
    \ZeroJoinSemi{\JoinOp,L,\varnothing}\dsep
    \JoinIrr{\JoinOp,L,\varnothing}{P}\dsep Q,S\in D_L(P)
    \vdash Q\JoinOp{}S\neq P
  }[JoinIrreducibleProperLowerJoinProper]
    \proofstepwidestar[1]{\ZeroJoinSemi{\JoinOp,L,\varnothing}}{\rA}
    \proofstepwidestar[2]{\JoinIrr{\JoinOp,L,\varnothing}{P}}{\rA}
    \proofstepwidestar[3]{Q,S\in D_L(P)}{\rA}
    \proofstepwidestar[3]{Q\in D_L(P)}{\rAEa{3}}
    \proofstepwidestar[3]{S\in D_L(P)}{\rAEb{3}}
    \proofstepwidestar[3]{Q\in L}{%
      \FormulaRefAuto{JoinIrreducibleProperLowerCarrier}[thm]{4}}
    \proofstepwidestar[3]{S\in L}{%
      \FormulaRefAuto{JoinIrreducibleProperLowerCarrier}[thm]{5}}
    \proofstepwidestar[3]{%
      Q\in L\land(Q,P)\in\JoinOrd{\JoinOp}{L}\land Q\neq P}{%
      \FormulaRefAuto{JoinIrreducibleProperLowerSet}[def]{4}}
    \proofstepwidestar[3]{Q\neq P}{\rAEb{\rAEb{8}}}
    \proofstepwidestar[3]{%
      S\in L\land(S,P)\in\JoinOrd{\JoinOp}{L}\land S\neq P}{%
      \FormulaRefAuto{JoinIrreducibleProperLowerSet}[def]{5}}
    \proofstepwidestar[3]{S\neq P}{\rAEb{\rAEb{10}}}
    \proofstepwidestar[12]{Q\JoinOp{}S=P}{\rA}
    \proofstepwidestar[2,3,12]{Q=P\lor S=P}{%
      \FormulaRefAuto{JoinIrreducible}[def]{2,6,7,12}}
    \proofstepwidestar[14]{Q=P}{\rA}
    \proofstepwidestar[3,14]{\bot}{\rBI{14,9}}
    \proofstepwidestar[16]{S=P}{\rA}
    \proofstepwidestar[3,16]{\bot}{\rBI{16,11}}
    \proofstepwidestar[2,3,12]{\bot}{\rOE{13,14,15,16,17}}
    \proofstepwidestar[2,3]{Q\JoinOp{}S\neq P}{\rCI{12,18}}
  \closeproofpartwideindR

  \proofpartwide{Schluss}
    \proofstepwidestar[1]{\ZeroJoinSemi{\JoinOp,L,\varnothing}}{\rA}
    \proofstepwidestar[2]{\JoinIrr{\JoinOp,L,\varnothing}{P}}{\rA}
    \proofstepwidestar[3]{Q,S\in D_L(P)}{\rA}
    \proofstepwidestar[1,2,3]{%
      Q\JoinOp{}S\in L\land
      (Q\JoinOp{}S,P)\in\JoinOrd{\JoinOp}{L}}{%
      \FormulaRefAuto{JoinIrreducibleProperLowerJoinBound}[thm]{1,2,3}}
    \proofstepwidestar[1,2,3]{Q\JoinOp{}S\in L}{\rAEa{4}}
    \proofstepwidestar[1,2,3]{%
      (Q\JoinOp{}S,P)\in\JoinOrd{\JoinOp}{L}}{\rAEb{4}}
    \proofstepwidestar[1,2,3]{Q\JoinOp{}S\neq P}{%
      \FormulaRefAuto{JoinIrreducibleProperLowerJoinProper}[thm]{1,2,3}}
    \proofstepwidestar[1,2,3]{Q\JoinOp{}S\in D_L(P)}{%
      \FormulaRefAuto{JoinIrreducibleProperLowerSet}[def]{5,6,7}}
  \closeproofpartwide
\end{tabproofsplitwide}

\FormulaThmDeltaK[Größtes echtes unteres Element eines Irreduziblen]{%
  \Finite L\dsep\ZeroJoinSemi{\JoinOp,L,\varnothing}\dsep
  \JoinIrr{\JoinOp,L,\varnothing}{P}\vdash
  \exists P_\ast\in D_L(P)\,\forall Q\in D_L(P)\,
    (Q,P_\ast)\in\JoinOrd{\JoinOp}{L}
}{JoinIrreducibleLargestProperLower}{%
  \DeltaRow{Mengen}{L\dsep P\dsep P_\ast\dsep Q\dsep S}
  \DeltaRow{Funktionen}{\JoinOp}
}
\begin{tabproofsplitwide}
  \proofpartwideindR[Ein maximales echtes unteres Element existiert]{%
    \begin{aligned}[t]
      &\Finite L\dsep\ZeroJoinSemi{\JoinOp,L,\varnothing}\dsep
        \JoinIrr{\JoinOp,L,\varnothing}{P}\vdash{}\\
      &\exists P_\ast\in D_L(P)\,\forall Q\in D_L(P)\,
        ((P_\ast,Q)\in\JoinOrd{\JoinOp}{L}\to Q=P_\ast)
    \end{aligned}
  }{%
    \Finite L\dsep\ZeroJoinSemi{\JoinOp,L,\varnothing}\dsep
    \JoinIrr{\JoinOp,L,\varnothing}{P}\vdash
    \exists P_\ast\in D_L(P)\,\forall Q\in D_L(P)\,
    ((P_\ast,Q)\in\JoinOrd{\JoinOp}{L}\to Q=P_\ast)
  }[JoinIrreducibleProperLowerMaximalExists]
    \proofstepwidestar[1]{\Finite L}{\rA}
    \proofstepwidestar[2]{\ZeroJoinSemi{\JoinOp,L,\varnothing}}{\rA}
    \proofstepwidestar[3]{\JoinIrr{\JoinOp,L,\varnothing}{P}}{\rA}
    \proofstepwidestar[1,2,3]{%
      \Finite{D_L(P)}\land D_L(P)\neq\varnothing\land D_L(P)\subseteq L}{%
      \FormulaRefAuto{JoinIrreducibleProperLowerFinite}[thm]{1,2,3}}
    \proofstepwidestar[1,2,3]{\Finite{D_L(P)}}{\rAEa{4}}
    \proofstepwidestar[1,2,3]{D_L(P)\neq\varnothing}{\rAEa{\rAEb{4}}}
    \proofstepwidestar[1,2,3]{D_L(P)\subseteq L}{\rAEb{\rAEb{4}}}
    \proofstepwidestar[2]{\Semi{\JoinOp,L}}{%
      \FormulaRefAuto{ZeroJoinBaseSemilatticeAxiom}[ax]{2}}
    \proofstepwidestar[2]{\PartOrd{\JoinOrd{\JoinOp}{L},L}}{%
      \FormulaRefAuto{JoinSemilatticeInducesOrder}[thm]{8}}
    \proofstepwidestar[1,2,3]{%
      \exists P_\ast\in D_L(P)\,\forall Q\in D_L(P)\,
      ((P_\ast,Q)\in\JoinOrd{\JoinOp}{L}\to Q=P_\ast)}{%
      \FormulaRefAuto{FinitePosetMaximalElement}[thm]{9,5,6,7}}
  \closeproofpartwideindR

  \proofpartwideindR[Ein maximales echtes unteres Element ist das größte]{%
    \begin{aligned}[t]
      &\ZeroJoinSemi{\JoinOp,L,\varnothing}\dsep
        \JoinIrr{\JoinOp,L,\varnothing}{P}\dsep P_\ast\in D_L(P)\dsep{}\\
      &\forall S\in D_L(P)
        ((P_\ast,S)\in\JoinOrd{\JoinOp}{L}\to S=P_\ast)
        \dsep Q\in D_L(P)\vdash
        (Q,P_\ast)\in\JoinOrd{\JoinOp}{L}
    \end{aligned}
  }{%
    \ZeroJoinSemi{\JoinOp,L,\varnothing}\dsep
    \JoinIrr{\JoinOp,L,\varnothing}{P}\dsep P_\ast\in D_L(P)\dsep
    \forall S\in D_L(P)
    ((P_\ast,S)\in\JoinOrd{\JoinOp}{L}\to S=P_\ast)
    \dsep Q\in D_L(P)\vdash (Q,P_\ast)\in\JoinOrd{\JoinOp}{L}
  }[JoinIrreducibleProperLowerMaximalDominates]
    \proofstepwidestar[1]{\ZeroJoinSemi{\JoinOp,L,\varnothing}}{\rA}
    \proofstepwidestar[2]{\JoinIrr{\JoinOp,L,\varnothing}{P}}{\rA}
    \proofstepwidestar[3]{P_\ast\in D_L(P)}{\rA}
    \proofstepwidestar[4]{%
      \forall S\in D_L(P)
      ((P_\ast,S)\in\JoinOrd{\JoinOp}{L}\to S=P_\ast)}{\rA}
    \proofstepwidestar[5]{Q\in D_L(P)}{\rA}
    \proofstepwidestar[1]{\Semi{\JoinOp,L}}{%
      \FormulaRefAuto{ZeroJoinBaseSemilatticeAxiom}[ax]{1}}
    \proofstepwidestar[3]{P_\ast\in L}{%
      \FormulaRefAuto{JoinIrreducibleProperLowerCarrier}[thm]{3}}
    \proofstepwidestar[5]{Q\in L}{%
      \FormulaRefAuto{JoinIrreducibleProperLowerCarrier}[thm]{5}}
    \proofstepwidestar[1,2,3,5]{P_\ast\JoinOp{}Q\in D_L(P)}{%
      \FormulaRefAuto{JoinIrreducibleProperLowerJoinClosed}[thm]{1,2,3,5}}
    \proofstepwidestar[1,3,5]{%
      (P_\ast,P_\ast\JoinOp{}Q)\in\JoinOrd{\JoinOp}{L}}{%
      \FormulaRefAuto{JoinUpperLeft}[thm]{6,7,8}}
    \proofstepwidestar[1,2,3,4,5]{P_\ast\JoinOp{}Q=P_\ast}{%
      \rRE{\rRE{9,\rUE{4}},10}}
    \proofstepwidestar[1,3,5]{%
      (Q,P_\ast\JoinOp{}Q)\in\JoinOrd{\JoinOp}{L}}{%
      \FormulaRefAuto{JoinUpperRight}[thm]{6,7,8}}
    \proofstepwidestar[1,2,3,4,5]{%
      (Q,P_\ast)\in\JoinOrd{\JoinOp}{L}}{\rIE{11,12}}
  \closeproofpartwideindR

  \proofpartwide{Schluss}
    \proofstepwidestar[1]{\Finite L}{\rA}
    \proofstepwidestar[2]{\ZeroJoinSemi{\JoinOp,L,\varnothing}}{\rA}
    \proofstepwidestar[3]{\JoinIrr{\JoinOp,L,\varnothing}{P}}{\rA}
    \proofstepwidestar[1,2,3]{%
      \exists P_\ast\in D_L(P)\,\forall Q\in D_L(P)\,
      ((P_\ast,Q)\in\JoinOrd{\JoinOp}{L}\to Q=P_\ast)}{%
      \FormulaRefAuto{JoinIrreducibleProperLowerMaximalExists}[thm]{1,2,3}}
    \proofstepwidestar[5]{%
      P_\ast\in D_L(P)\land
      \forall S\in D_L(P)
      ((P_\ast,S)\in\JoinOrd{\JoinOp}{L}\to S=P_\ast)}{\rA}
    \proofstepwidestar[5]{P_\ast\in D_L(P)}{\rAEa{5}}
    \proofstepwidestar[5]{%
      \forall S\in D_L(P)
      ((P_\ast,S)\in\JoinOrd{\JoinOp}{L}\to S=P_\ast)}{\rAEb{5}}
    \proofstepwidestar[2,3,5]{%
      Q\in D_L(P)\vdash(Q,P_\ast)\in\JoinOrd{\JoinOp}{L}}{%
      \FormulaRefAuto{JoinIrreducibleProperLowerMaximalDominates}[thm]{%
        2,3,6,7}}
    \proofstepwidestar[1,2,3]{%
      \exists P_\ast\in D_L(P)\,\forall Q\in D_L(P)
      (Q,P_\ast)\in\JoinOrd{\JoinOp}{L}}{%
      \rEE{4,5,\rEI{\rAI{6,8}}}}
  \closeproofpartwide
\end{tabproofsplitwide}

\FormulaThmDeltaK[Charakterisierung eines privaten Punkts]{%
  \begin{aligned}[t]
    &\ZeroJoinSemi{\JoinOp,L,\varnothing}\dsep\InterClosedFam L\dsep
      \forall Q,S\in L((Q,S)\in\JoinOrd{\JoinOp}{L}\leftrightarrow Q\subseteq S)\dsep{}\\
    &P_\ast\in D_L(P)\dsep
      \forall Q\in D_L(P)(Q,P_\ast)\in\JoinOrd{\JoinOp}{L}\dsep
      a\in P\setminus P_\ast\vdash{}\\
    &\forall Q\in L\,(a\in Q\leftrightarrow
      Q\in\PrincipalFilter{\JoinOrd{\JoinOp}{L},L}{P})
  \end{aligned}
}{JoinIrreduciblePrivatePointCharacterization}{%
  \DeltaRow{Mengen}{L\dsep P\dsep P_\ast\dsep Q\dsep S\dsep a}
  \DeltaRow{Funktionen}{\JoinOp}
}
\begin{tabproofwide}
  \proofstepwidestar[1]{\ZeroJoinSemi{\JoinOp,L,\varnothing}}{\rA}
  \proofstepwidestar[2]{\InterClosedFam L}{\rA}
  \proofstepwidestar[3]{\forall Q,S\in L
    ((Q,S)\in\JoinOrd{\JoinOp}{L}\leftrightarrow Q\subseteq S)}{\rA}
  \proofstepwidestar[4]{P_\ast\in D_L(P)}{\rA}
  \proofstepwidestar[5]{\forall Q\in D_L(P)
    (Q,P_\ast)\in\JoinOrd{\JoinOp}{L}}{\rA}
  \proofstepwidestar[6]{a\in P\setminus P_\ast}{\rA}
  \proofstepwidestar[1]{\Semi{\JoinOp,L}}{%
    \FormulaRefAuto{ZeroJoinBaseSemilatticeAxiom}[ax]{1}}
  \proofstepwidestar[4]{P_\ast\in L\land
    (P_\ast,P)\in\JoinOrd{\JoinOp}{L}\land P_\ast\neq P}{%
    \FormulaRefAuto{JoinIrreducibleProperLowerSet}[def]{4}}
  \proofstepwidestar[1,4]{P\in L}{%
    \FormulaRefAuto{JoinOrderSecondCarrier}[thm]{7,\rAEa{\rAEb{8}}}}
  \proofstepwidestar[6]{a\in P}{%
    \rAEa{\FormulaRefAuto{x \in A \setminus B \eqvdash
      x \in A \land x \notin B}{6}}}
  \proofstepwidestar[6]{a\notin P_\ast}{%
    \rAEb{\FormulaRefAuto{x \in A \setminus B \eqvdash
      x \in A \land x \notin B}{6}}}
  \proofstepwidestar[12]{Q\in L}{\rA}
  \proofstepwidestar[13]{a\in Q}{\rA}
  \proofstepwidestar[14]{P\nsubseteq Q}{\rA}
  \proofstepwidestar[2,9,12]{P\cap Q\in L}{%
    \FormulaRefAuto{IntersectionClosedFamily}[def]{2,9,12}}
  \proofstepwidestar[]{P\cap Q\subseteq P}{\FormulaRefAuto{A\cap B\subseteq A}}
  \proofstepwidestar[3,9,12]{(P\cap Q,P)\in\JoinOrd{\JoinOp}{L}}{%
    \FormulaRefAuto{P\leftrightarrow Q,Q\vdash P}{%
      \rUE{\rUE{3}},16}}
  \proofstepwidestar[14]{P\cap Q\neq P}{%
    \FormulaRefAuto{P\leftrightarrow Q,\neg P\vdash\neg Q}{%
      \FormulaRefAuto{A \subseteq B \eqvdash A \cap B = A},14}}
  \proofstepwidestar[3,9,12,14]{P\cap Q\in D_L(P)}{%
    \FormulaRefAuto{JoinIrreducibleProperLowerSet}[def]{15,17,18}}
  \proofstepwidestar[3,5,9,12,14]{%
    (P\cap Q,P_\ast)\in\JoinOrd{\JoinOp}{L}}{\rRE{19,\rUE{5}}}
  \proofstepwidestar[3,5,9,12,14]{P\cap Q\subseteq P_\ast}{%
    \FormulaRefAuto{P\leftrightarrow Q,P\vdash Q}{%
      \rUE{\rUE{3}},20}}
  \proofstepwidestar[10,13]{a\in P\cap Q}{%
    \FormulaRefAuto{x\in A\cap B\eqvdash x\in A\land x\in B}{%
      \rAI{10,13}}}
  \proofstepwidestar[3,5,9,10,12,13,14]{a\in P_\ast}{%
    \FormulaRefAuto{A\subseteq B,\,x\in A\vdash x\in B}{21,22}}
  \proofstepwidestar[3,5,6,9,10,12,13,14]{\bot}{\rBI{23,11}}
  \proofstepwidestar[1,2,3,4,5,6,12,13]{P\subseteq Q}{%
    \FormulaRefAuto{rDN}{\rCI{14,24}}}
  \proofstepwidestar[1,2,3,4,5,6,12]{a\in Q\to P\subseteq Q}{%
    \rRI{13,25}}
  \proofstepwidestar[27]{P\subseteq Q}{\rA}
  \proofstepwidestar[10,27]{a\in Q}{%
    \FormulaRefAuto{A\subseteq B,\,x\in A\vdash x\in B}{27,10}}
  \proofstepwidestar[10,12]{P\subseteq Q\to a\in Q}{\rRI{27,28}}
  \proofstepwidestar[1,2,3,4,5,6,12]{a\in Q\leftrightarrow P\subseteq Q}{%
    \FormulaRefAuto{P \leftrightarrow Q \eqvdash
      (P \rightarrow Q)\land(Q\rightarrow P)}{\rAI{26,29}}}
  \proofstepwidestar[12]{Q\in
    \PrincipalFilter{\JoinOrd{\JoinOp}{L},L}{P}\leftrightarrow
    (P,Q)\in\JoinOrd{\JoinOp}{L}}{%
    \FormulaRefAuto{PrincipalFilter}[def]{12}}
  \proofstepwidestar[3,9,12]{(P,Q)\in\JoinOrd{\JoinOp}{L}\leftrightarrow
    P\subseteq Q}{\rUE{\rUE{3}}}
  \proofstepwidestar[3,9,12]{Q\in
    \PrincipalFilter{\JoinOrd{\JoinOp}{L},L}{P}\leftrightarrow
    P\subseteq Q}{\rChain{31,32}}
  \proofstepwidestar[3,9,12]{P\subseteq Q\leftrightarrow Q\in
    \PrincipalFilter{\JoinOrd{\JoinOp}{L},L}{P}}{%
    \FormulaRefAuto{P\leftrightarrow Q\vdash Q\leftrightarrow P}{33}}
  \proofstepwidestar[1,2,3,4,5,6,12]{a\in Q\leftrightarrow
    Q\in\PrincipalFilter{\JoinOrd{\JoinOp}{L},L}{P}}{\rChain{30,34}}
  \proofstepwidestar[1,2,3,4,5,6]{\forall Q\in L\,(a\in Q\leftrightarrow
    Q\in\PrincipalFilter{\JoinOrd{\JoinOp}{L},L}{P})}{%
    \rUI{\rRI{12,35}}}
\end{tabproofwide}

\FormulaThmDeltaK[Privater Punkt eines \(\JoinOp\)-Irreduziblen]{%
  \begin{aligned}[t]
    &\Finite L\dsep\ZeroJoinSemi{\JoinOp,L,\varnothing}\dsep
      \InterClosedFam L\dsep{}\\
    &\forall Q,S\in L((Q,S)\in\JoinOrd{\JoinOp}{L}\leftrightarrow Q\subseteq S)
      \dsep\JoinIrr{\JoinOp,L,\varnothing}{P}\vdash{}\\
    &\exists a\in P\,\forall Q\in L\,(a\in Q\leftrightarrow
      Q\in\PrincipalFilter{\JoinOrd{\JoinOp}{L},L}{P})
  \end{aligned}
}{JoinIrreduciblePrivatePoint}{%
  \DeltaRow{Mengen}{L\dsep P\dsep P_\ast\dsep Q\dsep S\dsep a}
  \DeltaRow{Funktionen}{\JoinOp}
}
\begin{tabproofwide}
  \proofstepwidestar[1]{\Finite L}{\rA}
  \proofstepwidestar[2]{\ZeroJoinSemi{\JoinOp,L,\varnothing}}{\rA}
  \proofstepwidestar[3]{\InterClosedFam L}{\rA}
  \proofstepwidestar[4]{\forall Q,S\in L
    ((Q,S)\in\JoinOrd{\JoinOp}{L}\leftrightarrow Q\subseteq S)}{\rA}
  \proofstepwidestar[5]{\JoinIrr{\JoinOp,L,\varnothing}{P}}{\rA}
  \proofstepwidestar[1,2,5]{\exists P_\ast\in D_L(P)\,
    \forall Q\in D_L(P)(Q,P_\ast)\in\JoinOrd{\JoinOp}{L}}{%
    \FormulaRefAuto{JoinIrreducibleLargestProperLower}[thm]{1,2,5}}
  \proofstepwidestar[7]{P_\ast\in D_L(P)\land
    \forall Q\in D_L(P)(Q,P_\ast)\in\JoinOrd{\JoinOp}{L}}{\rA}
  \proofstepwidestar[7]{P_\ast\in L\land
    (P_\ast,P)\in\JoinOrd{\JoinOp}{L}\land P_\ast\neq P}{%
    \FormulaRefAuto{JoinIrreducibleProperLowerSet}[def]{\rAEa{7}}}
  \proofstepwidestar[4,7]{P_\ast\subseteq P}{%
    \FormulaRefAuto{P\leftrightarrow Q,P\vdash Q}{%
      \rUE{\rUE{4}},\rAEa{\rAEb{8}}}}
  \proofstepwidestar[4,7]{P_\ast\subsetneq P}{%
    \FormulaRefAuto{A\subset B\coloneq A\subseteq B\land A\neq B}[def]{%
      \rAI{9,\rAEb{\rAEb{8}}}}}
  \proofstepwidestar[4,7]{\exists a\in P\setminus P_\ast}{%
    \FormulaRefAuto{A\subset B\vdash \exists x\in B\setminus A}{10}}
  \proofstepwidestar[12]{a\in P\setminus P_\ast}{\rA}
  \proofstepwidestar[2,3,4,7,12]{\forall Q\in L\,(a\in Q\leftrightarrow
    Q\in\PrincipalFilter{\JoinOrd{\JoinOp}{L},L}{P})}{%
    \FormulaRefAuto{JoinIrreduciblePrivatePointCharacterization}[thm]{%
      2,3,4,\rAEa{7},\rAEb{7},12}}
  \proofstepwidestar[12]{a\in P}{%
    \rAEa{\FormulaRefAuto{x \in A \setminus B \eqvdash
      x \in A \land x \notin B}{12}}}
  \proofstepwidestar[1,2,3,4,5]{\exists a\in P\,\forall Q\in L\,
    (a\in Q\leftrightarrow
      Q\in\PrincipalFilter{\JoinOrd{\JoinOp}{L},L}{P})}{%
    \rEE{6,7,\rEE{11,12,\rEI{\rAI{14,13}}}}}
\end{tabproofwide}

\FormulaThmDeltaK[Komplementfaser der vermeidenden Teilfamilie]{%
  \begin{aligned}[t]
    &\FranklFam M\dsep H\subseteq\mathcal L_M\dsep a\in U_M\dsep{}\\
    &\forall Q\in\mathcal L_M\,(a\in Q\leftrightarrow Q\in H)
      \vdash
      \kappa_M^{-1}[H]=\FamAvoiding{M^\circ}{a}
  \end{aligned}
}{FranklComplementAvoidingFibre}{%
  \DeltaRow{Mengen}{M\dsep H\dsep Q\dsep X\dsep a}
  \DeltaRow{Funktionen}{\kappa_M}
}
\begin{tabproofwide}
  \proofstepwidestar[1]{\FranklFam M}{\rA}
  \proofstepwidestar[2]{H\subseteq\mathcal L_M}{\rA}
  \proofstepwidestar[3]{a\in U_M}{\rA}
  \proofstepwidestar[4]{\forall Q\in\mathcal L_M
    (a\in Q\leftrightarrow Q\in H)}{\rA}
  \proofstepwidestar[1]{\kappa_M\colon M^\circ\bij\mathcal L_M}{%
    \FormulaRefAuto{FranklComplementMapBijective}[thm]{1}}
  \proofstepwide[1,2]{X\in\kappa_M^{-1}[H]}{\leftrightarrow}{%
    X\in M^\circ\land\kappa_M(X)\in H}{%
    \FormulaRefAuto{F\colon A\to B\dsep C\subseteq B
      \vdash(x\in F^{-1}[C]\leftrightarrow(x\in A\land F(x)\in C))}{5,2}}
  \proofstepwide[1,2,4]{}{\leftrightarrow}{%
    X\in M^\circ\land a\in\kappa_M(X)}{%
    \rLRS{\FormulaRefAuto{P\leftrightarrow Q\vdash Q\leftrightarrow P}{%
      \rUE{4}},6}}
  \proofstepwide[1,2,4]{}{\leftrightarrow}{%
    X\in M^\circ\land a\in U_M\setminus X}{%
    \FormulaRefAuto{FranklComplementMap}[def]}
  \proofstepwide[1,2,3,4]{}{\leftrightarrow}{%
    X\in M^\circ\land a\notin X}{%
    \FormulaRefAuto{x\in A\setminus B\eqvdash x\in A\land x\notin B}{3}}
  \proofstepwide[1,2,3,4]{}{\leftrightarrow}{%
    X\in\FamAvoiding{M^\circ}{a}}{%
    \FormulaRefAuto{\FamAvoiding{M}{a}\coloneqq
      \{X\in M\mid a\notin X\}}}
  \proofstepwidestar[1,2,3,4]{%
    \kappa_M^{-1}[H]=\FamAvoiding{M^\circ}{a}}{%
    \FormulaRefAuto{\forall x\,(x\in A\leftrightarrow x\in B)
      \eqvdash A=B}{\rUI{10}}}
\end{tabproofwide}

\FormulaThmDeltaK[Komplementfaser der enthaltenden Teilfamilie]{%
  \begin{aligned}[t]
    &\FranklFam M\dsep H\subseteq\mathcal L_M\dsep a\in U_M\dsep{}\\
    &\forall Q\in\mathcal L_M\,(a\in Q\leftrightarrow Q\in H)
      \vdash
      \kappa_M^{-1}[\mathcal L_M\setminus H]=
      \FamContaining{M^\circ}{a}
  \end{aligned}
}{FranklComplementContainingFibre}{%
  \DeltaRow{Mengen}{M\dsep H\dsep Q\dsep X\dsep a}
  \DeltaRow{Funktionen}{\kappa_M}
}
\begin{tabproofwide}
  \proofstepwidestar[1]{\FranklFam M}{\rA}
  \proofstepwidestar[2]{H\subseteq\mathcal L_M}{\rA}
  \proofstepwidestar[3]{a\in U_M}{\rA}
  \proofstepwidestar[4]{\forall Q\in\mathcal L_M
    (a\in Q\leftrightarrow Q\in H)}{\rA}
  \proofstepwidestar[1]{\kappa_M\colon M^\circ\bij\mathcal L_M}{%
    \FormulaRefAuto{FranklComplementMapBijective}[thm]{1}}
  \proofstepwide[1]{X\in\kappa_M^{-1}[\mathcal L_M\setminus H]}
    {\leftrightarrow}{%
    X\in M^\circ\land\kappa_M(X)\in\mathcal L_M\setminus H}{%
    \FormulaRefAuto{F\colon A\to B\dsep C\subseteq B
      \vdash(x\in F^{-1}[C]\leftrightarrow(x\in A\land F(x)\in C))}{%
      5,\FormulaRefAuto{A\setminus B\subseteq A}}}
  \proofstepwide[1,4]{}{\leftrightarrow}{%
    X\in M^\circ\land a\notin\kappa_M(X)}{%
    \FormulaRefAuto{x\in A\setminus B\eqvdash x\in A\land x\notin B};
    \FormulaRefAuto{P\leftrightarrow Q\eqvdash\neg P\leftrightarrow\neg Q}[thm]{%
      \rUE{4}}}
  \proofstepwide[1,3,4]{}{\leftrightarrow}{%
    X\in M^\circ\land a\in X}{%
    \FormulaRefAuto{FranklComplementMap}[def];
    \FormulaRefAuto{x\in A\setminus B\eqvdash x\in A\land x\notin B}{3};
    \FormulaRefAuto{rDN}}
  \proofstepwide[1,3,4]{}{\leftrightarrow}{%
    X\in\FamContaining{M^\circ}{a}}{%
    \FormulaRefAuto{\FamContaining{M}{a}\coloneqq
      \{X\in M\mid a\in X\}}}
  \proofstepwidestar[1,2,3,4]{%
    \kappa_M^{-1}[\mathcal L_M\setminus H]=
    \FamContaining{M^\circ}{a}}{%
    \FormulaRefAuto{\forall x\,(x\in A\leftrightarrow x\in B)
      \eqvdash A=B}{\rUI{9}}}
\end{tabproofwide}

\FormulaThmDeltaK[Transport eines privaten Punkts]{%
  \begin{aligned}[t]
    &\FranklFam M\dsep H\subseteq\mathcal L_M\dsep a\in U_M\dsep{}\\
    &\forall Q\in\mathcal L_M\,(a\in Q\leftrightarrow Q\in H)
      \dsep\AtMostHalf{H}{\mathcal L_M}
      \vdash\HalfOcc{a}{M}
  \end{aligned}
}{FranklComplementPrivatePointTransport}{%
  \DeltaRow{Mengen}{M\dsep H\dsep Q\dsep X\dsep a}
  \DeltaRow{Funktionen}{\kappa_M\dsep F}
}
\begin{tabproofwide}
  \proofstepwidestar[1]{\FranklFam M}{\rA}
  \proofstepwidestar[2]{H\subseteq\mathcal L_M}{\rA}
  \proofstepwidestar[3]{a\in U_M}{\rA}
  \proofstepwidestar[4]{\forall Q\in\mathcal L_M
    (a\in Q\leftrightarrow Q\in H)}{\rA}
  \proofstepwidestar[5]{\AtMostHalf{H}{\mathcal L_M}}{\rA}
  \proofstepwidestar[1]{\kappa_M\colon M^\circ\bij\mathcal L_M}{%
    \FormulaRefAuto{FranklComplementMapBijective}[thm]{1}}
  \proofstepwidestar[1,2,3,4]{\kappa_M^{-1}[H]=
    \FamAvoiding{M^\circ}{a}}{%
    \FormulaRefAuto{FranklComplementAvoidingFibre}[thm]{1,2,3,4}}
  \proofstepwidestar[1,2,3,4]{%
    \kappa_M^{-1}[\mathcal L_M\setminus H]=
    \FamContaining{M^\circ}{a}}{%
    \FormulaRefAuto{FranklComplementContainingFibre}[thm]{1,2,3,4}}
  \proofstepwidestar[1,2,5]{%
    \exists F\colon\kappa_M^{-1}[H]
      \inj\kappa_M^{-1}[\mathcal L_M\setminus H]}{%
    \FormulaRefAuto{BijectionAtMostHalfTransport}[thm]{6,2,5}}
  \proofstepwidestar[1,2,3,4,5]{%
    \exists F\colon\FamAvoiding{M^\circ}{a}
      \inj\FamContaining{M^\circ}{a}}{\rIE{7,\rIE{8,9}}}
  \proofstepwidestar[1]{M\subseteq M^\circ}{%
    \rIE{\FormulaRefAuto{FranklComplementConstruction}[def]{1},
      \FormulaRefAuto{A\subseteq A\cup B}}}
  \proofstepwidestar[1]{\FamAvoiding{M}{a}\subseteq
    \FamAvoiding{M^\circ}{a}}{%
    \FormulaRefAuto{FamAvoidingMonotone}[thm]{11}}
  \proofstepwidestar[1]{\FamContaining{M^\circ}{a}=
    \FamContaining{M}{a}}{%
    \FormulaRefAuto{FranklAugmentedContainingUnchanged}[thm]{1}}
  \proofstepwidestar[14]{F\colon\FamAvoiding{M^\circ}{a}
    \inj\FamContaining{M^\circ}{a}}{\rA}
  \proofstepwidestar[12,14]{F\restriction_{\FamAvoiding{M}{a}}
    \colon\FamAvoiding{M}{a}\inj\FamContaining{M^\circ}{a}}{%
    \FormulaRefAuto{F\colon A\inj B\dsep C\subseteq A\vdash
      F\restriction_C\colon C\inj B}{14,12}}
  \proofstepwidestar[12,13,14]{\exists F\colon\FamAvoiding{M}{a}
    \inj\FamContaining{M}{a}}{\rEI{\rIE{13,15}}}
  \proofstepwidestar[1,2,3,4,5]{\exists F\colon\FamAvoiding{M}{a}
    \inj\FamContaining{M}{a}}{\rEE{10,14,16}}
  \proofstepwidestar[1,2,3,4,5]{\HalfOcc{a}{M}}{%
    \FormulaRefAuto{Halbhäufiges Element}[def]{3,17}}
\end{tabproofwide}

\FormulaDefDeltaK[Hauptfilter in der Komplementfamilie]{%
  H_M(P)\coloneqq
  \PrincipalFilter{\JoinOrd{\FamilyJoinOp{M}}{\mathcal L_M},\mathcal L_M}{P}%
}{FranklComplementPrincipalFilterAbbreviation}{%
  \DeltaRow{Mengen}{M\dsep P}
  \DeltaRow{Relationen}{\JoinOrd{\FamilyJoinOp{M}}{\mathcal L_M}}
  \DeltaRow{\textbf{Neues Symbol}}{}
  \DeltaPrem{Hauptfilter in der Komplementfamilie}{H_M(P)}
}

\FormulaThmDeltaK[Halbverbandsform impliziert Mengenform]{%
  \begin{aligned}[t]
    &\FranklFam M\dsep{}\\
    &\forall A\,\forall \JoinOp\,\forall o\,\Bigl(
      \Finite A\land\ZeroJoinSemi{\JoinOp,A,o}\land{}\\
    &\qquad
      \exists x\in A\,(x\neq o)
      \to\FranklSemi{\JoinOp,A,o}\Bigr)
      \vdash\Frankl M
  \end{aligned}
}{FranklSemilatticeImpliesSet}{%
  \DeltaRow{Mengen}{A\dsep M\dsep\mathcal L_M\dsep
    P\dsep Q\dsep X\dsep a\dsep o}
  \DeltaRow{Funktionen}{\JoinOp\dsep\FamilyJoinOp{M}\dsep F}
}

\begin{tabproofsplitwide}
  \proofpartwideindR[Anwendung auf die Komplementfamilie]{%
    \begin{aligned}[t]
      &\forall A\,\forall\JoinOp\,\forall o\,\Bigl(
        \Finite A\land\ZeroJoinSemi{\JoinOp,A,o}\land
        \exists x\in A\,(x\neq o)\\[-2pt]
      &\qquad\to\FranklSemi{\JoinOp,A,o}\Bigr)
        \dsep\FranklFam M\vdash
        \FranklSemi{\FamilyJoinOp{M},\mathcal L_M,\varnothing}
    \end{aligned}
  }{%
    \begin{aligned}[t]
      &\forall A\,\forall\JoinOp\,\forall o\,\Bigl(
        \Finite A\land\ZeroJoinSemi{\JoinOp,A,o}\land
        \exists x\in A\,(x\neq o)\\[-2pt]
      &\qquad\to\FranklSemi{\JoinOp,A,o}\Bigr)
        \dsep\FranklFam M\vdash
        \FranklSemi{\FamilyJoinOp{M},\mathcal L_M,\varnothing}
    \end{aligned}
  }[FranklComplementSemilatticeInstance]
    \proofstepwidestar[1]{%
      \forall A\,\forall\JoinOp\,\forall o\,\Bigl(
        \Finite A\land\ZeroJoinSemi{\JoinOp,A,o}\land
        \exists x\in A\,(x\neq o)
        \to\FranklSemi{\JoinOp,A,o}\Bigr)}{\rA}
    \proofstepwidestar[2]{\FranklFam M}{\rA}
    \proofstepwidestar[2]{\Finite{\mathcal L_M}}{%
      \FormulaRefAuto{FranklComplementFamilyFinite}[thm]{2}}
    \proofstepwidestar[2]{%
      \ZeroJoinSemi{\FamilyJoinOp{M},\mathcal L_M,\varnothing}}{%
      \FormulaRefAuto{FranklComplementFiniteLattice}[thm]{2}}
    \proofstepwidestar[2]{%
      \exists X\in\mathcal L_M\,(X\neq\varnothing)}{%
      \FormulaRefAuto{FranklComplementNontrivial}[thm]{2}}
    \proofstepwidestar[1]{%
      \begin{aligned}[t]
        \forall\JoinOp\,\forall o\,\Bigl(&
          \Finite{\mathcal L_M}\land
          \ZeroJoinSemi{\JoinOp,\mathcal L_M,o}\land{}\\[-2pt]
        &\exists x\in\mathcal L_M\,(x\neq o)
          \to\FranklSemi{\JoinOp,\mathcal L_M,o}\Bigr)
      \end{aligned}}{\rUE{1}}
    \proofstepwidestar[1]{%
      \begin{aligned}[t]
        \forall o\,\Bigl(&\Finite{\mathcal L_M}\land
          \ZeroJoinSemi{\FamilyJoinOp{M},\mathcal L_M,o}\land{}\\[-2pt]
        &\exists x\in\mathcal L_M\,(x\neq o)
          \to\FranklSemi{\FamilyJoinOp{M},\mathcal L_M,o}\Bigr)
      \end{aligned}}{\rUE{6}}
    \proofstepwidestar[1]{%
      \begin{aligned}[t]
        &\Finite{\mathcal L_M}\land
          \ZeroJoinSemi{\FamilyJoinOp{M},\mathcal L_M,\varnothing}\land
          \exists x\in\mathcal L_M\,(x\neq\varnothing)\\[-2pt]
        &\qquad\to
          \FranklSemi{\FamilyJoinOp{M},\mathcal L_M,\varnothing}
      \end{aligned}}{\rUE{7}}
    \proofstepwidestar[2]{%
      \begin{aligned}[t]
        &\Finite{\mathcal L_M}\land
          \ZeroJoinSemi{\FamilyJoinOp{M},\mathcal L_M,\varnothing}\land{}\\[-2pt]
        &\qquad\exists x\in\mathcal L_M\,(x\neq\varnothing)
      \end{aligned}}{\rAI{3,\rAI{4,5}}}
    \proofstepwidestar[1,2]{%
      \FranklSemi{\FamilyJoinOp{M},\mathcal L_M,\varnothing}}{\rRE{8,9}}
  \closeproofpartwideindR

  \proofpartwideindR[Punktweise Form der Komplementordnung]{%
    \begin{aligned}[t]
      \FranklFam M\vdash
      \forall Q,S\in\mathcal L_M\,\bigl(&
        (Q,S)\in\JoinOrd{\FamilyJoinOp{M}}{\mathcal L_M}\\[-2pt]
        &\leftrightarrow Q\subseteq S\bigr)
    \end{aligned}
  }{%
    \FranklFam M\vdash
    \forall Q,S\in\mathcal L_M\,
      ((Q,S)\in\JoinOrd{\FamilyJoinOp{M}}{\mathcal L_M}
      \leftrightarrow Q\subseteq S)
  }[FranklComplementOrderPointwise]
    \proofstepwidestar[1]{\FranklFam M}{\rA}
    \proofstepwidestar[1]{%
      \JoinOrd{\FamilyJoinOp{M}}{\mathcal L_M}=\mathcal R_M}{%
      \FormulaRefAuto{FranklComplementOrderRecovered}[thm]{1}}
    \proofstepwidestar[1]{%
      \mathcal R_M=\JoinOrd{\FamilyJoinOp{M}}{\mathcal L_M}}{%
      \FormulaRefAuto{a=b\vdash b=a}{2}}
    \proofstepwidestar[4]{Q\in\mathcal L_M}{\rA}
    \proofstepwidestar[5]{S\in\mathcal L_M}{\rA}
    \proofstepwidestar[6]{%
      (Q,S)\in\JoinOrd{\FamilyJoinOp{M}}{\mathcal L_M}}{\rA}
    \proofstepwidestar[1,6]{(Q,S)\in\mathcal R_M}{\rIE{2,6}}
    \proofstepwidestar[1,6]{Q\subseteq S}{%
      \FormulaRefAuto{FranklComplementConstruction}[def]{7}}
    \proofstepwidestar[1,4,5]{%
      (Q,S)\in\JoinOrd{\FamilyJoinOp{M}}{\mathcal L_M}
      \to Q\subseteq S}{\rRI{6,8}}
    \proofstepwidestar[10]{Q\subseteq S}{\rA}
    \proofstepwidestar[4,5,10]{(Q,S)\in\mathcal R_M}{%
      \FormulaRefAuto{FranklComplementConstruction}[def]{4,5,10}}
    \proofstepwidestar[1,4,5,10]{%
      (Q,S)\in\JoinOrd{\FamilyJoinOp{M}}{\mathcal L_M}}{\rIE{3,11}}
    \proofstepwidestar[1,4,5]{%
      Q\subseteq S\to
      (Q,S)\in\JoinOrd{\FamilyJoinOp{M}}{\mathcal L_M}}{\rRI{10,12}}
    \proofstepwidestar[1,4,5]{%
      (Q,S)\in\JoinOrd{\FamilyJoinOp{M}}{\mathcal L_M}
      \leftrightarrow Q\subseteq S}{%
      \FormulaRefAuto{P \leftrightarrow Q \eqvdash
        (P \rightarrow Q)\land(Q\rightarrow P)}{\rAI{9,13}}}
    \proofstepwidestar[1,4]{%
      \forall S\in\mathcal L_M\,
      ((Q,S)\in\JoinOrd{\FamilyJoinOp{M}}{\mathcal L_M}
      \leftrightarrow Q\subseteq S)}{\rUI{\rRI{5,14}}}
    \proofstepwidestar[1]{%
      \forall Q,S\in\mathcal L_M\,
      ((Q,S)\in\JoinOrd{\FamilyJoinOp{M}}{\mathcal L_M}
      \leftrightarrow Q\subseteq S)}{\rUI{\rRI{4,15}}}
  \closeproofpartwideindR

  \proofpartwideindR[Irreduzibler Hauptfilter]{%
    \begin{aligned}[t]
      &\FranklSemi{\FamilyJoinOp{M},\mathcal L_M,\varnothing}\vdash{}\\[-2pt]
      &\exists P\Bigl(
        \JoinIrr{\FamilyJoinOp{M},\mathcal L_M,\varnothing}{P}
        \land\AtMostHalf{H_M(P)}{\mathcal L_M}\Bigr)
    \end{aligned}
  }{%
    \begin{aligned}[t]
      &\FranklSemi{\FamilyJoinOp{M},\mathcal L_M,\varnothing}\vdash{}\\[-2pt]
      &\exists P\Bigl(
        \JoinIrr{\FamilyJoinOp{M},\mathcal L_M,\varnothing}{P}
        \land\AtMostHalf{H_M(P)}{\mathcal L_M}\Bigr)
    \end{aligned}
  }[FranklComplementIrreducibleFilter]
    \proofstepwidestar[1]{%
      \FranklSemi{\FamilyJoinOp{M},\mathcal L_M,\varnothing}}{\rA}
    \proofstepwidestar[1]{%
      \begin{aligned}[t]
        \exists P\Bigl(&
          \JoinIrr{\FamilyJoinOp{M},\mathcal L_M,\varnothing}{P}\land{}\\[-2pt]
        &\AtMostHalf{%
          \PrincipalFilter{\JoinOrd{\FamilyJoinOp{M}}{\mathcal L_M},
            \mathcal L_M}{P}}{\mathcal L_M}\Bigr)
      \end{aligned}}{%
      \FormulaRefAuto{FranklSemilatticeProperty}[def]{1}}
    \proofstepwidestar[3]{%
      \begin{aligned}[t]
        &\JoinIrr{\FamilyJoinOp{M},\mathcal L_M,\varnothing}{P}\land{}\\[-2pt]
        &\AtMostHalf{%
          \PrincipalFilter{\JoinOrd{\FamilyJoinOp{M}}{\mathcal L_M},
            \mathcal L_M}{P}}{\mathcal L_M}
      \end{aligned}}{\rA}
    \proofstepwidestar[3]{%
      \JoinIrr{\FamilyJoinOp{M},\mathcal L_M,\varnothing}{P}}{\rAEa{3}}
    \proofstepwidestar[3]{%
      \AtMostHalf{%
        \PrincipalFilter{\JoinOrd{\FamilyJoinOp{M}}{\mathcal L_M},
          \mathcal L_M}{P}}{\mathcal L_M}}{\rAEb{3}}
    \proofstepwidestar[]{%
      H_M(P)=\PrincipalFilter{%
        \JoinOrd{\FamilyJoinOp{M}}{\mathcal L_M},\mathcal L_M}{P}}{%
      \FormulaRefAuto{FranklComplementPrincipalFilterAbbreviation}[def]}
    \proofstepwidestar[]{%
      \PrincipalFilter{\JoinOrd{\FamilyJoinOp{M}}{\mathcal L_M},
        \mathcal L_M}{P}=H_M(P)}{%
      \FormulaRefAuto{a=b\vdash b=a}{6}}
    \proofstepwidestar[3]{\AtMostHalf{H_M(P)}{\mathcal L_M}}{\rIE{7,5}}
    \proofstepwidestar[3]{%
      \exists P\Bigl(
        \JoinIrr{\FamilyJoinOp{M},\mathcal L_M,\varnothing}{P}
        \land\AtMostHalf{H_M(P)}{\mathcal L_M}\Bigr)}{%
      \rEI{\rAI{4,8}}}
    \proofstepwidestar[1]{%
      \exists P\Bigl(
        \JoinIrr{\FamilyJoinOp{M},\mathcal L_M,\varnothing}{P}
        \land\AtMostHalf{H_M(P)}{\mathcal L_M}\Bigr)}{%
      \rEE{2,3,9}}
  \closeproofpartwideindR

  \proofpartwideindR[Privater Punkt des Hauptfilters]{%
    \begin{aligned}[t]
      &\FranklFam M\dsep
        \JoinIrr{\FamilyJoinOp{M},\mathcal L_M,\varnothing}{P}\vdash{}\\[-2pt]
      &\exists a\in U_M\,\forall Q\in\mathcal L_M\,
        (a\in Q\leftrightarrow Q\in H_M(P))
    \end{aligned}
  }{%
    \begin{aligned}[t]
      &\FranklFam M\dsep
        \JoinIrr{\FamilyJoinOp{M},\mathcal L_M,\varnothing}{P}\vdash{}\\[-2pt]
      &\exists a\in U_M\,\forall Q\in\mathcal L_M\,
        (a\in Q\leftrightarrow Q\in H_M(P))
    \end{aligned}
  }[FranklComplementPrivateFilterPoint]
    \proofstepwidestar[1]{\FranklFam M}{\rA}
    \proofstepwidestar[2]{%
      \JoinIrr{\FamilyJoinOp{M},\mathcal L_M,\varnothing}{P}}{\rA}
    \proofstepwidestar[1]{\Finite{\mathcal L_M}}{%
      \FormulaRefAuto{FranklComplementFamilyFinite}[thm]{1}}
    \proofstepwidestar[1]{%
      \ZeroJoinSemi{\FamilyJoinOp{M},\mathcal L_M,\varnothing}}{%
      \FormulaRefAuto{FranklComplementFiniteLattice}[thm]{1}}
    \proofstepwidestar[1]{\InterClosedFam{\mathcal L_M}}{%
      \FormulaRefAuto{FranklComplementFamilyIntersectionClosed}[thm]{1}}
    \proofstepwidestar[1]{%
      \forall Q,S\in\mathcal L_M\,
      ((Q,S)\in\JoinOrd{\FamilyJoinOp{M}}{\mathcal L_M}
      \leftrightarrow Q\subseteq S)}{%
      \FormulaRefAuto{FranklComplementOrderPointwise}[thm]{1}}
    \proofstepwidestar[1,2]{%
      \begin{aligned}[t]
        \exists a\in P\,\forall Q\in\mathcal L_M\,
        \bigl(a\in Q\leftrightarrow
          Q\in\PrincipalFilter{%
            \JoinOrd{\FamilyJoinOp{M}}{\mathcal L_M},
            \mathcal L_M}{P}\bigr)
      \end{aligned}}{%
      \FormulaRefAuto{JoinIrreduciblePrivatePoint}[thm]{3,4,5,6,2}}
    \proofstepwidestar[8]{%
      \begin{aligned}[t]
        a\in P\land\forall Q\in\mathcal L_M\,
        \bigl(a\in Q\leftrightarrow
          Q\in\PrincipalFilter{%
            \JoinOrd{\FamilyJoinOp{M}}{\mathcal L_M},
            \mathcal L_M}{P}\bigr)
      \end{aligned}}{\rA}
    \proofstepwidestar[8]{a\in P}{\rAEa{8}}
    \proofstepwidestar[8]{%
      \forall Q\in\mathcal L_M\,
      \bigl(a\in Q\leftrightarrow
        Q\in\PrincipalFilter{%
          \JoinOrd{\FamilyJoinOp{M}}{\mathcal L_M},
          \mathcal L_M}{P}\bigr)}{\rAEb{8}}
    \proofstepwidestar[2]{P\in\mathcal L_M}{%
      \FormulaRefAuto{JoinIrreducible}[def]{2}}
    \proofstepwidestar[1,2]{P\subseteq U_M}{%
      \FormulaRefAuto{FranklComplementMemberSubsetCarrier}[thm]{1,11}}
    \proofstepwidestar[1,2,8]{a\in U_M}{%
      \FormulaRefAuto{A\subseteq B,\,x\in A\vdash x\in B}{12,9}}
    \proofstepwidestar[]{%
      H_M(P)=\PrincipalFilter{%
        \JoinOrd{\FamilyJoinOp{M}}{\mathcal L_M},\mathcal L_M}{P}}{%
      \FormulaRefAuto{FranklComplementPrincipalFilterAbbreviation}[def]}
    \proofstepwidestar[]{%
      \PrincipalFilter{\JoinOrd{\FamilyJoinOp{M}}{\mathcal L_M},
        \mathcal L_M}{P}=H_M(P)}{%
      \FormulaRefAuto{a=b\vdash b=a}{14}}
    \proofstepwidestar[8]{%
      \forall Q\in\mathcal L_M\,(a\in Q\leftrightarrow Q\in H_M(P))}{%
      \rIE{15,10}}
    \proofstepwidestar[1,2,8]{%
      a\in U_M\land
      \forall Q\in\mathcal L_M\,(a\in Q\leftrightarrow Q\in H_M(P))}{%
      \rAI{13,16}}
    \proofstepwidestar[1,2]{%
      \exists a\in U_M\,\forall Q\in\mathcal L_M\,
        (a\in Q\leftrightarrow Q\in H_M(P))}{%
      \rEE{7,8,\rEI{17}}}
  \closeproofpartwideindR

  \proofpartwideindR[Transport des privaten Punkts]{%
    \begin{aligned}[t]
      &\FranklFam M\dsep\AtMostHalf{H_M(P)}{\mathcal L_M}\dsep
        a\in U_M\dsep{}\\[-2pt]
      &\forall Q\in\mathcal L_M\,(a\in Q\leftrightarrow Q\in H_M(P))
        \vdash\HalfOcc{a}{M}
    \end{aligned}
  }{%
    \begin{aligned}[t]
      &\FranklFam M\dsep\AtMostHalf{H_M(P)}{\mathcal L_M}\dsep
        a\in U_M\dsep{}\\[-2pt]
      &\forall Q\in\mathcal L_M\,(a\in Q\leftrightarrow Q\in H_M(P))
        \vdash\HalfOcc{a}{M}
    \end{aligned}
  }[FranklComplementPrivatePointHalfOccurrence]
    \proofstepwidestar[1]{\FranklFam M}{\rA}
    \proofstepwidestar[2]{\AtMostHalf{H_M(P)}{\mathcal L_M}}{\rA}
    \proofstepwidestar[3]{a\in U_M}{\rA}
    \proofstepwidestar[4]{%
      \forall Q\in\mathcal L_M\,(a\in Q\leftrightarrow Q\in H_M(P))}{\rA}
    \proofstepwidestar[]{%
      H_M(P)=\PrincipalFilter{%
        \JoinOrd{\FamilyJoinOp{M}}{\mathcal L_M},\mathcal L_M}{P}}{%
      \FormulaRefAuto{FranklComplementPrincipalFilterAbbreviation}[def]}
    \proofstepwidestar[]{%
      \PrincipalFilter{\JoinOrd{\FamilyJoinOp{M}}{\mathcal L_M},
        \mathcal L_M}{P}=H_M(P)}{%
      \FormulaRefAuto{a=b\vdash b=a}{5}}
    \proofstepwidestar[]{%
      \PrincipalFilter{\JoinOrd{\FamilyJoinOp{M}}{\mathcal L_M},
        \mathcal L_M}{P}\subseteq\mathcal L_M}{%
      \FormulaRefAuto{PrincipalFilterSubsetCarrier}[thm]}
    \proofstepwidestar[]{H_M(P)\subseteq\mathcal L_M}{\rIE{6,7}}
    \proofstepwidestar[1,2,3,4]{\HalfOcc{a}{M}}{%
      \FormulaRefAuto{FranklComplementPrivatePointTransport}[thm]{1,8,3,4,2}}
  \closeproofpartwideindR

  \proofpartwide{Schluss}
    \proofstepwidestar[1]{%
      \forall A\,\forall\JoinOp\,\forall o\,\Bigl(
        \Finite A\land\ZeroJoinSemi{\JoinOp,A,o}\land
        \exists x\in A\,(x\neq o)
        \to\FranklSemi{\JoinOp,A,o}\Bigr)}{\rA}
    \proofstepwidestar[2]{\FranklFam M}{\rA}
    \proofstepwidestar[1,2]{%
      \FranklSemi{\FamilyJoinOp{M},\mathcal L_M,\varnothing}}{%
      \FormulaRefAuto{FranklComplementSemilatticeInstance}[thm]{1,2}}
    \proofstepwidestar[1,2]{%
      \exists P\Bigl(
        \JoinIrr{\FamilyJoinOp{M},\mathcal L_M,\varnothing}{P}
        \land\AtMostHalf{H_M(P)}{\mathcal L_M}\Bigr)}{%
      \FormulaRefAuto{FranklComplementIrreducibleFilter}[thm]{3}}
    \proofstepwidestar[5]{%
      \JoinIrr{\FamilyJoinOp{M},\mathcal L_M,\varnothing}{P}
      \land\AtMostHalf{H_M(P)}{\mathcal L_M}}{\rA}
    \proofstepwidestar[5]{%
      \JoinIrr{\FamilyJoinOp{M},\mathcal L_M,\varnothing}{P}}{\rAEa{5}}
    \proofstepwidestar[5]{\AtMostHalf{H_M(P)}{\mathcal L_M}}{\rAEb{5}}
    \proofstepwidestar[2,5]{%
      \exists a\in U_M\,\forall Q\in\mathcal L_M\,
        (a\in Q\leftrightarrow Q\in H_M(P))}{%
      \FormulaRefAuto{FranklComplementPrivateFilterPoint}[thm]{2,6}}
    \proofstepwidestar[9]{%
      a\in U_M\land
      \forall Q\in\mathcal L_M\,(a\in Q\leftrightarrow Q\in H_M(P))}{\rA}
    \proofstepwidestar[9]{a\in U_M}{\rAEa{9}}
    \proofstepwidestar[9]{%
      \forall Q\in\mathcal L_M\,(a\in Q\leftrightarrow Q\in H_M(P))}{\rAEb{9}}
    \proofstepwidestar[2,5,9]{\HalfOcc{a}{M}}{%
      \FormulaRefAuto{FranklComplementPrivatePointHalfOccurrence}[thm]{%
        2,7,10,11}}
    \proofstepwidestar[2,5,9]{\Frankl M}{%
      \FormulaRefAuto{Frankl-Eigenschaft}[def]{\rEI{12}}}
    \proofstepwidestar[1,2]{\Frankl M}{%
      \rEE{4,5,\rEE{8,9,13}}}
  \closeproofpartwide
\end{tabproofsplitwide}

\section{Globaler Äquivalenzsatz}

\FormulaDefDeltaK[Globale Mengenform der Frankl-Vermutung]{%
  \FranklSetStatement\coloneqq
  \forall M\,(\FranklFam M\to\Frankl M)%
}{FranklGlobalSetStatement}{%
  \DeltaRow{Mengen}{M}
  \DeltaRow{\textbf{Neues Symbol}}{}
  \DeltaPrem{Globale Mengenform}{\FranklSetStatement}
}

\FormulaDefDeltaK[Globale Halbverbandsform der Frankl-Vermutung]{%
  \begin{aligned}[t]
    \FranklJoinStatement\coloneqq
    \forall A\,\forall\JoinOp\,\forall o\,\Bigl(&
      \Finite A\land\ZeroJoinSemi{\JoinOp,A,o}\land{}\\[-2pt]
      &\exists x\in A\,(x\neq o)
      \to\FranklSemi{\JoinOp,A,o}\Bigr)
  \end{aligned}
}{FranklGlobalJoinStatement}{%
  \DeltaRow{Mengen}{A\dsep x\dsep o}
  \DeltaRow{Funktionen}{\JoinOp}
  \DeltaRow{\textbf{Neues Symbol}}{}
  \DeltaPrem{Globale Halbverbandsform}{\FranklJoinStatement}
}

\FormulaThmDeltaK[Globale Äquivalenz der beiden Frankl-Formen]{%
  \FranklSetStatement\leftrightarrow\FranklJoinStatement
}{FranklSetSemilatticeEquivalence}{%
  \DeltaRow{Mengen}{A\dsep M\dsep x\dsep o}
  \DeltaRow{Funktionen}{\JoinOp}
}
Die Lemmon-Tabelle verwendet die beiden registrierten Kurzzeichen. Jede
Entfaltung wird ausdrücklich auf die zugehörige Definition gestützt.
\begin{tabproofsplitwide}
  \proofpartwideindR[Von der Mengenform zur Halbverbandsform]{%
    \FranklSetStatement\vdash\FranklJoinStatement
  }{%
    \FranklSetStatement\vdash\FranklJoinStatement
  }[FranklSetToJoinGlobal]
    \proofstepwidestar[1]{\FranklSetStatement}{\rA}
    \proofstepwidestar[1]{\forall M\,(\FranklFam M\to\Frankl M)}{%
      \FormulaRefAuto{FranklGlobalSetStatement}[def]{1}}
    \proofstepwidestar[3]{\Finite A\land
      \ZeroJoinSemi{\JoinOp,A,o}\land\exists x\in A\,(x\neq o)}{\rA}
    \proofstepwidestar[3]{\Finite A}{\rAEa{3}}
    \proofstepwidestar[3]{\ZeroJoinSemi{\JoinOp,A,o}}{\rAEa{\rAEb{3}}}
    \proofstepwidestar[3]{\exists x\in A\,(x\neq o)}{\rAEb{\rAEb{3}}}
    \proofstepwidestar[1,3]{\FranklSemi{\JoinOp,A,o}}{%
      \FormulaRefAuto{FranklSetImpliesSemilattice}[thm]{2,4,5,6}}
    \proofstepwidestar[1]{%
      \begin{aligned}[t]
        &\Finite A\land\ZeroJoinSemi{\JoinOp,A,o}\land
          \exists x\in A\,(x\neq o)\\[-2pt]
        &\qquad\to\FranklSemi{\JoinOp,A,o}
      \end{aligned}}{\rRI{3,7}}
    \proofstepwidestar[1]{%
      \begin{aligned}[t]
        \forall A\,\forall\JoinOp\,\forall o\,\Bigl(&
          \Finite A\land\ZeroJoinSemi{\JoinOp,A,o}\land{}\\[-2pt]
        &\exists x\in A\,(x\neq o)
          \to\FranklSemi{\JoinOp,A,o}\Bigr)
      \end{aligned}}{\rUI{\rUI{\rUI{8}}}}
    \proofstepwidestar[1]{\FranklJoinStatement}{%
      \FormulaRefAuto{FranklGlobalJoinStatement}[def]{9}}
  \closeproofpartwideindR

  \proofpartwideindR[Von der Halbverbandsform zur Mengenform]{%
    \FranklJoinStatement\vdash\FranklSetStatement
  }{%
    \FranklJoinStatement\vdash\FranklSetStatement
  }[FranklJoinToSetGlobal]
    \proofstepwidestar[1]{\FranklJoinStatement}{\rA}
    \proofstepwidestar[1]{%
      \begin{aligned}[t]
        \forall A\,\forall\JoinOp\,\forall o\,\Bigl(&
          \Finite A\land\ZeroJoinSemi{\JoinOp,A,o}\land{}\\[-2pt]
        &\exists x\in A\,(x\neq o)
          \to\FranklSemi{\JoinOp,A,o}\Bigr)
      \end{aligned}}{%
      \FormulaRefAuto{FranklGlobalJoinStatement}[def]{1}}
    \proofstepwidestar[3]{\FranklFam M}{\rA}
    \proofstepwidestar[1,3]{\Frankl M}{%
      \FormulaRefAuto{FranklSemilatticeImpliesSet}[thm]{3,2}}
    \proofstepwidestar[1]{\FranklFam M\to\Frankl M}{\rRI{3,4}}
    \proofstepwidestar[1]{\forall M\,(\FranklFam M\to\Frankl M)}{\rUI{5}}
    \proofstepwidestar[1]{\FranklSetStatement}{%
      \FormulaRefAuto{FranklGlobalSetStatement}[def]{6}}
  \closeproofpartwideindR

  \proofpartwide{Schluss}
    \proofstepwidestar[]{\FranklSetStatement\to\FranklJoinStatement}{%
      \FormulaRefAuto{FranklSetToJoinGlobal}[thm]}
    \proofstepwidestar[]{\FranklJoinStatement\to\FranklSetStatement}{%
      \FormulaRefAuto{FranklJoinToSetGlobal}[thm]}
    \proofstepwidestar[]{\FranklSetStatement\leftrightarrow
      \FranklJoinStatement}{%
      \FormulaRefAuto{P \leftrightarrow Q \eqvdash
        (P \rightarrow Q)\land(Q\rightarrow P)}{\rAI{1,2}}}
  \closeproofpartwide
\end{tabproofsplitwide}

Der Äquivalenzsatz beweist die Übersetzung, nicht die Frankl-Vermutung selbst:
\(\FranklSetStatement\Longleftrightarrow\FranklJoinStatement\).
Ein künftiger Beweis oder ein Gegenbeispiel lässt sich damit ohne Verlust
zwischen beiden Sprachen übertragen.

\section{Status und Literaturhinweise}

Die Frankl-Vermutung ist mit Stand vom 18.~Juli~2026 offen; siehe
\href{https://doi.org/10.23952/jano.8.2026.1.02}%
  {van der Hout--Roos, \emph{Some Results and Conjectures Related to
  Frankl's Union Closed Conjecture}, 2026}.
Äquivalente Fassungen behandelt
\href{https://doi.org/10.1016/0097-3165(92)90068-6}%
  {Poonen, \emph{Union-closed families}, 1992}.

Auch die allgemeine Gitterform der Frankl-Vermutung ist offen.
\href{https://arxiv.org/abs/2503.00277}%
  {Bouchard (2025)} formuliert diese Fassung und leitet notwendige Bedingungen
für einen kleinsten Gegenbeispielverband her; die Arbeit beweist die allgemeine
Gitterform nicht. Für untere semimodulare Gitter ist sie bewiesen; siehe
\href{https://doi.org/10.1007/s003730050008}%
  {Reinhold (2000)}.

\begin{remark}
Der allgemeine quantitative Ersatz ist schwächer als die Hälfte:
\href{https://doi.org/10.37236/12232}%
  {Alweiss--Huang--Sellke, \emph{Improved lower bound for the union-closed
  sets conjecture}, 2024} beweisen die Schranke \((3-\sqrt5)/2\).
Die Übersetzung ergibt daraus ein \(\JoinOp\)-irreduzibles \(j\) mit
\[
  \lvert\PrincipalFilter{\JoinOrd{\JoinOp}{A},A}{j}\rvert
  \leq \frac{\sqrt5-1}{2}\,\lvert A\rvert.
\]
Dieser externe Satz wird nur eingeordnet, nicht als intern bewiesenes Theorem
registriert.
\end{remark}

\end{document}
